% Generated mechanically from the frozen divisors.tex target.
% Do not edit this generated file; edit the bound locale source instead.

% OK, start here.
%
\raggedbottom

\setcounter{chapter}{30}
\chapter{除子}



\phantomsection
\label{divisors-section-phantom}


\section{引言}
\label{divisors-section-introduction}

\noindent
本章研究定义除子等问题所涉及的一些最基本问题。
基础参考文献是 \cite{EGA}。



\section{伴随点}
\label{divisors-section-associated}

\noindent
设 $R$ 为环，$M$ 为 $R$-模。
回忆，若素理想 $\mathfrak p \subset R$ 具有如下性质，
则称它是 $M$ 的{\it 伴随素理想}：$M$ 中存在某个元素，
其零化子为 $\mathfrak p$。
参见《代数》定义 \ref{algebra-definition-associated}。
下面给出概形上拟凝聚层的伴随点定义，
该定义取自 \cite[IV Definition 3.1.1]{EGA}。

\begin{definition}
\label{divisors-definition-associated}
设 $X$ 为概形。
设 $\mathcal{F}$ 为 $X$ 上的拟凝聚层。
\begin{enumerate}
\item 称 $x \in X$ 是 $\mathcal{F}$ 的{\it 伴随点}，
如果极大理想 $\mathfrak m_x$ 是 $\mathcal{O}_{X, x}$-模
$\mathcal{F}_x$ 的伴随素理想。
\item 以 $\text{Ass}(\mathcal{F})$ 或 $\text{Ass}_X(\mathcal{F})$
表示 $\mathcal{F}$ 的伴随点集。
\item 所谓 $X$ 的{\it 伴随点}，是指 $\mathcal{O}_X$ 的伴随点。
\end{enumerate}
\end{definition}

\noindent
当 $X$ 局部 Noether 且 $\mathcal{F}$ 为有限型时，
这些定义最为有用。
例如，$X$ 的某个不可约分支的泛点可能不是 $X$ 的伴随点，见
例 \ref{divisors-example-no-associated-prime}。
在非 Noether 情形下，使用弱伴随点有时更为方便，见
第 \ref{divisors-section-weakly-associated} 节。
下面把概形论概念与仿射开集上的代数概念联系起来；
注意，这一对应只有对局部 Noether 概形才完全成立。

\begin{lemma}
\label{divisors-lemma-associated-affine-open}
设 $X$ 为概形，$\mathcal{F}$ 为 $X$ 上的拟凝聚层。
设 $\Spec(A) = U \subset X$ 为仿射开集，并置
$M = \Gamma(U, \mathcal{F})$。
设 $x \in U$，并令 $\mathfrak p \subset A$ 为与之对应的素理想。
\begin{enumerate}
\item 若 $\mathfrak p$ 是 $M$ 的伴随素理想，则 $x$ 是
$\mathcal{F}$ 的伴随点。
\item 若 $\mathfrak p$ 有限生成，则逆命题也成立。
\end{enumerate}
特别地，若 $X$ 局部 Noether，则对上述所有配对
$(\mathfrak p, x)$ 都有等价关系
$$
\mathfrak p \in \text{Ass}(M) \Leftrightarrow x \in \text{Ass}(\mathcal{F})
$$
成立。
\end{lemma}

\begin{proof}
这由《代数》引理
\ref{algebra-lemma-associated-primes-localize} 推出。
也可如下直接论证。
假设 $\mathfrak p$ 是 $M$ 的伴随素理想。
则存在 $m \in M$，其零化子为 $\mathfrak p$。
由于局部化是正合的，可见
$\mathfrak pA_{\mathfrak p}$ 是
$m/1 \in M_{\mathfrak p}$ 的零化子。又因 $M_{\mathfrak p} = \mathcal{F}_x$
（《概形》引理 \ref{schemes-lemma-spec-sheaves}），
可知 $x$ 是 $\mathcal{F}$ 的伴随点。

\medskip\noindent
反之，假设 $x$ 是 $\mathcal{F}$ 的伴随点，且
$\mathfrak p$ 有限生成。
由 $x$ 是 $\mathcal{F}$ 的伴随点，
存在 $m' \in M_{\mathfrak p}$，其零化子为
$\mathfrak pA_{\mathfrak p}$。将其写成
$m' = m/f$，其中 $f \in A$ 且 $f \not \in \mathfrak p$。
设 $I$ 为 $m$ 的零化子。它是 $A$ 的理想，并满足
$IA_{\mathfrak p} = \mathfrak pA_{\mathfrak p}$。因此
$I \subset \mathfrak p$，且 $(\mathfrak p/I)_{\mathfrak p} = 0$。
由于 $\mathfrak p$ 有限生成，存在
$g \in A$ 且 $g \not \in \mathfrak p$，使得
$g(\mathfrak p/I) = 0$。因此 $gm$ 的零化子为
$\mathfrak p$，结论得证。

\medskip\noindent
若 $X$ 局部 Noether，则 $A$ 为 Noether 环
（《性质》引理 \ref{properties-lemma-locally-Noetherian}），
因而 $\mathfrak p$ 总是有限生成的。
\end{proof}

\begin{lemma}
\label{divisors-lemma-ass-support}
设 $X$ 为概形。
设 $\mathcal{F}$ 为拟凝聚 $\mathcal{O}_X$-模。
则 $\text{Ass}(\mathcal{F}) \subset \text{Supp}(\mathcal{F})$。
\end{lemma}

\begin{proof}
这由定义立即推出。
\end{proof}

\begin{lemma}
\label{divisors-lemma-ses-ass}
设 $X$ 为概形，并设
$0 \to \mathcal{F}_1 \to \mathcal{F}_2 \to \mathcal{F}_3 \to 0$
为 $X$ 上拟凝聚层的短正合列。则
$\text{Ass}(\mathcal{F}_2) \subset
\text{Ass}(\mathcal{F}_1) \cup \text{Ass}(\mathcal{F}_3)$，
并且
$\text{Ass}(\mathcal{F}_1) \subset \text{Ass}(\mathcal{F}_2)$。
\end{lemma}

\begin{proof}
对每个点 $x \in X$，茎的序列
$0 \to \mathcal{F}_{1, x} \to \mathcal{F}_{2, x} \to \mathcal{F}_{3, x} \to 0$
是 $\mathcal{O}_{X, x}$-模的短正合列。
因此本引理由《代数》引理 \ref{algebra-lemma-ass} 推出。
\end{proof}

\begin{lemma}
\label{divisors-lemma-finite-ass}
设 $X$ 为局部 Noether 概形。
设 $\mathcal{F}$ 为凝聚 $\mathcal{O}_X$-模。
则 $\text{Ass}(\mathcal{F}) \cap U$ 对每个拟紧开集
$U \subset X$ 均为有限集。
\end{lemma}

\begin{proof}
这是因为 Noether 环上的有限模只有有限多个伴随素理想；见
《代数》引理 \ref{algebra-lemma-finite-ass}。
为把概形陈述转化为代数陈述，使用如下事实：
$U$ 是有限多个仿射开集的并；每个这样的开集都是某个 Noether 环的谱
（《性质》引理 \ref{properties-lemma-locally-Noetherian}）；
$\mathcal{F}$ 对应于该环上的有限模
（《概形上同调》引理 \ref{coherent-lemma-coherent-Noetherian}）；
最后应用引理 \ref{divisors-lemma-associated-affine-open}。
\end{proof}

\begin{lemma}
\label{divisors-lemma-ass-zero}
设 $X$ 为局部 Noether 概形，$\mathcal{F}$ 为拟凝聚
$\mathcal{O}_X$-模。则
$$
\mathcal{F} = 0 \Leftrightarrow \text{Ass}(\mathcal{F}) = \emptyset.
$$
\end{lemma}

\begin{proof}
若 $\mathcal{F} = 0$，则由定义
$\text{Ass}(\mathcal{F}) = \emptyset$。
反之，若 $\text{Ass}(\mathcal{F}) = \emptyset$，则由
《代数》引理 \ref{algebra-lemma-ass-zero} 可知 $\mathcal{F} = 0$。
为把概形陈述转化为代数陈述，只需限制到任意仿射开集并应用
引理 \ref{divisors-lemma-associated-affine-open}。
\end{proof}

\begin{example}
\label{divisors-example-no-associated-prime}
设 $k$ 为域。环 $R = k[x_1, x_2, x_3, \ldots]/(x_i^2)$
是局部环，其极大理想 $\mathfrak m$ 局部幂零。
$R$ 中不存在零化子为 $\mathfrak m$ 的元素。
因此 $\text{Ass}(R) = \emptyset$，而 $X = \Spec(R)$
给出了一个没有伴随点的概形例子。
\end{example}

\begin{lemma}
\label{divisors-lemma-restriction-injective-open-contains-ass}
设 $X$ 为局部 Noether 概形，$\mathcal{F}$ 为拟凝聚
$\mathcal{O}_X$-模。若 $U \subset X$ 为开集且
$\text{Ass}(\mathcal{F}) \subset U$，则
$\Gamma(X, \mathcal{F}) \to \Gamma(U, \mathcal{F})$ 为单射。
\end{lemma}

\begin{proof}
设 $s \in \Gamma(X, \mathcal{F})$ 为限制到 $U$ 上等于零的截面。
令 $\mathcal{F}' \subset \mathcal{F}$ 为映射
$\mathcal{O}_X \to \mathcal{F}$ 的像，其中该映射由 $s$ 定义。于是
$\text{Supp}(\mathcal{F}') \cap U = \emptyset$。另一方面，由
引理 \ref{divisors-lemma-ses-ass} 有
$\text{Ass}(\mathcal{F}') \subset \text{Ass}(\mathcal{F})$。
又有
$\text{Ass}(\mathcal{F}') \subset \text{Supp}(\mathcal{F}')$
（引理 \ref{divisors-lemma-ass-support}），故
$\text{Ass}(\mathcal{F}') = \emptyset$。
最后由引理 \ref{divisors-lemma-ass-zero} 得 $\mathcal{F}' = 0$。
\end{proof}

\begin{lemma}
\label{divisors-lemma-minimal-support-in-ass}
设 $X$ 为局部 Noether 概形。
设 $\mathcal{F}$ 为拟凝聚 $\mathcal{O}_X$-模。
设 $x \in \text{Supp}(\mathcal{F})$ 为 $\mathcal{F}$ 的支撑中的点，
且它不是 $\text{Supp}(\mathcal{F})$ 中任何其他点的特殊化。
则 $x \in \text{Ass}(\mathcal{F})$。
特别地，$X$ 的每个不可约分支的泛点都是 $X$ 的伴随点。
\end{lemma}

\begin{proof}
由于 $x \in \text{Supp}(\mathcal{F})$，模 $\mathcal{F}_x$
非零。因此
$\text{Ass}(\mathcal{F}_x) \subset \Spec(\mathcal{O}_{X, x})$
由《代数》引理 \ref{algebra-lemma-ass-zero} 非空。
另一方面，由假设
$\text{Supp}(\mathcal{F}_x) = \{\mathfrak m_x\}$。
又因
$\text{Ass}(\mathcal{F}_x) \subset \text{Supp}(\mathcal{F}_x)$
（《代数》引理 \ref{algebra-lemma-ass-support}），
可知 $\mathfrak m_x$ 是 $\mathcal{F}_x$ 的伴随素理想，
结论得证。
\end{proof}

\noindent
下一个引理对应于《代数进阶》引理
\ref{more-algebra-lemma-check-injective-on-ass}。

\begin{lemma}
\label{divisors-lemma-check-injective-on-ass}
设 $X$ 为局部 Noether 概形，并设
$\varphi : \mathcal{F} \to \mathcal{G}$ 为拟凝聚
$\mathcal{O}_X$-模之间的映射。
假设对每个 $x \in X$，以下至少一种情形成立：
\begin{enumerate}
\item $\mathcal{F}_x \to \mathcal{G}_x$ 为单射；或
\item $x \not \in \text{Ass}(\mathcal{F})$。
\end{enumerate}
则 $\varphi$ 为单射。
\end{lemma}

\begin{proof}
这些假设蕴含 $\text{Ass}(\Ker(\varphi)) = \emptyset$，
故由引理 \ref{divisors-lemma-ass-zero} 有 $\Ker(\varphi) = 0$。
\end{proof}

\begin{lemma}
\label{divisors-lemma-check-isomorphism-via-depth-and-ass}
设 $X$ 为局部 Noether 概形，并设
$\varphi : \mathcal{F} \to \mathcal{G}$ 为拟凝聚
$\mathcal{O}_X$-模之间的映射。假设 $\mathcal{F}$ 凝聚，
并且对每个 $x \in X$，以下一种情形成立：
\begin{enumerate}
\item $\mathcal{F}_x \to \mathcal{G}_x$ 为同构；或
\item $\text{depth}(\mathcal{F}_x) \geq 2$ 且
$x \not \in \text{Ass}(\mathcal{G})$。
\end{enumerate}
则 $\varphi$ 为同构。
\end{lemma}

\begin{proof}
这只是把《代数进阶》引理
\ref{more-algebra-lemma-check-isomorphism-via-depth-and-ass}
改写成概形语言。
\end{proof}

\section{态射与伴随点}
\label{divisors-section-morphisms-associated}

\noindent
设 $f : X \to S$ 为概形的态射。
设 $\mathcal{F}$ 为 $\mathcal{O}_X$-模层。
若 $s \in S$ 是一点，则常用 $\mathcal{F}_s$ 表示下述
$\mathcal{O}_{X_s}$-模：它是把 $\mathcal{F}$ 沿态射
$i_s : X_s \to X$ 拉回所得的模。
这里 $X_s$ 是 $f$ 在 $s$ 上的概形论纤维。用公式表示即
$$
\mathcal{F}_s = i_s^*\mathcal{F}
$$
当然，这一记号会与 $\mathcal{F}$ 在点 $x \in X$ 处的茎所用的
既有记号冲突，这里指 $f = \text{id}_X$ 的情形。不过，这一记号通常很方便，
如下述引理的表述所示。

\begin{lemma}
\label{divisors-lemma-bourbaki}
设 $f : X \to S$ 为概形的态射。
设 $\mathcal{F}$ 为 $X$ 上在 $S$ 上平坦的拟凝聚层。
设 $\mathcal{G}$ 为 $S$ 上的拟凝聚层。则有
$$
\text{Ass}_X(\mathcal{F} \otimes_{\mathcal{O}_X} f^*\mathcal{G})
\supset
\bigcup\nolimits_{s \in \text{Ass}_S(\mathcal{G})}
\text{Ass}_{X_s}(\mathcal{F}_s)
$$
而当 $S$ 局部 Noether 时等号成立（记号 $\mathcal{F}_s$ 见上文）。
\end{lemma}

\begin{proof}
设 $x \in X$，并令 $s = f(x) \in S$。
置 $B = \mathcal{O}_{X, x}$、$A = \mathcal{O}_{S, s}$、
$N = \mathcal{F}_x$ 和 $M = \mathcal{G}_s$。
注意，$\mathcal{F} \otimes_{\mathcal{O}_X} f^*\mathcal{G}$ 在
$x$ 处的茎等于 $B$-模 $M \otimes_A N$。因此，
$x \in \text{Ass}_X(\mathcal{F} \otimes_{\mathcal{O}_X} f^*\mathcal{G})$
当且仅当 $\mathfrak m_B$ 属于 $\text{Ass}_B(M \otimes_A N)$。
类似地，$s \in \text{Ass}_S(\mathcal{G})$ 且
$x \in \text{Ass}_{X_s}(\mathcal{F}_s)$，当且仅当
$\mathfrak m_A \in \text{Ass}_A(M)$ 且
$\mathfrak m_B/\mathfrak m_A B \in
\text{Ass}_{B \otimes \kappa(\mathfrak m_A)}(N \otimes \kappa(\mathfrak m_A))$。
于是本引理由《代数》引理 \ref{algebra-lemma-bourbaki-fibres} 得出。
\end{proof}





\section{嵌入点}
\label{divisors-section-embedded}

\noindent
设 $R$ 为环，$M$ 为 $R$-模。
回忆，若素理想 $\mathfrak p \subset R$ 是 $M$ 的伴随素理想，
但在 $M$ 的伴随素理想中不是极小的，则称它为 $M$ 的
{\it 嵌入伴随素理想}。参见《代数》定义
\ref{algebra-definition-embedded-primes}。
下面给出概形上拟凝聚层的嵌入伴随点定义，
该定义取自 \cite[IV Definition 3.1.1]{EGA}。

\begin{definition}
\label{divisors-definition-embedded}
设 $X$ 为概形。
设 $\mathcal{F}$ 为 $X$ 上的拟凝聚层。
\begin{enumerate}
\item 称 $\mathcal{F}$ 的某个伴随点为{\it 嵌入伴随点}，
如果它在 $\mathcal{F}$ 的伴随点中不是极大的；换言之，
它是 $\mathcal{F}$ 的另一个伴随点的特化。
\item 设点 $x$ 所在的概形为 $X$。若 $x$ 是
$\mathcal{O}_X$ 的嵌入伴随点，则称其为{\it 嵌入点}。
\item $X$ 的一个{\it 嵌入分支}是不可约闭子集
$Z = \overline{\{x\}}$，其中 $x$ 是 $X$ 的嵌入点。
\end{enumerate}
\end{definition}

\noindent
在 Noether 情形下，当 $\mathcal{F}$ 凝聚时，有下述结论。

\begin{lemma}
\label{divisors-lemma-embedded}
设 $X$ 为局部 Noether 概形。
设 $\mathcal{F}$ 为凝聚 $\mathcal{O}_X$-模。则
\begin{enumerate}
\item $\text{Supp}(\mathcal{F})$ 的不可约分支的泛点是
$\mathcal{F}$ 的伴随点；
\item $\mathcal{F}$ 的伴随点是嵌入伴随点，当且仅当它不是
$\text{Supp}(\mathcal{F})$ 的某个不可约分支的泛点。
\end{enumerate}
特别地，$X$ 的嵌入点就是 $X$ 的伴随点中那些不是 $X$ 的
不可约分支泛点的点。
\end{lemma}

\begin{proof}
回忆，在此情形下 $Z = \text{Supp}(\mathcal{F})$ 是闭集，见
《态射》引理 \ref{morphisms-lemma-support-finite-type}；并且 $Z$
的不可约分支的泛点是 $\mathcal{F}$ 的伴随点，见引理
\ref{divisors-lemma-minimal-support-in-ass}。最后，由引理
\ref{divisors-lemma-ass-support} 有 $\text{Ass}(\mathcal{F}) \subset Z$。
结合这些结论以及 $Z$ 是 sober 拓扑空间这一事实（因而 $Z$ 的
每个点都是 $Z$ 的某个泛点的特化），便得到 (1) 和 (2)。
\end{proof}

\begin{lemma}
\label{divisors-lemma-S1-no-embedded}
设 $X$ 为局部 Noether 概形。
设 $\mathcal{F}$ 为 $X$ 上的凝聚层。则以下条件等价：
\begin{enumerate}
\item $\mathcal{F}$ 没有嵌入伴随点；
\item $\mathcal{F}$ 具有性质 $(S_1)$。
\end{enumerate}
\end{lemma}

\begin{proof}
这是《代数》引理 \ref{algebra-lemma-criterion-no-embedded-primes}
与上面的引理 \ref{divisors-lemma-associated-affine-open} 的结合。
\end{proof}

\begin{lemma}
\label{divisors-lemma-noetherian-dim-1-CM-no-embedded-points}
设 $X$ 为维数 $\leq 1$ 的局部 Noether 概形。以下条件等价：
\begin{enumerate}
\item $X$ 是 Cohen-Macaulay 概形；
\item $X$ 没有嵌入点。
\end{enumerate}
\end{lemma}

\begin{proof}
由引理 \ref{divisors-lemma-S1-no-embedded} 及各定义立即得到。
\end{proof}

\begin{lemma}
\label{divisors-lemma-scheme-theoretically-dense-contain-embedded-points}
设 $X$ 为局部 Noether 概形，并设 $U \subset X$ 为开子概形。
以下条件等价：
\begin{enumerate}
\item $U$ 在 $X$ 中概形论稠密
（《态射》定义 \ref{morphisms-definition-scheme-theoretically-dense}）；
\item $U$ 在 $X$ 中稠密，且 $U$ 包含 $X$ 的所有嵌入点；
\item $U$ 包含 $X$ 的所有伴随点。
\end{enumerate}
\end{lemma}

\begin{proof}
问题在 $X$ 上是局部的，故可假设 $X = \Spec(A)$，其中 $A$
为 Noether 环。此时 $U$ 拟紧（《性质》引理
\ref{properties-lemma-immersion-into-noetherian}），因而
$U = D(f_1) \cup \ldots \cup D(f_n)$
（《代数》引理 \ref{algebra-lemma-qc-open}）。在此情形下，
$U$ 在 $X$ 中概形论稠密，当且仅当
$A \to A_{f_1} \times \ldots \times A_{f_n}$ 是单射，见《态射》例
\ref{morphisms-example-scheme-theoretic-closure}。
条件 (3) 翻译为代数语言就是说：对每个伴随素理想
$\mathfrak p$，若它是 $A$ 的伴随素理想，则都存在某个 $i$
使得 $f_i \not \in \mathfrak p$。

\medskip\noindent
假设 (1)，即 $A \to A_{f_1} \times \ldots \times A_{f_n}$ 是单射。
设 $a \in A$ 的零化子是素理想 $\mathfrak p$。由于 $a$ 在某个
$A_{f_i}$ 中的像非零，便知对某个 $i$ 有
$f_i \not \in \mathfrak p$。故 (3) 成立。

\medskip\noindent
注意，$U$ 在 $\Spec(A)$ 中稠密，当且仅当它包含所有极小素理想。
由于 $A$ 的伴随素理想之集等于极小素理想之集与嵌入伴随素理想
之集的并（见《代数》命题
\ref{algebra-proposition-minimal-primes-associated-primes}），
可知 (2) 与 (3) 等价。

\medskip\noindent
假设 (3)。这表示每个伴随素理想 $\mathfrak p$，如果它是 $A$
的伴随素理想，便对应于某个 $A_{f_i}$ 的一个素理想，其中 $i$
可适当选取。因此
$A \to A_{f_1} \times \ldots \times A_{f_n}$ 是单射，因为由
《代数》引理 \ref{algebra-lemma-zero-at-ass-zero}，映射
$A \to \prod_{\mathfrak p \in \text{Ass}(A)} A_\mathfrak p$ 是单射。
由此可见 (3) 蕴含 (1)。
\end{proof}

\begin{lemma}
\label{divisors-lemma-remove-embedded-points}
设 $X$ 为局部 Noether 概形。
设 $\mathcal{F}$ 为 $X$ 上的凝聚层。满足下述条件的凝聚子层之集
$$
\{
\mathcal{K} \subset \mathcal{F}
\mid
\text{Supp}(\mathcal{K})\text{ is nowhere dense in }\text{Supp}(\mathcal{F})
\}
$$
有一个极大元 $\mathcal{K}$。令
$\mathcal{F}' = \mathcal{F}/\mathcal{K}$，则有
\begin{enumerate}
\item $\text{Supp}(\mathcal{F}') = \text{Supp}(\mathcal{F})$；
\item $\mathcal{F}'$ 没有嵌入伴随点；
\item 存在稠密开集 $U \subset X$，使得
$U \cap \text{Supp}(\mathcal{F})$ 在 $\text{Supp}(\mathcal{F})$
中稠密，且 $\mathcal{F}'|_U \cong \mathcal{F}|_U$。
\end{enumerate}
\end{lemma}

\begin{proof}
这由《代数》引理 \ref{algebra-lemma-remove-embedded-primes} 和
\ref{algebra-lemma-remove-embedded-primes-localize} 得出。
注意，$U$ 可取为 $\mathcal{F}$ 的嵌入伴随点之集的闭包的补集。
\end{proof}

\begin{lemma}
\label{divisors-lemma-no-embedded-points-endos}
设 $X$ 为局部 Noether 概形。
设 $\mathcal{F}$ 为没有嵌入伴随点的凝聚 $\mathcal{O}_X$-模。置
$$
\mathcal{I}
=
\Ker(\mathcal{O}_X
\longrightarrow
\SheafHom_{\mathcal{O}_X}(\mathcal{F}, \mathcal{F})).
$$
这是一个凝聚理想层，并定义了没有嵌入点的闭子概形
$Z \subset X$。此外，存在凝聚层 $\mathcal{G}$，它定义在 $Z$ 上，并且
(a) $\mathcal{F} = (Z \to X)_*\mathcal{G}$，
(b) $\mathcal{G}$ 没有嵌入伴随点，且
(c) $\text{Supp}(\mathcal{G}) = Z$（作为集合）。
\end{lemma}

\begin{proof}
其中一些断言已在《概形上同调》引理
\ref{coherent-lemma-coherent-support-closed} 的证明中见过，
其余断言由《代数》引理
\ref{algebra-lemma-no-embedded-primes-endos} 得出。
\end{proof}



\section{弱伴随点}
\label{divisors-section-weakly-associated}

\noindent
设 $R$ 为环，$M$ 为 $R$-模。
回忆，称素理想 $\mathfrak p \subset R$ 与 $M$ {\it 弱伴随}，
如果存在一个元素 $m$，它属于 $M$，并且 $\mathfrak p$
在包含 $m$ 的零化子的素理想中极小。参见《代数》定义
\ref{algebra-definition-weakly-associated}。
若 $R$ 是极大理想为 $\mathfrak m$ 的局部环，则
$\mathfrak m$ 与 $M$ 弱伴随，当且仅当存在元素 $m \in M$，
其零化子的根理想为 $\mathfrak m$，见《代数》引理
\ref{algebra-lemma-weakly-ass-local}。

\begin{definition}
\label{divisors-definition-weakly-associated}
设 $X$ 为概形。
设 $\mathcal{F}$ 为 $X$ 上的拟凝聚层。
\begin{enumerate}
\item 称 $x \in X$ 与 $\mathcal{F}$ {\it 弱伴随}，如果极大理想
$\mathfrak m_x$ 与 $\mathcal{O}_{X, x}$-模 $\mathcal{F}_x$ 弱伴随。
\item 用 $\text{WeakAss}(\mathcal{F})$ 表示 $\mathcal{F}$ 的
弱伴随点之集。
\item $X$ 的{\it 弱伴随点}是 $\mathcal{O}_X$ 的弱伴随点。
\end{enumerate}
\end{definition}

\noindent
在这种情形下，在任意仿射开集上，这一概念恰好对应于上面定义的
弱伴随素理想。精确表述如下。

\begin{lemma}
\label{divisors-lemma-weakly-associated-affine-open}
设 $X$ 为概形，$\mathcal{F}$ 为 $X$ 上的拟凝聚层。
设 $\Spec(A) = U \subset X$ 为仿射开集，并置
$M = \Gamma(U, \mathcal{F})$。
设 $x \in U$，并设 $\mathfrak p \subset A$ 为对应的素理想。
以下条件等价：
\begin{enumerate}
\item $\mathfrak p$ 与 $M$ 弱伴随；
\item $x$ 与 $\mathcal{F}$ 弱伴随。
\end{enumerate}
\end{lemma}

\begin{proof}
这由《代数》引理 \ref{algebra-lemma-weakly-ass-local} 得出。
\end{proof}

\begin{lemma}
\label{divisors-lemma-weakly-ass-support}
设 $X$ 为概形。
设 $\mathcal{F}$ 为拟凝聚 $\mathcal{O}_X$-模。则
$$
\text{Ass}(\mathcal{F}) \subset \text{WeakAss}(\mathcal{F}) \subset
\text{Supp}(\mathcal{F}).
$$
\end{lemma}

\begin{proof}
这由定义立即得到。
\end{proof}

\begin{lemma}
\label{divisors-lemma-ses-weakly-ass}
设 $X$ 为概形。设
$0 \to \mathcal{F}_1 \to \mathcal{F}_2 \to \mathcal{F}_3 \to 0$
为 $X$ 上拟凝聚层的短正合列。则
$\text{WeakAss}(\mathcal{F}_2) \subset
\text{WeakAss}(\mathcal{F}_1) \cup \text{WeakAss}(\mathcal{F}_3)$，且
$\text{WeakAss}(\mathcal{F}_1) \subset \text{WeakAss}(\mathcal{F}_2)$。
\end{lemma}

\begin{proof}
对每个点 $x \in X$，茎的序列
$0 \to \mathcal{F}_{1, x} \to \mathcal{F}_{2, x} \to \mathcal{F}_{3, x} \to 0$
是 $\mathcal{O}_{X, x}$-模的短正合列。因此，本引理由《代数》引理
\ref{algebra-lemma-weakly-ass} 得出。
\end{proof}

\begin{lemma}
\label{divisors-lemma-weakly-ass-zero}
设 $X$ 为概形。
设 $\mathcal{F}$ 为拟凝聚 $\mathcal{O}_X$-模。则
$$
\mathcal{F} = (0) \Leftrightarrow \text{WeakAss}(\mathcal{F}) = \emptyset
$$
\end{lemma}

\begin{proof}
由引理 \ref{divisors-lemma-weakly-associated-affine-open} 和《代数》引理
\ref{algebra-lemma-weakly-ass-zero} 得出。
\end{proof}

\begin{lemma}
\label{divisors-lemma-restriction-injective-open-contains-weakly-ass}
设 $X$ 为概形，$\mathcal{F}$ 为拟凝聚 $\mathcal{O}_X$-模。
若 $U \subset X$ 是开集，且
$\text{WeakAss}(\mathcal{F}) \subset U$，则映射
$\Gamma(X, \mathcal{F}) \to \Gamma(U, \mathcal{F})$ 是单射。
\end{lemma}

\begin{proof}
设 $s \in \Gamma(X, \mathcal{F})$ 为一个限制到 $U$ 上为零的截面。
令 $\mathcal{F}' \subset \mathcal{F}$ 为映射
$\mathcal{O}_X \to \mathcal{F}$ 的像，该映射由 $s$ 定义。于是
$\text{Supp}(\mathcal{F}') \cap U = \emptyset$。另一方面，由引理
\ref{divisors-lemma-ses-weakly-ass}，有
$\text{WeakAss}(\mathcal{F}') \subset \text{WeakAss}(\mathcal{F})$。
又有 $\text{WeakAss}(\mathcal{F}') \subset \text{Supp}(\mathcal{F}')$
（引理 \ref{divisors-lemma-weakly-ass-support}），故
$\text{WeakAss}(\mathcal{F}') = \emptyset$。
再由引理 \ref{divisors-lemma-weakly-ass-zero} 得 $\mathcal{F}' = 0$。
\end{proof}

\begin{lemma}
\label{divisors-lemma-minimal-support-in-weakly-ass}
设 $X$ 为概形。
设 $\mathcal{F}$ 为拟凝聚 $\mathcal{O}_X$-模。
设 $x \in \text{Supp}(\mathcal{F})$ 是 $\mathcal{F}$ 的支撑中的一点，
并且它不是 $\text{Supp}(\mathcal{F})$ 中另一个点的特化。则
$x \in \text{WeakAss}(\mathcal{F})$。
特别地，$X$ 的任一不可约分支的泛点都与 $\mathcal{O}_X$ 弱伴随。
\end{lemma}

\begin{proof}
由于 $x \in \text{Supp}(\mathcal{F})$，模 $\mathcal{F}_x$ 非零。
因此，由《代数》引理 \ref{algebra-lemma-weakly-ass-zero}，集合
$\text{WeakAss}(\mathcal{F}_x) \subset \Spec(\mathcal{O}_{X, x})$
非空。另一方面，根据假设，
$\text{Supp}(\mathcal{F}_x) = \{\mathfrak m_x\}$。
又因为
$\text{WeakAss}(\mathcal{F}_x) \subset \text{Supp}(\mathcal{F}_x)$
（《代数》引理 \ref{algebra-lemma-weakly-ass-support}），可见
$\mathfrak m_x$ 与 $\mathcal{F}_x$ 弱伴随，结论成立。
\end{proof}

\begin{lemma}
\label{divisors-lemma-ass-weakly-ass}
设 $X$ 为概形。
设 $\mathcal{F}$ 为拟凝聚 $\mathcal{O}_X$-模。
若 $\mathfrak m_x$ 是 $\mathcal{O}_{X, x}$ 的有限生成理想，则
$$
x \in \text{Ass}(\mathcal{F}) \Leftrightarrow
x \in \text{WeakAss}(\mathcal{F}).
$$
特别地，若 $X$ 局部 Noether，则
$\text{Ass}(\mathcal{F}) = \text{WeakAss}(\mathcal{F})$。
\end{lemma}

\begin{proof}
见《代数》引理 \ref{algebra-lemma-ass-weakly-ass}。
\end{proof}

\begin{lemma}
\label{divisors-lemma-weakass-pushforward}
设 $f : X \to S$ 为拟紧拟分离的概形态射。
设 $\mathcal{F}$ 为拟凝聚 $\mathcal{O}_X$-模。
设 $s \in S$ 是一个不属于 $f$ 的像的点。则
$s$ 不与 $f_*\mathcal{F}$ 弱伴随。
\end{lemma}

\begin{proof}
考虑基变换 $f' : X' \to \Spec(\mathcal{O}_{S, s})$；它由 $f$
沿态射 $g : \Spec(\mathcal{O}_{S, s}) \to S$ 得到。记另一个投影为
$g' : X' \to X$。于是
$$
(f_*\mathcal{F})_s = (g^*f_*\mathcal{F})_s = (f'_*(g')^*\mathcal{F})_s
$$
第一个等式成立，是因为 $g$ 在 $s$ 处的局部环上诱导同构；
第二个等式由平坦基变换得到（《概形上同调》引理
\ref{coherent-lemma-flat-base-change-cohomology}）。当然，
$s \in \Spec(\mathcal{O}_{S, s})$ 不属于 $f'$ 的像。
因此可假设 $S$ 是局部环 $(A, \mathfrak m)$ 的谱，并且 $s$
对应于 $\mathfrak m$。由《概形》引理
\ref{schemes-lemma-push-forward-quasi-coherent}，层
$f_*\mathcal{F}$ 拟凝聚，设它对应于 $A$-模 $M$。
由于 $s$ 不属于 $f$ 的像，可见
$X = \bigcup_{a \in \mathfrak m} f^{-1}D(a)$ 是一个开覆盖。
因为 $X$ 拟紧，可找到 $a_1, \ldots, a_n \in \mathfrak m$，使得
$X = f^{-1}D(a_1) \cup \ldots \cup f^{-1}D(a_n)$。由此
$$
M \to M_{a_1} \oplus \ldots \oplus M_{a_r}
$$
是单射。因此，对任意非零元素 $m$，若把它视作茎 $M_\mathfrak p$ 的元素，
存在一个 $i$，使得 $a_i^n m$ 对所有 $n \geq 0$ 都非零。
故 $\mathfrak m$ 不与 $M$ 弱伴随。
\end{proof}

\begin{lemma}
\label{divisors-lemma-check-injective-on-weakass}
设 $X$ 为概形。设 $\varphi : \mathcal{F} \to \mathcal{G}$ 为拟凝聚
$\mathcal{O}_X$-模之间的映射。假设对每个 $x \in X$，
以下至少一种情形成立：
\begin{enumerate}
\item $\mathcal{F}_x \to \mathcal{G}_x$ 是单射；
\item $x \not \in \text{WeakAss}(\mathcal{F})$。
\end{enumerate}
则 $\varphi$ 是单射。
\end{lemma}

\begin{proof}
这些假设蕴含 $\text{WeakAss}(\Ker(\varphi)) = \emptyset$，
故由引理 \ref{divisors-lemma-weakly-ass-zero} 有 $\Ker(\varphi) = 0$。
\end{proof}

\begin{lemma}
\label{divisors-lemma-depth-2-hartog}
设 $X$ 为局部 Noether 概形，$\mathcal{F}$ 为凝聚
$\mathcal{O}_X$-模。设 $j : U \to X$ 为开子概形，并且对
$x \in X \setminus U$ 有 $\text{depth}(\mathcal{F}_x) \geq 2$。则
$$
\mathcal{F} \longrightarrow j_*(\mathcal{F}|_U)
$$
是同构，从而
$\Gamma(X, \mathcal{F}) \to \Gamma(U, \mathcal{F})$ 也是同构。
\end{lemma}

\begin{proof}
我们断言，引理 \ref{divisors-lemma-check-isomorphism-via-depth-and-ass}
可用于本引理中展示的映射。
设 $x \in X$。若 $x \in U$，则该映射在茎上是同构，因为
$j_*(\mathcal{F}|_U)|_U = \mathcal{F}|_U$。
若 $x \in X \setminus U$，则
$x \not \in \text{Ass}(j_*(\mathcal{F}|_U))$
（引理 \ref{divisors-lemma-weakass-pushforward} 和
\ref{divisors-lemma-weakly-ass-support}）。再结合假设
$\text{depth}(\mathcal{F}_x) \geq 2$，证明即告完成。
\end{proof}

\begin{lemma}
\label{divisors-lemma-weakass-reduced}
设 $X$ 为约化概形。则 $X$ 的弱伴随点恰好是 $X$ 的
不可约分支的泛点。
\end{lemma}

\begin{proof}
由《代数》引理 \ref{algebra-lemma-reduced-weakly-ass-minimal} 得出。
\end{proof}



\section{态射与弱伴随点}
\label{divisors-section-morphisms-weakly-associated}

\begin{lemma}
\label{divisors-lemma-weakly-ass-reverse-functorial}
设 $f : X \to S$ 为概形的仿射态射。
设 $\mathcal{F}$ 为拟凝聚 $\mathcal{O}_X$-模。则
$$
\text{WeakAss}_S(f_*\mathcal{F}) \subset f(\text{WeakAss}_X(\mathcal{F}))
$$
\end{lemma}

\begin{proof}
可以假设 $X$ 和 $S$ 仿射，故 $X \to S$ 来自环映射
$A \to B$。于是 $\mathcal{F} = \widetilde M$，这里所用的是某个
$B$-模 $M$。由引理 \ref{divisors-lemma-weakly-associated-affine-open}，
$\mathcal{F}$ 的弱伴随点恰好对应于 $M$ 的弱伴随素理想。
类似地，$f_*\mathcal{F}$ 的弱伴随点恰好对应于把 $M$ 看作
$A$-模时的弱伴随素理想。因此，本引理由《代数》引理
\ref{algebra-lemma-weakly-ass-reverse-functorial} 得出。
\end{proof}

\begin{lemma}
\label{divisors-lemma-ass-functorial-equal}
设 $f : X \to S$ 为概形的仿射态射。
设 $\mathcal{F}$ 为拟凝聚 $\mathcal{O}_X$-模。
若 $X$ 局部 Noether，则
$$
f(\text{Ass}_X(\mathcal{F})) =
\text{Ass}_S(f_*\mathcal{F}) =
\text{WeakAss}_S(f_*\mathcal{F}) =
f(\text{WeakAss}_X(\mathcal{F}))
$$
\end{lemma}

\begin{proof}
可以假设 $X$ 和 $S$ 仿射，故 $X \to S$ 来自环映射
$A \to B$。由于 $X$ 局部 Noether，环 $B$ 是 Noether 环，见
《性质》引理 \ref{properties-lemma-locally-Noetherian}。
写 $\mathcal{F} = \widetilde M$，这里所用的是某个 $B$-模 $M$。
由引理 \ref{divisors-lemma-associated-affine-open}，$\mathcal{F}$ 的伴随点
恰好对应于 $M$ 的伴随素理想；而把 $M$ 看作 $A$-模时的任一
伴随素理想，都是 $f_*\mathcal{F}$ 的一个伴随点。因此，由《代数》引理
\ref{algebra-lemma-ass-functorial-Noetherian} 得到包含关系
$$
f(\text{Ass}_X(\mathcal{F})) \subset \text{Ass}_S(f_*\mathcal{F})
$$
由引理 \ref{divisors-lemma-weakly-ass-support}，有包含关系
$$
\text{Ass}_S(f_*\mathcal{F}) \subset \text{WeakAss}_S(f_*\mathcal{F})
$$
由引理 \ref{divisors-lemma-weakly-ass-reverse-functorial}，有包含关系
$$
\text{WeakAss}_S(f_*\mathcal{F}) \subset f(\text{WeakAss}_X(\mathcal{F}))
$$
由引理 \ref{divisors-lemma-ass-weakly-ass}，最外侧的两个集合相等，
故处处均为等号。
\end{proof}

\begin{lemma}
\label{divisors-lemma-weakly-associated-finite}
设 $f : X \to S$ 为概形的有限态射。
设 $\mathcal{F}$ 为拟凝聚 $\mathcal{O}_X$-模。则
$\text{WeakAss}(f_*\mathcal{F}) = f(\text{WeakAss}(\mathcal{F}))$。
\end{lemma}

\begin{proof}
可以假设 $X$ 和 $S$ 仿射，故 $X \to S$ 来自有限环映射
$A \to B$。写 $\mathcal{F} = \widetilde M$，这里所用的是某个
$B$-模 $M$。由引理 \ref{divisors-lemma-weakly-associated-affine-open}，
$\mathcal{F}$ 的弱伴随点恰好对应于 $M$ 的弱伴随素理想。
类似地，$f_*\mathcal{F}$ 的弱伴随点恰好对应于把 $M$ 看作
$A$-模时的弱伴随素理想。因此，本引理由《代数》引理
\ref{algebra-lemma-weakly-ass-finite-ring-map} 得出。
\end{proof}

\begin{lemma}
\label{divisors-lemma-weakly-ass-pullback}
设 $f : X \to S$ 为概形的态射，$\mathcal{G}$ 为拟凝聚
$\mathcal{O}_S$-模。设 $x \in X$ 且 $s = f(x)$。
若 $f$ 在 $x$ 处平坦，点 $x$ 是纤维 $X_s$ 的泛点，并且
$s \in \text{WeakAss}_S(\mathcal{G})$，则
$x \in \text{WeakAss}(f^*\mathcal{G})$。
\end{lemma}

\begin{proof}
置 $A = \mathcal{O}_{S, s}$、$B = \mathcal{O}_{X, x}$ 和
$M = \mathcal{G}_s$。设 $m \in M$ 为一个元素，其零化子
$I = \{a \in A \mid am = 0\}$ 的根理想为 $\mathfrak m_A$。
$m \otimes 1$ 的零化子为 $I B$，因为 $A \to B$ 忠实平坦。
因此，只须证明 $\sqrt{I B} = \mathfrak m_B$。这由以下两点得出：
$B/\mathfrak m_AB$ 的极大理想局部幂零（见《代数》引理
\ref{algebra-lemma-minimal-prime-reduced-ring}），以及假设
$\sqrt{I} = \mathfrak m_A$。略去一些细节。
\end{proof}

\begin{lemma}
\label{divisors-lemma-weakly-ass-change-fields}
设 $K/k$ 为域扩张，$X$ 为 $k$ 上的概形。
设 $\mathcal{F}$ 为拟凝聚 $\mathcal{O}_X$-模。
设 $y \in X_K$ 的像为 $x \in X$。若 $y$ 是拉回
$\mathcal{F}_K$ 的弱伴随点，则 $x$ 是 $\mathcal{F}$ 的弱伴随点。
\end{lemma}

\begin{proof}
这只是把《代数》引理 \ref{algebra-lemma-weakly-ass-change-fields}
改写成概形语言。
\end{proof}

\noindent
下面是一个简单引理，它表明直像常常具有至少为 2 的深度。

\begin{lemma}
\label{divisors-lemma-depth-pushforward}
设 $f : X \to S$ 为拟紧拟分离的概形态射。
设 $\mathcal{F}$ 为拟凝聚 $\mathcal{O}_X$-模。
设 $s \in S$。
\begin{enumerate}
\item 若 $s \not \in f(X)$，则 $s$ 不与 $f_*\mathcal{F}$ 弱伴随。
\item 若 $s \not \in f(X)$ 且 $\mathcal{O}_{S, s}$ 是 Noether 环，
则 $s$ 不是 $f_*\mathcal{F}$ 的伴随点。
\item 若 $s \not \in f(X)$，$(f_*\mathcal{F})_s$ 是有限
$\mathcal{O}_{S, s}$-模，并且 $\mathcal{O}_{S, s}$ 是 Noether 环，
则 $\text{depth}((f_*\mathcal{F})_s) \geq 2$。
\item 若 $\mathcal{F}$ 在 $S$ 上平坦，并且 $a \in \mathfrak m_s$
是非零因子，则 $a$ 在 $(f_*\mathcal{F})_s$ 上是非零因子。
\item 若 $\mathcal{F}$ 在 $S$ 上平坦，并且
$a, b \in \mathfrak m_s$ 是正则序列，则 $a$ 在
$(f_*\mathcal{F})_s$ 上是非零因子，且 $b$ 在
$(f_*\mathcal{F})_s/a(f_*\mathcal{F})_s$ 上是非零因子。
\item 若 $\mathcal{F}$ 在 $S$ 上平坦，并且 $(f_*\mathcal{F})_s$
是有限 $\mathcal{O}_{S, s}$-模，则
$\text{depth}((f_*\mathcal{F})_s) \geq
\min(2, \text{depth}(\mathcal{O}_{S, s}))$。
\end{enumerate}
\end{lemma}

\begin{proof}
(1) 是引理 \ref{divisors-lemma-weakass-pushforward}。
(2) 由 (1) 和引理 \ref{divisors-lemma-ass-weakly-ass} 得出。

{\small
\medskip\noindent
证明 (3)。要证明深度 $\geq 2$，只须证明
$\Hom_{\mathcal{O}_{S, s}}(\kappa(s), (f_*\mathcal{F})_s) = 0$ 及
$\Ext^1_{\mathcal{O}_{S, s}}(\kappa(s), (f_*\mathcal{F})_s) = 0$，
见《代数》引理 \ref{algebra-lemma-depth-ext}。\par}
利用正合列
$0 \to \mathfrak m_s \to \mathcal{O}_{S, s} \to \kappa(s) \to 0$，
只须证明映射
$$
\Hom_{\mathcal{O}_{S, s}}(\mathcal{O}_{S, s}, (f_*\mathcal{F})_s)
\to
\Hom_{\mathcal{O}_{S, s}}(\mathfrak m_s, (f_*\mathcal{F})_s)
$$
是同构。由平坦基变换（《概形上同调》引理
\ref{coherent-lemma-flat-base-change-cohomology}），可以把 $S$ 替换为
$\Spec(\mathcal{O}_{S, s})$，并把 $X$ 替换为
$\Spec(\mathcal{O}_{S, s}) \times_S X$。
用 $\mathfrak m \subset \mathcal{O}_S$ 表示 $s$ 的理想层。于是
$$
\Hom_{\mathcal{O}_{S, s}}(\mathfrak m_s, (f_*\mathcal{F})_s) =
\Hom_{\mathcal{O}_S}(\mathfrak m, f_*\mathcal{F}) =
\Hom_{\mathcal{O}_X}(f^*\mathfrak m, \mathcal{F})
$$
第一个等式成立，是因为 $S$ 是以 $s$ 为闭点的局部概形；
第二个等式由拟凝聚模上 $f^*, f_*$ 的伴随性得到。然而，由于
$s \not \in f(X)$，可见 $f^*\mathfrak m = \mathcal{O}_X$。
沿上述论证逆推，便得到所需等式。

\medskip\noindent
为证明 (4)、(5) 和 (6)，利用平坦基变换（《概形上同调》引理
\ref{coherent-lemma-flat-base-change-cohomology}），可约化到 $S$ 是
$\mathcal{O}_{S, s}$ 的谱的情形。此时非零因子
$a \in \mathcal{O}_{S, s}$ 确定短正合列
$$
0 \to \mathcal{O}_S \xrightarrow{a} \mathcal{O}_S \to
\mathcal{O}_S/a \mathcal{O}_S \to 0
$$
由于 $\mathcal{F}$ 在 $S$ 上平坦，得到正合列
$$
0 \to \mathcal{F} \xrightarrow{a} \mathcal{F} \to
\mathcal{F}/a\mathcal{F} \to 0
$$
取直像，得到正合列
$$
0 \to f_*\mathcal{F} \xrightarrow{a} f_*\mathcal{F} \to
f_*(\mathcal{F}/a\mathcal{F})
$$
这证明了 (4)，并且说明
$f_*\mathcal{F}/ af_*\mathcal{F} \subset f_*(\mathcal{F}/a\mathcal{F})$。
若 $b$ 在 $\mathcal{O}_{S, s}/a\mathcal{O}_{S, s}$ 上是非零因子，
则完全相同的论证说明
$b : \mathcal{F}/a\mathcal{F} \to \mathcal{F}/a\mathcal{F}$ 是单射。
取直像可知
$$
b : f_*(\mathcal{F}/a\mathcal{F}) \to f_*(\mathcal{F}/a\mathcal{F})
$$
是单射，因而
$b : f_*\mathcal{F}/ af_*\mathcal{F} \to f_*\mathcal{F}/ af_*\mathcal{F}$
也为单射。这证明了 (5)。(6) 由 (4)、(5) 及各定义得出。
\end{proof}


\section{相对伴随点集}
\label{divisors-section-relative-assassin}

\noindent
设 $A \to B$ 为环映射，$N$ 为 $B$-模。回忆，若素理想
$\mathfrak q \subset B$ 满足如下条件，则称它属于 $N$ 在
$B/A$ 上的相对伴随素理想集：$\mathfrak q$ 是
$N \otimes_A \kappa(\mathfrak p)$ 的伴随素理想。这里
$\mathfrak p = A \cap \mathfrak q$。见《代数》定义
\ref{algebra-definition-relative-assassin}。
下面定义概形态射上拟凝聚层的相对伴随点集。

\begin{definition}
\label{divisors-definition-relative-assassin}
设 $f : X \to S$ 为概形的态射。
设 $\mathcal{F}$ 为拟凝聚 $\mathcal{O}_X$-模。
$\mathcal{F}$ 在 $X$ 中、$S$ 上的{\it 相对伴随点集}是集合
$$
\text{Ass}_{X/S}(\mathcal{F}) =
\bigcup\nolimits_{s \in S} \text{Ass}_{X_s}(\mathcal{F}_s)
$$
其中 $\mathcal{F}_s = (X_s \to X)^*\mathcal{F}$ 是 $\mathcal{F}$
在 $f$ 于 $s$ 处的纤维上的限制。
\end{definition}

\noindent
这里仍有一个注意事项：最好只在 $f$ 的纤维局部 Noether 且
$\mathcal{F}$ 为有限型时使用这一概念。一般情形下，或许应当使用
相对弱伴随点集（定义见下一节）。下面把概形论概念与仿射开集上的
代数概念联系起来；注意，这一对应只有在概形态射的纤维局部
Noether 时才完全成立。

\begin{lemma}
\label{divisors-lemma-relative-assassin-affine-open}
设 $f : X \to S$ 为概形的态射。
设 $\mathcal{F}$ 为 $X$ 上的拟凝聚层。
设 $U \subset X$ 和 $V \subset S$ 为仿射开集，且
$f(U) \subset V$。写 $U = \Spec(A)$、$V = \Spec(R)$，并置
$M = \Gamma(U, \mathcal{F})$。
设 $x \in U$，并设 $\mathfrak p \subset A$ 为对应的素理想。则
$$
\mathfrak p \in \text{Ass}_{A/R}(M) \Rightarrow
x \in \text{Ass}_{X/S}(\mathcal{F})
$$
若所有纤维 $X_s$（即 $f$ 的纤维）都局部 Noether，则等价式
$\mathfrak p \in \text{Ass}_{A/R}(M) \Leftrightarrow
x \in \text{Ass}_{X/S}(\mathcal{F})$
对上述每一对 $(\mathfrak p, x)$ 都成立。
\end{lemma}

\begin{proof}
集合 $\text{Ass}_{A/R}(M)$ 定义于《代数》定义
\ref{algebra-definition-relative-assassin}。选取一对
$(\mathfrak p, x)$，并令 $s = f(x)$。
设 $\mathfrak r \subset R$ 为位于 $\mathfrak p$ 之下的素理想，
即对应于 $s$ 的素理想。设
$\mathfrak p' \subset A \otimes_R \kappa(\mathfrak r)$ 为逆像等于
$\mathfrak p$ 的素理想；换言之，它对应于把 $x$ 看作其纤维
$X_s$ 的点时的那个点。于是，
$\mathfrak p \in \text{Ass}_{A/R}(M)$ 当且仅当 $\mathfrak p'$ 是
$M \otimes_R \kappa(\mathfrak r)$ 的伴随素理想，见《代数》引理
\ref{algebra-lemma-compare-relative-assassins}。
注意，环 $A \otimes_R \kappa(\mathfrak r)$ 对应于 $U_s$，而模
$M \otimes_R \kappa(\mathfrak r)$ 对应于拟凝聚层
$\mathcal{F}_s|_{U_s}$。因此，由引理
\ref{divisors-lemma-associated-affine-open}，$x$ 是 $\mathcal{F}_s$ 的伴随点。
若 $\mathfrak p'$ 有限生成，则逆命题也成立；这便说明了最后一句。
\end{proof}

\begin{lemma}
\label{divisors-lemma-base-change-relative-assassin}
设 $f : X \to S$ 为概形的态射。
设 $\mathcal{F}$ 为拟凝聚 $\mathcal{O}_X$-模。
设 $g : S' \to S$ 为概形的态射。考虑基变换图
$$
\xymatrix{
X' \ar[d] \ar[r]_{g'} & X \ar[d] \\
S' \ar[r]^g & S
}
$$
并置 $\mathcal{F}' = (g')^*\mathcal{F}$。设点 $x' \in X'$ 的像
分别为 $x \in X$、$s' \in S'$ 和 $s \in S$。假设 $f$ 局部有限型。
则 $x' \in \text{Ass}_{X'/S'}(\mathcal{F}')$ 当且仅当
$x \in \text{Ass}_{X/S}(\mathcal{F})$，并且 $x'$ 对应于
$\Spec(\kappa(s') \otimes_{\kappa(s)} \kappa(x))$ 的一个不可约分支的泛点。
\end{lemma}

\begin{proof}
考虑纤维态射 $X'_{s'} \to X_s$。由于
$X_{s'} = X_s \times_{\Spec(\kappa(s))} \Spec(\kappa(s'))$，
这是平坦态射。此外，$\mathcal{F}'_{s'}$ 是 $\mathcal{F}_s$ 沿此态射的拉回。
由于 $X_s$ 在 Noether 概形 $\Spec(\kappa(s))$ 上局部有限型，
故 $X_s$ 局部 Noether，见《态射》引理
\ref{morphisms-lemma-finite-type-noetherian}。因此可应用引理
\ref{divisors-lemma-bourbaki}，得到
$$
\text{Ass}_{X'_{s'}}(\mathcal{F}'_{s'}) =
\bigcup\nolimits_{x \in \text{Ass}(\mathcal{F}_s)} \text{Ass}((X'_{s'})_x).
$$
于是，为证明本引理，只须说明纤维 $(X'_{s'})_x$ 的伴随点就是其泛点；
这里该纤维是态射 $X'_{s'} \to X_s$ 在 $x$ 上的纤维。注意，作为概形有
$(X'_{s'})_x = \Spec(\kappa(s') \otimes_{\kappa(s)} \kappa(x))$。
由《代数》引理 \ref{algebra-lemma-tensor-fields-CM}，环
$\kappa(s') \otimes_{\kappa(s)} \kappa(x)$ 是 Noether
Cohen-Macaulay 环。因此，它的伴随素理想就是其极小素理想；见
《代数》命题 \ref{algebra-proposition-minimal-primes-associated-primes}
（极小素理想是伴随素理想）和《代数》引理
\ref{algebra-lemma-criterion-no-embedded-primes}（没有嵌入素理想）。
\end{proof}

\begin{remark}
\label{divisors-remark-base-change-relative-assassin}
采用引理 \ref{divisors-lemma-base-change-relative-assassin} 的记号和假设，
总有
$(g')^{-1}(\text{Ass}_{X/S}(\mathcal{F})) \supset
\text{Ass}_{X'/S'}(\mathcal{F}')$。
若态射 $S' \to S$ 局部拟有限，则事实上有
$$
(g')^{-1}(\text{Ass}_{X/S}(\mathcal{F}))
=
\text{Ass}_{X'/S'}(\mathcal{F}')
$$
因为此时域扩张 $\kappa(s')/\kappa(s)$ 总是有限的。事实上，
对任意满足下述条件的态射 $g : S' \to S$，这一结论仍成立：
所有域扩张 $\kappa(s')/\kappa(s)$ 都是代数扩张。这是因为此时
$\kappa(s') \otimes_{\kappa(s)} \kappa(x)$ 的所有素理想都是
极大（且极小）素理想，见《代数》引理
\ref{algebra-lemma-integral-over-field}。
\end{remark}




\section{相对弱伴随点集}
\label{divisors-section-relative-weak-assassin}

\begin{definition}
\label{divisors-definition-relative-weak-assassin}
设 $f : X \to S$ 为概形的态射。
设 $\mathcal{F}$ 为拟凝聚 $\mathcal{O}_X$-模。
$\mathcal{F}$ 在 $X$ 中、$S$ 上的{\it 相对弱伴随点集}是集合
$$
\text{WeakAss}_{X/S}(\mathcal{F}) =
\bigcup\nolimits_{s \in S} \text{WeakAss}(\mathcal{F}_s)
$$
其中 $\mathcal{F}_s = (X_s \to X)^*\mathcal{F}$ 是 $\mathcal{F}$
在 $f$ 于 $s$ 处的纤维上的限制。
\end{definition}

\begin{lemma}
\label{divisors-lemma-relative-weak-assassin-assassin-finite-type}
设 $f : X \to S$ 为局部有限型的概形态射。
设 $\mathcal{F}$ 为拟凝聚 $\mathcal{O}_X$-模。则
$\text{WeakAss}_{X/S}(\mathcal{F}) = \text{Ass}_{X/S}(\mathcal{F})$。
\end{lemma}

\begin{proof}
这是因为 $f$ 的纤维都是局部 Noether 概形，并且在局部 Noether
概形上，伴随点与弱伴随点一致，见引理
\ref{divisors-lemma-ass-weakly-ass}。
\end{proof}

\begin{lemma}
\label{divisors-lemma-relative-weak-assassin-finite}
设 $f : X \to S$ 为概形的态射。
设 $i : Z \to X$ 为有限态射。
设 $\mathcal{F}$ 为拟凝聚 $\mathcal{O}_Z$-模。则
$\text{WeakAss}_{X/S}(i_*\mathcal{F}) =
i(\text{WeakAss}_{Z/S}(\mathcal{F}))$。
\end{lemma}

\begin{proof}
设 $i_s : Z_s \to X_s$ 为纤维间的诱导态射。于是
$(i_*\mathcal{F})_s = i_{s, *}(\mathcal{F}_s)$；这是《概形上同调》引理
\ref{coherent-lemma-affine-base-change} 与 $i$ 是仿射态射这一事实的结果。
因此可应用引理
\ref{divisors-lemma-weakly-associated-finite} 得出结论。
\end{proof}





\section{Fitting 理想}
\label{divisors-section-fitting-ideals}

\noindent
本节继续《代数进阶》第
\ref{more-algebra-section-fitting-ideals} 节中的讨论。
设 $S$ 为概形，$\mathcal{F}$ 为有限型拟凝聚 $\mathcal{O}_S$-模。
在这种情形下，可以构造 Fitting 理想
$$
0 = \text{Fit}_{-1}(\mathcal{F}) \subset \text{Fit}_0(\mathcal{F}) \subset
\text{Fit}_1(\mathcal{F}) \subset \ldots \subset \mathcal{O}_S
$$
它们是由下述性质刻画的拟凝聚理想序列：对每个仿射开集
$U = \Spec(A)$，若它是 $S$ 的开集且 $\mathcal{F}|_U$ 对应于 $A$-模 $M$，则
$\text{Fit}_i(\mathcal{F})|_U$ 对应于理想
$\text{Fit}_i(M) \subset A$。这一定义良好，并给出拟凝聚理想层，
因为若 $f \in A$，则第 $i$ 个 Fitting 理想（即 $M_f$ 在 $A_f$
上的相应理想）等于 $\text{Fit}_i(M) A_f$，见《代数进阶》引理
\ref{more-algebra-lemma-fitting-ideal-basics}。

\medskip\noindent
也可以用 $\mathcal{F}$ 的局部展示来构造 Fitting 理想。具体地，
若 $U \subset X$ 为开集，并且
$$
\bigoplus\nolimits_{i \in I} \mathcal{O}_U \to
\mathcal{O}_U^{\oplus n} \to \mathcal{F}|_U \to 0
$$
是 $\mathcal{F}$ 在 $U$ 上的一个展示，则
$\text{Fit}_r(\mathcal{F})|_U$ 由定义该展示中第一个箭头的矩阵的
所有 $(n - r) \times (n - r)$ 子式生成。这与上述构造相容，
因为环上模的 Fitting 理想本来就是这样定义的。略去一些细节。

\begin{lemma}
\label{divisors-lemma-base-change-fitting-ideal}
设 $f : T \to S$ 为概形的态射。
设 $\mathcal{F}$ 为有限型拟凝聚 $\mathcal{O}_S$-模。则
$f^{-1}\text{Fit}_i(\mathcal{F}) \cdot \mathcal{O}_T =
\text{Fit}_i(f^*\mathcal{F})$。
\end{lemma}

\begin{proof}
由《代数进阶》引理 \ref{more-algebra-lemma-fitting-ideal-basics}
的 (3) 立即得到。
\end{proof}

\begin{lemma}
\label{divisors-lemma-fitting-ideal-of-finitely-presented}
设 $S$ 为概形。
设 $\mathcal{F}$ 为有限展示 $\mathcal{O}_S$-模。则
$\text{Fit}_r(\mathcal{F})$ 是有限型拟凝聚理想。
\end{lemma}

\begin{proof}
由《代数进阶》引理 \ref{more-algebra-lemma-fitting-ideal-basics}
的 (4) 立即得到。
\end{proof}

\begin{lemma}
\label{divisors-lemma-on-subscheme-cut-out-by-Fit-0}
设 $S$ 为概形。
设 $\mathcal{F}$ 为有限型拟凝聚 $\mathcal{O}_S$-模。
设 $Z_0 \subset S$ 为由 $\text{Fit}_0(\mathcal{F})$ 截出的闭子概形。
设 $Z \subset S$ 为 $\mathcal{F}$ 的概形论支撑。则
\begin{enumerate}
\item 作为闭子概形，有 $Z \subset Z_0 \subset S$；
\item 作为闭子集，有 $Z = Z_0 = \text{Supp}(\mathcal{F})$；
\item 存在有限型拟凝聚 $\mathcal{O}_{Z_0}$-模 $\mathcal{G}_0$，使得
$$
(Z_0 \to X)_*\mathcal{G}_0 = \mathcal{F}.
$$
\end{enumerate}
\end{lemma}

\begin{proof}
回忆，$Z$ 局部由 $\mathcal{F}$ 的零化子截出，见《态射》定义
\ref{morphisms-definition-scheme-theoretic-support}（该定义使用《态射》引理
\ref{morphisms-lemma-scheme-theoretic-support} 来定义 $Z$）。因此，
由《代数进阶》引理 \ref{more-algebra-lemma-fitting-ideal-basics} 的 (6)，
可见在概形论意义下 $Z \subset Z_0$。另一方面，由《态射》引理
\ref{morphisms-lemma-scheme-theoretic-support}，集合论上有
$Z = \text{Supp}(\mathcal{F})$；又由《代数进阶》引理
\ref{more-algebra-lemma-fitting-ideal-basics} 的 (7)，集合论上有
$Z_0 = Z$。最后，为得到 (3) 中的 $\mathcal{G}_0$，可以使用
《态射》引理 \ref{morphisms-lemma-scheme-theoretic-support} 给出的
$\mathcal{G}$（它定义在 $Z$ 上），并置
$\mathcal{G}_0 = (Z \to Z_0)_*\mathcal{G}$；也可以使用《态射》引理
\ref{morphisms-lemma-i-star-equivalence}，以及由《代数进阶》引理
\ref{more-algebra-lemma-fitting-ideal-basics} 的 (6) 得到的事实：
$\text{Fit}_0(\mathcal{F})$ 零化 $\mathcal{F}$。
\end{proof}

\begin{lemma}
\label{divisors-lemma-fitting-ideal-generate-locally}
设 $S$ 为概形，$\mathcal{F}$ 为有限型拟凝聚 $\mathcal{O}_S$-模，
并设 $s \in S$。则 $\mathcal{F}$ 能由 $r$ 个元素生成，
并且可在 $s$ 的一个邻域中如此生成，当且仅当
$\text{Fit}_r(\mathcal{F})_s = \mathcal{O}_{S, s}$。
\end{lemma}

\begin{proof}
由《代数进阶》引理
\ref{more-algebra-lemma-fitting-ideal-generate-locally} 立即得到。
\end{proof}

\begin{lemma}
\label{divisors-lemma-fitting-ideal-finite-locally-free}
设 $S$ 为概形，$\mathcal{F}$ 为有限型拟凝聚 $\mathcal{O}_S$-模，
并设 $r \geq 0$。以下条件等价：
\begin{enumerate}
\item $\mathcal{F}$ 是秩为 $r$ 的有限局部自由模；
\item $\text{Fit}_{r - 1}(\mathcal{F}) = 0$ 且
$\text{Fit}_r(\mathcal{F}) = \mathcal{O}_S$；
\item $\text{Fit}_k(\mathcal{F}) = 0$ 对 $k < r$ 成立，而
$\text{Fit}_k(\mathcal{F}) = \mathcal{O}_S$ 对 $k \geq r$ 成立。
\end{enumerate}
\end{lemma}

\begin{proof}
由《代数进阶》引理
\ref{more-algebra-lemma-fitting-ideal-finite-locally-free} 立即得到。
\end{proof}

\begin{lemma}
\label{divisors-lemma-locally-free-rank-r-pullback}
设 $S$ 为概形，$\mathcal{F}$ 为有限型拟凝聚 $\mathcal{O}_S$-模。
由 $\mathcal{F}$ 的 Fitting 理想定义的闭子概形
$$
S = Z_{-1} \supset Z_0 \supset Z_1 \supset Z_2 \ldots
$$
具有下述性质：
\begin{enumerate}
\item 交集 $\bigcap Z_r$ 为空。
\item 由规则
$$
T \longmapsto
\left\{
\begin{matrix}
\{*\} & \text{若 }\mathcal{F}_T\text{ is locally generated by }
\leq r\text{ sections} \\
\emptyset & \text{otherwise}
\end{matrix}
\right.
$$
定义的函子 $(\Sch/S)^{opp} \to \textit{Sets}$ 由开子概形
$S \setminus Z_r$ 表示。
\item 由规则
$$
T \longmapsto
\left\{
\begin{matrix}
\{*\} & \text{若 }\mathcal{F}_T\text{ locally free rank }r\\
\emptyset & \text{otherwise}
\end{matrix}
\right.
$$
定义的函子 $F_r : (\Sch/S)^{opp} \to \textit{Sets}$ 由局部闭子概形
$Z_{r - 1} \setminus Z_r$ 表示；它是 $S$ 的局部闭子概形。
\end{enumerate}
若 $\mathcal{F}$ 有限展示，则态射 $Z_r \to S$、
$S \setminus Z_r \to S$ 和 $Z_{r - 1} \setminus Z_r \to S$
都是有限展示的。
\end{lemma}

\begin{proof}
(1) 成立，是因为在每个仿射开集 $U$ 上，都存在整数 $n$，使得
$\text{Fit}_n(\mathcal{F})|_U = \mathcal{O}_U$。事实上，可取 $n$
为 $\mathcal{F}$ 在 $U$ 上的生成元个数；见《代数进阶》第
\ref{more-algebra-section-fitting-ideals} 节。

\medskip\noindent
对任意态射 $g : T \to S$，由引理 \ref{divisors-lemma-base-change-fitting-ideal}
和 \ref{divisors-lemma-fitting-ideal-generate-locally} 可见，$\mathcal{F}_T$
局部由 $\leq r$ 个截面生成，当且仅当
$\text{Fit}_r(\mathcal{F}) \cdot \mathcal{O}_T = \mathcal{O}_T$。
这证明了 (2)。

\medskip\noindent
对任意态射 $g : T \to S$，由引理 \ref{divisors-lemma-base-change-fitting-ideal}
和 \ref{divisors-lemma-fitting-ideal-finite-locally-free} 可见，$\mathcal{F}_T$
是秩为 $r$ 的自由模，当且仅当
$\text{Fit}_r(\mathcal{F}) \cdot \mathcal{O}_T = \mathcal{O}_T$ 且
$\text{Fit}_{r - 1}(\mathcal{F}) \cdot \mathcal{O}_T = 0$。
这证明了 (3)。

\medskip\noindent
假设 $\mathcal{F}$ 有限展示。则每个态射 $Z_r \to S$ 都是有限展示的，
因为 $\text{Fit}_r(\mathcal{F})$ 有限型（引理
\ref{divisors-lemma-fitting-ideal-of-finitely-presented} 和《态射》引理
\ref{morphisms-lemma-closed-immersion-finite-presentation}）。这说明
$Z_{r - 1} \setminus Z_r$ 是 $Z_r$ 中的逆紧开集（《性质》引理
\ref{properties-lemma-quasi-coherent-finite-type-ideals}），因而态射
$Z_{r - 1} \setminus Z_r \to Z_r$ 也是有限展示的。
\end{proof}

\noindent
尽管有引理 \ref{divisors-lemma-locally-free-rank-r-pullback}，若 $\mathcal{F}$
只是有限型拟凝聚模，下述引理并不成立。确实，此时分层仍然存在，
但一般而言它并不表示函子 $F_{flat}$。

\begin{lemma}
\label{divisors-lemma-finite-presentation-module}
设 $S$ 为概形，$\mathcal{F}$ 为有限展示 $\mathcal{O}_S$-模。
设 $S = Z_{-1} \supset Z_0 \supset Z_1 \supset \ldots$ 如引理
\ref{divisors-lemma-locally-free-rank-r-pullback} 所述。置
$S_r = Z_{r - 1} \setminus Z_r$。则
$S' = \coprod_{r \geq 0} S_r$ 表示函子
$$
F_{flat} : \Sch/S \longrightarrow \textit{Sets},\quad\quad
T \longmapsto
\left\{
\begin{matrix}
\{*\} & \text{若 }\mathcal{F}_T\text{ flat over }T\\
\emptyset & \text{otherwise}
\end{matrix}
\right.
$$
此外，$\mathcal{F}|_{S_r}$ 是秩为 $r$ 的局部自由模，且态射
$S_r \to S$ 和 $S' \to S$ 都是有限展示的。
\end{lemma}

\begin{proof}
设 $g : T \to S$ 是概形态射，并且拉回
$\mathcal{F}_T = g^*\mathcal{F}$ 平坦。则 $\mathcal{F}_T$ 是有限展示
平坦 $\mathcal{O}_T$-模。因此，由《性质》引理
\ref{properties-lemma-finite-locally-free}，$\mathcal{F}_T$ 有限局部自由。
于是 $T = \coprod_{r \geq 0} T_r$，其中 $\mathcal{F}_T|_{T_r}$
局部自由且秩为 $r$。这说明在 $\Sch/S$ 上的 Zariski 层范畴中
$$
F_{flat} = \coprod\nolimits_{r \geq 0} F_r
$$
其中 $F_r$ 如引理 \ref{divisors-lemma-locally-free-rank-r-pullback} 所述。
因此，$F_{flat}$ 由
$\coprod_{r \geq 0} (Z_{r - 1} \setminus Z_r)$ 表示，其中 $Z_r$
如引理 \ref{divisors-lemma-locally-free-rank-r-pullback} 所述。
其他断言也由该引理得出。
\end{proof}

\begin{example}
\label{divisors-example-not-fp-MB}
设 $R = \prod_{n \in \mathbf{N}} \mathbf{F}_2$。设 $I \subset R$
为由这样的元素 $a = (a_n)_{n \in \mathbf{N}}$ 构成的理想：
它们只有有限多个分量非零。设 $\mathfrak m$ 为包含 $I$ 的极大理想。
则 $M = R/\mathfrak m$ 是有限平坦 $R$-模，因为 $R$ 绝对平坦
（《代数进阶》引理
\ref{more-algebra-lemma-product-fields-absolutely-flat}）。
置 $S = \Spec(R)$ 和 $\mathcal{F} = \widetilde{M}$。
引理 \ref{divisors-lemma-locally-free-rank-r-pullback} 中的闭子概形为
$S = Z_{-1}$、$Z_0 = \Spec(R/\mathfrak m)$，以及
$Z_i = \emptyset$（当 $i > 0$ 时）。但是，$\text{id} : S \to S$ 不通过
$(S \setminus Z_0) \amalg Z_0$ 分解，因为 $\mathfrak m$ 是 $S$
的非孤立点。因此，引理 \ref{divisors-lemma-finite-presentation-module}
对有限型模不成立。
\end{example}







\section{态射的奇异点集}
\label{divisors-section-singular-locus-morphism}

\noindent
设 $f : X \to S$ 为概形的有限型态射。令点集 $U$ 由 $f$ 光滑的点组成；
它是 $X$ 的开集（由《态射》定义 \ref{morphisms-definition-smooth}）。
在许多情形下，拥有一个典范闭子概形
$\text{Sing}(f) \subset X$ 很有用：它的补集为 $U$，并且它的构造
与任意基变换交换。

\medskip\noindent
若 $f$ 有限展示，一种选择是考虑由仿射局部意义下“严格标准”的
函数截出的闭子概形 $Z$；“严格标准”取《环映射的光滑化》定义
\ref{smoothing-definition-strictly-standard} 的意义。
由《环映射的光滑化》引理
\ref{smoothing-lemma-strictly-standard-base-change}，若
$f' : X' \to S'$ 是 $f$ 沿态射 $S' \to S$ 的基变换，则
$Z' \subset S' \times_S Z$，其中 $Z'$ 是 $X'$ 中由仿射局部严格标准
函数截出的闭子概形。然而，等号是否成立并不清楚。
严格标准元这一概念在 Popescu 定理一章中很有用。据我们所知，
由这些元素定义的闭子概形在文献中并未使用\footnote{若 $f$ 是
局部完全交态射（《态射进阶》定义
\ref{more-morphisms-definition-lci}），则由局部严格标准元截出的闭子概形
正是应当考察的对象。}。

\medskip\noindent
若 $f$ 平坦且有限展示，并且 $f$ 的所有纤维都等维且维数为 $d$，
则可用第 $d$ 个 Fitting 理想（即 $\Omega_{X/S}$ 的相应理想）得到一个
良好的闭子概形。对任意有限型态射，$X$ 的各闭子概形若由
$\Omega_{X/S}$ 的 Fitting 理想定义，则给出 $X$ 的一个分层；
该分层按照 $\Omega_{X/S}$ 的秩划分，其构造与基变换交换。
这一分层可能很有帮助，并与纤维的
嵌入维数有关，见《簇》第 \ref{varieties-section-embedding-dimension} 节。

\begin{lemma}
\label{divisors-lemma-base-change-and-fitting-ideal-omega}
设 $f : X \to S$ 为局部有限型的概形态射。
设 $X = Z_{-1} \supset Z_0 \supset Z_1 \supset \ldots$ 为由
$\Omega_{X/S}$ 的 Fitting 理想定义的闭子概形。则 $Z_i$ 的构造
与任意基变换交换。
\end{lemma}

\begin{proof}
注意，$\Omega_{X/S}$ 是有限型拟凝聚 $\mathcal{O}_X$-模
（《态射》引理 \ref{morphisms-lemma-finite-type-differentials}），
故其 Fitting 理想有定义。若 $f' : X' \to S'$ 是 $f$ 沿
$g : S' \to S$ 的基变换，则
$\Omega_{X'/S'} = (g')^*\Omega_{X/S}$，其中 $g' : X' \to X$ 为投影
（《态射》引理 \ref{morphisms-lemma-base-change-differentials}）。因此
$(g')^{-1}\text{Fit}_i(\Omega_{X/S}) \cdot \mathcal{O}_{X'} =
\text{Fit}_i(\Omega_{X'/S'})$。这表示在概形论意义下
$$
Z'_i = (g')^{-1}(Z_i) = Z_i \times_X X'
$$
而这正是本引理断言的含义。
\end{proof}

\noindent
第 $0$ 个 Fitting 理想（关于 $\Omega$）截出态射的“分歧点集”。

\begin{lemma}
\label{divisors-lemma-zero-fitting-ideal-omega-unramified}
设 $f : X \to S$ 为局部有限型的概形态射。
闭子概形 $Z \subset X$ 由第 $0$ 个 Fitting 理想（关于
$\Omega_{X/S}$）截出；它恰好是 $f$ 不为非分歧的点之集。
\end{lemma}

\begin{proof}
由引理 \ref{divisors-lemma-on-subscheme-cut-out-by-Fit-0}，$Z$ 的补集恰好是
$\Omega_{X/S}$ 为零的点集。由《态射》引理
\ref{morphisms-lemma-unramified-omega-zero}，这恰好是 $f$ 非分歧的点集。
\end{proof}

\begin{lemma}
\label{divisors-lemma-d-fitting-ideal-omega-smooth}
设 $f : X \to S$ 为概形的态射，并设 $d \geq 0$ 为整数。假设
\begin{enumerate}
\item $f$ 平坦；
\item $f$ 局部有限展示；
\item $f$ 的每个非空纤维都等维且维数为 $d$。
\end{enumerate}
设 $Z \subset X$ 为第 $d$ 个 Fitting 理想（关于 $\Omega_{X/S}$）
截出的闭子概形。则 $Z$ 恰好是 $f$ 不光滑的点之集。
\end{lemma}

\begin{proof}
由引理 \ref{divisors-lemma-locally-free-rank-r-pullback}，$Z$ 的补集恰好是
$\Omega_{X/S}$ 可由至多 $d$ 个元素生成的点集。因此，本引理由
《态射》引理 \ref{morphisms-lemma-smooth-at-point} 得出。
\end{proof}







\section{无挠模}
\label{divisors-section-torsion-free}

\noindent
本节是《代数进阶》第
\ref{more-algebra-section-torsion-flat} 节关于拟凝聚模的对应版本。

\begin{lemma}
\label{divisors-lemma-torsion-sections}
设 $X$ 为泛点是 $\eta$ 的整概形，$\mathcal{F}$ 为拟凝聚
$\mathcal{O}_X$-模。设 $U \subset X$ 为非空开集，且
$s \in \mathcal{F}(U)$。则下列条件等价：
\begin{enumerate}
\item 对某个 $x \in U$，$s$ 在 $\mathcal{F}_x$ 中的像是挠元；
\item 对所有 $x \in U$，$s$ 在 $\mathcal{F}_x$ 中的像都是挠元；
\item $s$ 在 $\mathcal{F}_\eta$ 中的像为零；
\item $s$ 在 $j_*\mathcal{F}_\eta$ 中的像为零，其中
$j : \eta \to X$ 是包含态射。
\end{enumerate}
\end{lemma}

\begin{proof}
略。
\end{proof}

\begin{definition}
\label{divisors-definition-torsion}
设 $X$ 为整概形，$\mathcal{F}$ 为拟凝聚 $\mathcal{O}_X$-模。
\begin{enumerate}
\item 若 $\mathcal{F}$ 的一个局部截面满足引理
\ref{divisors-lemma-torsion-sections} 中的等价条件，则称它是{\it 挠截面}。
\item 称 $\mathcal{F}$ 是{\it 无挠的}，如果 $\mathcal{F}$ 的每个
挠截面都是 $0$。
\end{enumerate}
\end{definition}

\noindent
下面给出按例必须有的引理，将此定义与通常的代数定义作比较。

\begin{lemma}
\label{divisors-lemma-check-torsion-on-affines}
设 $X$ 为整概形，$\mathcal{F}$ 为拟凝聚 $\mathcal{O}_X$-模。
则下列条件等价：
\begin{enumerate}
\item $\mathcal{F}$ 无挠；
\item 对每个仿射开集 $U \subset X$，$\mathcal{F}(U)$ 都是无挠
$\mathcal{O}(U)$-模。
\end{enumerate}
\end{lemma}

\begin{proof}
略。
\end{proof}

\begin{lemma}
\label{divisors-lemma-torsion}
设 $X$ 为整概形，$\mathcal{F}$ 为拟凝聚 $\mathcal{O}_X$-模。
$\mathcal{F}$ 的挠截面构成一个拟凝聚 $\mathcal{O}_X$-子模
$\mathcal{F}_{tors} \subset \mathcal{F}$。
商模 $\mathcal{F}/\mathcal{F}_{tors}$ 无挠。
\end{lemma}

\begin{proof}
略。代数版本见《代数进阶》引理 \ref{more-algebra-lemma-torsion}。
\end{proof}

\begin{lemma}
\label{divisors-lemma-flat-torsion-free}
设 $X$ 为整概形。任何平坦的拟凝聚 $\mathcal{O}_X$-模都是无挠的。
\end{lemma}

\begin{proof}
略。参见《代数进阶》引理
\ref{more-algebra-lemma-flat-torsion-free}。
\end{proof}

\begin{lemma}
\label{divisors-lemma-flat-pullback-torsion}
设 $f : X \to Y$ 为整概形之间的平坦态射，$\mathcal{G}$ 为无挠的
拟凝聚 $\mathcal{O}_Y$-模。则 $f^*\mathcal{G}$ 是无挠的
$\mathcal{O}_X$-模。
\end{lemma}

\begin{proof}
略。代数版本见《代数进阶》引理
\ref{more-algebra-lemma-flat-pullback-torsion}。
\end{proof}

\begin{lemma}
\label{divisors-lemma-flat-over-integral-integral-fibre}
设 $f : X \to Y$ 为概形的平坦态射。若 $Y$ 是整概形，且 $f$ 的泛纤维
是整概形，则 $X$ 是整概形。
\end{lemma}

\begin{proof}
其代数版本如下：设 $A$ 为整环，分式域为 $K$；设 $B$ 为平坦
$A$-代数，且 $B \otimes_A K$ 是整环。则 $B$ 是整环。这是因为由
《代数进阶》引理 \ref{more-algebra-lemma-flat-torsion-free}，
$B$ 无挠，因而 $B \subset B \otimes_A K$。
\end{proof}

\begin{lemma}
\label{divisors-lemma-check-torsion}
设 $X$ 为整概形，$\mathcal{F}$ 为拟凝聚 $\mathcal{O}_X$-模。
则 $\mathcal{F}$ 无挠，当且仅当 $\mathcal{F}_x$ 是无挠的
$\mathcal{O}_{X, x}$-模；这一条件须对所有 $x \in X$ 成立。
\end{lemma}

\begin{proof}
略。参见《代数进阶》引理
\ref{more-algebra-lemma-check-torsion}。
\end{proof}

\begin{lemma}
\label{divisors-lemma-extension-torsion-free}
设 $X$ 为整概形，并设
$0 \to \mathcal{F} \to \mathcal{F}' \to \mathcal{F}'' \to 0$
为拟凝聚 $\mathcal{O}_X$-模的短正合列。若 $\mathcal{F}$ 和
$\mathcal{F}''$ 无挠，则 $\mathcal{F}'$ 无挠。
\end{lemma}

\begin{proof}
略。代数版本见《代数进阶》引理
\ref{more-algebra-lemma-extension-torsion-free}。
\end{proof}

\begin{lemma}
\label{divisors-lemma-torsion-free-finite-noetherian-domain}
设 $X$ 为泛点是 $\eta$ 的局部诺特整概形，$\mathcal{F}$ 为非零凝聚
$\mathcal{O}_X$-模。则下列条件等价：
\begin{enumerate}
\item $\mathcal{F}$ 无挠；
\item $\eta$ 是 $\mathcal{F}$ 的唯一伴随素理想；
\item $\eta$ 属于 $\mathcal{F}$ 的支撑，且 $\mathcal{F}$ 具有性质
$(S_1)$；
\item $\eta$ 属于 $\mathcal{F}$ 的支撑，且 $\mathcal{F}$
没有嵌入伴随素理想。
\end{enumerate}
\end{lemma}

\begin{proof}
这是把《代数进阶》引理
\ref{more-algebra-lemma-torsion-free-finite-noetherian-domain}
翻译成概形语言的结果。略去该翻译。
\end{proof}

\begin{lemma}
\label{divisors-lemma-torsion-free-over-regular-dim-1}
设 $X$ 为维数 $\leq 1$ 的整正则概形，$\mathcal{F}$ 为凝聚
$\mathcal{O}_X$-模。则下列条件等价：
\begin{enumerate}
\item $\mathcal{F}$ 无挠；
\item $\mathcal{F}$ 有限局部自由。
\end{enumerate}
\end{lemma}

\begin{proof}
有限局部自由模显然无挠。反过来，我们将说明：若 $\mathcal{F}$ 无挠，
则 $\mathcal{F}_x$ 是自由 $\mathcal{O}_{X, x}$-模；这一结论对所有
$x \in X$ 成立。由《代数》引理
\ref{algebra-lemma-finite-projective} 以及 $\mathcal{F}$ 凝聚这一事实，
这就足够了。若 $\dim(\mathcal{O}_{X, x}) = 0$，则
$\mathcal{O}_{X, x}$ 是域，结论显然。若
$\dim(\mathcal{O}_{X, x}) = 1$，则 $\mathcal{O}_{X, x}$ 是离散赋值环
（《代数》引理 \ref{algebra-lemma-characterize-dvr}），并且
$\mathcal{F}_x$ 无挠。因此，由《代数进阶》引理
\ref{more-algebra-lemma-dedekind-torsion-free-flat}，
$\mathcal{F}_x$ 是自由的。
\end{proof}

\begin{lemma}
\label{divisors-lemma-hom-into-torsion-free}
设 $X$ 为整概形，$\mathcal{F}$、$\mathcal{G}$ 为拟凝聚
$\mathcal{O}_X$-模。若 $\mathcal{G}$ 无挠，且 $\mathcal{F}$ 有限展示，
则 $\SheafHom_{\mathcal{O}_X}(\mathcal{F}, \mathcal{G})$ 无挠。
\end{lemma}

\begin{proof}
该断言有意义，因为
$\SheafHom_{\mathcal{O}_X}(\mathcal{F}, \mathcal{G})$
由《概形》第 \ref{schemes-section-quasi-coherent} 节可知是拟凝聚的。
断言的证明见《代数进阶》引理
\ref{more-algebra-lemma-hom-into-torsion-free}。略去一些细节。
\end{proof}

\begin{lemma}
\label{divisors-lemma-isom-depth-2-torsion-free}
设 $X$ 为整的局部诺特概形，并设
$\varphi : \mathcal{F} \to \mathcal{G}$ 为拟凝聚
$\mathcal{O}_X$-模的映射。假设 $\mathcal{F}$ 凝聚、$\mathcal{G}$ 无挠，
且对每个 $x \in X$，下列情形之一成立：
\begin{enumerate}
\item $\mathcal{F}_x \to \mathcal{G}_x$ 是同构；或
\item $\text{depth}(\mathcal{F}_x) \geq 2$。
\end{enumerate}
则 $\varphi$ 是同构。
\end{lemma}

\begin{proof}
这是把《代数进阶》引理
\ref{more-algebra-lemma-isom-depth-2-torsion-free}
翻译成概形语言的结果。
\end{proof}










\section{模的秩}
\label{divisors-section-rank}

\noindent
本节是《代数进阶》第 \ref{more-algebra-section-rank} 节关于拟凝聚模的
对应版本。回忆，若 $X$ 是泛点为 $\eta$ 的整概形，则局部环
$\mathcal{O}_{X, \eta}$ 等于剩余域 $\kappa(\eta)$，见《态射》引理
\ref{morphisms-lemma-integral-scheme-rational-functions}。

\begin{definition}
\label{divisors-definition-rank}
设 $X$ 为泛点是 $\eta$ 的整概形，$\mathcal{F}$ 为拟凝聚
$\mathcal{O}_X$-模。$\mathcal{F}$ 的{\it 秩}定义为
$\dim_{\kappa(\eta)} \mathcal{F}_\eta$。
\end{definition}

\noindent
当两个定义都适用时，这与局部自由模之秩的定义并不冲突。若概形 $X$
不是不可约的，或者泛点处的局部环 $\mathcal{O}_{X, \eta}$ 不是域，则给
拟凝聚模定义秩看来不是好主意\footnote{若 $X$ 不可约，且
$\mathcal{F}_\eta$ 在 $\mathcal{O}_{X, \eta}$ 上自由，则采用该自由模的
秩似乎很自然。须注意，有些文献用 $\mathcal{F}_\eta$ 的长度定义秩。}

\begin{lemma}
\label{divisors-lemma-check-rank-on-affines}
设 $X$ 为整概形，$\mathcal{F}$ 为拟凝聚 $\mathcal{O}_X$-模，并设
$r$ 为一个基数。则下列条件等价：
\begin{enumerate}
\item $\mathcal{F}$ 的秩为 $r$；
\item 对所有非空仿射开集 $U \subset X$，$\mathcal{F}(U)$ 都是
$\mathcal{O}_X(U)$-模，其秩为 $r$；
\item 对某个非空仿射开集 $U \subset X$，$\mathcal{F}(U)$ 是
$\mathcal{O}_X(U)$-模，其秩为 $r$。
\end{enumerate}
\end{lemma}

\begin{proof}
略。
\end{proof}

\begin{lemma}
\label{divisors-lemma-dominant-pullback-rank}
设 $f : X \to Y$ 为整概形之间的支配态射，$\mathcal{G}$ 为拟凝聚
$\mathcal{O}_Y$-模。则 $\mathcal{G}$ 在 $Y$ 上的秩等于
$f^*\mathcal{G}$ 在 $X$ 上的秩。
\end{lemma}

\begin{proof}
使用引理 \ref{divisors-lemma-check-rank-on-affines}、《态射》引理
\ref{morphisms-lemma-dominant-between-integral} 以及《代数进阶》引理
\ref{more-algebra-lemma-dominant-pullback-rank}。
\end{proof}

\begin{lemma}
\label{divisors-lemma-rank-additive}
设 $X$ 为整概形，并设
$0 \to \mathcal{F} \to \mathcal{F}' \to \mathcal{F}'' \to 0$
为拟凝聚 $\mathcal{O}_X$-模的短正合列。则 $\mathcal{F}'$ 的秩等于
$\mathcal{F}$ 与 $\mathcal{F}''$ 的秩之和。
\end{lemma}

\begin{proof}
略。参见《代数进阶》引理 \ref{more-algebra-lemma-rank-additive}。
\end{proof}

\begin{lemma}
\label{divisors-lemma-rank-tensor}
设 $X$ 为整概形，$\mathcal{F}$ 和 $\mathcal{G}$ 为拟凝聚
$\mathcal{O}_X$-模。
\begin{enumerate}
\item $\mathcal{F} \otimes_{\mathcal{O}_X} \mathcal{G}$ 的秩等于
$\mathcal{F}$ 与 $\mathcal{G}$ 的秩之积。
\item 若 $\mathcal{F}$ 有限展示，则
$\SheafHom_{\mathcal{O}_X}(\mathcal{F}, \mathcal{G})$ 的秩等于
$\mathcal{F}$ 与 $\mathcal{G}$ 的秩之积。
\end{enumerate}
\end{lemma}

\begin{proof}
略。注意，(2) 是有意义的，因为
$\SheafHom_{\mathcal{O}_X}(\mathcal{F}, \mathcal{G})$
由《概形》第 \ref{schemes-section-quasi-coherent} 节的 (10)
可知是拟凝聚的。本引理的代数版本是《代数进阶》引理
\ref{more-algebra-lemma-rank-tensor}。
\end{proof}

\begin{lemma}
\label{divisors-lemma-finite-module-rank}
设 $X$ 为整概形，$\mathcal{F}$ 为有限型拟凝聚 $\mathcal{O}_X$-模。
则
\begin{enumerate}
\item $\mathcal{F}$ 具有有限秩 $r \geq 0$；
\item 存在非空开集 $U \subset X$，使得 $\mathcal{F}|_U$ 是秩为
$r$ 的自由模。
\end{enumerate}
\end{lemma}

\begin{proof}
代数版本见《代数进阶》引理
\ref{more-algebra-lemma-finite-module-rank}。
\end{proof}










\section{自反模}
\label{divisors-section-reflexive}

\noindent
本节是《代数进阶》第 \ref{more-algebra-section-reflexive} 节关于局部诺特
概形上凝聚模的对应版本。之所以考察凝聚模，是因为
$\SheafHom_{\mathcal{O}_X}(\mathcal{F}, \mathcal{G})$ 对任意一对凝聚
$\mathcal{O}_X$-模 $\mathcal{F}, \mathcal{G}$ 都是凝聚的，
见《模》引理
\ref{modules-lemma-internal-hom-locally-kernel-direct-sum}。

\begin{definition}
\label{divisors-definition-reflexive}
设 $X$ 为整的局部诺特概形，$\mathcal{F}$ 为凝聚
$\mathcal{O}_X$-模。$\mathcal{F}$ 的{\it 自反包}是
$\mathcal{O}_X$-模
$$
\mathcal{F}^{**} = \SheafHom_{\mathcal{O}_X}(
\SheafHom_{\mathcal{O}_X}(\mathcal{F}, \mathcal{O}_X), \mathcal{O}_X)
$$
称 $\mathcal{F}$ 是{\it 自反的}，如果自然映射
$j : \mathcal{F} \longrightarrow \mathcal{F}^{**}$ 是同构。
\end{definition}

\noindent
由引理 \ref{divisors-lemma-dual-reflexive}，自反包是自反
$\mathcal{O}_X$-模。在更一般的情形中也可以用同一定义来定义自反模，
但这看来并不十分有用。下面给出按例必须有的引理，将此定义与通常的
代数定义作比较。

\begin{lemma}
\label{divisors-lemma-check-reflexive-on-affines}
设 $X$ 为整的局部诺特概形，$\mathcal{F}$ 为凝聚
$\mathcal{O}_X$-模。则下列条件等价：
\begin{enumerate}
\item $\mathcal{F}$ 自反；
\item 对仿射开集 $U \subset X$，$\mathcal{F}(U)$ 是自反
$\mathcal{O}(U)$-模。
\end{enumerate}
\end{lemma}

\begin{proof}
略。
\end{proof}

\begin{remark}
\label{divisors-remark-different-reflexive}
若 $X$ 是域上的有限型概形，则有时会使用另一种自反模的概念
（例如见 \cite[第 5 页底部及定义 1.1.9]{HL}）。这种概念使用
$R\SheafHom$ 取值于对偶化复形 $\omega_X^\bullet$，而不是取值于
$\mathcal{O}_X$ 的版本；它也许应当另取名称，因为当 $X$ 不是
Gorenstein 概形时，两者可能不同。例如，若
$X = \Spec(k[t^3, t^4, t^5])$，则计算表明，对偶化层 $\omega_X$
按我们的定义并不自反，但按另一种定义是自反的，因为
$\omega_X \to \SheafHom(\SheafHom(\omega_X, \omega_X), \omega_X)$
是同构。
\end{remark}

\begin{lemma}
\label{divisors-lemma-reflexive-torsion-free}
设 $X$ 为整的局部诺特概形，$\mathcal{F}$ 为凝聚
$\mathcal{O}_X$-模。
\begin{enumerate}
\item 若 $\mathcal{F}$ 自反，则 $\mathcal{F}$ 无挠。
\item 映射 $j : \mathcal{F} \longrightarrow \mathcal{F}^{**}$ 是单射，
当且仅当 $\mathcal{F}$ 无挠。
\end{enumerate}
\end{lemma}

\begin{proof}
略。参见《代数进阶》引理
\ref{more-algebra-lemma-reflexive-torsion-free}。
\end{proof}

\begin{lemma}
\label{divisors-lemma-check-reflexive}
设 $X$ 为整的局部诺特概形，$\mathcal{F}$ 为凝聚
$\mathcal{O}_X$-模。则下列条件等价：
\begin{enumerate}
\item $\mathcal{F}$ 自反；
\item $\mathcal{F}_x$ 是自反 $\mathcal{O}_{X, x}$-模，这对所有
$x \in X$ 都成立；
\item $\mathcal{F}_x$ 是自反 $\mathcal{O}_{X, x}$-模，这对所有闭点
$x \in X$ 都成立。
\end{enumerate}
\end{lemma}

\begin{proof}
由《模》引理 \ref{modules-lemma-stalk-internal-hom}，(1) 与 (2)
等价。由于 $X$ 的每个点都特化到某个闭点（《性质》引理
\ref{properties-lemma-locally-Noetherian-closed-point}），
(2) 与 (3) 等价。
\end{proof}

\begin{lemma}
\label{divisors-lemma-flat-pullback-reflexive}
设 $f : X \to Y$ 为整的局部诺特概形之间的平坦态射，$\mathcal{G}$
为凝聚自反 $\mathcal{O}_Y$-模。则 $f^*\mathcal{G}$ 是凝聚自反
$\mathcal{O}_X$-模。
\end{lemma}

\begin{proof}
略。代数版本见《代数进阶》引理
\ref{more-algebra-lemma-flat-pullback-reflexive}。
\end{proof}

\begin{lemma}
\label{divisors-lemma-sequence-reflexive}
设 $X$ 为整的局部诺特概形，并设
$0 \to \mathcal{F} \to \mathcal{F}' \to \mathcal{F}''$
为凝聚 $\mathcal{O}_X$-模的正合列。若 $\mathcal{F}'$ 自反且
$\mathcal{F}''$ 无挠，则 $\mathcal{F}$ 自反。
\end{lemma}

\begin{proof}
略。参见《代数进阶》引理
\ref{more-algebra-lemma-sequence-reflexive}。
\end{proof}

\begin{lemma}
\label{divisors-lemma-dual-reflexive}
设 $X$ 为整的局部诺特概形，$\mathcal{F}$、$\mathcal{G}$ 为凝聚
$\mathcal{O}_X$-模。若 $\mathcal{G}$ 自反，则
$\SheafHom_{\mathcal{O}_X}(\mathcal{F}, \mathcal{G})$ 自反。
\end{lemma}

\begin{proof}
该断言有意义，因为
$\SheafHom_{\mathcal{O}_X}(\mathcal{F}, \mathcal{G})$
由《概形的上同调》引理
\ref{coherent-lemma-tensor-hom-coherent} 可知是凝聚的。断言的证明见
《代数进阶》引理 \ref{more-algebra-lemma-dual-reflexive}。
略去一些细节。
\end{proof}

\begin{remark}
\label{divisors-remark-tensor}
设 $X$ 为整的局部诺特概形。由引理 \ref{divisors-lemma-dual-reflexive}，
我们知道自反包 $\mathcal{F}^{**}$ 对凝聚 $\mathcal{O}_X$-模而言
是凝聚自反的。考虑范畴 $\mathcal{C}$，其对象为凝聚自反
$\mathcal{O}_X$-模。取自反包给出包含函子
$\mathcal{C} \to \textit{Coh}(\mathcal{O}_X)$ 的左伴随。
注意，$\mathcal{C}$ 是具有核和余核的加性范畴。事实上，给定
$\varphi : \mathcal{F} \to \mathcal{G}$，并设它是 $\mathcal{C}$ 中的态射，
通常的核 $\Ker(\varphi)$ 是自反的（引理
\ref{divisors-lemma-sequence-reflexive}），而通常余核的自反包
$\Coker(\varphi)^{**}$ 是 $\mathcal{C}$ 中的余核。此外，
$\mathcal{C}$ 继承了张量积
$$
\mathcal{F} \otimes_\mathcal{C} \mathcal{G} =
(\mathcal{F} \otimes_{\mathcal{O}_X} \mathcal{G})^{**}
$$
它是结合且对称的。还存在内 Hom：对任意三个对象
$\mathcal{F}, \mathcal{G}, \mathcal{H}$（它们属于 $\mathcal{C}$），有恒等式
$$
\Hom_\mathcal{C}(\mathcal{F} \otimes_\mathcal{C} \mathcal{G}, \mathcal{H}) =
\Hom_\mathcal{C}(\mathcal{F},
\SheafHom_{\mathcal{O}_X}(\mathcal{G}, \mathcal{H}))
$$
见《模》引理 \ref{modules-lemma-internal-hom}。在 $\mathcal{C}$ 中，
每个对象 $\mathcal{F}$ 都有一个{\it 对偶对象}
$\SheafHom_{\mathcal{O}_X}(\mathcal{F}, \mathcal{O}_X)$。
若不对 $X$ 加进一步条件，则可能有
$$
\SheafHom_{\mathcal{O}_X}(\mathcal{F}, \mathcal{G}) \not \cong
\SheafHom_{\mathcal{O}_X}(\mathcal{F}, \mathcal{O}_X)
\otimes_\mathcal{C} \mathcal{G}
\quad\text{且}\quad
\mathcal{F} \otimes_\mathcal{C}
\SheafHom_{\mathcal{O}_X}(\mathcal{F}, \mathcal{O}_X)
\not \cong \mathcal{O}_X
$$
其中 $\mathcal{F}, \mathcal{G}$ 的秩为 $1$，且它们是 $\mathcal{C}$ 中的对象。
为构造一个例子，令 $X = \Spec(R)$，其中 $R$ 如《代数进阶》例
\ref{more-algebra-example-ring-not-S2} 所述，并令 $\mathcal{F}, \mathcal{G}$
为与 $M$ 对应的模。略去计算。
\end{remark}

\begin{lemma}
\label{divisors-lemma-reflexive-depth-2}
设 $X$ 为整的局部诺特概形，$\mathcal{F}$ 为凝聚
$\mathcal{O}_X$-模。则下列条件等价：
\begin{enumerate}
\item $\mathcal{F}$ 自反；
\item 对每个 $x \in X$，下列情形之一成立：
\begin{enumerate}
\item $\mathcal{F}_x$ 是自反 $\mathcal{O}_{X, x}$-模；或
\item $\text{depth}(\mathcal{F}_x) \geq 2$。
\end{enumerate}
\end{enumerate}
\end{lemma}

\begin{proof}
略。参见《代数进阶》引理
\ref{more-algebra-lemma-reflexive-depth-2}。
\end{proof}

\begin{lemma}
\label{divisors-lemma-reflexive-S2}
设 $X$ 为整的局部诺特概形，$\mathcal{F}$ 为凝聚自反
$\mathcal{O}_X$-模，并设 $x \in X$。
\begin{enumerate}
\item 若 $\text{depth}(\mathcal{O}_{X, x}) \geq 2$，则
$\text{depth}(\mathcal{F}_x) \geq 2$。
\item 若 $X$ 具有 $(S_2)$ 性质，则 $\mathcal{F}$ 具有 $(S_2)$ 性质。
\end{enumerate}
\end{lemma}

\begin{proof}
略。参见《代数进阶》引理
\ref{more-algebra-lemma-reflexive-S2}。
\end{proof}

\begin{lemma}
\label{divisors-lemma-reflexive-S2-extend}
设 $X$ 为整的局部诺特概形，$j : U \to X$ 为补集是 $Z$ 的开子概形。
假设 $\mathcal{O}_{X, z}$ 的深度 $\geq 2$，且这一条件对所有
$z \in Z$ 成立。则 $j^*$ 与 $j_*$ 给出凝聚自反 $\mathcal{O}_X$-模范畴
同凝聚自反 $\mathcal{O}_U$-模范畴之间的等价。
\end{lemma}

\begin{proof}
设 $\mathcal{F}$ 为凝聚自反 $\mathcal{O}_X$-模。对 $z \in Z$，
由引理 \ref{divisors-lemma-reflexive-S2}，茎 $\mathcal{F}_z$ 的深度
$\geq 2$。因此，由引理 \ref{divisors-lemma-depth-2-hartog}，
$\mathcal{F} \to j_*j^*\mathcal{F}$ 是同构。反之，设 $\mathcal{G}$
为凝聚自反 $\mathcal{O}_U$-模。只须证明 $j_*\mathcal{G}$ 是凝聚自反
$\mathcal{O}_X$-模。为此，可以假设 $X$ 仿射。由《性质》引理
\ref{properties-lemma-lift-finite-presentation}，存在凝聚
$\mathcal{O}_X$-模 $\mathcal{F}$，使得
$\mathcal{G} = j^*\mathcal{F}$。取 $\mathcal{F}$ 的自反包后，
可以假设 $\mathcal{F}$ 自反
（见上文讨论，特别是引理 \ref{divisors-lemma-dual-reflexive}）。
由上述论证，$j_*\mathcal{G} = j_*j^*\mathcal{F} = \mathcal{F}$，
这正是所需结论。
\end{proof}

\noindent
若概形正规，则自反等价于无挠并满足 $(S_2)$。

\begin{lemma}
\label{divisors-lemma-reflexive-over-normal}
设 $X$ 为整的局部诺特正规概形，$\mathcal{F}$ 为凝聚
$\mathcal{O}_X$-模。则下列条件等价：
\begin{enumerate}
\item $\mathcal{F}$ 自反；
\item $\mathcal{F}$ 无挠并具有性质 $(S_2)$；
\item 存在开子概形 $j : U \to X$，使得
\begin{enumerate}
\item $X \setminus U$ 的每个不可约分支的余维都 $\geq 2$
（余维在 $X$ 中计算）；
\item $j^*\mathcal{F}$ 有限局部自由；
\item $\mathcal{F} = j_*j^*\mathcal{F}$。
\end{enumerate}
\end{enumerate}
\end{lemma}

\begin{proof}
使用引理 \ref{divisors-lemma-check-reflexive-on-affines}，(1) 与 (2) 的等价性
由《代数进阶》引理
\ref{more-algebra-lemma-reflexive-over-normal} 得出。设
$U \subset X$ 如 (3) 所述。由《性质》引理
\ref{properties-lemma-criterion-normal}，有
$\text{depth}(\mathcal{O}_{X, x}) \geq 2$，其中 $x \not \in U$。
由于有限局部自由模自反，
引理 \ref{divisors-lemma-reflexive-S2-extend} 表明 (3) 蕴含 (1)。

\medskip\noindent
假设 (1)。设 $U \subset X$ 为使
$j^*\mathcal{F} = \mathcal{F}|_U$ 有限局部自由的极大开子概形。
于是 (3)(b) 成立。设 $x \in X$。若 $\mathcal{F}_x$ 是自由
$\mathcal{O}_{X, x}$-模，则由《模》引理
\ref{modules-lemma-finite-presentation-stalk-free} 有 $x \in U$。
若 $\dim(\mathcal{O}_{X, x}) \leq 1$，则
$\mathcal{O}_{X, x}$ 是域或离散赋值环（《性质》引理
\ref{properties-lemma-criterion-normal}），从而 $\mathcal{F}_x$ 是自由的
（《代数进阶》引理
\ref{more-algebra-lemma-dedekind-torsion-free-flat}）。因此
$x \not \in U \Rightarrow \dim(\mathcal{O}_{X, x}) \geq 2$。
此时《性质》引理
\ref{properties-lemma-codimension-local-ring} 给出 (3)(a)。
再次使用前述《性质》引理
\ref{properties-lemma-criterion-normal}，还可知
$\text{depth}(\mathcal{O}_{X, x}) \geq 2$，其中 $x \not \in U$；
于是 (3)(c) 由引理
\ref{divisors-lemma-reflexive-S2-extend} 得出。
\end{proof}

\begin{lemma}
\label{divisors-lemma-describe-reflexive-hull}
设 $X$ 为泛点是 $\eta$ 的整局部诺特正规概形，$\mathcal{F}$、
$\mathcal{G}$ 为凝聚 $\mathcal{O}_X$-模。设
$T : \mathcal{G}_\eta \to \mathcal{F}_\eta$ 为线性映射。则 $T$
可延拓为映射 $\mathcal{G} \to \mathcal{F}^{**}$，且该映射为
$\mathcal{O}_X$-模映射，当且仅当
\begin{itemize}
\item[(*)] 对每个 $x \in X$，若 $\dim(\mathcal{O}_{X, x}) = 1$，则有
$$
T\left(\Im(\mathcal{G}_x \to \mathcal{G}_\eta)\right) \subset
\Im(\mathcal{F}_x \to \mathcal{F}_\eta).
$$
\end{itemize}
\end{lemma}

\begin{proof}
由于 $\mathcal{F}^{**}$ 无挠，且
$\mathcal{F}_\eta = \mathcal{F}^{**}_\eta$，所以延拓若存在就唯一。
因此，只须在 $X$ 的一个开覆盖的各成员上证明本引理；也就是说，可以假设
$X$ 仿射。在这种情形下，我们问的是如下代数问题：设 $R$ 为分式域是
$K$ 的诺特正规整环，$M$、$N$ 为有限 $R$-模，并设
$T : M \otimes_R K \to N \otimes_R K$ 为 $K$-线性映射。$T$
何时能延拓为映射 $N \to M^{**}$？由《代数进阶》引理
\ref{more-algebra-lemma-describe-reflexive-hull}，这当且仅当
$N_\mathfrak p$ 映入 $(M/M_{tors})_\mathfrak p$，而这对每个高为
$1$ 的素理想 $\mathfrak p$（属于 $R$）都成立。这恰好是本引理的
条件 $(*)$。
\end{proof}

\begin{lemma}
\label{divisors-lemma-reflexive-over-regular-dim-2}
设 $X$ 为维数 $\leq 2$ 的正则概形，$\mathcal{F}$ 为凝聚
$\mathcal{O}_X$-模。则下列条件等价：
\begin{enumerate}
\item $\mathcal{F}$ 自反；
\item $\mathcal{F}$ 有限局部自由。
\end{enumerate}
\end{lemma}

\begin{proof}
有限局部自由模显然自反。反过来，我们将说明：若 $\mathcal{F}$ 自反，
则 $\mathcal{F}_x$ 是自由 $\mathcal{O}_{X, x}$-模，这对所有
$x \in X$ 都成立。由《代数》引理
\ref{algebra-lemma-finite-projective} 以及 $\mathcal{F}$ 凝聚这一事实，
这就足够了。若 $\dim(\mathcal{O}_{X, x}) = 0$，则
$\mathcal{O}_{X, x}$ 是域，结论显然。若
$\dim(\mathcal{O}_{X, x}) = 1$，则 $\mathcal{O}_{X, x}$ 是离散赋值环
（《代数》引理 \ref{algebra-lemma-characterize-dvr}），并且
$\mathcal{F}_x$ 无挠。因此，由《代数进阶》引理
\ref{more-algebra-lemma-dedekind-torsion-free-flat}，
$\mathcal{F}_x$ 是自由的。若 $\dim(\mathcal{O}_{X, x}) = 2$，
则 $\mathcal{O}_{X, x}$ 是维数为 $2$ 的正则局部环。由《代数进阶》引理
\ref{more-algebra-lemma-reflexive-over-normal}，
$\mathcal{F}_x$ 的深度 $\geq 2$。因此，由《代数》引理
\ref{algebra-lemma-regular-mcm-free}，$\mathcal{F}_x$ 是自由的。
\end{proof}






\section{有效 Cartier 除子}
\label{divisors-section-effective-Cartier-divisors}

\noindent
在定义其他类型的除子之前，我们先定义有效 Cartier 除子。

\begin{definition}
\label{divisors-definition-effective-Cartier-divisor}
设 $S$ 为概形。
\begin{enumerate}
\item $S$ 的一个{\it 局部主闭子概形}，是其理想层可局部由单个元素
生成的闭子概形。
\item $S$ 上的一个{\it 有效 Cartier 除子}，是闭子概形
$D \subset S$，其理想层 $\mathcal{I}_D \subset \mathcal{O}_S$
是可逆 $\mathcal{O}_S$-模。
\end{enumerate}
\end{definition}

\noindent
因此，有效 Cartier 除子是局部主闭子概形，但反之并非总是成立。
在最强的可能意义下，有效 Cartier 除子是纯余维 $1$ 的闭子概形。
具体而言，它们局部由单个非零因子截出。特别地，它们无处稠密。

\begin{lemma}
\label{divisors-lemma-characterize-effective-Cartier-divisor}
设 $S$ 为概形，$D \subset S$ 为闭子概形。则下列条件等价：
\begin{enumerate}
\item 子概形 $D$ 是 $S$ 上的有效 Cartier 除子。
\item 对每个 $x \in D$，都存在仿射开邻域
$\Spec(A) = U \subset S$，它是 $x$ 的邻域，并且
$U \cap D = \Spec(A/(f))$，其中 $f \in A$ 是非零因子。
\end{enumerate}
\end{lemma}

\begin{proof}
假设 (1)。对每个 $x \in D$，都存在仿射开邻域
$\Spec(A) = U \subset S$，它是 $x$ 的邻域，并且
$\mathcal{I}_D|_U \cong \mathcal{O}_U$。换言之，存在截面
$f \in \Gamma(U, \mathcal{I}_D)$，它自由生成限制
$\mathcal{I}_D|_U$。因此 $f \in A$，且乘法映射
$f : A \to A$ 是单射。又因为 $\mathcal{I}_D$ 拟凝聚，所以
$D \cap U = \Spec(A/(f))$。

\medskip\noindent
假设 (2)。设 $x \in D$。根据假设，存在仿射开邻域
$\Spec(A) = U \subset S$，它是 $x$ 的邻域，并且
$U \cap D = \Spec(A/(f))$，其中 $f \in A$ 是非零因子。
于是 $\mathcal{I}_D|_U \cong \mathcal{O}_U$，因为它等于
$\widetilde{(f)} \cong \widetilde{A} \cong \mathcal{O}_U$。
当然，$\mathcal{I}_D$ 在开子概形 $S \setminus D$ 上的限制同构于
$\mathcal{O}_{S \setminus D}$。因此 $\mathcal{I}_D$ 是可逆
$\mathcal{O}_S$-模。
\end{proof}

\begin{lemma}
\label{divisors-lemma-complement-locally-principal-closed-subscheme}
设 $S$ 为概形，$Z \subset S$ 为局部主闭子概形，并设
$U = S \setminus Z$。则 $U \to S$ 是仿射态射。
\end{lemma}

\begin{proof}
该问题在 $S$ 上是局部的，见《态射》引理
\ref{morphisms-lemma-characterize-affine}。因此可以假设
$S = \Spec(A)$ 且 $Z = V(f)$，其中 $f \in A$。此时
$U = D(f) = \Spec(A_f)$ 仿射，故 $U \to S$ 仿射。
\end{proof}

\begin{lemma}
\label{divisors-lemma-complement-effective-Cartier-divisor}
设 $S$ 为概形，$D \subset S$ 为有效 Cartier 除子，并设
$U = S \setminus D$。则 $U \to S$ 是仿射态射，且 $U$ 在 $S$ 中
概形论稠密。
\end{lemma}

\begin{proof}
仿射性由引理
\ref{divisors-lemma-complement-locally-principal-closed-subscheme} 给出。
稠密性问题在 $S$ 上是局部的，见《态射》引理
\ref{morphisms-lemma-characterize-scheme-theoretically-dense}。
因此可以假设 $S = \Spec(A)$，而 $D$ 对应于非零因子 $f \in A$，
见引理 \ref{divisors-lemma-characterize-effective-Cartier-divisor}。于是
$A \subset A_f$，这说明 $U \subset S$ 概形论稠密，见《态射》例
\ref{morphisms-example-scheme-theoretic-closure}。
\end{proof}

\begin{lemma}
\label{divisors-lemma-effective-Cartier-makes-dimension-drop}
设 $S$ 为概形，$D \subset S$ 为有效 Cartier 除子，并设
$s \in D$。若 $\dim_s(S) < \infty$，则
$\dim_s(D) < \dim_s(S)$。
\end{lemma}

\begin{proof}
假设 $\dim_s(S) < \infty$。设 $U = \Spec(A) \subset S$ 为 $s$
的仿射开邻域，使得 $\dim(U) = \dim_s(S)$，并使得
$D = V(f)$，其中 $f \in A$ 是非零因子（见引理
\ref{divisors-lemma-characterize-effective-Cartier-divisor}）。回忆，
$\dim(U)$ 是环 $A$ 的 Krull 维数，而 $\dim(U \cap D)$ 是环
$A/(f)$ 的 Krull 维数。于是 $f$ 不属于 $A$ 的任何极小素理想。
因此，把 $A/(f)$ 中任一极大素理想链看作 $A$ 中的素理想链后，
都可以添加一个极小素理想来延长它。
\end{proof}

\begin{definition}
\label{divisors-definition-sum-effective-Cartier-divisors}
设 $S$ 为概形。给定有效 Cartier 除子 $D_1$、$D_2$，它们定义在 $S$ 上，
令 $D = D_1 + D_2$ 为 $S$ 的闭子概形，它对应于拟凝聚理想层
$\mathcal{I}_{D_1}\mathcal{I}_{D_2} \subset \mathcal{O}_S$。
称它为{\it 有效 Cartier 除子 $D_1$ 与 $D_2$ 的和}。
\end{definition}

\noindent
显然，可以定义和 $\sum n_iD_i$，只要给定有限多个有效 Cartier
除子 $D_i$（定义在 $X$ 上）和非负整数 $n_i$。

\begin{lemma}
\label{divisors-lemma-sum-effective-Cartier-divisors}
两个有效 Cartier 除子的和仍是有效 Cartier 除子。
\end{lemma}

\begin{proof}
略。局部地，$f_1, f_2 \in A$ 都是非零因子，因而
$f_1f_2 \in A$ 也是非零因子。
\end{proof}

\begin{lemma}
\label{divisors-lemma-difference-effective-Cartier-divisors}
设 $X$ 为概形，$D, D'$ 为 $X$ 上两个有效 Cartier 除子。
若 $D \subset D'$（作为 $X$ 的闭子概形），则存在有效 Cartier
除子 $D''$，使得 $D' = D + D''$。
\end{lemma}

\begin{proof}
略。
\end{proof}

\begin{lemma}
\label{divisors-lemma-sum-closed-subschemes-effective-Cartier}
设 $X$ 为概形，$Z, Y$ 为 $X$ 的两个闭子概形，其理想层分别为
$\mathcal{I}$ 和 $\mathcal{J}$。若 $\mathcal{I}\mathcal{J}$
定义有效 Cartier 除子 $D \subset X$，则 $Z$ 和 $Y$ 都是
有效 Cartier 除子，且 $D = Z + Y$。
\end{lemma}

\begin{proof}
应用引理 \ref{divisors-lemma-characterize-effective-Cartier-divisor}，得到如下
代数情形：$A$ 是环，$I, J \subset A$ 是理想，$f \in A$ 是非零因子，
且 $IJ = (f)$。因此，结论由《代数》引理
\ref{algebra-lemma-product-ideals-principal} 得出。
\end{proof}

\begin{lemma}
\label{divisors-lemma-sum-effective-Cartier-divisors-union}
设 $X$ 为概形，$D, D' \subset X$ 为有效 Cartier 除子，并假设
概形论交 $D \cap D'$ 是 $D'$ 上的有效 Cartier 除子。则
$D + D'$ 是 $D$ 与 $D'$ 的概形论并。
\end{lemma}

\begin{proof}
概形论交与并的定义见《态射》定义
\ref{morphisms-definition-scheme-theoretic-intersection-union}。
为证明本引理，局部工作（使用引理
\ref{divisors-lemma-characterize-effective-Cartier-divisor}）后得到如下代数问题：
给定环 $A$ 与非零因子 $f_1, f_2 \in A$，且 $f_1$ 在
$A/f_2A$ 中的像是非零因子，证明
$f_1A \cap f_2A = f_1f_2A$。略去这个直接的论证。
\end{proof}

\noindent
回忆，我们已在《概形》定义
\ref{schemes-definition-inverse-image-closed-subscheme} 中定义了闭子概形
在任意概形态射下的逆像。

\begin{lemma}
\label{divisors-lemma-pullback-locally-principal}
设 $f : S' \to S$ 为概形态射，$Z \subset S$ 为局部主闭子概形。
则逆像 $f^{-1}(Z)$ 是 $S'$ 的局部主闭子概形。
\end{lemma}

\begin{proof}
略。
\end{proof}

\begin{definition}
\label{divisors-definition-pullback-effective-Cartier-divisor}
设 $f : S' \to S$ 为概形态射，$D \subset S$ 为有效 Cartier 除子。
称{\it $D$ 沿 $f$ 的拉回有定义}，如果闭子概形
$f^{-1}(D) \subset S'$ 是有效 Cartier 除子。此时将它记作 $f^*D$ 或
$f^{-1}(D)$，并称之为{\it 有效 Cartier 除子的拉回}。
\end{definition}

\noindent
$f^{-1}(D)$ 是有效 Cartier 除子这一条件在实践中经常满足。
下面给出一个示例性引理。

\begin{lemma}
\label{divisors-lemma-pullback-effective-Cartier-defined}
设 $f : X \to Y$ 为概形态射，$D \subset Y$ 为有效 Cartier 除子。
在下列每种情形中，$D$ 沿 $f$ 的拉回都有定义：
\begin{enumerate}
\item $f(x) \not \in D$ 对每个弱伴随点 $x$（属于 $X$）都成立；
\item $X$、$Y$ 均为整概形，且 $f$ 支配；
\item $X$ 约化，并且 $f(\xi) \not \in D$ 对任一泛点 $\xi$
（属于 $X$ 的任一不可约分支）都成立；
\item $X$ 局部诺特，并且 $f(x) \not \in D$ 对任一伴随点 $x$
（属于 $X$）都成立；
\item $X$ 局部诺特且没有嵌入点，并且 $f(\xi) \not \in D$ 对任一
泛点 $\xi$（属于 $X$ 的某个不可约分支）都成立；
\item $f$ 平坦；
\item 可按需在此添加更多情形。
\end{enumerate}
\end{lemma}

\begin{proof}
该问题在 $X$ 上是局部的，因此可约化到如下情形：
$X = \Spec(A)$，$Y = \Spec(R)$，$f$ 由
$\varphi : R \to A$ 给出，并且 $D = \Spec(R/(t))$，其中
$t \in R$ 是非零因子。每种情形中的目标都是证明
$\varphi(t) \in A$ 是非零因子。

\medskip\noindent
情形 (1) 由《代数》引理
\ref{algebra-lemma-weakly-ass-zero-divisors} 得出。由引理
\ref{divisors-lemma-ass-weakly-ass}，情形 (4) 是 (1) 的特例。情形 (5)
由 (4) 和定义得出。由引理 \ref{divisors-lemma-weakass-reduced}，情形 (3)
是 (1) 的特例。情形 (2) 是 (3) 的特例。若 $R \to A$ 平坦，则
$t : R \to R$ 是单射这一事实表明 $t : A \to A$ 是单射。
这证明了 (6)。
\end{proof}

\begin{lemma}
\label{divisors-lemma-pullback-effective-Cartier-divisors-additive}
设 $f : S' \to S$ 为概形态射，$D_1$、$D_2$ 为 $S$ 上的
有效 Cartier 除子。若 $D_1$ 与 $D_2$ 的拉回都有定义，则
$D = D_1 + D_2$ 的拉回有定义，且
$f^*D = f^*D_1 + f^*D_2$。
\end{lemma}

\begin{proof}
略。
\end{proof}





\section{有效 Cartier 除子与可逆层}
\label{divisors-section-effective-Cartier-invertible}

\noindent
由于有效 Cartier 除子具有可逆理想层（定义
\ref{divisors-definition-effective-Cartier-divisor}），下述定义是有意义的。

\begin{definition}
\label{divisors-definition-invertible-sheaf-effective-Cartier-divisor}
设 $S$ 为概形，$D \subset S$ 为有效 Cartier 除子，其理想层为
$\mathcal{I}_D$。
\begin{enumerate}
\item {\it 可逆层 $\mathcal{O}_S(D)$（与 $D$ 伴随）}定义为
$$
\mathcal{O}_S(D) =
\SheafHom_{\mathcal{O}_S}(\mathcal{I}_D, \mathcal{O}_S) =
\mathcal{I}_D^{\otimes -1}.
$$
\item {\it 典范截面}通常记作 $1$ 或 $1_D$；它是
$\mathcal{O}_S(D)$ 的整体截面，对应于包含映射
$\mathcal{I}_D \to \mathcal{O}_S$。
\item 记
$\mathcal{O}_S(-D) = \mathcal{O}_S(D)^{\otimes -1} = \mathcal{I}_D$。
\item 给定第二个有效 Cartier 除子 $D' \subset S$，定义
$\mathcal{O}_S(D - D') =
\mathcal{O}_S(D) \otimes_{\mathcal{O}_S} \mathcal{O}_S(-D')$。
\end{enumerate}
\end{definition}

\noindent
作几点说明。下文将看到，赋值 $D \mapsto \mathcal{O}_S(D)$ 把有效
Cartier 除子的加法（定义
\ref{divisors-definition-sum-effective-Cartier-divisors}）变为 $S$ 的 Picard 群
中的加法（引理
\ref{divisors-lemma-invertible-sheaf-sum-effective-Cartier-divisors}）。
然而，上述定义中的表达式 $D - D'$ 没有任何几何含义。更确切地说，
可以把 $S$ 上的有效 Cartier 除子之集看作交换幺半群
$\text{EffCart}(S)$，其零元为空的有效 Cartier 除子。于是赋值
$(D, D') \mapsto \mathcal{O}_S(D - D')$ 定义群同态
$$
\text{EffCart}(S)^{gp} \longrightarrow \Pic(S)
$$
其中左端是 $\text{EffCart}(S)$ 的群完备化。换言之，当写
$\mathcal{O}_S(D - D')$ 时，可以把 $D - D'$ 看作
$\text{EffCart}(S)^{gp}$ 的元素。

\begin{lemma}
\label{divisors-lemma-conormal-effective-Cartier-divisor}
设 $S$ 为概形，$D \subset S$ 为有效 Cartier 除子。则余法层为
$\mathcal{C}_{D/S} = \mathcal{I}_D|_D =
\mathcal{O}_S(-D)|_D$，而法层为
$\mathcal{N}_{D/S} = \mathcal{O}_S(D)|_D$。
\end{lemma}

\begin{proof}
这由《态射》引理 \ref{morphisms-lemma-affine-conormal} 得出。
\end{proof}

\begin{lemma}
\label{divisors-lemma-ses-add-divisor}
设 $X$ 为概形，$D, D' \subset X$ 为有效 Cartier 除子，且
$D \subset D'$；再设 $C \subset X$ 为满足 $D' = D + C$ 的有效
Cartier 除子（其存在性由引理
\ref{divisors-lemma-difference-effective-Cartier-divisors} 保证）。则有短正合列
$$
0 \to \mathcal{O}_X(-D)|_C \to \mathcal{O}_{D'} \to \mathcal{O}_D \to 0
$$
它由 $\mathcal{O}_X$-模组成。
\end{lemma}

\begin{proof}
在本引理的陈述和证明中，我们使用《态射》引理
\ref{morphisms-lemma-i-star-equivalence} 的等价，把 $X$ 的闭子概形
上的拟凝聚模视为 $X$ 上的拟凝聚模。设 $\mathcal{I}$ 为 $D$ 在
$D'$ 中的理想层，并设 $\mathcal{I}_D = \mathcal{O}_X(-D)$ 为
$D$ 在 $X$ 中的理想层。于是有典范满射
$\mathcal{I}_D \to \mathcal{I}$。因此，只须证明该映射的核等于
$\mathcal{I}_C\mathcal{I}_D$（采用显然的记号）。这可局部检验，
故可假设 $X = \Spec(A)$、$D = \Spec(A/fA)$ 且
$C = \Spec(A/gA)$，其中 $f, g \in A$ 为非零因子
（引理 \ref{divisors-lemma-characterize-effective-Cartier-divisor}）。
于是 $\mathcal{I}_D$ 对应于 $fA$，非零因子 $fg \in A$ 生成
$D'$ 的理想，而 $\mathcal{I}$ 对应于 $I = fA/fgA$。结论随即清楚。
\end{proof}

\begin{lemma}
\label{divisors-lemma-invertible-sheaf-sum-effective-Cartier-divisors}
设 $S$ 为概形，$D_1$、$D_2$ 为 $S$ 上的有效 Cartier 除子，并设
$D = D_1 + D_2$。则存在唯一同构
$$
\mathcal{O}_S(D_1) \otimes_{\mathcal{O}_S} \mathcal{O}_S(D_2)
\longrightarrow
\mathcal{O}_S(D)
$$
它把 $1_{D_1} \otimes 1_{D_2}$ 映到 $1_D$。
\end{lemma}

\begin{proof}
略。
\end{proof}

\begin{lemma}
\label{divisors-lemma-pullback-effective-Cartier-divisors}
设 $f : S' \to S$ 为概形态射，$D$ 为 $S$ 上的有效 Cartier 除子。
若 $D$ 的拉回有定义，则
$f^*\mathcal{O}_S(D) = \mathcal{O}_{S'}(f^*D)$，并且典范截面
$1_D$ 拉回为典范截面 $1_{f^*D}$。
\end{lemma}

\begin{proof}
略。
\end{proof}

\begin{definition}
\label{divisors-definition-regular-section}
设 $(X, \mathcal{O}_X)$ 为局部赋环空间，$\mathcal{L}$ 为 $X$ 上的
可逆层。若整体截面 $s \in \Gamma(X, \mathcal{L})$ 所给出的映射
$\mathcal{O}_X \to \mathcal{L}$（$f \mapsto fs$）是单射，则称它为
{\it 正则截面}。
\end{definition}

\begin{lemma}
\label{divisors-lemma-regular-section-structure-sheaf}
设 $X$ 为局部赋环空间，$f \in \Gamma(X, \mathcal{O}_X)$。
则下列条件等价：
\begin{enumerate}
\item $f$ 是正则截面；
\item 对任意 $x \in X$，像 $f \in \mathcal{O}_{X, x}$ 是非零因子。
\end{enumerate}
若 $X$ 是概形，则它们还等价于
\begin{enumerate}
\item[(3)] 对任意仿射开集 $\Spec(A) = U \subset X$，像 $f \in A$
是非零因子；
\item[(4)] 存在仿射开覆盖 $X = \bigcup \Spec(A_i)$，使得 $f$ 在
$A_i$ 中的像对所有 $i$ 都是非零因子。
\end{enumerate}
\end{lemma}

\begin{proof}
略。
\end{proof}

\noindent
注意，整体截面 $s$ 属于某个可逆 $\mathcal{O}_X$-模
$\mathcal{L}$；它可以视为 $\mathcal{O}_X$-模映射
$s : \mathcal{O}_X \to \mathcal{L}$。因此，其对偶是映射
$s : \mathcal{L}^{\otimes -1} \to \mathcal{O}_X$。
（对偶可逆层的定义见《模》定义
\ref{modules-definition-powers}。）

\begin{definition}
\label{divisors-definition-zero-scheme-s}
设 $X$ 为概形，$\mathcal{L}$ 为可逆层，并设
$s \in \Gamma(X, \mathcal{L})$ 为整体截面。$s$ 的{\it 零点概形}
是闭子概形 $Z(s) \subset X$，它由拟凝聚理想层
$\mathcal{I} \subset \mathcal{O}_X$ 定义；该理想层是映射
$s : \mathcal{L}^{\otimes -1} \to \mathcal{O}_X$ 的像。
\end{definition}

\begin{lemma}
\label{divisors-lemma-zero-scheme}
设 $X$ 为概形，$\mathcal{L}$ 为可逆层，并设
$s \in \Gamma(X, \mathcal{L})$。
\begin{enumerate}
\item 考虑闭浸入 $i : Z \to X$，要求
$i^*s \in \Gamma(Z, i^*\mathcal{L})$ 为零，并按包含关系排序。
零点概形 $Z(s)$ 是该有序集的极大元。
\item 对任意概形态射 $f : Y \to X$，有
$f^*s = 0$（在 $\Gamma(Y, f^*\mathcal{L})$ 中），当且仅当 $f$
经由 $Z(s)$ 分解。
\item 零点概形 $Z(s)$ 是局部主闭子概形。
\item 零点概形 $Z(s)$ 是有效 Cartier 除子，当且仅当 $s$ 是
$\mathcal{L}$ 的正则截面。
\end{enumerate}
\end{lemma}

\begin{proof}
略。
\end{proof}

\begin{lemma}
\label{divisors-lemma-characterize-OD}
\begin{slogan}
概形上的有效 Cartier 除子等同于带有指定正则整体截面的可逆层。
\end{slogan}
设 $X$ 为概形。
\begin{enumerate}
\item 若 $D \subset X$ 是有效 Cartier 除子，则典范截面 $1_D$
（属于 $\mathcal{O}_X(D)$）是正则的。
\item 反之，若 $s$ 是可逆层 $\mathcal{L}$ 的正则截面，则存在唯一
有效 Cartier 除子 $D = Z(s) \subset X$ 和唯一同构
$\mathcal{O}_X(D) \to \mathcal{L}$，它把 $1_D$ 映到 $s$。
\end{enumerate}
构造 $D \mapsto (\mathcal{O}_X(D), 1_D)$ 和
$(\mathcal{L}, s) \mapsto Z(s)$ 给出互逆的映射
$$
\left\{
\begin{matrix}
\text{effective Cartier divisors on }X
\end{matrix}
\right\}
\leftrightarrow
\left\{
\begin{matrix}
\text{isomorphism classes of pairs }(\mathcal{L}, s)\\
\text{consisting of an invertible }
\mathcal{O}_X\text{-模}\\
\mathcal{L}\text{ and a regular global section }s
\end{matrix}
\right\}
$$
\end{lemma}

\begin{proof}
略。
\end{proof}

\begin{remark}
\label{divisors-remark-ses-regular-section}
设 $X$ 为概形，$\mathcal{L}$ 为可逆 $\mathcal{O}_X$-模，$s$ 为
$\mathcal{L}$ 的正则截面。则零点概形 $D = Z(s)$ 是 $X$ 上的
有效 Cartier 除子，并有短正合列
$$
0 \to \mathcal{O}_X \to \mathcal{L} \to i_*(\mathcal{L}|_D) \to 0
\quad\text{且}\quad
0 \to \mathcal{L}^{\otimes -1} \to \mathcal{O}_X \to i_*\mathcal{O}_D \to 0.
$$
给定有效 Cartier 除子 $D \subset X$，使用引理
\ref{divisors-lemma-characterize-OD} 和
\ref{divisors-lemma-conormal-effective-Cartier-divisor}，得到
$$
0 \to \mathcal{O}_X \to \mathcal{O}_X(D) \to i_*(\mathcal{N}_{D/X}) \to 0
\quad\text{且}\quad
0 \to \mathcal{O}_X(-D) \to \mathcal{O}_X \to i_*(\mathcal{O}_D) \to 0
$$
\end{remark}






\section{诺特概形上的有效 Cartier 除子}
\label{divisors-section-Noetherian-effective-Cartier}

\noindent
在局部诺特的情形中，关于有效 Cartier 除子和正则截面的许多讨论
都有所简化。

\begin{lemma}
\label{divisors-lemma-regular-section-associated-points}
设 $X$ 为局部诺特概形，$\mathcal{L}$ 为可逆
$\mathcal{O}_X$-模，并设 $s \in \Gamma(X, \mathcal{L})$。则 $s$
是正则截面，当且仅当 $s$ 在 $X$ 的伴随点处不消失。
\end{lemma}

\begin{proof}
略。提示：约化到仿射情形，并令 $\mathcal{L}$ 平凡，然后使用引理
\ref{divisors-lemma-regular-section-structure-sheaf} 和《代数》引理
\ref{algebra-lemma-ass-zero-divisors}。
\end{proof}

\begin{lemma}
\label{divisors-lemma-effective-Cartier-in-points}
设 $X$ 为局部诺特概形，$D \subset X$ 为闭子概形，对应于拟凝聚理想层
$\mathcal{I} \subset \mathcal{O}_X$。
\begin{enumerate}
\item 若对每个 $x \in D$，理想
$\mathcal{I}_x \subset \mathcal{O}_{X, x}$ 都可以由一个元素生成，
则 $D$ 是局部主的。
\item 若对每个 $x \in D$，理想
$\mathcal{I}_x \subset \mathcal{O}_{X, x}$ 都可以由一个非零因子生成，
则 $D$ 是有效 Cartier 除子。
\end{enumerate}
\end{lemma}

\begin{proof}
设 $\Spec(A)$ 为点 $x \in D$ 的仿射邻域。设
$\mathfrak p \subset A$ 为对应于 $x$ 的素理想。设 $I \subset A$
为定义 $D$ 在 $\Spec(A)$ 上之迹的理想。因为 $A$ 是诺特的
（由于 $X$ 局部诺特），理想 $I$ 可由有限多个元素生成，记为
$I = (f_1, \ldots, f_r)$。在 (1) 的假设下，有
$I_\mathfrak p = (f)$，其中 $f \in A_\mathfrak p$。于是有
$f_i = g_i f$，其中 $g_i \in A_\mathfrak p$。写
$g_i = a_i/h_i$ 且 $f = f'/h$，其中
$a_i, h_i, f', h \in A$，并且 $h_i, h \not \in \mathfrak p$。
于是 $I_{h_1 \ldots h_r h} \subset A_{h_1 \ldots h_r h}$ 是主理想，
因为它由 $f'$ 生成。这证明了 (1)。对于 (2)，可以假设
$I = (f)$。假设意味着 $f$ 在 $A_\mathfrak p$ 中的像是非零因子。
由《代数》引理
\ref{algebra-lemma-regular-sequence-in-neighbourhood}，$f$ 在 $x$
的某个邻域上是非零因子。这证明了 (2)。
\end{proof}

\begin{lemma}
\label{divisors-lemma-effective-Cartier-codimension-1}
设 $X$ 为局部诺特概形。
\begin{enumerate}
\item 设 $D \subset X$ 为局部主闭子概形，$\xi \in D$ 为 $D$
某个不可约分支的泛点。则 $\dim(\mathcal{O}_{X, \xi}) \leq 1$。
\item 设 $D \subset X$ 为有效 Cartier 除子，$\xi \in D$ 为 $D$
某个不可约分支的泛点。则 $\dim(\mathcal{O}_{X, \xi}) = 1$。
\end{enumerate}
\end{lemma}

\begin{proof}
证明 (1)。根据假设，可以设 $X = \Spec(A)$ 且
$D = \Spec(A/(f))$，其中 $A$ 为诺特环，$f \in A$。设 $\xi$
对应于素理想 $\mathfrak p \subset A$。$\xi$ 是 $D$ 某个不可约分支
的泛点这一假设，意味着 $\mathfrak p$ 是包含 $(f)$ 的极小素理想。
于是由《代数》引理 \ref{algebra-lemma-minimal-over-1}，
$\dim(A_\mathfrak p) \leq 1$。

\medskip\noindent
证明 (2)。由 (1) 可知 $\dim(\mathcal{O}_{X, \xi}) \leq 1$。
另一方面，由引理
\ref{divisors-lemma-characterize-effective-Cartier-divisor}，局部方程 $f$
在 $A_\mathfrak p$ 中是非零因子，这说明维数至少为 $1$
（因为由初等的《代数》引理
\ref{algebra-lemma-Zariski-topology}，$A_\mathfrak p$ 中必有一个
不包含 $f$ 的素理想）。
\end{proof}

\begin{lemma}
\label{divisors-lemma-integral-effective-Cartier-divisor-dvr}
设 $X$ 为诺特概形，$D \subset X$ 为整闭子概形，同时也是有效
Cartier 除子。则 $X$ 在 $D$ 的泛点处的局部环是离散赋值环。
\end{lemma}

\begin{proof}
由引理 \ref{divisors-lemma-characterize-effective-Cartier-divisor}，可以假设
$X = \Spec(A)$ 且 $D = \Spec(A/(f))$，其中 $A$ 为诺特环，
$f \in A$ 为非零因子。$D$ 为整概形这一假设意味着 $(f)$ 是素理想。
因此，$X$ 在泛点处的局部环为 $A_{(f)}$；这是诺特局部环，其极大理想
由一个非零因子生成。于是由《代数》引理
\ref{algebra-lemma-characterize-dvr}，它是离散赋值环。
\end{proof}

\begin{lemma}
\label{divisors-lemma-effective-Cartier-divisor-Sk}
设 $X$ 为局部诺特概形，$D \subset X$ 为有效 Cartier 除子。
若 $X$ 具有 $(S_k)$ 性质，则 $D$ 具有 $(S_{k - 1})$ 性质。
\end{lemma}

\begin{proof}
设 $x \in D$。则
$\mathcal{O}_{D, x} = \mathcal{O}_{X, x}/(f)$，其中
$f \in \mathcal{O}_{X, x}$ 是非零因子。根据假设，有
$\text{depth}(\mathcal{O}_{X, x}) \geq \min(\dim(\mathcal{O}_{X, x}), k)$。
由《代数》引理 \ref{algebra-lemma-depth-drops-by-one}，有
$\text{depth}(\mathcal{O}_{D, x}) = \text{depth}(\mathcal{O}_{X, x}) - 1$；
由《代数》引理 \ref{algebra-lemma-one-equation}，有
$\dim(\mathcal{O}_{D, x}) = \dim(\mathcal{O}_{X, x}) - 1$。
因此
$\text{depth}(\mathcal{O}_{D, x}) \geq \min(\dim(\mathcal{O}_{D, x}), k - 1)$，
这正是所需结论。
\end{proof}

\begin{lemma}
\label{divisors-lemma-normal-effective-Cartier-divisor-S1}
设 $X$ 为局部诺特正规概形，$D \subset X$ 为有效 Cartier 除子。
则 $D$ 具有 $(S_1)$ 性质。
\end{lemma}

\begin{proof}
由《性质》引理 \ref{properties-lemma-criterion-normal}，可知 $X$
具有 $(S_2)$ 性质。因此，由引理
\ref{divisors-lemma-effective-Cartier-divisor-Sk} 得出结论。
\end{proof}

\begin{lemma}
\label{divisors-lemma-weil-divisor-is-cartier-UFD}
设 $X$ 为局部诺特概形，$D \subset X$ 为整闭子概形。假设
\begin{enumerate}
\item $D$ 的余维为 $1$（余维在 $X$ 中计算）；
\item $\mathcal{O}_{X, x}$ 是 UFD，且这对所有 $x \in D$ 都成立。
\end{enumerate}
则 $D$ 是有效 Cartier 除子。
\end{lemma}

\begin{proof}
设 $x \in D$，并令 $A = \mathcal{O}_{X, x}$。设
$\mathfrak p \subset A$ 对应于 $D$ 的泛点。根据 (1)，
$A_\mathfrak p$ 的维数为 $1$。因此 $\mathfrak p$ 是高为 $1$
的素理想。由于 $A$ 是 UFD，这意味着
$\mathfrak p = (f)$，其中 $f \in A$。当然，$f$ 是非零因子，
于是结论由引理
\ref{divisors-lemma-effective-Cartier-in-points} 得出。
\end{proof}

\begin{lemma}
\label{divisors-lemma-codim-1-part}
设 $X$ 为诺特概形，$Z \subset X$ 为闭子概形。假设存在
\begin{enumerate}
\item 一族整有效 Cartier 除子 $D_i \subset X$，$i \in I$；
\item 闭子集 $Z' \subset X$，其所有不可约分支的余维都
$\geq 2$（余维在 $X$ 中计算；《拓扑》定义
\ref{topology-definition-codimension}）。
\end{enumerate}
并且集合论上有 $Z \subset Z' \cup \bigcup_{i \in I} D_i$。
则存在整数 $a_i \geq 0$，其中 $a_i = 0$ 对几乎所有
$i \in I$ 都成立，使得有效 Cartier 除子
$$
D = \sum a_i D_i
$$
包含在 $Z$ 中，并且包含态射 $D \to Z$ 在余维 $2$ 之外
是同构（余维在 $X$ 中计算；其含义是：存在开集 $U \subset Z$，使
$D \cap U \to Z \cap U$ 为同构，并且 $Z \setminus U$ 的每个
不可约分支的余维都 $\geq 2$，余维在 $X$ 中计算）。当 $Z$ 在 $X$
中无处稠密时，$D_i$、$i \in I$ 和 $Z'$ 的存在性得到保证，只要
$\mathcal{O}_{X, x}$ 是 UFD 且这对所有 $x \in Z$ 都成立，或者
$X$ 正则。
\end{lemma}

\begin{proof}
设 $\xi_i \in D_i$ 为泛点，并令
$\mathcal{O}_i = \mathcal{O}_{X, \xi_i}$ 为局部环；由引理
\ref{divisors-lemma-integral-effective-Cartier-divisor-dvr}，它是离散赋值环。
设 $a_i \geq 0$ 为
$\mathcal{I}_{Z, \xi_i} \subset \mathcal{O}_i$ 中元素的最小赋值。
当然，$a_i > 0$ 仅可能在 $D_i$ 是 $Z$ 的不可约分支时成立，因此
$a_i > 0$ 只对有限多个 $i \in I$ 成立。我们断言有效 Cartier
除子 $D = \sum a_i D_i$ 符合要求。

\medskip\noindent
具体地，设 $x \in X$，并令 $A = \mathcal{O}_{X, x}$。设
$D_1, \ldots, D_n$ 为两两不同的除子 $D_i$，满足
$x \in D_i$ 且 $a_i > 0$。对 $1 \leq i \leq n$，令
$f_i \in A$ 为 $D_i$ 的局部方程。则 $f_i$ 是 $A$ 的素元，且
$\mathcal{O}_i = A_{(f_i)}$。设
$I = \mathcal{I}_{Z, x} \subset A$ 为 $Z$ 的理想层之茎。
根据 $a_i$ 的选择，有
$I A_{(f_i)} = f_i^{a_i}A_{(f_i)}$。我们断言
$I \subset (\prod_{i = 1, \ldots, n} f_i^{a_i})$。

\medskip\noindent
证明该断言。局部化映射
$\varphi : A/(f_i) \to A_{(f_i)}/f_iA_{(f_i)}$ 是单射，因为素理想
$(f_i)$ 是极大理想 $f_iA_{(f_i)}$ 的逆像。对 $n$ 归纳可知
$\varphi_n : A/(f_i^n)\to A_{(f_i)}/f_i^nA_{(f_i)}$ 也为单射。
由于 $\varphi_{a_i}(I) = 0$，有 $I \subset (f_i^{a_i})$。
因此，对任意 $x \in I$，可写成 $x = f_1^{a_1}x_1$，其中
$x_1 \in A$。由于 $D_1, \ldots, D_n$ 两两不同，$f_i$ 在
$A_{(f_j)}$ 中是单位，只要 $i \not = j$。比较 $x$ 与 $x_1$ 在
$A_{(f_i)}$ 中的像，对 $n \geq i > 1$ 仍有
$x_1 \in (f_i^{a_i})$。重复上述过程，归纳地写出
$x_i = f_{i + 1}^{a_{i + 1}}x_{i + 1}$，这对任意
$n > i \geq 1$ 成立。总之，有
$x \in (\prod_{i = 1, \ldots n} f_i^{a_i})$，这对任意
$x \in I$ 成立，证毕。

\medskip\noindent
该断言表明 $\mathcal{I}_Z \subset \mathcal{I}_D$，即
$D \subset Z$。此外，还可看出 $D$ 与 $Z$ 在 $\xi_i$ 处相同，
这证明 $D \to Z$ 在余维 $2$ 之外是同构（余维在 $X$ 中计算）。

\medskip\noindent
最后几个断言可如下证明。正则局部环是 UFD（《代数进阶》引理
\ref{more-algebra-lemma-regular-local-UFD}），故只须处理 UFD 情形。
在该情形中，令 $D_i$ 为 $Z$ 的余维为 $1$ 的不可约分支
（余维在 $X$ 中计算）。由引理
\ref{divisors-lemma-weil-divisor-is-cartier-UFD}，每个 $D_i$ 都是有效
Cartier 除子。
\end{proof}

\begin{lemma}
\label{divisors-lemma-codimension-1-is-effective-Cartier}
设 $Z \subset X$ 为诺特概形的闭子概形。假设
\begin{enumerate}
\item $Z$ 没有嵌入点；
\item $Z$ 的每个不可约分支的余维均为 $1$（余维在 $X$ 中计算）；
\item 每个局部环 $\mathcal{O}_{X, x}$（$x \in Z$）都是
UFD，或者 $X$ 正则。
\end{enumerate}
则 $Z$ 是有效 Cartier 除子。
\end{lemma}

\begin{proof}
设 $D = \sum a_i D_i$ 如引理 \ref{divisors-lemma-codim-1-part} 中所述，
其中 $D_i \subset Z$ 是 $Z$ 的不可约分支。
若 $D \to Z$ 不是同构，则 $\mathcal{O}_Z \to \mathcal{O}_D$
有一个非零核，其支集位于余维 $\geq 2$ 处。这将意味着
$Z$ 存在嵌入点，与假设 (1) 矛盾。因此如所需有 $D \cong Z$。
\end{proof}

\begin{lemma}
\label{divisors-lemma-UFD-one-equation-CM}
设 $R$ 为诺特 UFD。设 $I \subset R$ 为一个理想，
使得 $R/I$ 没有嵌入素理想，且包含 $I$
的每个极小素理想的高均为 $1$。
则 $I = (f)$，其中 $f \in R$。
\end{lemma}

\begin{proof}
由引理 \ref{divisors-lemma-codimension-1-is-effective-Cartier}，
理想层 $\tilde I$ 在 $\Spec(R)$ 上可逆。
由《代数进阶》引理 \ref{more-algebra-lemma-UFD-Pic-trivial}，
它由一个元生成。
\end{proof}

\begin{lemma}
\label{divisors-lemma-effective-Cartier-divisor-is-a-sum}
设 $X$ 为诺特概形，$D \subset X$ 为有效 Cartier 除子。
假设存在整有效 Cartier 除子 $D_i \subset X$，使得集合论上有
$D \subset \bigcup D_i$。则 $D = \sum a_i D_i$，其中 $a_i \geq 0$。
$D_i$ 的存在性得到保证，只要 $\mathcal{O}_{X, x}$ 是 UFD，
且这对所有 $x \in D$ 都成立，或者 $X$ 正则。
\end{lemma}

\begin{proof}
如引理 \ref{divisors-lemma-codim-1-part} 中选取 $a_i$，并令
$D' = \sum a_i D_i$。则 $D' \to D$ 是有效 Cartier 除子的包含，
且在余维 $2$ 之外是同构（余维在 $X$ 上计算）。取 $x \in X$。
令 $A = \mathcal{O}_{X, x}$，并设 $f, f' \in A$ 分别生成
$D, D'$ 之理想的非零因子（在 $A$ 中）。则 $f = gf'$，
其中 $g \in A$。此外，对每个素理想 $\mathfrak p$，若其高
$\leq 1$（在 $A$ 中计算），则 $g$ 在 $A_\mathfrak p$ 中的像是单位。这意味着
$g$ 是单位，因为包含 $(g)$ 的极小素理想的高为 $1$
（《代数》引理 \ref{algebra-lemma-minimal-over-1}）。
\end{proof}

\begin{lemma}
\label{divisors-lemma-quasi-projective-Noetherian-pic-effective-Cartier}
\begin{slogan}
在射影概形上，每个线丛都有正则亚纯截面。
\end{slogan}
设 $X$ 为具有丰富可逆层的诺特概形。
则每个可逆 $\mathcal{O}_X$-模都同构于
$$
\mathcal{O}_X(D - D') =
\mathcal{O}_X(D) \otimes_{\mathcal{O}_X} \mathcal{O}_X(D')^{\otimes -1}
$$
其中 $D, D'$ 是 $X$ 中的某些有效 Cartier 除子。此外，给定
有限子集 $E \subset X$，可以选取 $D, D'$ 使得
$E \cap D = \emptyset$ 且 $E \cap D' = \emptyset$。若
$X$ 拟仿射，则可选取 $D' = \emptyset$。
\end{lemma}

\begin{proof}
设 $x_1, \ldots, x_n$ 为 $X$ 的伴随点
（引理 \ref{divisors-lemma-finite-ass}）。

\medskip\noindent
若 $X$ 拟仿射，而 $\mathcal{N}$ 是任意可逆
$\mathcal{O}_X$-模，则可以选取一个截面 $t$ 属于
$\mathcal{N}$，
使其在 $E \cup \{x_1, \ldots, x_n\}$ 的任一点处都不消失，见《性质》引理
\ref{properties-lemma-quasi-affine-invertible-nonvanishing-section}。
由引理 \ref{divisors-lemma-regular-section-associated-points}，$t$ 是
$\mathcal{N}$ 的正则截面。因此 $\mathcal{N} \cong \mathcal{O}_X(D)$，其中
$D = Z(t)$ 是对应于 $t$ 的有效 Cartier 除子，见引理
\ref{divisors-lemma-characterize-OD}。由构造可知 $E \cap D = \emptyset$，
因而该情形已证。

\medskip\noindent
回到一般情形，设 $\mathcal{L}$ 为 $X$ 上的丰富可逆层。
存在 $m > 0$ 以及截面 $s \in \Gamma(X, \mathcal{L}^{\otimes m})$，
使得 $X_s$ 为仿射概形，并且
$E \cup \{x_1, \ldots, x_n\} \subset X_s$
（《性质》引理 \ref{properties-lemma-ample-finite-set-in-principal-affine}）。

\medskip\noindent
设 $\mathcal{N}$ 为任意可逆 $\mathcal{O}_X$-模。
由拟仿射情形，可以找到截面 $t \in \mathcal{N}(X_s)$，
它在 $E \cup \{x_1, \ldots, x_n\}$ 的任一点处均不消失。
由《性质》引理 \ref{properties-lemma-invert-s-sections}，
对某个 $e \geq 0$，截面 $s^e|_{X_s} t$ 可以延拓为
一个整体截面 $\tau$，它是 $\mathcal{L}^{\otimes me} \otimes \mathcal{N}$ 的截面。
因此 $\mathcal{L}^{\otimes me} \otimes \mathcal{N}$ 与
$\mathcal{L}^{\otimes me}$ 都是可逆层，并且各自有一个整体截面，
该截面在 $E \cup \{x_1, \ldots, x_n\}$ 的任一点处均不消失。
由引理 \ref{divisors-lemma-regular-section-associated-points}，这些都是正则截面。
故有 $\mathcal{L}^{\otimes me} \otimes \mathcal{N} \cong \mathcal{O}_X(D)$
以及 $\mathcal{L}^{\otimes me} \cong \mathcal{O}_X(D')$，其中 $D$ 与 $D'$
是某些有效 Cartier 除子，见引理
\ref{divisors-lemma-characterize-OD}。由构造可知 $E \cap D = \emptyset$ 且
$E \cap D' = \emptyset$，证毕。
\end{proof}

\begin{lemma}
\label{divisors-lemma-wedge-product-ses}
设 $X$ 为整正则概形，其维数为 $2$。
设 $i : D \to X$ 为一个有效 Cartier 除子的浸入。
设 $\mathcal{F} \to \mathcal{F}' \to i_*\mathcal{G} \to 0$
为凝聚 $\mathcal{O}_X$-模的正合列。假设
\begin{enumerate}
\item $\mathcal{F}, \mathcal{F}'$ 是秩为 $r$ 的局部自由模，这在
$X$ 的一个非空开集上成立；
\item $D$ 是整概形；
\item $\mathcal{G}$ 是有限局部自由 $\mathcal{O}_D$-模，秩为 $s$。
\end{enumerate}
则 $\mathcal{L} = (\wedge^r\mathcal{F})^{**}$ 与
$\mathcal{L}' = (\wedge^r \mathcal{F}')^{**}$ 是可逆
$\mathcal{O}_X$-模，且 $\mathcal{L}' \cong \mathcal{L}(k D)$，其中
$k \in \{0, \ldots, \min(s, r)\}$。
\end{lemma}

\begin{proof}
第一个断言由引理 \ref{divisors-lemma-reflexive-over-regular-dim-2} 得出，
因为假设 (1) 意味着 $\mathcal{L}$ 与 $\mathcal{L}'$ 的秩都为 $1$。
取 $\wedge^r$ 和取双对偶都是函子，因此得到一个典范映射
$\sigma : \mathcal{L} \to \mathcal{L}'$。它在 (1) 中的非空开集上是同构，
因而非零。为了完成证明，只须证明：将 $\sigma$ 视为
$\mathcal{L}' \otimes \mathcal{L}^{\otimes -1}$ 的整体截面时，
它在 $X$ 的任意余维点处都不消失，除了
$D$ 的泛点；而在该点处，其消失阶至多为 $\min(s, r)$。

\medskip\noindent
翻译成代数语言，我们得到如下问题：
设 $(A, \mathfrak m, \kappa)$ 为离散赋值环，分式域为 $K$。
设 $M \to M' \to N \to 0$ 为有限 $A$-模的正合列，满足
$\dim_K(M \otimes K) = \dim_K(M' \otimes K) = r$，并且
$N \cong \kappa^{\oplus s}$。证明诱导映射
$L = \wedge^r(M)^{**} \to L' = \wedge^r(M')^{**}$ 的消失阶
至多为 $\min(s, r)$。我们将使用 $A$ 上模的结构定理，参见
《代数进阶》引理
\ref{more-algebra-lemma-generalized-valuation-ring-modules} 或
\ref{more-algebra-lemma-modules-PID}。
有限 $A$-模对其一个挠子模取商，不会改变其双对偶。
因此，可以将 $M$ 替换为 $M/M_{tors}$，并将 $M'$
替换为 $M'/\Im(M_{tors} \to M')$，从而假设 $M$ 无挠。
此时 $M \to M'$ 是单射，且 $M'_{tors} \to N$ 是单射。
因此，可以将 $M'$ 替换为 $M'/M'_{tors}$，并将 $N$
替换为 $N/M'_{tors}$。于是问题归约到 $M$ 和 $M'$ 都是秩为 $r$
的自由模且 $N \cong \kappa^{\oplus s}$ 的情形。此时 $\sigma$
是 $M \to M'$ 的行列式，并且其消失阶为 $s$；例如这可由
《代数》引理 \ref{algebra-lemma-order-vanishing-determinant} 得出。
\end{proof}

\section{仿射开集的补集}
\label{divisors-section-complement-affine-open}

\noindent
本节讨论如下结果：簇中一个仿射开集的补集具有纯余维 $1$。

\begin{lemma}
\label{divisors-lemma-affine-punctured-spec}
设 $(A, \mathfrak m)$ 为诺特局部环。
记 $U = \Spec(A) \setminus \{\mathfrak m\}$；这是 $A$ 的穿孔谱，并且它是仿射的，
当且仅当 $\dim(A) \leq 1$。
\end{lemma}

\begin{proof}
若 $\dim(A) = 0$，则 $U$ 为空，因而仿射（它等于 $0$ 环的谱）。
若 $\dim(A) = 1$，则可以选择元素 $f \in \mathfrak m$，使其不包含于
$A$ 的有限多个极小素理想中的任何一个
（《代数》引理 \ref{algebra-lemma-Noetherian-irreducible-components} 与
\ref{algebra-lemma-silly}）。于是 $U = \Spec(A_f)$ 是仿射的。

\medskip\noindent
逆命题更有意思。我们将给出一个稍微非标准的证明，并在下面的评注中讨论
标准论证。假设 $U = \Spec(B)$ 是仿射的。由于取约化不影响仿射性与维数，
可以用 $A$ 对其幂零元理想所取的商替换原环，从而假设 $A$ 约化。
令 $Q = B/A$，将其视为 $A$-模。$Q$ 的支集是
$\{\mathfrak m\}$，因为有 $A_\mathfrak p = B_\mathfrak p$，这对每个非极大素理想
$\mathfrak p$（属于 $A$）都成立。
可以假设 $\dim(A) \geq 1$；于是如上可选取 $f \in \mathfrak m$，
使其不包含于 $A$ 的任何极小理想。由于 $A$ 约化，这意味着 $f$ 在 $A$
中是非零因子。同一论证表明 $f : B \to B$ 是单射，因为 $\Spec(B)$
是 $\Spec(A)$ 的开集。注意，由《代数》引理
\ref{algebra-lemma-one-equation}，有
$\dim(A/fA) = \dim(A) - 1$。
对乘以 $f$ 的映射应用蛇形引理，其中该映射作用于短正合列
$0 \to A \to B \to Q \to 0$，
得到
$$
0 \to Q[f] \to A/fA \to B/fB \to Q/fQ \to 0
$$
其中 $Q[f] = \Ker(f : Q \to Q)$。
这说明 $Q[f]$ 是有限 $A$-模。由于 $Q[f]$ 的支集是
$\{\mathfrak m\}$，可知
$l = \text{length}_A(Q[f]) < \infty$
（《代数》引理 \ref{algebra-lemma-support-point}）。
令 $l_n = \text{length}_A(Q[f^n])$。正合列
$$
0 \to Q[f^n] \to Q[f^{n + 1}] \xrightarrow{f^n} Q[f]
$$
归纳地表明 $l_n < \infty$，且 $l_n \leq l_{n + 1}$。
考察正合列
$$
0 \to Q[f] \to Q[f^{n + 1}] \xrightarrow{f} Q[f^n] \to Q/fQ
$$
可知 $Q[f^n]$ 在 $Q/fQ$ 中之像的长度为
$l_n - l_{n + 1} + l \leq l$。由于 $Q = \bigcup Q[f^n]$，
可知 $Q/fQ$ 的长度至多为 $l$，即有界。
因此 $Q/fQ$ 是有限 $A$-模。故 $A/fA \to B/fB$ 是有限环同态，
特别地，它在谱上诱导闭映射
（《代数》引理 \ref{algebra-lemma-integral-going-up} 与
\ref{algebra-lemma-going-up-closed}）。
另一方面，$\Spec(B/fB)$ 是 $\Spec(A/fA)$ 的穿孔谱。
除非 $\Spec(B/fB) = \emptyset$，否则就产生矛盾；而该等式意味着
$\dim(A/fA) = 0$，正是所需结论。
\end{proof}

\begin{remark}
\label{divisors-remark-affine-punctured-spectrum-standard-proof}
设 $(A, \mathfrak m)$ 为维数 $\geq 2$ 的诺特局部正规整环，并设
$U$ 为 $A$ 的穿孔谱，则 $\Gamma(U, \mathcal{O}_U) = A$。
Hartogs 定理的这一代数版本来自如下事实：
$A = \bigcap_{\text{高度}(\mathfrak p) = 1} A_\mathfrak p$，
我们已在《代数》引理
\ref{algebra-lemma-normal-domain-intersection-localizations-height-1}
中见过它。因此在该情形中，$U$ 不可能仿射（否则会迫使 $\mathfrak m$
成为 $U$ 的一个点）。这常被用作引理
\ref{divisors-lemma-affine-punctured-spec} 之证明的出发点。
为把一般诺特局部环的情形约化到该情形，首先作完备化
（以得到 Nagata 局部环），然后对某个适当的极小素理想将 $A$
替换为 $A/\mathfrak q$，再作正规化。这些步骤都不改变维数，并由此得到矛盾。
也可以跳过完备化步骤，但此时正规化一般不是诺特整环。不过，它仍然是
同维数的 Krull 整环（这利用 Krull--Akizuki 定理证明），因而可以应用
同一论证。
\end{remark}

\begin{remark}
\label{divisors-remark-affine-puctured-spectrum-general}
尚不清楚如何刻画穿孔谱为仿射的非诺特局部环 $(A, \mathfrak m)$。
这样的环有一个有限生成理想 $I$，满足
$\mathfrak m = \sqrt{I}$。当然，若可以取 $I$ 由 $1$ 个元素生成，
则 $A$ 的穿孔谱是仿射的；这给出许多非诺特例子。
反过来，由引理 \ref{divisors-lemma-affine-punctured-spec} 之证明中的论证可知，
这样的环不可能有非零因子 $f \in \mathfrak m$ 满足
$H^0_I(A/fA) = 0$（所以 $A$ 不可能有长度为 $2$ 的正则列）。
此外，对任何环 $A'$，若它是局部环同态 $A \to A'$ 的目标环，只要
$\mathfrak m_{A'} = \sqrt{\mathfrak mA'}$，同一结论也成立。
\end{remark}

\begin{lemma}
\label{divisors-lemma-complement-affine-open-immersion}
\begin{reference}
\cite[EGA IV, Corollaire 21.12.7]{EGA4}
\end{reference}
设 $X$ 为局部诺特概形。设 $U \subset X$ 为开子概形，并且包含态射
$U \to X$ 是仿射的。对每个泛点 $\xi$，若它来自
$X \setminus U$ 的某个不可约分支，则局部环 $\mathcal{O}_{X, \xi}$ 的维数
$\leq 1$。若 $U$ 稠密，或 $\xi$ 位于 $U$ 的闭包中，则
$\dim(\mathcal{O}_{X, \xi}) = 1$。
\end{lemma}

\begin{proof}
由于 $\xi$ 是 $X \setminus U$ 的泛点，可知
$$
U_\xi = U \times_X \Spec(\mathcal{O}_{X, \xi}) \subset
\Spec(\mathcal{O}_{X, \xi})
$$
是 $\mathcal{O}_{X, \xi}$ 的穿孔谱（提示：使用《概形》引理
\ref{schemes-lemma-specialize-points}）。
因为 $U \to X$ 是仿射的，故
$U_\xi \to \Spec(\mathcal{O}_{X, \xi})$ 是仿射的
（《态射》引理 \ref{morphisms-lemma-base-change-affine}），
从而 $U_\xi$ 是仿射的。因此由引理
\ref{divisors-lemma-affine-punctured-spec}，有
$\dim(\mathcal{O}_{X, \xi}) \leq 1$。
若 $\xi \in \overline{U}$，则存在特化 $\eta \to \xi$，其中
$\eta \in U$（只需取泛点 $\eta$，它属于 $\overline{U}$ 的一个包含
$\xi$ 的不可约分支；由于 $\overline{U}$ 局部诺特，因而局部只有有限多个
不可约分支，可知 $\eta \in U$）。于是
$\eta \in \Spec(\mathcal{O}_{X, \xi})$，故维数不可能为 $0$。
\end{proof}

\begin{lemma}
\label{divisors-lemma-complement-affine-open}
设 $X$ 为分离的局部诺特概形。设 $U \subset X$ 为仿射开集。
对每个泛点 $\xi$，若它来自 $X \setminus U$ 的某个不可约分支，则局部环
$\mathcal{O}_{X, \xi}$ 的维数均 $\leq 1$。若 $U$ 稠密，或
$\xi$ 位于 $U$ 的闭包中，则
$\dim(\mathcal{O}_{X, \xi}) = 1$。
\end{lemma}

\begin{proof}
这由引理 \ref{divisors-lemma-complement-affine-open-immersion} 得出，因为
由《态射》引理 \ref{morphisms-lemma-affine-permanence}，
态射 $U \to X$ 是仿射的。
\end{proof}

\noindent
下面的引理有时可以用来构造有效 Cartier 除子。

\begin{lemma}
\label{divisors-lemma-complement-open-affine-effective-cartier-divisor}
设 $X$ 为诺特分离概形。设 $U \subset X$ 为稠密仿射开集。
若 $\mathcal{O}_{X, x}$ 是 UFD，且这对所有 $x \in X \setminus U$ 都成立，
则存在有效 Cartier 除子 $D \subset X$，使得 $U = X \setminus D$。
\end{lemma}

\begin{proof}
由于 $X$ 诺特，补集 $X \setminus U$ 只有有限多个不可约分支
$D_1, \ldots, D_r$
（将《性质》引理 \ref{properties-lemma-Noetherian-irreducible-components}
应用于 $X \setminus U$ 上的既约诱导子概形结构）。
由引理 \ref{divisors-lemma-complement-affine-open}
（以及《性质》引理 \ref{properties-lemma-codimension-local-ring}），
每个 $D_i \subset X$ 的余维均为 $1$。
于是由引理 \ref{divisors-lemma-weil-divisor-is-cartier-UFD}，
$D_i$ 是有效 Cartier 除子。因此可以取
$D = D_1 + \ldots + D_r$。
\end{proof}

\begin{lemma}
\label{divisors-lemma-complement-open-affine-effective-cartier-divisor-bis}
设 $X$ 为具有仿射对角的诺特概形。设 $U \subset X$ 为稠密仿射开集。
若 $\mathcal{O}_{X, x}$ 是 UFD，且这对所有 $x \in X \setminus U$ 都成立，
则存在有效 Cartier 除子 $D \subset X$，使得 $U = X \setminus D$。
\end{lemma}

\begin{proof}
由于 $X$ 诺特，补集 $X \setminus U$ 只有有限多个不可约分支
$D_1, \ldots, D_r$
（将《性质》引理 \ref{properties-lemma-Noetherian-irreducible-components}
应用于 $X \setminus U$ 上的既约诱导子概形结构）。
将 $D_i$ 视为 $X$ 的既约闭子概形。
设 $X = \bigcup_{j \in J} X_j$ 为 $X$ 的仿射开覆盖。对
所有 $j$（其属于 $J$），令 $U_j = U \cap X_j$。由于 $X$ 具有仿射对角，
概形
$$
U_j = X \times_{(X \times X)} (U \times X_j)
$$
是仿射的。因此，由于 $X_j$ 分离，根据引理
\ref{divisors-lemma-complement-open-affine-effective-cartier-divisor}
及其证明，对所有 $j \in J$ 与 $1 \leq i \leq r$，交
$D_i \cap X_j$ 或为空，或为 $X_j$ 中的有效 Cartier 除子。
因此 $D_i \subset X$ 是有效 Cartier 除子（因为这是局部性质）。
故可以取 $D = D_1 + \ldots + D_r$。
\end{proof}

\begin{lemma}
\label{divisors-lemma-regular-ample-family}
设 $X$ 为具有仿射对角的拟紧正则概形。则 $X$ 具有一个可逆模的丰富族
（《态射》定义
\ref{morphisms-definition-family-ample-invertible-modules}）。
\end{lemma}

\begin{proof}
注意 $X$ 是整概形的有限不交并
（《性质》引理
\ref{properties-lemma-regular-normal} 与
\ref{properties-lemma-normal-Noetherian}）。
因此可以进一步假设 $X$ 是整的、诺特的、正则的，并且具有仿射对角。
设 $x \in X$。选择仿射开邻域 $U \subset X$，它包含 $x$。
由于 $X$ 是整概形，$U$ 在 $X$ 中稠密。
由《代数进阶》引理 \ref{more-algebra-lemma-regular-local-UFD}，
$X$ 的局部环都是 UFD。因此由引理
\ref{divisors-lemma-complement-open-affine-effective-cartier-divisor-bis}，
可以找到有效 Cartier 除子 $D \subset X$，其补集为 $U$。
于是典范截面 $s = 1_D$ 属于 $\mathcal{L} = \mathcal{O}_X(D)$，
见定义 \ref{divisors-definition-invertible-sheaf-effective-Cartier-divisor}，
恰在 $D$ 上消失，故 $U = X_s$。
因此《态射》定义
\ref{morphisms-definition-family-ample-invertible-modules}
中的两个条件都成立，证毕。
\end{proof}

\section{范数}
\label{divisors-section-norms}

\noindent
设 $\pi : X \to Y$ 为概形的有限态射，并设 $d \geq 1$ 为整数。
如果存在一个乘法映射，则称存在{\it 次数为 $d$ 的 $\pi$ 的范数}
\footnote{这是非标准记号。}
$$
\text{Norm}_\pi : \pi_*\mathcal{O}_X \to \mathcal{O}_Y
$$
它是层的映射，并且满足
\begin{enumerate}
\item 复合
$\mathcal{O}_Y \xrightarrow{\pi^\sharp} \pi_*\mathcal{O}_X
\xrightarrow{\text{Norm}_\pi} \mathcal{O}_Y$ 等于 $g \mapsto g^d$；
\item 对开集 $V \subset Y$，若 $f \in \mathcal{O}_X(\pi^{-1}V)$
在 $x \in \pi^{-1}(V)$ 处为零，则 $\text{Norm}_\pi(f)$
在 $\pi(x)$ 处为零。
\end{enumerate}
注意条件 (1) 迫使 $\pi$ 满射。由于 $\text{Norm}_\pi$ 是乘法的，
它把单位映到单位；因此，给定 $y \in Y$，若 $f$ 是 $X$ 上的正则函数，
它在每个 $x \in X$ 处有定义但不消失，并且 $\pi(x) = y$，则
$\text{Norm}_\pi(f)$ 有定义且在 $y$ 处不消失。该结论不需要条件 (2)；
事实上，本节中的构造只需要条件 (1)，只有某些消失性质
（特别是在引理 \ref{divisors-lemma-norm-ample} 的证明中使用）才需要性质 (2)。

\begin{lemma}
\label{divisors-lemma-finite-trivialize-invertible-upstairs}
设 $\pi : X \to Y$ 为概形的有限态射。设 $\mathcal{L}$ 为可逆
$\mathcal{O}_X$-模。设 $y \in Y$。则存在开邻域
$V \subset Y$，它包含 $y$，并使得 $\mathcal{L}|_{\pi^{-1}(V)}$ 平凡。
\end{lemma}

\begin{proof}
显然可以假设 $Y$ 仿射，因而 $X$ 也仿射。由于 $\pi$ 有限，
纤维 $\pi^{-1}(\{y\})$（位于 $y$ 上）是有限的。
因为 $X$ 仿射，可以选取 $s \in \Gamma(X, \mathcal{L})$，使其在
$\pi^{-1}(\{y\})$ 的任何点处都不消失。这由《性质》引理
\ref{properties-lemma-quasi-affine-invertible-nonvanishing-section}
得出，不过我们也给出直接论证。具体地，可以选取闭点的有限集
$E \subset X$，使每个 $x \in \pi^{-1}(\{y\})$ 都特化到
$E$ 的某个点。对 $x \in E$，用 $i_x : x \to X$ 表示闭浸入。
于是
$\mathcal{L} \to \bigoplus_{x \in E} i_{x, *}i_x^*\mathcal{L}$
是拟凝聚 $\mathcal{O}_X$-模的满射，故映射
$$
\Gamma(X, \mathcal{L}) \to
\bigoplus\nolimits_{x \in E} \mathcal{L}_x/\mathfrak m_x\mathcal{L}_x
$$
是满射（因为在拟凝聚 $\mathcal{O}_X$-模范畴上取整体截面是正合函子，
见《概形》引理 \ref{schemes-lemma-equivalence-quasi-coherent}）。
因此，可以找到 $s \in \Gamma(X, \mathcal{L})$，使其在特化到
$E$ 中某点的任何点处都不消失。于是 $X_s \subset X$ 是
$\pi^{-1}(\{y\})$ 的开邻域。由于 $\pi$ 有限，因而是闭态射，
可知存在开邻域 $V \subset Y$，它包含 $y$，且其逆像包含于 $X_s$，正是所需。
\end{proof}

\begin{lemma}
\label{divisors-lemma-norm-invertible}
设 $\pi : X \to Y$ 为概形的有限态射。若存在次数为 $d$ 的 $\pi$ 的范数，
则存在阿贝尔群同态
$$
\text{Norm}_\pi : \Pic(X) \to \Pic(Y)
$$
使得 $\text{Norm}_\pi(\pi^*\mathcal{N}) \cong \mathcal{N}^{\otimes d}$
对所有可逆 $\mathcal{O}_Y$-模 $\mathcal{N}$ 都成立。
\end{lemma}

\begin{proof}
我们将使用可逆 $\mathcal{O}_X$-模的同构类与
$H^1(X, \mathcal{O}_X^*)$ 的元素之间的对应，而不再另作说明；
该对应见《上同调》引理 \ref{cohomology-lemma-h1-invertible}。
现说明如何对可逆 $\mathcal{O}_X$-模 $\mathcal{L}$ 取范数。
由引理 \ref{divisors-lemma-finite-trivialize-invertible-upstairs}，存在开覆盖
$Y = \bigcup V_j$，使得 $\mathcal{L}|_{\pi^{-1}V_j}$ 平凡。
选择一个生成截面 $s_j \in \mathcal{L}(\pi^{-1}V_j)$，对每个 $j$ 都如此。
在交集 $\pi^{-1}V_j \cap \pi^{-1}V_{j'}$ 上，可以写成
$$
s_j = u_{jj'} s_{j'}
$$
其中存在唯一的
$u_{jj'} \in \mathcal{O}^*_X(\pi^{-1}V_j \cap \pi^{-1}V_{j'})$。
于是可以考察元素
$$
v_{jj'} = \text{Norm}_\pi(u_{jj'}) \in \mathcal{O}_Y^*(V_j \cap V_{j'})
$$
这些元素满足余循环条件（因为 $u_{jj'}$ 满足该条件，且
$\text{Norm}_\pi$ 是乘法的），因此定义了一个可逆
$\mathcal{O}_Y$-模。我们略去如下验证：该构造是良定义的，在 Picard 群上
可加，并且满足
$\text{Norm}_\pi(\pi^*\mathcal{N}) \cong \mathcal{N}^{\otimes d}$，
这对所有可逆 $\mathcal{O}_Y$-模 $\mathcal{N}$ 都成立。
\end{proof}

\begin{lemma}
\label{divisors-lemma-norm-map-invertible}
设 $\pi : X \to Y$ 为概形的有限态射。假设存在次数为 $d$ 的 $\pi$ 的范数。
对任意 $\mathcal{O}_X$-线性映射
$\varphi : \mathcal{L} \to \mathcal{L}'$，若它连接两个可逆
$\mathcal{O}_X$-模，则存在 $\mathcal{O}_Y$-线性映射
$$
\text{Norm}_\pi(\varphi) :
\text{Norm}_\pi(\mathcal{L})
\longrightarrow
\text{Norm}_\pi(\mathcal{L}')
$$
其中 $\text{Norm}_\pi(\mathcal{L})$、$\text{Norm}_\pi(\mathcal{L}')$
如引理 \ref{divisors-lemma-norm-invertible} 所述。此外，对 $y \in Y$，下列条件等价：
\begin{enumerate}
\item $\varphi$ 在某个点 $x \in X$ 处为零，且 $\pi(x) = y$；
\item $\text{Norm}_\pi(\varphi)$ 在 $y$ 处为零。
\end{enumerate}
\end{lemma}

\begin{proof}
选择开覆盖 $Y = \bigcup V_j$，使 $\mathcal{L}$ 和 $\mathcal{L}'$
在开集 $\pi^{-1}V_j$ 上平凡。由引理
\ref{divisors-lemma-finite-trivialize-invertible-upstairs}，这可以做到。
选择生成截面 $s_j$ 与 $s'_j$；它们分别是 $\mathcal{L}$ 与
$\mathcal{L}'$ 在开集 $\pi^{-1}V_j$ 上的生成截面。于是
$\varphi(s_j) = f_js'_j$，其中
$f_j \in \mathcal{O}_X(\pi^{-1}V_j)$。
把 $\text{Norm}_\pi(\varphi)$ 定义为乘以
$\text{Norm}_\pi(f_j)$；此定义在 $V_j$ 上给出。对引理
\ref{divisors-lemma-norm-invertible} 之证明中
用来构造 $\text{Norm}_\pi(\mathcal{L})$、$\text{Norm}_\pi(\mathcal{L}')$
的余循环作简单计算，即可表明这定义了引理所述映射。
最后的断言来自次数为 $d$ 的范数映射之定义中的条件 (2)。
略去一些细节。
\end{proof}

\begin{lemma}
\label{divisors-lemma-norm-ample}
设 $\pi : X \to Y$ 为概形的有限态射。假设 $X$ 具有丰富可逆层，
并且存在次数为 $d$ 的 $\pi$ 的范数。则 $Y$ 具有丰富可逆层。
\end{lemma}

\begin{proof}
设 $\mathcal{L}$ 为假设所给的 $X$ 上的丰富可逆层。我们将证明
$\mathcal{N} = \text{Norm}_\pi(\mathcal{L})$ 在 $Y$ 上丰富。

\medskip\noindent
由于 $X$ 拟紧（《性质》定义
\ref{properties-definition-ample}），而 $X \to Y$ 满射
（因为 $\text{Norm}_\pi$ 存在），可知 $Y$ 拟紧。
取一点 $y \in Y$。为完成证明，我们将证明：存在截面 $t$，它属于
$\mathcal{N}$ 的某个正张量幂，在 $y$ 处不消失，并且 $Y_t$ 仿射。
为此，选择仿射开邻域 $V \subset Y$，它包含 $y$。
选择 $n \gg 0$ 以及截面
$s \in \Gamma(X, \mathcal{L}^{\otimes n})$，使得
$$
\pi^{-1}(\{y\}) \subset X_s \subset \pi^{-1}V
$$
这可由《性质》引理
\ref{properties-lemma-ample-finite-set-in-principal-affine} 做到。
于是 $t = \text{Norm}_\pi(s)$ 是 $\mathcal{N}^{\otimes n}$ 的截面，
它在 $x$ 处不消失，并且 $Y_t \subset V$，见引理
\ref{divisors-lemma-norm-map-invertible}。再由《性质》引理
\ref{properties-lemma-affine-cap-s-open}，$Y_t$ 是仿射的。
\end{proof}

\begin{lemma}
\label{divisors-lemma-norm-quasi-affine}
设 $\pi : X \to Y$ 为概形的有限态射。假设 $X$ 拟仿射，并且存在
次数为 $d$ 的 $\pi$ 的范数。则 $Y$ 拟仿射。
\end{lemma}

\begin{proof}
由《性质》引理 \ref{properties-lemma-quasi-affine-O-ample}，
$\mathcal{O}_X$ 是 $X$ 上的丰富可逆层。
引理 \ref{divisors-lemma-norm-ample} 的证明表明
$\text{Norm}_\pi(\mathcal{O}_X) = \mathcal{O}_Y$
是丰富可逆 $\mathcal{O}_Y$-模。因此《性质》引理
\ref{properties-lemma-quasi-affine-O-ample} 表明 $Y$ 拟仿射。
\end{proof}

\begin{lemma}
\label{divisors-lemma-finite-locally-free-has-norm}
设 $\pi : X \to Y$ 为次数 $d \geq 1$ 的有限局部自由态射。
则存在次数为 $d$ 的典范范数，其构造与任意基变换交换。
\end{lemma}

\begin{proof}
设 $V \subset Y$ 为仿射开集，使得 $(\pi_*\mathcal{O}_X)|_V$
是秩为 $d$ 的有限自由模。选择一组基，得到同构
$$
\mathcal{O}_V^{\oplus d} \cong (\pi_*\mathcal{O}_X)|_V
$$
对每个 $f \in \pi_*\mathcal{O}_X(V) = \mathcal{O}_X(\pi^{-1}(V))$，
乘以 $f$ 定义了所示自由向量丛的 $\mathcal{O}_V$-线性自同态 $m_f$。
因此得到一个 $d \times d$ 矩阵
$M_f \in \text{Mat}(d \times d, \mathcal{O}_Y(V))$，并可定义
$$
\text{Norm}_\pi(f) = \det(M_f)
$$
由于矩阵的行列式与所选基无关，可知该定义良好；这也意味着该构造可以
粘合成所需的整体映射。构造显然与基变换相容。

\medskip\noindent
性质 (1) 来自如下事实：一个 $d \times d$ 对角矩阵的对角元为
$g, g, \ldots, g$ 时，其行列式为 $g^d$。
为证明性质 (2)，可以作基变换，并假设 $Y$ 是域 $k$ 的谱。
于是 $X = \Spec(A)$，其中 $A$ 为 $k$-代数，并且
$\dim_k(A) = d$。若存在 $x \in X$ 使 $f \in A$ 在 $x$ 处消失，
则存在到某个域的映射 $A \to \kappa$，使 $f$ 在 $\kappa$ 中的像为零。
于是 $f : A \to A$ 不可能满射，因而
$\det(f : A \to A) = 0$，正是所需。
\end{proof}

\begin{lemma}
\label{divisors-lemma-norm-in-normal-case}
设 $\pi : X \to Y$ 为有限满态射，其中 $X$ 与 $Y$ 都是整概形，
且 $Y$ 正规。则存在次数为 $[R(X) : R(Y)]$ 的 $\pi$ 的范数。
\end{lemma}

\begin{proof}
设 $\Spec(B) \subset Y$ 为仿射开子集，并设
$\Spec(A) \subset X$ 为其逆像。则 $A$ 与 $B$ 都是整环。
设 $K$ 为 $A$ 的分式域，$L$ 为 $B$ 的分式域。图示如下：
$$
\xymatrix{
L \ar[r] & K \\
B \ar[u] \ar[r] & A \ar[u]
}
$$
由于 $K/L$ 是有限扩张，存在范数映射
$\text{Norm}_{K/L} : K^* \to L^*$，其次数为 $d = [K : L]$；如引理
\ref{divisors-lemma-finite-locally-free-has-norm} 的证明所述，它把
$f \in K$ 映到 $\det_L(f : K \to K)$。
注意 $f : K \to K$ 的特征多项式是 $f$ 在 $L$ 上的极小多项式的一个幂；
特别地，$\text{Norm}_{K/L}(f)$ 是 $f$ 在 $L$ 上之极小多项式的常数项的一个幂。
因此，由《代数》引理
\ref{algebra-lemma-minimal-polynomial-normal-domain}，
$\text{Norm}_{K/L}$ 把 $A$ 映入 $B$。
这在仿射开集的截面上确定了一个相容映射系，从而确定整体范数映射
$\text{Norm}_\pi$，其次数为 $d$。

\medskip\noindent
性质 (1) 由构造立即得出。
为证明性质 (2)，设 $f \in A$ 包含于素理想 $\mathfrak p \subset A$。
设 $f^m + b_1 f^{m - 1} + \ldots + b_m$ 为 $f$ 在 $L$ 上的极小多项式。
由《代数》引理 \ref{algebra-lemma-minimal-polynomial-normal-domain}，
有 $b_i \in B$。因此 $b_0 \in B \cap \mathfrak p$。
由于 $\text{Norm}_{K/L}(f) = b_0^{d/m}$（见上文），可知范数在
$\mathfrak p$ 的像点处消失。
\end{proof}

\begin{lemma}
\label{divisors-lemma-Frobenius-gives-norm-for-reduction}
设 $X$ 为诺特概形。设 $p$ 为满足 $p\mathcal{O}_X = 0$ 的素数。
则对某个 $e > 0$，存在次数为 $p^e$ 的范数，它对应态射 $X_{red} \to X$，
其中 $X_{red}$ 是 $X$ 的约化。
\end{lemma}

\begin{proof}
设 $A$ 为满足 $pA = 0$ 的诺特环。设 $I \subset A$ 为幂零元的理想。
则有 $I^n = 0$，其中 $n$ 为某个整数（《代数》引理
\ref{algebra-lemma-Noetherian-power}）。
选择 $e$ 使 $p^e \geq n$。于是
$$
A/I \longrightarrow A,\quad
f \bmod I \longmapsto f^{p^e}
$$
是良定义的。这给出次数为 $p^e$ 的范数，它对应
$\Spec(A/I) \to \Spec(A)$。
现在，若 $X$ 由某些仿射概形 $\Spec(A_i)$ 粘合而成，则对某个
$e \gg 0$，这些映射粘合成 $X_{red} \to X$ 的范数映射。略去细节。
\end{proof}

\begin{proposition}
\label{divisors-proposition-push-down-ample}
设 $\pi : X \to Y$ 为概形的有限满态射。假设 $X$ 具有丰富可逆
$\mathcal{O}_X$-模。若
\begin{enumerate}
\item $\pi$ 有限局部自由；或
\item $Y$ 为整正规概形；或
\item $Y$ 诺特，$p\mathcal{O}_Y = 0$，并且 $X = Y_{red}$，
\end{enumerate}
则 $Y$ 具有丰富可逆 $\mathcal{O}_Y$-模。
\end{proposition}

\begin{proof}
情形 (1) 由引理 \ref{divisors-lemma-finite-locally-free-has-norm} 与
\ref{divisors-lemma-norm-ample} 结合得出。
情形 (3) 由引理 \ref{divisors-lemma-Frobenius-gives-norm-for-reduction} 与
\ref{divisors-lemma-norm-ample} 结合得出。
在情形 (2) 中，首先替换 $X$；所取替代者是 $X$ 的一个支配 $Y$
的不可约分支（把它视为 $X$ 的既约闭子概形）。
然后可以应用引理 \ref{divisors-lemma-norm-in-normal-case}。
\end{proof}

\begin{lemma}
\label{divisors-lemma-push-down-quasi-affine}
设 $\pi : X \to Y$ 为概形的有限满态射。假设 $X$ 拟仿射。若下列任一条件成立：
\begin{enumerate}
\item $\pi$ 有限局部自由；或
\item $Y$ 为整正规概形，
\end{enumerate}
则 $Y$ 拟仿射。
\end{lemma}

\begin{proof}
情形 (1) 由引理
\ref{divisors-lemma-finite-locally-free-has-norm} 与 \ref{divisors-lemma-norm-quasi-affine}
结合得出。
在情形 (2) 中，首先替换 $X$；所取替代者是 $X$ 的一个支配 $Y$
的不可约分支（把它视为 $X$ 的既约闭子概形）。
然后可以应用引理 \ref{divisors-lemma-norm-in-normal-case}。
\end{proof}

\section{相对有效 Cartier 除子}
\label{divisors-section-effective-Cartier-morphisms}

\noindent
下面的引理表明，在基上平坦的有效 Cartier 除子确实是基上的一个
“有效 Cartier 除子族”。例如，它在任意纤维上的限制都是有效
Cartier 除子。

\begin{lemma}
\label{divisors-lemma-relative-Cartier}
设 $f : X \to S$ 为概形的态射。设 $D \subset X$ 为闭子概形。
假设
\begin{enumerate}
\item $D$ 是有效 Cartier 除子；并且
\item $D \to S$ 是平坦态射。
\end{enumerate}
则对概形的每个态射 $g : S' \to S$，拉回
$(g')^{-1}D$ 都是 $X' = S' \times_S X$ 上的有效 Cartier 除子，
其中 $g' : X' \to X$ 是投影。
\end{lemma}

\begin{proof}
使用引理 \ref{divisors-lemma-characterize-effective-Cartier-divisor}，
可把它翻译成如下代数陈述。设 $A \to B$ 为环同态，且 $h \in B$。
假设 $h$ 是非零因子，并且 $B/hB$ 在 $A$ 上平坦。则
$$
0 \to B \xrightarrow{h} B \to B/hB \to 0
$$
是 $A$-模的短正合列，其中 $B/hB$ 在 $A$ 上平坦。
由《代数》引理 \ref{algebra-lemma-flat-tor-zero}，该列与任意模在
$A$ 上张量后仍然正合，特别地，也可以与任意 $A$-代数 $A'$ 张量。
\end{proof}

\noindent
该引理是下述定义的动机。

\begin{definition}
\label{divisors-definition-relative-effective-Cartier-divisor}
设 $f : X \to S$ 为概形的态射。$X/S$ 上的一个{\it 相对有效 Cartier 除子}
是有效 Cartier 除子 $D \subset X$，使得 $D \to S$ 是概形的平坦态射。
\end{definition}

\noindent
提醒读者，这可能是非标准记号。特别地，在
\cite[IV, Section 21.15]{EGA} 中，只有当 $X \to S$ 平坦且局部有限表示时，
才讨论相对除子的概念。我们的定义稍微更一般。不过，事实证明，在许多情形
（但并非总是如此）下，若 $x \in D$，则 $X \to S$ 在 $x$ 处平坦。

\begin{lemma}
\label{divisors-lemma-sum-relative-effective-Cartier-divisor}
设 $f : X \to S$ 为概形的态射。若 $D_1, D_2 \subset X$ 是
$X/S$ 上的相对有效 Cartier 除子，则 $D_1 + D_2$ 也是
（定义 \ref{divisors-definition-sum-effective-Cartier-divisors}）。
\end{lemma}

\begin{proof}
这翻译成如下代数事实：
设 $A \to B$ 为环同态，且 $h_1, h_2 \in B$。
假设 $h_i$ 都是非零因子，并且 $B/h_iB$ 在 $A$ 上平坦。
则 $h_1h_2$ 是非零因子，并且 $B/h_1h_2B$ 在 $A$ 上平坦。
理由是存在短正合列
$$
0 \to B/h_1B \to B/h_1h_2B \to B/h_2B \to 0
$$
其中第一个箭头由乘以 $h_2$ 给出。由于两端的两项在 $A$ 上平坦，
中间项也平坦，见《代数》引理 \ref{algebra-lemma-flat-ses}。
\end{proof}

\begin{lemma}
\label{divisors-lemma-difference-relative-effective-Cartier-divisor}
设 $f : X \to S$ 为概形的态射。若 $D_1, D_2 \subset X$ 是
$X/S$ 上的相对有效 Cartier 除子，且作为闭子概形有 $D_1 \subset D_2$，
则有效 Cartier 除子 $D$，它满足 $D_2 = D_1 + D$
（引理 \ref{divisors-lemma-difference-effective-Cartier-divisors}）
是 $X/S$ 上的相对有效 Cartier 除子。
\end{lemma}

\begin{proof}
这翻译成如下代数事实：
设 $A \to B$ 为环同态，且 $h_1, h_2 \in B$。
假设 $h_i$ 都是非零因子，$B/h_iB$ 在 $A$ 上平坦，并且
$(h_2) \subset (h_1)$。于是可以写成 $h_2 = h h_1$，
其中 $h \in B$ 是非零因子。得到短正合列
$$
0 \to B/hB \to B/h_2B \to B/h_1B \to 0
$$
其中第一个箭头由乘以 $h_1$ 给出。由于右边两项在 $A$ 上平坦，
中间项也平坦，见《代数》引理 \ref{algebra-lemma-flat-ses}。
\end{proof}

\begin{lemma}
\label{divisors-lemma-flat-at-x}
设 $f : X \to S$ 为概形的态射。设 $D \subset X$ 为
$X/S$ 上的相对有效 Cartier 除子。若 $x \in D$ 且
$\mathcal{O}_{X, x}$ 诺特，则 $f$ 在 $x$ 处平坦。
\end{lemma}

\begin{proof}
令 $A = \mathcal{O}_{S, f(x)}$ 且 $B = \mathcal{O}_{X, x}$。
设 $h \in B$ 为生成 $D$ 之理想的元素。则 $h$ 是 $B$ 中的非零因子，
并且 $B/hB$ 是平坦局部 $A$-代数。设 $I \subset A$ 为有限生成理想。
考察交换图
$$
\xymatrix{
0 \ar[r] &
B \ar[r]_h &
B \ar[r] &
B/hB \ar[r] & 0 \\
0 \ar[r] &
B \otimes_A I \ar[r]^h \ar[u] &
B \otimes_A I \ar[r] \ar[u] &
B/hB \otimes_A I \ar[r] \ar[u] & 0
}
$$
由于 $B/hB$ 在 $A$ 上平坦，下面一列是短正合列，见《代数》引理
\ref{algebra-lemma-flat-tor-zero}。
由于 $B/hB$ 在 $A$ 上平坦，右侧竖直箭头是单射，见《代数》引理
\ref{algebra-lemma-flat}。
因此，乘以 $h$ 在中间竖直箭头的核 $K$ 上是满射。
由 Nakayama 引理，见《代数》引理 \ref{algebra-lemma-NAK}，
可知 $K= 0$。因此 $B$ 在 $A$ 上平坦，见《代数》引理
\ref{algebra-lemma-flat}。
\end{proof}

\noindent
下面的引理依赖平坦点集为开的代数版本。其概形论版本见
《态射进阶》第 \ref{more-morphisms-section-open-flat} 节。

\begin{lemma}
\label{divisors-lemma-flat-relative-Cartier-divisor}
设 $f : X \to S$ 为概形的态射。设 $D \subset X$ 为相对有效
Cartier 除子。若 $f$ 局部有限表示，则存在开子概形 $U \subset X$，
使得 $D \subset U$ 且 $f|_U : U \to S$ 平坦。
\end{lemma}

\begin{proof}
取 $x \in D$。只须找到开邻域 $U \subset X$，它包含 $x$，并使得
$f|_U$ 平坦。因此，该引理约化到如下情形：
$X = \Spec(B)$ 与 $S = \Spec(A)$ 都仿射，并且 $D$ 由非零因子
$h \in B$ 给出。根据假设，$B$ 是有限表示 $A$-代数，而
$B/hB$ 是平坦 $A$-代数。我们将使用绝对诺特逼近。

\medskip\noindent
写成 $B = A[x_1, \ldots, x_n]/(g_1, \ldots, g_m)$。
假设 $h$ 是 $h' \in A[x_1, \ldots, x_n]$ 的像。
选择一个在 $\mathbf{Z}$ 上有限型的子代数 $A_0 \subset A$，使多项式
$h', g_1, \ldots, g_m$ 的所有系数都属于 $A_0$。于是可以令
$B_0 = A_0[x_1, \ldots, x_n]/(g_1, \ldots, g_m)$，并令 $h_0$
为 $h'$ 在 $B_0$ 中的像。此时
$B = B_0 \otimes_{A_0} A$，且
$B/hB = B_0/h_0B_0 \otimes_{A_0} A$。由《代数》引理
\ref{algebra-lemma-flat-finite-presentation-limit-flat}，扩大 $A_0$ 后，
可以假设 $B_0/h_0B_0$ 在 $A_0$ 上平坦。
令 $K_0 = \Ker(h_0 : B_0 \to B_0)$。
由于 $B_0$ 在 $\mathbf{Z}$ 上有限型，可知 $K_0$ 是有限生成理想。
设 $A_1 \subset A$ 为有限型 $\mathbf{Z}$-子代数，它包含 $A_0$，
并用 $B_1$、$h_1$、$K_1$ 表示在 $A_1$ 上的相应对象。
由《代数进阶》引理
\ref{more-algebra-lemma-base-change-H1-regular}，映射
$K_0 \otimes_{A_0} A_1 \to K_1$ 是满射。另一方面，根据假设，
$h : B \to B$ 的核为零。因此，$K_0$ 的每个元素在 $K_1$ 中的像
都为零，这对充分大的子环 $A_1 \subset A$ 成立。由于 $K_0$ 有限生成，
可知 $K_1 = 0$，只须适当选择 $A_1$。

\medskip\noindent
令 $f_1 : X_1 \to S_1$ 等于 $\Spec$，它取自环同态 $A_1 \to B_1$。
令 $D_1 = \Spec(B_1/h_1B_1)$。由于
$B = B_1 \otimes_{A_1} A$，即 $X = X_1 \times_{S_1} S$，
现在只须对 $X_1 \to S_1$ 与相对有效 Cartier 除子 $D_1$
证明该引理，见《态射》引理
\ref{morphisms-lemma-base-change-module-flat}。
因此，问题已约化到 $A$ 为诺特环的情形。在该情形中，由引理
\ref{divisors-lemma-flat-at-x}，环同态 $A \to B$ 在每个素理想
$\mathfrak q$ 处平坦，这些素理想属于 $V(h)$。结合此情形下平坦点集为开的事实，见《代数》定理
\ref{algebra-theorem-openness-flatness}，即得结论。
\end{proof}

\noindent
还有下面的引理（其思想显然归于 Michael Artin，见 \cite{Nobile}），
它完全不需要任何有限性假设。

\begin{lemma}
\label{divisors-lemma-michael-artin}
设 $f : X \to S$ 为概形的态射。设 $D \subset X$ 为
$X/S$ 上的相对有效 Cartier 除子。若 $f$ 在 $X \setminus D$
的所有点处平坦，则 $f$ 平坦。
\end{lemma}

\begin{proof}
这翻译成如下代数事实：
设 $A \to B$ 为环同态，且 $h \in B$。
假设 $h$ 是非零因子，$B/hB$ 在 $A$ 上平坦，并且局部化
$B_h$ 在 $A$ 上平坦。则 $B$ 在 $A$ 上平坦。
理由是存在短正合列
$$
0 \to B \to B_h \to \colim_n (1/h^n)B/B \to 0
$$
并且第二项与第三项在 $A$ 上平坦，这意味着 $B$ 在 $A$ 上平坦
（见《代数》引理 \ref{algebra-lemma-flat-ses}）。
注意，平坦模的滤过余极限是平坦的
（见《代数》引理 \ref{algebra-lemma-colimit-flat}），
而对 $n$ 归纳可知，每个 $(1/h^n)B/B \cong B/h^nB$ 都在 $A$ 上平坦，
因为它嵌入短正合列
$$
0 \to B/h^{n - 1}B \xrightarrow{h} B/h^nB \to B/hB \to 0
$$
略去一些细节。
\end{proof}

\begin{example}
\label{divisors-example-relative-cartier-ambient-space-not-flat}
下面给出一个相对有效 Cartier 除子 $D$ 的例子，其中环境概形在 $D$
的一个邻域上不平坦。具体地，令 $A = k[t]$，并令
$$
B = k[t, x, y, x^{-1}y, x^{-2}y, \ldots]/(ty, tx^{-1}y, tx^{-2}y, \ldots)
$$
则 $B$ 在 $A$ 上不平坦，但 $B/xB \cong A$ 在 $A$ 上平坦。
此外，$x$ 是非零因子，因而在 $\Spec(B)$ 中定义了一个位于
$\Spec(A)$ 上的相对有效 Cartier 除子。
\end{example}

\noindent
若环境概形在基上平坦且局部有限表示，则可以用纤维来刻画相对有效
Cartier 除子。关于该引理稍有不同的表述，另见《态射进阶》引理
\ref{more-morphisms-lemma-slice-given-element}。

\begin{lemma}
\label{divisors-lemma-fibre-Cartier}
设 $\varphi : X \to S$ 为局部有限表示的平坦态射。设
$Z \subset X$ 为闭子概形。设 $x \in Z$，其像为 $s \in S$。
\begin{enumerate}
\item 若 $Z_s \subset X_s$ 在 $x$ 的一个邻域中是 Cartier 除子，
则存在开集 $U \subset X$ 以及相对有效 Cartier 除子 $D \subset U$，
使得 $Z \cap U \subset D$ 且 $Z_s \cap U = D_s$。
\item 若 $Z_s \subset X_s$ 在 $x$ 的一个邻域中是 Cartier 除子，
态射 $Z \to X$ 有限表示，并且 $Z \to S$ 在 $x$ 处平坦，
则可以选择 $U$ 与 $D$ 使得 $Z \cap U = D$。
\item 若 $Z_s \subset X_s$ 在 $x$ 的一个邻域中是 Cartier 除子，
并且 $Z$ 是 $X$ 的局部主闭子概形，这在 $x$ 的一个邻域中成立，
则可以选择 $U$ 与 $D$ 使得 $Z \cap U = D$。
\end{enumerate}
特别地，若 $Z \to S$ 局部有限表示且平坦，并且所有纤维
$Z_s \subset X_s$ 都是有效 Cartier 除子，则 $Z$ 是相对有效
Cartier 除子。类似地，若 $Z$ 是 $X$ 的局部主闭子概形，且所有纤维
$Z_s \subset X_s$ 都是有效 Cartier 除子，则 $Z$ 是相对有效
Cartier 除子。
\end{lemma}

\begin{proof}
选择仿射开邻域 $\Spec(A)$，它包含 $s$；以及仿射开邻域
$\Spec(B)$，它包含 $x$，并使得
$\varphi(\Spec(B)) \subset \Spec(A)$。
设 $\mathfrak p \subset A$ 为对应于 $s$ 的素理想。
设 $\mathfrak q \subset B$ 为对应于 $x$ 的素理想。
设 $I \subset B$ 为对应于 $Z$ 的理想。
由引理的初始假设可知 $A \to B$ 平坦且有限表示。
(1) 中的假设意味着，缩小 $\Spec(B)$ 后，可以假设
$I(B \otimes_A \kappa(\mathfrak p))$ 由一个元素生成，并且该元素在
$B \otimes_A \kappa(\mathfrak p)$ 中是非零因子。
设 $f \in I$ 映到这个生成元。我们断言，在可逆某个
$g \in B$（其中 $g \not \in \mathfrak q$）后，闭子概形
$D = V(f) \subset \Spec(B_g)$ 是相对有效 Cartier 除子。

\medskip\noindent
由《代数》引理
\ref{algebra-lemma-flat-finite-presentation-limit-flat}，
可以找到诺特环之间的平坦有限型环同态 $A_0 \to B_0$、一个元素
$f_0 \in B_0$、一个环同态 $A_0 \to A$ 以及同构
$A \otimes_{A_0} B_0 \cong B$。若
$\mathfrak p_0 = A_0 \cap \mathfrak p$，则可知
$$
B \otimes_A \kappa(\mathfrak p) =
\left(B_0 \otimes_{A_0} \kappa(\mathfrak p_0)\right)
\otimes_{\kappa(\mathfrak p_0))} \kappa(\mathfrak p)
$$
因此，$f_0$ 在 $B_0 \otimes_{A_0} \kappa(\mathfrak p_0)$ 中是非零因子。
由《代数》引理 \ref{algebra-lemma-grothendieck}，可知 $f_0$ 在
$(B_0)_{\mathfrak q_0}$ 中是非零因子，其中
$\mathfrak q_0 = B_0 \cap \mathfrak q$，并且
$(B_0/f_0B_0)_{\mathfrak q_0}$ 在 $A_0$ 上平坦。
因此，由《代数》引理
\ref{algebra-lemma-regular-sequence-in-neighbourhood}
与《代数》定理 \ref{algebra-theorem-openness-flatness}，
存在 $g_0 \in B_0$，其中 $g_0 \not \in \mathfrak q_0$，使得
$f_0$ 在 $(B_0)_{g_0}$ 中是非零因子，并且
$(B_0/f_0B_0)_{g_0}$ 在 $A_0$ 上平坦。因此可知
$D_0 = V(f_0) \subset \Spec((B_0)_{g_0})$ 是相对有效 Cartier 除子。
由于该性质在基变换下保持，见引理
\ref{divisors-lemma-relative-Cartier}，可得上面所述的断言，其中 $g$ 等于
$g_0$ 在 $B$ 中的像。

\medskip\noindent
至此已证明 (1)。为证明 (2)，考察闭浸入 $Z \to D$。
满环同态
$u : \mathcal{O}_{D, x} \to \mathcal{O}_{Z, x}$
是本质有限表示的平坦局部 $\mathcal{O}_{S, s}$-代数之间的映射，
并且模掉 $\mathfrak m_s$ 后成为同构。因此它是同构，见《代数》引理
\ref{algebra-lemma-mod-injective-general}。
由此可知 $Z \to D$ 在 $x$ 的一个邻域中是同构，见《代数》引理
\ref{algebra-lemma-local-isomorphism}。
为证明 (3)，必要时缩小 $U$ 后，可以假设 $D$ 的理想由单个非零因子
$f$ 生成，而 $Z$ 的理想由元素 $g$ 生成。于是 $f = gh$。
但是 $g|_{U_s}$ 与 $f|_{U_s}$ 在 $x$ 的一个邻域中截出同一个
有效 Cartier 除子。因此 $h|_{X_s}$ 是 $\mathcal{O}_{X_s, x}$ 中的单位，
从而 $h$ 是 $\mathcal{O}_{X, x}$ 中的单位，进而 $h$ 在 $x$
的一个开邻域中是单位。也就是说，有 $Z \cap U = D$，这是在缩小 $U$ 后成立。

\medskip\noindent
引理最后的断言立即由 (2)、(3) 以及如下事实得出：
$Z \to S$ 局部有限表示，当且仅当 $Z \to X$ 有限表示，见《态射》引理
\ref{morphisms-lemma-composition-finite-presentation} 与
\ref{morphisms-lemma-finite-presentation-permanence}。
\end{proof}

\section{浸入的法锥}
\label{divisors-section-normal-cone}

\noindent
设 $i : Z \to X$ 为闭浸入。设
$\mathcal{I} \subset \mathcal{O}_X$ 为相应的拟凝聚理想层。
考察分次 $\mathcal{O}_X$-代数的拟凝聚层
$\bigoplus_{n \geq 0} \mathcal{I}^n/\mathcal{I}^{n + 1}$。
由于每个层 $\mathcal{I}^n/\mathcal{I}^{n + 1}$ 都被
$\mathcal{I}$ 零化，由《态射》引理
\ref{morphisms-lemma-i-star-equivalence}，这个分次代数对应于一个
分次 $\mathcal{O}_Z$-代数的拟凝聚层。该拟凝聚分次
$\mathcal{O}_Z$-代数称为{\it $Z$ 在 $X$ 中的余法代数}，并且按照
《态射》第 \ref{morphisms-section-closed-immersions-quasi-coherent} 节
所述的记号滥用，常简记为
$\bigoplus_{n \geq 0} \mathcal{I}^n/\mathcal{I}^{n + 1}$。

\medskip\noindent
设 $f : Z \to X$ 为浸入。把 $f$ 的余法代数定义为闭浸入
$i : Z \to X \setminus \partial Z$ 的余法代数，其中
$\partial Z = \overline{Z} \setminus Z$。它常记为
$\bigoplus_{n \geq 0} \mathcal{I}^n/\mathcal{I}^{n + 1}$，
其中 $\mathcal{I}$ 是闭浸入
$i : Z \to X \setminus \partial Z$ 的理想层。

\begin{definition}
\label{divisors-definition-conormal-sheaf}
设 $f : Z \to X$ 为浸入。所谓{\it 余法代数
$\mathcal{C}_{Z/X, *}$，即 $Z$ 在 $X$ 中的余法代数}，或所谓
{\it $f$ 的余法代数}，是上面所述的分次 $\mathcal{O}_Z$-代数拟凝聚层
$\bigoplus_{n \geq 0} \mathcal{I}^n/\mathcal{I}^{n + 1}$。
\end{definition}

\noindent
因此 $\mathcal{C}_{Z/X, 1} = \mathcal{C}_{Z/X}$ 是该浸入的余法层。
此外，$\mathcal{C}_{Z/X, 0} = \mathcal{O}_Z$，而
$\mathcal{C}_{Z/X, n}$ 是由下述性质刻画的拟凝聚 $\mathcal{O}_Z$-模：
\begin{equation}
\label{divisors-equation-conormal-in-degree-n}
i_*\mathcal{C}_{Z/X, n} = \mathcal{I}^n/\mathcal{I}^{n + 1}
\end{equation}
其中 $i : Z \to X \setminus \partial Z$，且 $\mathcal{I}$ 如上是
$i$ 的理想层。最后，注意存在拟凝聚分次 $\mathcal{O}_Z$-代数的典范满射
\begin{equation}
\label{divisors-equation-conormal-algebra-quotient}
\text{Sym}^*(\mathcal{C}_{Z/X}) \longrightarrow \mathcal{C}_{Z/X, *}
\end{equation}
它在次数 $0$ 与 $1$ 上是同构。

\begin{lemma}
\label{divisors-lemma-affine-conormal-sheaf}
设 $i : Z \to X$ 为浸入。$i$ 的余法代数具有下列性质：
\begin{enumerate}
\item 设 $U \subset X$ 为任意开集，使得 $i(Z)$ 是 $U$ 的闭子集。
设 $\mathcal{I} \subset \mathcal{O}_U$ 为对应于闭子概形
$i(Z) \subset U$ 的理想层。则
$$
\mathcal{C}_{Z/X, *} =
i^*\left(\bigoplus\nolimits_{n \geq 0} \mathcal{I}^n\right) =
i^{-1}\left(
\bigoplus\nolimits_{n \geq 0} \mathcal{I}^n/\mathcal{I}^{n + 1}
\right)
$$
\item 对任意仿射开集 $\Spec(R) = U \subset X$，若
$Z \cap U = \Spec(R/I)$，则存在典范同构
$\Gamma(Z \cap U, \mathcal{C}_{Z/X, *}) = \bigoplus_{n \geq 0} I^n/I^{n + 1}$。
\end{enumerate}
\end{lemma}

\begin{proof}
大多由定义显然可得。注意，给定环 $R$ 以及一个理想 $I$（属于 $R$），有
$I^n/I^{n + 1} = I^n \otimes_R R/I$。略去细节。
\end{proof}

\begin{lemma}
\label{divisors-lemma-conormal-algebra-functorial}
设
$$
\xymatrix{
Z \ar[r]_i \ar[d]_f & X \ar[d]^g \\
Z' \ar[r]^{i'} & X'
}
$$
为概形范畴中的交换图。假设 $i$、$i'$ 都是浸入。存在分次
$\mathcal{O}_Z$-代数的典范映射
$$
f^*\mathcal{C}_{Z'/X', *}
\longrightarrow
\mathcal{C}_{Z/X, *}
$$
它由下列性质刻画：对每一对仿射开集
$(\Spec(R) = U \subset X, \Spec(R') = U' \subset X')$，若
$g(U) \subset U'$，且
$Z \cap U = \Spec(R/I)$、$Z' \cap U' = \Spec(R'/I')$，
则诱导映射
$$
\Gamma(Z' \cap U', \mathcal{C}_{Z'/X', *}) =
\bigoplus\nolimits (I')^n/(I')^{n + 1}
\longrightarrow
\bigoplus\nolimits_{n \geq 0} I^n/I^{n + 1} =
\Gamma(Z \cap U, \mathcal{C}_{Z/X, *})
$$
是由环同态 $f^\sharp : R' \to R$ 诱导的映射，该同态满足
$f^\sharp(I') \subset I$。
\end{lemma}

\begin{proof}
令 $\partial Z' = \overline{Z'} \setminus Z'$ 且
$\partial Z = \overline{Z} \setminus Z$。它们分别是 $X'$ 与 $X$
的闭子集。替换 $X'$，使用 $X' \setminus \partial Z'$；并替换 $X$，
使用 $X \setminus \big(g^{-1}(\partial Z') \cup \partial Z\big)$，
可知可以假设 $i$ 与 $i'$ 都是闭浸入。

\medskip\noindent
{\small
$g \circ i$ 经过 $i'$ 分解这一事实意味着，$g^*\mathcal{I}'$ 映入
$\mathcal{I}$，这是在典范映射 $g^*\mathcal{I}' \to \mathcal{O}_X$ 下成立，见《概形》引理
\ref{schemes-lemma-characterize-closed-subspace} 与
\ref{schemes-lemma-restrict-map-to-closed}。
因此得到拟凝聚层的诱导映射
$g^*((\mathcal{I}')^n/(\mathcal{I}')^{n + 1}) \to
\mathcal{I}^n/\mathcal{I}^{n + 1}$。
再由 $i$ 拉回，得到
$i^*g^*((\mathcal{I}')^n/(\mathcal{I}')^{n + 1}) \to
i^*(\mathcal{I}^n/\mathcal{I}^{n + 1})$。
注意
$i^*(\mathcal{I}^n/\mathcal{I}^{n + 1}) = \mathcal{C}_{Z/X, n}$。
另一方面，
$i^*g^*((\mathcal{I}')^n/(\mathcal{I}')^{n + 1}) =
f^*(i')^*((\mathcal{I}')^n/(\mathcal{I}')^{n + 1}) =
f^*\mathcal{C}_{Z'/X', n}$。
这给出所需映射。\par}

\medskip\noindent
检验该映射局部地由给定映射
$(I')^n/(I')^{n + 1} \to I^n/I^{n + 1}$ 描述，只须展开定义，
故略去。另一个观察是：给定任意 $x \in i(Z)$，确实存在仿射开邻域
$U$、$U'$，满足 $f(U) \subset U'$，并使 $Z \cap U$ 以及
$U' \cap Z'$ 都是闭的，且 $x \in U$。略去证明。因此，引理中的要求
确实刻画了该映射（也本可以用来定义它）。
\end{proof}

\begin{lemma}
\label{divisors-lemma-conormal-algebra-functorial-flat}
设
$$
\xymatrix{
Z \ar[r]_i \ar[d]_f & X \ar[d]^g \\
Z' \ar[r]^{i'} & X'
}
$$
{\small 为概形范畴中的纤维积图，其中 $i$、$i'$ 都是浸入。则引理
\ref{divisors-lemma-conormal-algebra-functorial} 的典范映射
$f^*\mathcal{C}_{Z'/X', *} \to \mathcal{C}_{Z/X, *}$ 是满射。
若 $g$ 平坦，则它是同构。\par}
\end{lemma}

\begin{proof}
设 $R' \to R$ 为环同态，且 $I' \subset R'$ 为理想。
令 $I = I'R$。则
$(I')^n/(I')^{n + 1} \otimes_{R'} R \to I^n/I^{n + 1}$
是满射。若 $R' \to R$ 平坦，则
$I^n = (I')^n \otimes_{R'} R$，从而该映射是同构。
\end{proof}

\begin{definition}
\label{divisors-definition-normal-cone}
设 $i : Z \to X$ 为概形的浸入。所谓{\it 法锥 $C_ZX$，即 $Z$ 在 $X$
中的法锥}，是
$$
C_ZX = \underline{\Spec}_Z(\mathcal{C}_{Z/X, *})
$$
见《构造》定义 \ref{constructions-definition-cone} 与
\ref{constructions-definition-abstract-cone}。$Z$ 在 $X$ 中的{\it 法丛}
是向量丛
$$
N_ZX = \underline{\Spec}_Z(\text{Sym}(\mathcal{C}_{Z/X}))
$$
见《构造》定义 \ref{constructions-definition-vector-bundle} 与
\ref{constructions-definition-abstract-vector-bundle}。
\end{definition}

\noindent
因此 $C_ZX \to Z$ 是 $Z$ 上的锥，$N_ZX \to Z$ 是 $Z$ 上的向量丛
（回忆：在我们的术语中，这并不意味着余法层是有限局部自由层）。
此外，分次代数的典范满射
(\ref{divisors-equation-conormal-algebra-quotient}) 定义了 $Z$ 上锥的典范闭浸入
\begin{equation}
\label{divisors-equation-normal-cone-in-normal-bundle}
C_ZX \longrightarrow N_ZX
\end{equation}

\section{正则理想层}
\label{divisors-section-regular-ideal-sheaves}

\noindent
本节把有效 Cartier 除子的概念推广到更高余维。回忆：元素列
$f_1, \ldots, f_r$（属于环 $R$）称为{\it 正则列}，如果对每个
$i = 1, \ldots, r$，元素 $f_i$ 在
$R/(f_1, \ldots, f_{i - 1})$ 上都是非零因子，并且
$R/(f_1, \ldots, f_r) \not = 0$，见《代数》定义
\ref{algebra-definition-regular-sequence}。
还可以施加三个密切相关的较弱条件。第一个条件是假设
$f_1, \ldots, f_r$ 是{\it Koszul-正则列}，即
$H_i(K_\bullet(f_1, \ldots, f_r)) = 0$ 对所有 $i > 0$ 成立，见《代数进阶》定义
\ref{more-algebra-definition-koszul-regular-sequence}。
称该列为{\it $H_1$-正则列}，如果
$H_1(K_\bullet(f_1, \ldots, f_r)) = 0$。另一个可以施加的条件是：令
$J = (f_1, \ldots, f_r)$，映射
$$
R/J[T_1, \ldots, T_r]
\longrightarrow
\bigoplus\nolimits_{n \geq 0}
J^n/J^{n + 1}
$$
把 $T_i$ 映到 $f_i \bmod J^2$，并且它是同构。在这种情形下，称
$f_1, \ldots, f_r$ 为{\it 拟正则列}，见《代数》定义
\ref{algebra-definition-quasi-regular-sequence}。
给定 $R$-模 $M$，还有 $M$-正则列与 $M$-拟正则列的概念。

\medskip\noindent
可以如下把这些概念推广到带环空间。设 $X$ 为带环空间，并设
$f_1, \ldots, f_r \in \Gamma(X, \mathcal{O}_X)$。
称 $f_1, \ldots, f_r$ 为{\it 正则列}，如果对每个
$i = 1, \ldots, r$，映射
\begin{equation}
\label{divisors-equation-map-regular}
f_i :
\mathcal{O}_X/(f_1, \ldots, f_{i - 1})
\longrightarrow
\mathcal{O}_X/(f_1, \ldots, f_{i - 1})
\end{equation}
都是层的单射。称 $f_1, \ldots, f_r$ 为{\it Koszul-正则列}，如果
Koszul 复形
\begin{equation}
\label{divisors-equation-koszul}
K_\bullet(\mathcal{O}_X, f_\bullet),
\end{equation}
见《模》定义 \ref{modules-definition-koszul-complex}，
在次数 $> 0$ 上无环。称 $f_1, \ldots, f_r$ 为
{\it $H_1$-正则列}，如果 Koszul 复形
$K_\bullet(\mathcal{O}_X, f_\bullet)$ 在次数 $1$ 上正合。
最后，称 $f_1, \ldots, f_r$ 为{\it 拟正则列}，如果映射
\begin{equation}
\label{divisors-equation-map-quasi-regular}
\mathcal{O}_X/\mathcal{J}[T_1, \ldots, T_r]
\longrightarrow
\bigoplus\nolimits_{d \geq 0} \mathcal{J}^d/\mathcal{J}^{d + 1}
\end{equation}
是层的同构，其中 $\mathcal{J} \subset \mathcal{O}_X$ 是由
$f_1, \ldots, f_r$ 生成的理想层。（还有 $\mathcal{F}$-正则列与
$\mathcal{F}$-拟正则列的概念，其中给定的 $\mathcal{O}_X$-模为
$\mathcal{F}$；若以后需要，我们将在那里引入。）

\begin{lemma}
\label{divisors-lemma-types-regular-sequences-implications}
设 $X$ 为带环空间。设
$f_1, \ldots, f_r \in \Gamma(X, \mathcal{O}_X)$。
有下列蕴含：
$f_1, \ldots, f_r$ 是正则列 $\Rightarrow$
$f_1, \ldots, f_r$ 是 Koszul-正则列 $\Rightarrow$
$f_1, \ldots, f_r$ 是 $H_1$-正则列 $\Rightarrow$
$f_1, \ldots, f_r$ 是拟正则列。
\end{lemma}

\begin{proof}
由于正合性可以在茎上检验，列 $f_1, \ldots, f_r$ 是正则列，
当且仅当映射
$$
f_i :
\mathcal{O}_{X, x}/(f_1, \ldots, f_{i - 1})
\longrightarrow
\mathcal{O}_{X, x}/(f_1, \ldots, f_{i - 1})
$$
对所有 $x \in X$ 都是单射。其他类型的正则性也都可以逐茎检验
（略去细节）。因此，这些蕴含由《代数进阶》引理
\ref{more-algebra-lemma-regular-koszul-regular}、
\ref{more-algebra-lemma-koszul-regular-H1-regular} 与
\ref{more-algebra-lemma-H1-regular-quasi-regular} 得出。
\end{proof}

\begin{definition}
\label{divisors-definition-regular-ideal-sheaf}
\begin{reference}
Koszul-正则理想层的概念由
\cite[Expose VII, Definition 1.4]{SGA6} 引入，其中称为正则理想层。
\end{reference}
设 $X$ 为带环空间。设 $\mathcal{J} \subset \mathcal{O}_X$ 为理想层。
\begin{enumerate}
\item 称 $\mathcal{J}$ {\it 正则}，如果对每个
$x \in \text{Supp}(\mathcal{O}_X/\mathcal{J})$，都存在开邻域
$x \in U \subset X$ 以及正则列
$f_1, \ldots, f_r \in \mathcal{O}_X(U)$，使得 $\mathcal{J}|_U$
由 $f_1, \ldots, f_r$ 生成。
\item 称 $\mathcal{J}$ {\it Koszul-正则}，如果对每个
$x \in \text{Supp}(\mathcal{O}_X/\mathcal{J})$，都存在开邻域
$x \in U \subset X$ 以及 Koszul-正则列
$f_1, \ldots, f_r \in \mathcal{O}_X(U)$，使得 $\mathcal{J}|_U$
由 $f_1, \ldots, f_r$ 生成。
\item 称 $\mathcal{J}$ {\it $H_1$-正则}，如果对每个
$x \in \text{Supp}(\mathcal{O}_X/\mathcal{J})$，都存在开邻域
$x \in U \subset X$ 以及 $H_1$-正则列
$f_1, \ldots, f_r \in \mathcal{O}_X(U)$，使得 $\mathcal{J}|_U$
由 $f_1, \ldots, f_r$ 生成。
\item 称 $\mathcal{J}$ {\it 拟正则}，如果对每个
$x \in \text{Supp}(\mathcal{O}_X/\mathcal{J})$，都存在开邻域
$x \in U \subset X$ 以及拟正则列
$f_1, \ldots, f_r \in \mathcal{O}_X(U)$，使得 $\mathcal{J}|_U$
由 $f_1, \ldots, f_r$ 生成。
\end{enumerate}
\end{definition}

\noindent
该概念的许多性质立即由环中正则列与拟正则列的相应概念得出。

\begin{lemma}
\label{divisors-lemma-regular-quasi-regular-scheme}
设 $X$ 为带环空间。设 $\mathcal{J}$ 为理想层。则有下列蕴含：
$\mathcal{J}$ 正则 $\Rightarrow$
$\mathcal{J}$ Koszul-正则 $\Rightarrow$
$\mathcal{J}$ 是 $H_1$-正则 $\Rightarrow$
$\mathcal{J}$ 拟正则。
\end{lemma}

\begin{proof}
该引理立即约化为引理
\ref{divisors-lemma-types-regular-sequences-implications}。
\end{proof}

\begin{lemma}
\label{divisors-lemma-quasi-regular-ideal}
设 $X$ 为局部带环空间。设 $\mathcal{J} \subset \mathcal{O}_X$ 为理想层。
则 $\mathcal{J}$ 拟正则，当且仅当下列条件成立：
\begin{enumerate}
\item $\mathcal{J}$ 是有限型 $\mathcal{O}_X$-模；
\item $\mathcal{J}/\mathcal{J}^2$ 是有限局部自由
$\mathcal{O}_X/\mathcal{J}$-模；并且
\item 典范映射
$$
\text{Sym}^n_{\mathcal{O}_X/\mathcal{J}}(\mathcal{J}/\mathcal{J}^2)
\longrightarrow
\mathcal{J}^n/\mathcal{J}^{n + 1}
$$
对所有 $n \geq 0$ 都是同构。
\end{enumerate}
\end{lemma}

\begin{proof}
显然，若 $U \subset X$ 是开集，并且 $\mathcal{J}|_U$ 由拟正则列
$f_1, \ldots, f_r \in \mathcal{O}_X(U)$ 生成，则 $\mathcal{J}|_U$
有限型，$\mathcal{J}|_U/\mathcal{J}^2|_U$ 自由且具有基
$f_1, \ldots, f_r$，而 (3) 中的映射是同构，因为它们是
(\ref{divisors-equation-map-quasi-regular}) 中次数 $n$ 部分的无坐标表述。
因此显然，拟正则蕴含条件 (1)、(2)、(3)。

\medskip\noindent
反过来，假设 (1)、(2)、(3) 成立。取一点
$x \in \text{Supp}(\mathcal{O}_X/\mathcal{J})$。则存在邻域
$U \subset X$，它包含 $x$，使得 $\mathcal{J}|_U/\mathcal{J}^2|_U$
是秩为 $r$、底环为 $\mathcal{O}_U/\mathcal{J}|_U$ 的自由模。
必要时缩小 $U$ 后，可以假设存在
$f_1, \ldots, f_r \in \mathcal{J}(U)$，它们映为
$\mathcal{J}|_U/\mathcal{J}^2|_U$ 作为
$\mathcal{O}_U/\mathcal{J}|_U$-模的一组基。
特别地，可知 $f_1, \ldots, f_r$ 在
$\mathcal{J}_x/\mathcal{J}^2_x$ 中的像生成该模。因此，由 Nakayama 引理
（《代数》引理 \ref{algebra-lemma-NAK}），可知
$f_1, \ldots, f_r$ 生成茎 $\mathcal{J}_x$。
于是，由于 $\mathcal{J}$ 有限型，根据《模》引理
\ref{modules-lemma-finite-type-surjective-on-stalk}，缩小 $U$ 后，
可以假设 $f_1, \ldots, f_r$ 生成 $\mathcal{J}$。
最后，由 (3) 与同构
$\mathcal{J}|_U/\mathcal{J}^2|_U = \bigoplus \mathcal{O}_U/\mathcal{J}|_U f_i$，显然
$f_1, \ldots, f_r \in \mathcal{O}_X(U)$ 是拟正则列。
\end{proof}

\begin{lemma}
\label{divisors-lemma-generate-regular-ideal}
设 $(X, \mathcal{O}_X)$ 为局部带环空间。设
$\mathcal{J} \subset \mathcal{O}_X$ 为理想层。设 $x \in X$，并设
$f_1, \ldots, f_r \in \mathcal{J}_x$ 的像给出
$\kappa(x)$-向量空间 $\mathcal{J}_x/\mathfrak m_x\mathcal{J}_x$ 的一组基。
\begin{enumerate}
\item 若 $\mathcal{J}$ 拟正则，则存在一个开邻域，使得
$f_1, \ldots, f_r \in \mathcal{O}_X(U)$ 构成生成
$\mathcal{J}|_U$ 的拟正则列。
\item 若 $\mathcal{J}$ 是 $H_1$-正则，则存在一个开邻域，使得
$f_1, \ldots, f_r \in \mathcal{O}_X(U)$ 构成 $H_1$-正则列并生成
$\mathcal{J}|_U$。
\item 若 $\mathcal{J}$ Koszul-正则，则存在一个开邻域，使得
$f_1, \ldots, f_r \in \mathcal{O}_X(U)$ 构成生成
$\mathcal{J}|_U$ 的 Koszul-正则列。
\end{enumerate}
\end{lemma}

\begin{proof}
首先假设 $\mathcal{J}$ 拟正则。可以选择开邻域
$U \subset X$（它包含 $x$）以及拟正则列
$g_1, \ldots, g_s \in \mathcal{O}_X(U)$，后者生成 $\mathcal{J}|_U$。
注意，这意味着 $\mathcal{J}/\mathcal{J}^2$ 是秩为 $s$、底环为
$\mathcal{O}_U/\mathcal{J}|_U$ 的自由模
（见引理 \ref{divisors-lemma-quasi-regular-ideal} 及其证明），因而 $r = s$。
可以缩小 $U$，并假设
$f_1, \ldots, f_r \in \mathcal{J}(U)$。于是可以写成
$$
f_i = \sum a_{ij} g_j
$$
其中 $a_{ij} \in \mathcal{O}_X(U)$。根据假设，矩阵
$A = (a_{ij})$ 映为 $\kappa(x)$ 上的可逆矩阵。
因此，再次缩小 $U$ 后，可以假设 $(a_{ij})$ 可逆。于是可知
$f_1, \ldots, f_r$ 给出 $(\mathcal{J}/\mathcal{J}^2)|_U$ 的一组基，
这证明 $f_1, \ldots, f_r$ 是 $U$ 上的拟正则列。

\medskip\noindent
注意，为证明 (2)、(3)，由于其假设比 (1) 的假设更强，可以直接假设
$f_1, \ldots, f_r \in \mathcal{J}(U)$ 且
$f_i = \sum a_{ij}g_j$，其中 $(a_{ij})$ 如上可逆，而现在
$g_1, \ldots, g_r$ 是 $H_1$-正则列或 Koszul-正则列。
由于 $f_1, \ldots, f_r$ 上的 Koszul 复形同构于
$g_1, \ldots, g_r$ 上的 Koszul 复形，此同构由矩阵 $(a_{ij})$ 给出
（见《代数进阶》引理 \ref{more-algebra-lemma-change-basis}），
可知 $f_1, \ldots, f_r$ 按所需是 $H_1$-正则或 Koszul-正则的。
\end{proof}

\begin{lemma}
\label{divisors-lemma-regular-ideal-sheaf-quasi-coherent}
概形上的任何正则、Koszul-正则、$H_1$-正则或拟正则理想层，
都是有限型拟凝聚理想层。
\end{lemma}

\begin{proof}
这是因为这样的理想层局部地由有限多个截面生成。概形上任何局部地由
截面生成的理想层都是拟凝聚的，见《概形》引理
\ref{schemes-lemma-closed-subspace-scheme}。
\end{proof}

\begin{lemma}
\label{divisors-lemma-regular-ideal-sheaf-scheme}
设 $X$ 为概形。设 $\mathcal{J}$ 为理想层。则 $\mathcal{J}$ 正则
（相应地，Koszul-正则、$H_1$-正则、拟正则），当且仅当对每个
$x \in \text{Supp}(\mathcal{O}_X/\mathcal{J})$，都存在仿射开邻域
$x \in U \subset X$，$U = \Spec(A)$，使得
$\mathcal{J}|_U = \widetilde{I}$，并且 $I$ 由正则列
（相应地，Koszul-正则列、$H_1$-正则列、拟正则列）
$f_1, \ldots, f_r \in A$ 生成。
\end{lemma}

\begin{proof}
根据假设，可以找到开邻域 $U$（它包含 $x$），使得 $\mathcal{J}$ 在其上由正则列
（相应地，Koszul-正则列、$H_1$-正则列、拟正则列）
$f_1, \ldots, f_r \in \mathcal{O}_X(U)$ 生成。缩小 $U$ 后，
可以假设 $U$ 仿射，记为 $U = \Spec(A)$。
由于由引理 \ref{divisors-lemma-regular-ideal-sheaf-quasi-coherent}，
$\mathcal{J}$ 拟凝聚，可知有 $\mathcal{J}|_U = \widetilde{I}$，其中
$I \subset A$ 为某个理想。现在可以使用如下事实：
$$
\widetilde{\ } : \text{Mod}_A \longrightarrow \QCoh(\mathcal{O}_U)
$$
是保持正合性的范畴等价。例如，函数 $f_i$ 生成 $\mathcal{J}$
这一事实，意味着把 $f_i$ 视为 $A$ 的元素时，它们生成 $I$。
映射 (\ref{divisors-equation-map-regular}) 为单射
（相应地，(\ref{divisors-equation-koszul}) 正合，(\ref{divisors-equation-koszul})
在次数 $1$ 上正合，(\ref{divisors-equation-map-quasi-regular}) 为同构）
这一事实，蕴含映射
$A/(f_1, \ldots, f_{i - 1}) \to A/(f_1, \ldots, f_{i - 1})$
（相应地，复形 $K_\bullet(A, f_1, \ldots, f_r)$，映射
$A/I[T_1, \ldots, T_r] \to \bigoplus I^n/I^{n + 1}$）
具有相应性质。因此，$f_1, \ldots, f_r \in A$ 是正则列
（相应地，Koszul-正则列、$H_1$-正则列、拟正则列），所在的环为 $A$。
\end{proof}

\begin{lemma}
\label{divisors-lemma-Noetherian-scheme-regular-ideal}
设 $X$ 为局部诺特概形。设 $\mathcal{J} \subset \mathcal{O}_X$
为拟凝聚理想层。设 $x$ 为 $\mathcal{O}_X/\mathcal{J}$ 之支集中的点。
下列条件等价：
\begin{enumerate}
\item $\mathcal{J}_x$ 在 $\mathcal{O}_{X, x}$ 中由正则列生成；
\item $\mathcal{J}_x$ 在 $\mathcal{O}_{X, x}$ 中由 Koszul-正则列生成；
\item $\mathcal{J}_x$ 由 $H_1$-正则列生成，该列位于
$\mathcal{O}_{X, x}$ 中；
\item $\mathcal{J}_x$ 在 $\mathcal{O}_{X, x}$ 中由拟正则列生成；
\item 存在仿射邻域 $U = \Spec(A)$，它包含 $x$，使得
$\mathcal{J}|_U = \widetilde{I}$，并且 $I$ 在 $A$ 中由正则列生成；
\item 存在仿射邻域 $U = \Spec(A)$，它包含 $x$，使得
$\mathcal{J}|_U = \widetilde{I}$，并且 $I$ 在 $A$ 中由 Koszul-正则列生成；
\item 存在仿射邻域 $U = \Spec(A)$，它包含 $x$，使得
$\mathcal{J}|_U = \widetilde{I}$，并且 $I$ 由 $H_1$-正则列生成，
该列位于 $A$ 中；
\item 存在仿射邻域 $U = \Spec(A)$，它包含 $x$，使得
$\mathcal{J}|_U = \widetilde{I}$，并且 $I$ 在 $A$ 中由拟正则列生成；
\item 存在邻域 $U$（它包含 $x$），使得 $\mathcal{J}|_U$ 正则；
\item 存在邻域 $U$（它包含 $x$），使得 $\mathcal{J}|_U$ Koszul-正则；
\item 存在邻域 $U$（它包含 $x$），使得 $\mathcal{J}|_U$ 是 $H_1$-正则；
\item 存在邻域 $U$（它包含 $x$），使得 $\mathcal{J}|_U$ 拟正则。
\end{enumerate}
特别地，在局部诺特概形上，正则、Koszul-正则、$H_1$-正则与拟正则
理想层的概念全都一致。
\end{lemma}

\begin{proof}
由引理 \ref{divisors-lemma-regular-ideal-sheaf-scheme} 可知
(5) $\Leftrightarrow$ (9)、(6) $\Leftrightarrow$ (10)、
(7) $\Leftrightarrow$ (11)，以及 (8) $\Leftrightarrow$ (12)。
显然，(5) $\Rightarrow$ (1)、(6) $\Rightarrow$ (2)、
(7) $\Rightarrow$ (3)，以及 (8) $\Rightarrow$ (4)。
由《代数》引理 \ref{algebra-lemma-regular-sequence-in-neighbourhood}，
有 (1) $\Rightarrow$ (5)。
由引理 \ref{divisors-lemma-regular-quasi-regular-scheme}，有
(9) $\Rightarrow$ (10) $\Rightarrow$ (11) $\Rightarrow$ (12)。
最后，由《代数》引理 \ref{algebra-lemma-quasi-regular-regular}，
有 (4) $\Rightarrow$ (1)。现在全部 12 个陈述都等价。
\end{proof}

\section{正则浸入}
\label{divisors-section-regular-immersions}

\noindent
设 $i : Z \to X$ 为概形的浸入。按定义，这意味着存在开子概形
$U \subset X$，使得 $Z$ 可与 $U$ 的一个闭子概形等同。令
$\mathcal{I} \subset \mathcal{O}_U$ 为相应的拟凝聚理想层。假设
$U' \subset X$ 是另一个这样的开子概形，并记
$\mathcal{I}' \subset \mathcal{O}_{U'}$ 为相应的拟凝聚理想层。则
$\mathcal{I}|_{U \cap U'} = \mathcal{I}'|_{U \cap U'}$。
此外，$\mathcal{O}_U/\mathcal{I}$ 的支集是 $Z$；它包含于
$U \cap U'$，并且也是 $\mathcal{O}_{U'}/\mathcal{I}'$ 的支集。
因此，由定义 \ref{divisors-definition-regular-ideal-sheaf} 可知，
$\mathcal{I}$ 是正则理想，当且仅当 $\mathcal{I}'$ 是正则理想。
对 Koszul-正则、$H_1$-正则或拟正则也同样成立。

\begin{definition}
\label{divisors-definition-regular-immersion}
\begin{reference}
Koszul-正则浸入的概念由
\cite[Expose VII, Definition 1.4]{SGA6} 引入，其中称为正则浸入。
\end{reference}
设 $i : Z \to X$ 为概形的浸入。选择开子概形
$U \subset X$，使得 $i$ 把 $Z$ 与 $U$ 的一个闭子概形等同，并记
$\mathcal{I} \subset \mathcal{O}_U$ 为相应的拟凝聚理想层。
\begin{enumerate}
\item 称 $i$ 为{\it 正则浸入}，如果 $\mathcal{I}$ 正则。
\item 称 $i$ 为{\it Koszul-正则浸入}，如果 $\mathcal{I}$ Koszul-正则。
\item 称 $i$ 为{\it $H_1$-正则浸入}，如果 $\mathcal{I}$ 是 $H_1$-正则。
\item 称 $i$ 为{\it 拟正则浸入}，如果 $\mathcal{I}$ 拟正则。
\end{enumerate}
\end{definition}

\noindent
上面的讨论表明，这一定义不依赖于 $U$ 的选择。这些条件按强弱递减的
顺序列出，见引理 \ref{divisors-lemma-regular-quasi-regular-immersion}。
Koszul-正则闭浸入经平滑局部化后是正则浸入，见引理
\ref{divisors-lemma-koszul-regular-smooth-locally-regular}。
在局部 Noether 情形，四种概念全都一致，见引理
\ref{divisors-lemma-Noetherian-scheme-regular-ideal}。

\begin{lemma}
\label{divisors-lemma-regular-quasi-regular-immersion}
设 $i : Z \to X$ 为概形的浸入。则有下列蕴含：
$i$ 正则 $\Rightarrow$
$i$ Koszul-正则 $\Rightarrow$
$i$ 是 $H_1$-正则 $\Rightarrow$
$i$ 拟正则。
\end{lemma}

\begin{proof}
该引理立即约化为引理
\ref{divisors-lemma-regular-quasi-regular-scheme}。
\end{proof}

\begin{lemma}
\label{divisors-lemma-regular-immersion-noetherian}
设 $i : Z \to X$ 为概形的浸入。假设 $X$ 局部 Noether。则
$i$ 正则 $\Leftrightarrow$
$i$ Koszul-正则 $\Leftrightarrow$
$i$ 是 $H_1$-正则 $\Leftrightarrow$
$i$ 拟正则。
\end{lemma}

\begin{proof}
这立即由引理
\ref{divisors-lemma-regular-quasi-regular-immersion}
与引理
\ref{divisors-lemma-Noetherian-scheme-regular-ideal}
得出。
\end{proof}

\begin{lemma}
\label{divisors-lemma-flat-base-change-regular-immersion}
设 $i : Z \to X$ 为正则浸入（相应地，Koszul-正则、
$H_1$-正则、拟正则浸入）。设 $X' \to X$ 为平坦态射。则基变换
$i' : Z \times_X X' \to X'$ 是正则浸入（相应地，Koszul-正则、
$H_1$-正则、拟正则浸入）。
\end{lemma}

\begin{proof}
借助引理
\ref{divisors-lemma-regular-ideal-sheaf-scheme}，
这转化为《代数》引理 \ref{algebra-lemma-flat-increases-depth} 与
\ref{algebra-lemma-flat-base-change-quasi-regular}
以及《代数进阶》引理
\ref{more-algebra-lemma-koszul-regular-flat-base-change}
中的代数陈述。
\end{proof}

\begin{lemma}
\label{divisors-lemma-quasi-regular-immersion}
设 $i : Z \to X$ 为概形的浸入。则 $i$ 是拟正则浸入，当且仅当
下列条件成立：
\begin{enumerate}
\item $i$ 局部有限表示；
\item 余法层 $\mathcal{C}_{Z/X}$ 有限局部自由；并且
\item 映射 (\ref{divisors-equation-conormal-algebra-quotient}) 是同构。
\end{enumerate}
\end{lemma}

\begin{proof}
开浸入局部有限表示。因此，可以把 $X$ 替换为开子概形
$U \subset X$，使得 $i$ 把 $Z$ 与 $U$ 的一个闭子概形等同；换言之，可以假设
$i$ 是闭浸入。令 $\mathcal{I} \subset \mathcal{O}_X$ 为相应的
拟凝聚理想层。回忆《态射》引理
\ref{morphisms-lemma-closed-immersion-finite-presentation}：
$\mathcal{I}$ 有限型，当且仅当 $i$ 局部有限表示。因此，所述等价性
由引理 \ref{divisors-lemma-quasi-regular-ideal} 和展开定义得出。
\end{proof}

\begin{lemma}
\label{divisors-lemma-transitivity-conormal-quasi-regular}
设 $Z \to Y \to X$ 为概形的浸入。假设 $Z \to Y$ 是
$H_1$-正则的。则《态射》引理
\ref{morphisms-lemma-transitivity-conormal} 的典范序列
$$
0 \to i^*\mathcal{C}_{Y/X} \to
\mathcal{C}_{Z/X} \to
\mathcal{C}_{Z/Y} \to 0
$$
正合且局部分裂。
\end{lemma}

\begin{proof}
由于 $\mathcal{C}_{Z/Y}$ 有限局部自由（见引理
\ref{divisors-lemma-quasi-regular-immersion}
与引理
\ref{divisors-lemma-regular-quasi-regular-scheme}），
只需证明该序列正合。由《态射》引理
\ref{morphisms-lemma-transitivity-conormal} 中已经证明的结论，
只需说明第一个映射是单射。仿射局部地工作，可把问题约化如下：
假设有环 $A$ 及理想 $I \subset J \subset A$。假设
$J/I \subset A/I$ 由一个 $H_1$-正则列生成。这是否蕴含
$I/I^2 \otimes_A A/J \to J/J^2$ 为单射？
注意 $I/I^2 \otimes_A A/J = I/IJ$。因此，需要证明
$I \cap J^2 = IJ$。这正是《代数进阶》引理
\ref{more-algebra-lemma-conormal-sequence-H1-regular} 的结论。
\end{proof}

\noindent
拟正则浸入的复合未必拟正则，见《代数》评注
\ref{algebra-remark-join-quasi-regular-sequences}。
其他类型的正则浸入在复合下保持。

\begin{lemma}
\label{divisors-lemma-composition-regular-immersion}
设 $i : Z \to Y$ 与 $j : Y \to X$ 为概形的浸入。
\begin{enumerate}
\item 若 $i$ 与 $j$ 都是正则浸入，则 $j \circ i$ 也是。
\item 若 $i$ 与 $j$ 都是 Koszul-正则浸入，则 $j \circ i$ 也是。
\item 若 $i$ 与 $j$ 都是 $H_1$-正则浸入，则 $j \circ i$ 也是。
\item 若 $i$ 是 $H_1$-正则浸入而 $j$ 是拟正则浸入，则
$j \circ i$ 是拟正则浸入。
\end{enumerate}
\end{lemma}

\begin{proof}
(1) 的代数版本是《代数》引理
\ref{algebra-lemma-join-regular-sequences}。
(2) 的代数版本是《代数进阶》引理
\ref{more-algebra-lemma-join-koszul-regular-sequences}。
(3) 的代数版本是《代数进阶》引理
\ref{more-algebra-lemma-join-H1-regular-sequences}。
(4) 的代数版本是《代数进阶》引理
\ref{more-algebra-lemma-join-quasi-regular-H1-regular}。
\end{proof}

\begin{lemma}
\label{divisors-lemma-permanence-regular-immersion}
设 $i : Z \to Y$ 与 $j : Y \to X$ 为概形的浸入。假设
$j$ 局部有限表示，并且《态射》引理
\ref{morphisms-lemma-transitivity-conormal} 的序列
$$
0 \to i^*\mathcal{C}_{Y/X} \to
\mathcal{C}_{Z/X} \to
\mathcal{C}_{Z/Y} \to 0
$$
正合且局部分裂。
\begin{enumerate}
\item 若 $j \circ i$ 是拟正则浸入，则 $i$ 也是。
\item 若 $j \circ i$ 是 $H_1$-正则浸入，则 $i$ 也是。
\item 若 $j$ 与 $j \circ i$ 都是 Koszul-正则浸入，则 $i$ 也是。
\end{enumerate}
\end{lemma}

\begin{proof}
缩小 $Y$ 与 $X$ 后，可以假设 $i$ 与 $j$ 都是闭浸入。记
$\mathcal{I} \subset \mathcal{O}_X$ 为 $Y$ 的理想层，并记
$\mathcal{J} \subset \mathcal{O}_X$ 为 $Z$ 的理想层。
余法序列是 $0 \to \mathcal{I}/\mathcal{I}\mathcal{J}
\to \mathcal{J}/\mathcal{J}^2 \to
\mathcal{J}/(\mathcal{I} + \mathcal{J}^2) \to 0$。
设 $z \in Z$，并置 $y = i(z)$、$x = j(y) = j(i(z))$。
选择 $f_1, \ldots, f_n \in \mathcal{I}_x$，其像构成
$\mathcal{I}_x/\mathfrak m_z\mathcal{I}_x$ 的一组基。把它扩充为
$f_1, \ldots, f_n, g_1, \ldots, g_m \in \mathcal{J}_x$，使其像构成
$\mathcal{J}_x/\mathfrak m_z\mathcal{J}_x$ 的一组基。这是可能的，
因为已经假设余法层的序列在 $z$ 的某个邻域中分裂，从而
$\mathcal{I}_x/\mathfrak m_x\mathcal{I}_x \to
\mathcal{J}_x/\mathfrak m_x\mathcal{J}_x$ 为单射。

\medskip\noindent
{\small (1) 的证明。由引理
\ref{divisors-lemma-generate-regular-ideal}，可以找到仿射开邻域
$U$（它包含 $x$），使得 $f_1, \ldots, f_n, g_1, \ldots, g_m$ 构成生成
$\mathcal{J}$ 的拟正则列。由于 $i$ 局部有限表示，还可以假设
$\mathcal{I}$ 由 $f_1, \ldots, f_n$ 生成（使用《态射》引理
\ref{morphisms-lemma-closed-immersion-finite-presentation}、
Nakayama 引理以及《模》引理
\ref{modules-lemma-finite-type-surjective-on-stalk}）。
因此，由《代数》引理
\ref{algebra-lemma-truncate-quasi-regular} 可知，
$g_1, \ldots, g_m$ 在 $Y \cap U$ 上诱导一个切出 $Z$ 的拟正则列。\par}

\medskip\noindent
(2) 的证明。与 (1) 的证明完全相同，只需改用《代数进阶》引理
\ref{more-algebra-lemma-truncate-H1-regular}。

\medskip\noindent
(3) 的证明。将引理
\ref{divisors-lemma-generate-regular-ideal}
应用两次，可找到仿射开邻域 $U$（它包含 $x$），使得
$f_1, \ldots, f_n$ 构成生成 $\mathcal{I}$ 的 Koszul-正则列，并且
$f_1, \ldots, f_n, g_1, \ldots, g_m$ 构成生成 $\mathcal{J}$ 的
Koszul-正则列。因此，由《代数进阶》引理
\ref{more-algebra-lemma-truncate-koszul-regular} 可知，
$g_1, \ldots, g_m$ 在 $Y \cap U$ 上诱导一个切出 $Z$ 的
Koszul-正则列。
\end{proof}

\begin{lemma}
\label{divisors-lemma-extra-permanence-regular-immersion-noetherian}
设 $i : Z \to Y$ 与 $j : Y \to X$ 为概形的浸入。取 $z \in Z$，
并记 $y \in Y$、$x \in X$ 为相应的点。假设 $X$ 局部 Noether。
下列条件等价：
\begin{enumerate}
\item $i$ 在 $z$ 的一个邻域中是正则浸入，并且 $j$ 在
$y$ 的一个邻域中是正则浸入；
\item $i$ 与 $j \circ i$ 在 $z$ 的一个邻域中都是正则浸入；
\item $j \circ i$ 在 $z$ 的一个邻域中是正则浸入，并且余法序列
$$
0 \to i^*\mathcal{C}_{Y/X} \to
\mathcal{C}_{Z/X} \to
\mathcal{C}_{Z/Y} \to 0
$$
在 $z$ 的一个邻域中分裂正合。
\end{enumerate}
\end{lemma}

\begin{proof}
由于 $X$（因而 $Y$）局部 Noether，四种正则浸入的概念全都一致；
此外，可以在局部环层面检验一个态射是否为正则浸入，见引理
\ref{divisors-lemma-Noetherian-scheme-regular-ideal}。
蕴含 (1) $\Rightarrow$ (2) 是引理
\ref{divisors-lemma-composition-regular-immersion}。
蕴含 (2) $\Rightarrow$ (3) 是引理
\ref{divisors-lemma-transitivity-conormal-quasi-regular}。
因此，只需证明 (3) 蕴含 (1)。

\medskip\noindent
假设 (3)。置 $A = \mathcal{O}_{X, x}$。记 $I \subset A$ 为满射
$\mathcal{O}_{X, x} \to \mathcal{O}_{Y, y}$ 的核，并记
$J \subset A$ 为满射
$\mathcal{O}_{X, x} \to \mathcal{O}_{Z, z}$ 的核。
注意，生成 $J$ 的任意极小元素列（位于 $A$ 中）都是拟正则列，因而是正则列，
见引理 \ref{divisors-lemma-generate-regular-ideal}。由假设，余法序列
$$
0 \to I/IJ \to J/J^2 \to J/(I + J^2) \to 0
$$
作为 $A/J$-模的序列分裂正合。因此，可以选择极小生成系
$f_1, \ldots, f_n, g_1, \ldots, g_m$，它生成 $J$，并使得
$f_1, \ldots, f_n \in I$ 是 $I$ 的一个极小生成系。如上所述，
$f_1, \ldots, f_n, g_1, \ldots, g_m$ 是 $A$ 中的正则列。
由正则列的定义直接可知，$f_1, \ldots, f_n$ 是 $A$ 中的正则列，
而 $\overline{g}_1, \ldots, \overline{g}_m$ 是 $A/I$ 中的正则列。
因此，$j$ 在 $y$ 处是正则浸入，而 $i$ 在 $z$ 处是正则浸入。
\end{proof}

\begin{remark}
\label{divisors-remark-not-always-extra-permanence}
在引理
\ref{divisors-lemma-extra-permanence-regular-immersion-noetherian}
的情形中，(1)、(2)、(3) 与下述条件{\bf 不}等价：
“$j \circ i$ 与 $j$ 分别在 $z$ 与 $y$ 处为正则浸入”。
一个例子是 $X = \mathbf{A}^1_k = \Spec(k[x])$、
$Y = \Spec(k[x]/(x^2))$ 与 $Z = \Spec(k[x]/(x))$。
\end{remark}

\begin{lemma}
\label{divisors-lemma-koszul-regular-smooth-locally-regular}
设 $i : Z \to X$ 为 Koszul 正则闭浸入。则存在满的平滑态射
$X' \to X$，使得基变换
$i' : Z \times_X X' \to X'$（它来自 $i$）是正则浸入。
\end{lemma}

\begin{proof}
可以假设 $X$ 仿射，并且 $Z$ 的理想由 Koszul-正则列生成：为此，
用某个适当仿射开覆盖的成员替换 $X$（仿射开集取引理
\ref{divisors-lemma-regular-ideal-sheaf-scheme} 中的那些）。仿射情形是
《代数进阶》引理
\ref{more-algebra-lemma-Koszul-regular-flat-locally-regular}。
\end{proof}

\begin{lemma}
\label{divisors-lemma-immersion-regular-regular-immersion}
设 $i : Z \to X$ 为浸入。若 $Z$ 与 $X$ 都是正则概形，则 $i$ 是
正则浸入。
\end{lemma}

\begin{proof}
设 $z \in Z$。由引理 \ref{divisors-lemma-Noetherian-scheme-regular-ideal}，
只需证明映射
$\mathcal{O}_{X, z} \to \mathcal{O}_{Z, z}$
的核由正则列生成。这由《代数》引理
\ref{algebra-lemma-regular-quotient-regular} 与
\ref{algebra-lemma-regular-ring-CM}
得出。
\end{proof}





\section{相对正则浸入}
\label{divisors-section-relative-regular-immersion}

\noindent
本节考察正则浸入的基变换性质。下述引理对正则浸入或 Koszul 浸入
并不成立，见《例》引理
\ref{examples-lemma-base-change-regular-sequence}。

\begin{lemma}
\label{divisors-lemma-relative-regular-immersion}
设 $f : X \to S$ 为概形的态射。设 $i : Z \subset X$ 为浸入。假设
\begin{enumerate}
\item $i$ 是 $H_1$-正则浸入（相应地，拟正则浸入）；并且
\item $Z \to S$ 是平坦态射。
\end{enumerate}
则对概形的每个态射 $g : S' \to S$，基变换
$Z' = S' \times_S Z \to X' = S' \times_S X$
是 $H_1$-正则浸入（相应地，拟正则浸入）。
\end{lemma}

\begin{proof}
展开定义并使用引理
\ref{divisors-lemma-regular-ideal-sheaf-scheme}，
可把它转化为《代数进阶》引理
\ref{more-algebra-lemma-relative-regular-immersion-algebra}。
\end{proof}

\noindent
这个引理引出了下面的定义。

\begin{definition}
\label{divisors-definition-relative-H1-regular-immersion}
设 $f : X \to S$ 为概形的态射。设 $i : Z \to X$ 为浸入。
\begin{enumerate}
\item 称 $i$ 为{\it 相对拟正则浸入}，如果 $Z \to S$ 平坦且
$i$ 是拟正则浸入。
\item 称 $i$ 为{\it 相对 $H_1$-正则浸入}，如果 $Z \to S$ 平坦且
$i$ 是 $H_1$-正则浸入。
\end{enumerate}
\end{definition}

\noindent
提醒读者，这可能不是标准记号。引理
\ref{divisors-lemma-relative-regular-immersion}
保证相对拟正则浸入（相应地，$H_1$-正则浸入）在任意基变换下保持。
相对 $H_1$-正则浸入是相对拟正则浸入，见引理
\ref{divisors-lemma-regular-quasi-regular-immersion}。请参见引理
\ref{divisors-lemma-flat-relative-H1-regular}
（或引理
\ref{divisors-lemma-relative-regular-immersion-flat-in-neighbourhood}）：
它表明，若 $Z \to X$ 是相对 $H_1$-正则浸入（或拟正则浸入），
且环境概形在 $S$ 上（平坦并）局部有限表示，则 $Z \to X$
实际上是正则浸入，并且在任意基变换后仍然如此。

\begin{lemma}
\label{divisors-lemma-quasi-regular-immersion-flat-at-x}
设 $f : X \to S$ 为概形的态射。设 $Z \to X$ 为相对拟正则浸入。
若 $x \in Z$ 且 $\mathcal{O}_{X, x}$ 是 Noether 环，则 $f$ 在 $x$ 处平坦。
\end{lemma}

\begin{proof}
设 $f_1, \ldots, f_r \in \mathcal{O}_{X, x}$ 为切出
$Z$ 在 $x$ 处之理想的拟正则列。由《代数》引理
\ref{algebra-lemma-quasi-regular-regular} 可知
$f_1, \ldots, f_r$ 是正则列。因此，$f_r$ 在
$\mathcal{O}_{X, x}/(f_1, \ldots, f_{r - 1})$ 上是非零因子，
并且其商是平坦 $\mathcal{O}_{S, f(x)}$-模。由引理
\ref{divisors-lemma-flat-at-x} 可知
$\mathcal{O}_{X, x}/(f_1, \ldots, f_{r - 1})$
是平坦 $\mathcal{O}_{S, f(x)}$-模。继续归纳，可得
$\mathcal{O}_{X, x}$ 是平坦 $\mathcal{O}_{S, s}$-模。
\end{proof}

\begin{lemma}
\label{divisors-lemma-relative-regular-immersion-flat-in-neighbourhood}
设 $X \to S$ 为概形的态射。设 $Z \to X$ 为浸入。假设
\begin{enumerate}
\item $X \to S$ 平坦且局部有限表示；
\item $Z \to X$ 是相对拟正则浸入。
\end{enumerate}
则 $Z \to X$ 是正则浸入，并且在任意基变换后仍然如此。
\end{lemma}

\begin{proof}
取 $x \in Z$，其像为 $s \in S$。为证明此结论，只需找到 $x$ 的
一个包含于 $U$ 的仿射邻域，使结论在该仿射开集上成立。因此，可以
假设 $X$ 仿射，并且存在拟正则列
$f_1, \ldots, f_r \in \Gamma(X, \mathcal{O}_X)$，使得
$Z = V(f_1, \ldots, f_r)$。由《代数进阶》引理
\ref{more-algebra-lemma-relative-regular-immersion-algebra}，
列 $f_1|_{X_s}, \ldots, f_r|_{X_s}$ 是
$\Gamma(X_s, \mathcal{O}_{X_s})$ 中的拟正则列。由于 $X_s$
Noether，这蕴含在必要时稍微缩小 $X$ 后，
$f_1|_{X_s}, \ldots, f_r|_{X_s}$ 是正则列；见《代数》引理
\ref{algebra-lemma-quasi-regular-regular} 与
\ref{algebra-lemma-regular-sequence-in-neighbourhood}。
由引理
\ref{divisors-lemma-fibre-Cartier} 可知，
$Z_1 = V(f_1) \subset X$ 是相对有效 Cartier 除子；这里必要时再次
稍微缩小 $X$。再次应用同一引理，
这次应用于 $Z_2 = V(f_1, f_2) \subset Z_1$，可知
$Z_2 \subset Z_1$ 是相对有效 Cartier 除子。依此类推，直到到达
$Z = Z_n = V(f_1, \ldots, f_n)$。由于相对有效 Cartier 除子的性质
在任意基变换下保持，见引理
\ref{divisors-lemma-relative-Cartier}，引理的最后一个陈述也随之成立。
\end{proof}

\begin{remark}
\label{divisors-remark-relative-regular-immersion-elements}
相对拟正则浸入的余维若为常数，则在基变换后不变。事实上，若有环映射
$A \to B$ 以及拟正则列 $f_1, \ldots, f_r \in B$，使得
$B/(f_1, \ldots, f_r)$ 在 $A$ 上平坦，则对任意环映射
$A \to A'$，由《代数进阶》引理
\ref{more-algebra-lemma-relative-regular-immersion-algebra}，
$f_1 \otimes 1, \ldots, f_r \otimes 1$ 在
$B' = B \otimes_A A'$ 中是拟正则列
（该引理也用于上面引理
\ref{divisors-lemma-relative-regular-immersion} 的证明）。
现在，引理
\ref{divisors-lemma-relative-regular-immersion-flat-in-neighbourhood}
的证明表明，若 $A \to B$ 平坦且局部有限表示，则对每个素理想
$\mathfrak q' \subset B'$，列
$f_1 \otimes 1, \ldots, f_r \otimes 1$ 甚至是局部环
$B'_{\mathfrak q'}$ 中的正则列。
\end{remark}

\begin{lemma}
\label{divisors-lemma-flat-relative-H1-regular}
设 $X \to S$ 为概形的态射。设 $Z \to X$ 为相对 $H_1$-正则浸入。
假设 $X \to S$ 局部有限表示。则
\begin{enumerate}
\item 存在开子概形 $U \subset X$，使得 $Z \subset U$ 且
$U \to S$ 平坦；并且
\item $Z \to X$ 是正则浸入，并且在任意基变换后仍然如此。
\end{enumerate}
\end{lemma}

\begin{proof}
取 $x \in Z$。为证明 (1)，只需找到开邻域
$U \subset X$，它包含 $x$，使得 $U \to S$ 平坦。因此，引理约化到如下情形：
$X = \Spec(B)$ 与 $S = \Spec(A)$ 仿射，并且 $Z$ 由
$H_1$-正则列 $f_1, \ldots, f_r \in B$ 给出。由假设，
$B$ 是有限表示 $A$-代数，并且
$B/(f_1, \ldots, f_r)B$ 是平坦 $A$-代数。下面使用绝对 Noether 逼近。

\medskip\noindent
写成 $B = A[x_1, \ldots, x_n]/(g_1, \ldots, g_m)$。假设
$f_i$ 是 $f_i' \in A[x_1, \ldots, x_n]$ 的像。选择有限型
$\mathbf{Z}$-子代数 $A_0 \subset A$，使得多项式
$f_1', \ldots, f_r', g_1, \ldots, g_m$ 的所有系数都属于 $A_0$。
置 $B_0 = A_0[x_1, \ldots, x_n]/(g_1, \ldots, g_m)$，并记
$f_{i, 0}$ 为 $f_i'$ 在 $B_0$ 中的像。则
$B = B_0 \otimes_{A_0} A$，并且
$$
B/(f_1, \ldots, f_r) =
B_0/(f_{0, 1}, \ldots, f_{0, r}) \otimes_{A_0} A.
$$
由《代数》引理
\ref{algebra-lemma-flat-finite-presentation-limit-flat}，扩大 $A_0$ 后，
可以假设 $B_0/(f_{0, 1}, \ldots, f_{0, r})$ 在 $A_0$ 上平坦。
此时 Koszul 同调群
$H_1(K_\bullet(B_0, f_{0, 1}, \ldots, f_{0, r}))$ 未必为零。
另一方面，由于 $B_0$ Noether，它是有限生成 $B_0$-模。令
$\xi_1, \ldots, \xi_n \in H_1(K_\bullet(B_0, f_{0, 1}, \ldots, f_{0, r}))$
为生成元。设 $A_0 \subset A_1 \subset A$ 为一个更大的有限型
$\mathbf{Z}$-子代数，它包含在 $A$ 中。记 $f_{1, i}$ 为 $f_{0, i}$ 在
$B_1 = B_0 \otimes_{A_0} A_1$ 中的像。由《代数进阶》引理
\ref{more-algebra-lemma-base-change-H1-regular}，映射
$$
H_1(K_\bullet(B_0, f_{0, 1}, \ldots, f_{0, r})) \otimes_{A_0} A_1
\longrightarrow
H_1(K_\bullet(B_1, f_{1, 1}, \ldots, f_{1, r}))
$$
是满射。此外，显然，对上述所有 $A_1$ 取法，复形
$K_\bullet(B_1, f_{1, 1}, \ldots, f_{1, r})$ 的余极限是复形
$K_\bullet(B, f_1, \ldots, f_r)$，后者在次数 $1$ 上无环。因此，由
《代数》引理 \ref{algebra-lemma-directed-colimit-exact}，
$$
\colim_{A_0 \subset A_1 \subset A}
H_1(K_\bullet(B_1, f_{1, 1}, \ldots, f_{1, r}))
= 0
$$
于是可选择 $A_1$，使得 $\xi_1, \ldots, \xi_n$ 全部映到
$H_1(K_\bullet(B_1, f_{1, 1}, \ldots, f_{1, r}))$ 中的零。
换言之，Koszul 同调群
$H_1(K_\bullet(B_1, f_{1, 1}, \ldots, f_{1, r}))$
为零。

\medskip\noindent
考察仿射概形的态射 $X_1 \to S_1$，它等于 $\Spec$，取自环映射
$A_1 \to B_1$，并令
$Z_1 = \Spec(B_1/(f_{1, 1}, \ldots, f_{1, r}))$。
由于 $B = B_1 \otimes_{A_1} A$，即 $X = X_1 \times_{S_1} S$，
并且类似地 $Z = Z_1 \times_S S_1$，现在只需对
$X_1 \to S_1$ 与相对 $H_1$-正则浸入 $Z_1 \to X_1$ 证明 (1)，
见《态射》引理
\ref{morphisms-lemma-base-change-module-flat}。
于是已约化到 $X \to S$ 为 Noether 概形之间的有限型态射的情形。
在此情形，由引理
\ref{divisors-lemma-quasi-regular-immersion-flat-at-x}，$X \to S$ 在
$Z$ 的每个点处平坦。结合此时平坦点集为开集这一事实，见《代数》定理
\ref{algebra-theorem-openness-flatness}，可知 (1) 成立。然后 (2) 由应用引理
\ref{divisors-lemma-relative-regular-immersion-flat-in-neighbourhood}
得出。
\end{proof}

\noindent
如果环境概形在基上平坦且局部有限表示，则可以用纤维刻画相对拟正则浸入。

\begin{lemma}
\label{divisors-lemma-fibre-quasi-regular}
设 $\varphi : X \to S$ 为局部有限表示的平坦态射。设
$T \subset X$ 为闭子概形。设 $x \in T$，其像为 $s \in S$。
\begin{enumerate}
\item 若 $T_s \subset X_s$ 在 $x$ 的一个邻域中是拟正则浸入，
则存在开集 $U \subset X$ 以及相对拟正则浸入 $Z \subset U$，
使得 $Z_s = T_s \cap U_s$ 且 $T \cap U \subset Z$。
\item 若 $T_s \subset X_s$ 在 $x$ 的一个邻域中是拟正则浸入，
态射 $T \to X$ 有限表示，并且 $T \to S$ 在 $x$ 处平坦，则可以选择
(1) 中的 $U$ 与 $Z$，使得 $T \cap U = Z$。
\item 若 $T_s \subset X_s$ 在 $x$ 的一个邻域中是拟正则浸入，
并且 $T$ 由 $c$ 个方程在 $x$ 的一个邻域中切出，其中
$c = \dim_x(X_s) - \dim_x(T_s)$，则可以选择 (1) 中的 $U$ 与 $Z$，
使得 $T \cap U = Z$。
\end{enumerate}
每种情形中，由引理
\ref{divisors-lemma-relative-regular-immersion-flat-in-neighbourhood}，
$Z \to U$ 都是正则浸入。特别地，若 $T \to S$ 局部有限表示且平坦，
并且所有纤维 $T_s \subset X_s$ 都是拟正则浸入，则 $T \to X$
是相对拟正则浸入。
\end{lemma}

\begin{proof}
选择仿射开邻域 $\Spec(A)$（它包含 $s$）与仿射开邻域
$\Spec(B)$（它包含 $x$），使得
$\varphi(\Spec(B)) \subset \Spec(A)$。令
$\mathfrak p \subset A$ 为对应于 $s$ 的素理想，令
$\mathfrak q \subset B$ 为对应于 $x$ 的素理想，并令
$I \subset B$ 为对应于 $T$ 的理想。由引理的初始假设可知
$A \to B$ 平坦且有限表示。(1) 中的假设意味着，在缩小
$\Spec(B)$ 后，可以假设 $I(B \otimes_A \kappa(\mathfrak p))$
由一个元素拟正则列生成。必要时把 $B$ 在某个
$g \in B$（且 $g \not \in \mathfrak q$）处局部化后，可以假设存在
$f_1, \ldots, f_r \in I$，它们在
$B \otimes_A \kappa(\mathfrak p)$ 中的像构成一个拟正则列，并生成
$I(B \otimes_A \kappa(\mathfrak p))$。由《代数》引理
\ref{algebra-lemma-quasi-regular-regular} 与
\ref{algebra-lemma-regular-sequence-in-neighbourhood}，再次局部化后，
可以假设 $f_1, \ldots, f_r \in I$ 在
$B \otimes_A \kappa(\mathfrak p)$ 中构成正则列。由引理
\ref{divisors-lemma-fibre-Cartier}，
$Z_1 = V(f_1) \subset \Spec(B)$ 是相对有效 Cartier 除子；这里必要时
再次局部化 $B$。
再次应用同一引理，这次应用于
$Z_2 = V(f_1, f_2) \subset Z_1$，可知 $Z_2 \subset Z_1$
是相对有效 Cartier 除子。依此类推，直到到达
$Z = Z_n = V(f_1, \ldots, f_n)$。于是
$Z \to \Spec(B)$ 是正则浸入，并且 $Z$ 在 $S$ 上平坦；特别地，
$Z \to \Spec(B)$ 在 $\Spec(A)$ 上是相对拟正则浸入。
这证明了 (1)。

\medskip\noindent
为证明 (2)，考察闭浸入 $Z \to D$。满环映射
$u : \mathcal{O}_{D, x} \to \mathcal{O}_{Z, x}$
是本质有限表示的平坦局部 $\mathcal{O}_{S, s}$-代数之间的映射，
并且模掉 $\mathfrak m_s$ 后成为同构。因此，由《代数》引理
\ref{algebra-lemma-mod-injective-general}，它是同构。
由此，$Z \to D$ 在 $x$ 的一个邻域中是同构，见《代数》引理
\ref{algebra-lemma-local-isomorphism}。

\medskip\noindent
为证明 (3)，必要时缩小 $U$ 后，可以假设 $Z$ 的理想由正则列
$f_1, \ldots, f_r$ 生成（见上面对 $Z$ 的构造），并且 $T$ 的理想由
$g_1, \ldots, g_c$ 生成。断言 $c = r$。事实上，
\begin{align*}
\dim_x(X_s) & = \dim(\mathcal{O}_{X_s, x}) +
\text{trdeg}_{\kappa(s)}(\kappa(x)), \\
\dim_x(T_s) & = \dim(\mathcal{O}_{T_s, x}) +
\text{trdeg}_{\kappa(s)}(\kappa(x)), \\
\dim(\mathcal{O}_{X_s, x}) & = \dim(\mathcal{O}_{T_s, x}) + r
\end{align*}
前两个等式来自《代数》引理
\ref{algebra-lemma-dimension-at-a-point-finite-type-field}，第二个等式
来自将《代数》引理
\ref{algebra-lemma-one-equation}
应用 $r$ 次。由于 $T \subset Z$，可知
$f_i = \sum b_{ij} g_j$。但是，$Z$ 与 $T$ 的理想切出
$X_s$ 的同一个拟正则闭子概形，而这一点在 $x$ 的一个邻域中成立。因此，矩阵
$(b_{ij}) \bmod \mathfrak m_x$ 可逆（略去一些细节）。于是
$(b_{ij})$ 在 $x$ 的一个开邻域中可逆。换言之，
$T \cap U = Z$，这是在缩小 $U$ 后成立。

\medskip\noindent
引理的最后几个陈述立即由 (2) 得出，并结合如下事实：
$Z \to S$ 局部有限表示，当且仅当 $Z \to X$ 有限表示；见《态射》引理
\ref{morphisms-lemma-composition-finite-presentation} 与
\ref{morphisms-lemma-finite-presentation-permanence}。
\end{proof}

\noindent
下述引理加强了《态射》引理
\ref{morphisms-lemma-section-smooth-morphism}。

\begin{lemma}
\label{divisors-lemma-section-smooth-regular-immersion}
设 $f : X \to S$ 为概形的平滑态射。设
$\sigma : S \to X$ 为 $f$ 的截面。则 $\sigma$ 是正则浸入。
\end{lemma}

\begin{proof}
由《概形》引理
\ref{schemes-lemma-semi-diagonal}，态射 $\sigma$ 是浸入。
用 $X$ 的一个开子概形替换它，该开子概形是 $\sigma(S)$ 的邻域；
于是可以假设 $\sigma$ 是闭浸入。
令 $T = \sigma(S)$ 为 $X$ 的相应闭子概形。由于 $T \to S$ 是同构，
它平坦且有限表示。平滑态射也平坦且局部有限表示，见《态射》引理
\ref{morphisms-lemma-smooth-flat} 与
\ref{morphisms-lemma-smooth-locally-finite-presentation}。
因此，根据引理
\ref{divisors-lemma-fibre-quasi-regular}，只需证明
$T_s \subset X_s$ 是拟正则闭子概形。这立即由《态射》引理
\ref{morphisms-lemma-section-smooth-morphism} 得出，但也可以如下直接说明。
设 $k$ 为域，并设 $A$ 为平滑 $k$-代数。设
$\mathfrak m \subset A$ 为剩余域等于 $k$ 的极大理想。则
$\mathfrak m$ 由拟正则列生成；这里必要时先将 $A$ 替换为某个
$A_g$，其中 $g \in A$、$g \not \in \mathfrak m$。在《代数》引理
\ref{algebra-lemma-characterize-smooth-over-field} 中已证明
$A_{\mathfrak m}$ 是正则局部环，故
$\mathfrak mA_{\mathfrak m}$ 由正则列生成。这的确蕴含
$\mathfrak m$ 由正则列生成（在必要时将 $A$ 替换为某个
$A_g$，其中 $g \in A$、$g \not \in \mathfrak m$），见《代数》引理
\ref{algebra-lemma-regular-sequence-in-neighbourhood}。
\end{proof}

\noindent
下述引理有某种逆，见引理
\ref{divisors-lemma-push-regular-immersion-thru-smooth}。

\begin{lemma}
\label{divisors-lemma-lift-regular-immersion-to-smooth}
设
$$
\xymatrix{
Y \ar[rd]_j \ar[rr]_i & & X \ar[ld] \\
& S
}
$$
为概形态射的交换图。假设 $X \to S$ 平滑，并且 $i$、$j$ 是浸入。
若 $j$ 是正则浸入（相应地，Koszul-正则、$H_1$-正则、拟正则浸入），
则 $i$ 也是。
\end{lemma}

\begin{proof}
可以把 $i$ 写成复合
$$
Y \to Y \times_S X \to X
$$
由引理
\ref{divisors-lemma-section-smooth-regular-immersion}，第一个箭头是正则浸入。
第二个箭头是 $Y \to S$ 的平坦基变换，因而是正则浸入（相应地，
Koszul-正则、$H_1$-正则、拟正则浸入），见引理
\ref{divisors-lemma-flat-base-change-regular-immersion}。
应用引理
\ref{divisors-lemma-composition-regular-immersion}
即得结论。
\end{proof}

\begin{lemma}
\label{divisors-lemma-immersion-lci-into-smooth-regular-immersion}
设
$$
\xymatrix{
Y \ar[rd] \ar[rr]_i & & X \ar[ld] \\
& S
}
$$
为概形态射的交换图。假设 $Y \to S$ 是 syntomic 态射，
$X \to S$ 平滑，并且 $i$ 是浸入。则 $i$ 是正则浸入。
\end{lemma}

\begin{proof}
用 $X$ 的一个开子概形替换它，该开子概形是 $i(Y)$ 的邻域；
于是可以假设 $i$ 是闭浸入。
令 $T = i(Y)$ 为 $X$ 的相应闭子概形。由于 $T \cong Y$，态射
$T \to S$ 平坦且有限表示（《态射》引理
\ref{morphisms-lemma-syntomic-locally-finite-presentation} 与
\ref{morphisms-lemma-syntomic-flat}）。
平滑态射也平坦且局部有限表示（《态射》引理
\ref{morphisms-lemma-smooth-flat} 与
\ref{morphisms-lemma-smooth-locally-finite-presentation}）。
因此，根据引理
\ref{divisors-lemma-fibre-quasi-regular}，只需证明
$T_s \subset X_s$ 是拟正则闭子概形。由于 $X_s$ 在域上局部有限型，
它是 Noether 概形（《态射》引理
\ref{morphisms-lemma-finite-type-noetherian}）。因此，可以逐点检验
$T_s \subset X_s$ 是拟正则浸入，见引理
\ref{divisors-lemma-Noetherian-scheme-regular-ideal}。取 $t \in T_s$。
由《态射》引理
\ref{morphisms-lemma-local-complete-intersection}，局部环
$\mathcal{O}_{T_s, t}$ 是 $\kappa(s)$ 上的局部完全交。局部环
$\mathcal{O}_{X_s, t}$ 正则，见《代数》引理
\ref{algebra-lemma-characterize-smooth-over-field}。由《代数》引理
\ref{algebra-lemma-lci-local} 可知，满射
$\mathcal{O}_{X_s, t} \to \mathcal{O}_{T_s, t}$
的核由正则列生成，这正是所需结论。
\end{proof}

\begin{lemma}
\label{divisors-lemma-immersion-smooth-into-smooth-regular-immersion}
设
$$
\xymatrix{
Y \ar[rd] \ar[rr]_i & & X \ar[ld] \\
& S
}
$$
为概形态射的交换图。假设 $Y \to S$ 平滑、$X \to S$ 平滑，并且
$i$ 是浸入。则 $i$ 是正则浸入。
\end{lemma}

\begin{proof}
这是引理
\ref{divisors-lemma-immersion-lci-into-smooth-regular-immersion}
的特例，因为平滑态射是 syntomic 态射，见《态射》引理
\ref{morphisms-lemma-smooth-syntomic}。
\end{proof}

\begin{lemma}
\label{divisors-lemma-push-regular-immersion-thru-smooth}
设
$$
\xymatrix{
Y \ar[rd]_j \ar[rr]_i & & X \ar[ld] \\
& S
}
$$
为概形态射的交换图。假设 $X \to S$ 平滑，并且 $i$ 与 $j$ 是浸入。
若 $i$ 是 Koszul-正则浸入（相应地，$H_1$-正则、拟正则浸入），
则 $j$ 也是。
\end{lemma}

\begin{proof}
以下不再另行说明地使用引理
\ref{divisors-lemma-regular-quasi-regular-immersion}。
任取 $y \in Y$。置 $x = i(y)$，并置 $s = j(y)$。
只需在用 $X$ 与 $S$ 的开邻域 $U$ 与 $V$（分别包含 $x$ 与 $s$）
替换它们，并用 $Y$ 中 $y$ 的一个开邻域（包含于
$i^{-1}(U) \cap j^{-1}(V)$）替换它后证明结论。

\medskip\noindent
先对 $X = \mathbf{A}^n_S$ 证明结论。用 $S$ 的仿射开子集 $V$
替换它，并用 $Y$ 的开子集 $j^{-1}(V)$ 替换它后，可以假设
$j$ 是闭浸入且 $S$ 仿射。
写成 $S = \Spec(A)$。于是 $j : Y \to S$ 把 $Y$ 同构到闭子概形
$\Spec(A/I)$，其中 $I \subset A$ 为某个理想。映射
$i : Y = \Spec(A/I) \to
\mathbf{A}^n_S = \Spec(A[x_1, \ldots, x_n])$
对应于 $A$-代数同态
$i^\sharp : A[x_1, \ldots, x_n] \to A/I$。
选择 $a_i \in A$，使其像为 $i^\sharp(x_i)$，此像位于 $A/I$ 中。
注意，闭浸入 $i$ 的理想是
$$
J = (x_1 - a_1, \ldots, x_n - a_n) + IA[x_1, \ldots, x_n].
$$
置 $K = (x_1 - a_1, \ldots, x_n - a_n)$。断言序列
$$
0 \to K/KJ \to J/J^2 \to J/(K + J^2) \to 0
$$
分裂正合。为说明这一点，注意 $K/K^2$ 以 $x_i - a_i$ 为基，
并且在环 $A[x_1, \ldots, x_n]/K \cong A$ 上自由。
因此，$K/KJ$ 在环 $A[x_1, \ldots, x_n]/J \cong A/I$ 上自由，
基仍相同。另一方面，取微分给出映射
$$
\text{d}_{A[x_1, \ldots, x_n]/A} :
J/J^2
\longrightarrow
\Omega_{A[x_1, \ldots, x_n]/A} \otimes_{A[x_1, \ldots, x_n]}
A[x_1, \ldots, x_n]/J
$$
它把生成元 $x_i - a_i$ 映到右边自由模的基元素 $\text{d}x_i$。
断言随之成立。此外，注意
$x_1 - a_1, \ldots, x_n - a_n$ 是
$A[x_1, \ldots, x_n]$ 中的正则列，其商环为
$A[x_1, \ldots, x_n]/(x_1 - a_1, \ldots, x_n - a_n) \cong A$。
于是，闭浸入 $i$ 有分解
$$
Y \to V(x_1 - a_1, \ldots, x_n - a_n) \to \mathbf{A}^n_S
$$
其中复合是 Koszul-正则浸入（相应地，$H_1$-正则、拟正则浸入），
第二个箭头是正则浸入，并且相应的余法序列分裂。现在由引理
\ref{divisors-lemma-permanence-regular-immersion}
得出结论。

\medskip\noindent
接着，在 $i$ 是 $H_1$-正则的情形证明结论。也就是说，按证明第一段
那样缩小后，可以假设 $Y$、$X$ 与 $S$ 都仿射。此时，可以选择闭浸入
$h : X \to \mathbf{A}^n_S$，它在 $S$ 上，其中 $n$ 为某个整数。由引理
\ref{divisors-lemma-immersion-smooth-into-smooth-regular-immersion}，
$h$ 是正则浸入。因此，$h \circ i$ 是 $H_1$-正则浸入，见引理
\ref{divisors-lemma-composition-regular-immersion}（注意，这一步在
“拟正则情形”中不成立）。于是约化到上面已经证明的
$X = \mathbf{A}^n_S$ 且 $S$ 仿射的情形。

\medskip\noindent
最后，假设 $i$ 拟正则。按证明第一段那样缩小后，可以使用《态射》引理
\ref{morphisms-lemma-smooth-etale-over-affine-space}，把 $f$ 分解为
$X \to \mathbf{A}^n_S \to S$，其中第一个态射
$X \to \mathbf{A}^n_S$ 是 \'etale 的。这把问题约化到两个情形：
(a) $X = \mathbf{A}^n_S$，以及 (b) $f$ 是 \'etale 的。
情形 (a) 已在证明第二段处理。情形 (b) 由下一段处理。

\medskip\noindent
假设 $f$ 是 \'etale 的。缩小后，可以假设 $X$、$Y$ 与 $S$ 仿射，
$i$ 与 $j$ 是闭浸入（略去一个小细节）。记
$S = \Spec(A)$、$X = \Spec(B)$ 且
$Y = \Spec(B/J) = \Spec(A/I)$。进一步缩小后，可以假设 $J$
由拟正则列生成。环映射 $A \to B$ 是 \'etale 的，因而形式 \'etale
（《代数》引理 \ref{algebra-lemma-formally-etale-etale}）。
因此，由《代数》引理
\ref{algebra-lemma-formally-etale-lift-infinitesimal}，
$\bigoplus I^n/I^{n + 1} \cong \bigoplus J^n/J^{n + 1}$。
由于 $J$ 由拟正则列生成，$I$ 也如此。这就完成了证明。
\end{proof}









\section{亚纯函数与亚纯截面}
\label{divisors-section-meromorphic-functions}

\noindent
本节只包含一般定义和一些初等结果。若要了解可能遇到的陷阱，见
\cite{misconceptions}\footnote{危险，威尔·罗宾逊！}。

\medskip\noindent
设 $(X, \mathcal{O}_X)$ 为局部带环空间。对任意开集
$U \subset X$，已经定义了集合
$\mathcal{S}(U) \subset \mathcal{O}_X(U)$，即 $\mathcal{O}_X$ 在 $U$ 上的
正则截面集合，见定义
\ref{divisors-definition-regular-section}。正则截面限制到更小开集后仍然正则。
因此，$\mathcal{S} : U \mapsto \mathcal{S}(U)$ 是
$\mathcal{O}_X$ 的子层（集合层）。有时记
$\mathcal{S} = \mathcal{S}_X$，以表明其对 $X$ 的依赖。此外，
$\mathcal{S}(U)$ 是环 $\mathcal{O}_X(U)$ 的乘法子集；这对每个 $U$
都成立。因此，可以考察环的预层
$$
U \longmapsto \mathcal{S}(U)^{-1} \mathcal{O}_X(U),
$$
见《模》引理
\ref{modules-lemma-simple-invert}。

\begin{definition}
\label{divisors-definition-sheaf-meromorphic-functions}
设 $(X, \mathcal{O}_X)$ 为局部带环空间。$X$ 上的{\it 亚纯函数层}
是与上面所示预层相伴的层{\it $\mathcal{K}_X$}。$X$ 上的
{\it 亚纯函数}是 $\mathcal{K}_X$ 的整体截面。
\end{definition}

\noindent
由于每个 $\mathcal{S}(U)$ 的每个元素都是
$\mathcal{O}_X(U)$ 上的非零因子，可知环层的自然映射
$\mathcal{O}_X \to \mathcal{K}_X$ 是单射。

\begin{example}
\label{divisors-example-no-change}
设
$A = \mathbf{C}[x, \{y_\alpha\}_{\alpha \in \mathbf{C}}]/
((x - \alpha)y_\alpha, y_\alpha y_\beta)$。$A$ 的每个元素都可唯一写成
$f(x) + \sum \lambda_\alpha y_\alpha$，其中
$f(x) \in \mathbf{C}[x]$ 且 $\lambda_\alpha \in \mathbf{C}$。
令 $X = \Spec(A)$。在此情形，$\mathcal{O}_X = \mathcal{K}_X$，
因为在任意仿射开集 $D(f)$ 上，环 $A_f$ 中的每个非零因子都是单位
（略去证明）。
\end{example}

\noindent
设 $(X, \mathcal{O}_X)$ 为局部带环空间。设 $\mathcal{F}$ 为
$\mathcal{O}_X$-模层。考察预层
$U \mapsto \mathcal{S}(U)^{-1}\mathcal{F}(U)$。它的层化是层
$\mathcal{F} \otimes_{\mathcal{O}_X} \mathcal{K}_X$，见《模》引理
\ref{modules-lemma-simple-invert-module}。

\begin{definition}
\label{divisors-definition-meromorphic-section}
设 $X$ 为局部带环空间。设 $\mathcal{F}$ 为 $\mathcal{O}_X$-模层。
\begin{enumerate}
\item 记 $\mathcal{K}_X(\mathcal{F})$ 为如下
$\mathcal{K}_X$-模层：它是预层
$U \mapsto \mathcal{S}(U)^{-1}\mathcal{F}(U)$ 的层化。等价地，
$\mathcal{K}_X(\mathcal{F}) =
\mathcal{F} \otimes_{\mathcal{O}_X} \mathcal{K}_X$（见上文）。
\item $\mathcal{F}$ 的{\it 亚纯截面}是
$\mathcal{K}_X(\mathcal{F})$ 的整体截面。
\end{enumerate}
\end{definition}

\noindent
特别地，对任意点 $x \in X$，有
$$
\mathcal{K}_X(\mathcal{F})_x
=
\mathcal{F}_x \otimes_{\mathcal{O}_{X, x}} \mathcal{K}_{X, x}
=
\mathcal{S}_x^{-1}\mathcal{F}_x
$$
不过必须谨慎，因为 $\mathcal{S}_x$ 未必等于局部环
$\mathcal{O}_{X, x}$ 中的非零因子集合。具体地，总有到全商环的单射
$$
\mathcal{K}_{X, x} \longrightarrow Q(\mathcal{O}_{X, x})
$$
它为满射，当且仅当 $\mathcal{S}_x$ 是
$\mathcal{O}_{X, x}$ 中的非零因子集合。亚纯截面层一般不是拟凝聚模，
但它们与拟凝聚模有一些共同性质。

\begin{definition}
\label{divisors-definition-pullback-meromorphic-sections}
设 $f : (X, \mathcal{O}_X) \to (Y, \mathcal{O}_Y)$ 为局部带环空间的态射。
称{\it 亚纯函数对 $f$ 的拉回有定义}，如果对每一对开集
$U \subset X$、$V \subset Y$，只要
$f(U) \subset V$，并对任意截面 $s \in \Gamma(V, \mathcal{S}_Y)$，
拉回 $f^\sharp(s) \in \Gamma(U, \mathcal{O}_X)$ 都属于
$\Gamma(U, \mathcal{S}_X)$。
\end{definition}

\noindent
在这种情形，存在诱导映射
$f^\sharp : f^{-1}\mathcal{K}_Y \to \mathcal{K}_X$；换言之，得到
带环空间态射的交换图
$$
\xymatrix{
(X, \mathcal{K}_X) \ar[r] \ar[d]^f &
(X, \mathcal{O}_X) \ar[d]^f \\
(Y, \mathcal{K}_Y) \ar[r] &
(Y, \mathcal{O}_Y)
}
$$
有时记 $f^*(s) = f^\sharp(s)$，其中
$s \in \Gamma(Y, \mathcal{K}_Y)$。

\begin{lemma}
\label{divisors-lemma-pullback-meromorphic-sections-defined}
设 $f : X \to Y$ 为概形的态射。在下列每种情形中，亚纯函数拉回有定义：
\begin{enumerate}
\item $X$ 的每个弱伴随点都映到 $Y$ 的某个不可约分支的泛点；
\item $X$、$Y$ 整，且 $f$ 支配；
\item $X$ 整，且 $X$ 的泛点映到 $Y$ 的某个不可约分支的泛点；
\item $X$ 约化，且 $X$ 的每个不可约分支的每个泛点都映到
$Y$ 的某个不可约分支的泛点；
\item $X$ 局部 Noether，且 $X$ 的任意伴随点都映到
$Y$ 的某个不可约分支的泛点；
\item $X$ 局部 Noether、没有嵌入点，并且 $X$ 的每个不可约分支的
泛点都映到 $Y$ 的某个不可约分支的泛点；以及
\item $f$ 平坦。
\end{enumerate}
\end{lemma}

\begin{proof}
问题在 $X$ 与 $Y$ 上都是局部的。因此，约化到
$X = \Spec(A)$、$Y = \Spec(R)$，并且 $f$ 由环映射
$\varphi : R \to A$ 给出的情形。由引理
\ref{divisors-lemma-regular-section-structure-sheaf}
中对结构层正则截面的刻画，只需证明 $R \to A$ 把非零因子映到
非零因子。设 $t \in R$ 为非零因子。

\medskip\noindent
若 $R \to A$ 平坦，则 $t : R \to R$ 为单射这一事实表明
$t : A \to A$ 为单射。这证明了 (7)。

\medskip\noindent
在其他情形中，注意 $t$ 不属于 $R$ 的任何极小素理想
（因为环的极小素理想中的每个元素都是零因子）。
因此，在情形 (1)，$\varphi(t)$ 不属于 $A$ 的任何弱伴随素理想。
此情形遂由《代数》引理
\ref{algebra-lemma-weakly-ass-zero-divisors}
得出。由引理 \ref{divisors-lemma-ass-weakly-ass}，情形 (5) 是 (1) 的特例。
情形 (6) 由 (5) 和定义得出。由引理
\ref{divisors-lemma-weakass-reduced}，情形 (4) 是 (1) 的特例。
情形 (2)、(3) 是 (4) 的特例。
\end{proof}


\begin{lemma}
\label{divisors-lemma-meromorphic-weakass-finite}
设 $X$ 为满足下列条件的概形：
\begin{enumerate}
\item[(a)] $X$ 的每个弱伴随点都是 $X$ 的某个不可约分支的泛点；并且
\item[(b)] 每个拟紧开集只有有限多个不可约分支。
\end{enumerate}
令 $X^0$ 为 $X$ 的不可约分支的泛点集合。则有
$$
\mathcal{K}_X =
\bigoplus\nolimits_{\eta \in X^0} j_{\eta, *}\mathcal{O}_{X, \eta} =
\prod\nolimits_{\eta \in X^0} j_{\eta, *}\mathcal{O}_{X, \eta}
$$
其中 $j_\eta : \Spec(\mathcal{O}_{X, \eta}) \to X$ 是《概形》第
\ref{schemes-section-points} 节的典范映射。此外：
\begin{enumerate}
\item $\mathcal{K}_X$ 是拟凝聚 $\mathcal{O}_X$-代数层；
\item 对每个拟凝聚 $\mathcal{O}_X$-模 $\mathcal{F}$，$\mathcal{F}$ 的亚纯截面层
$$
\mathcal{K}_X(\mathcal{F}) =
\bigoplus\nolimits_{\eta \in X^0} j_{\eta, *}\mathcal{F}_\eta =
\prod\nolimits_{\eta \in X^0} j_{\eta, *}\mathcal{F}_\eta
$$
是拟凝聚的；
\item $\mathcal{S}_x \subset \mathcal{O}_{X, x}$ 是非零因子集合，
对任意 $x \in X$；
\item $\mathcal{K}_{X, x}$ 是 $\mathcal{O}_{X, x}$ 的全商环，
对任意 $x \in X$；
\item $\mathcal{K}_X(U)$ 等于 $\mathcal{O}_X(U)$ 的全商环，
对任意仿射开集 $U \subset X$；
\item $X$ 的有理函数环（《态射》定义
\ref{morphisms-definition-rational-function}）就是 $X$ 的亚纯函数环；
写成公式即 $R(X) = \Gamma(X, \mathcal{K}_X)$。
\end{enumerate}
\end{lemma}

\begin{proof}
注意，局部有限的模层直和等于其直积；例如，这可以在茎上检验。又因
$\mathcal{K}_X(\mathcal{F}) =
\mathcal{F} \otimes_{\mathcal{O}_X} \mathcal{K}_X$，
可知 (2) 由其他陈述推出。还要注意，当 $X^0$ 有限时，(6) 由引理的
初始陈述与《态射》引理
\ref{morphisms-lemma-integral-scheme-rational-functions}
得出；一般情形的 (6) 再由此粘合而得（略去论证）。

\medskip\noindent
令
$j : Y = \coprod\nolimits_{\eta \in X^0} \Spec(\mathcal{O}_{X, \eta}) \to X$
为态射 $j_\eta$ 的余积。需要证明
$\mathcal{K}_X = j_*\mathcal{O}_Y$。首先注意，
$\mathcal{K}_Y = \mathcal{O}_Y$，因为 $Y$ 是 $0$ 维局部环之谱的不交并：
在零维局部环中，每个非零因子都是单位。其次注意，由引理
\ref{divisors-lemma-pullback-meromorphic-sections-defined}，$j$ 的亚纯函数拉回
有定义。这给出映射
$$
\mathcal{K}_X \longrightarrow j_*\mathcal{O}_Y.
$$
设 $\Spec(A) = U \subset X$ 为仿射开集。则 $A$ 是只有有限多个
极小素理想 $\mathfrak q_1, \ldots, \mathfrak q_t$ 的环，并且
$A$ 的每个弱伴随素理想都是某个 $\mathfrak q_i$。由《代数》引理
\ref{algebra-lemma-total-ring-fractions-no-embedded-points}
与 \ref{algebra-lemma-weakly-ass-zero-divisors}，得到
$Q(A) = \prod A_{\mathfrak q_i}$。换言之，预层的取值
$U \mapsto \mathcal{S}(U)^{-1}\mathcal{O}_X(U)$ 就与
$j_*\mathcal{O}_Y(U)$ 在仿射开集 $U$ 上一致。因此，上示映射是同构，
从而证明引理陈述中
第一个上示等式。

\medskip\noindent
最后，证明 (1)、(3)、(4)、(5)。在证明
$\mathcal{K}_X = j_*\mathcal{O}_Y$ 的过程中，已经看到了 (5)。
态射 $j$ 因我们关于 $X$ 的不可约分支集合局部有限这一假设而拟紧。
因此，$j$ 拟紧且拟分离（因为 $Y$ 分离）。由《概形》引理
\ref{schemes-lemma-push-forward-quasi-coherent}，
$j_*\mathcal{O}_Y$ 拟凝聚。这证明了 (1)。设 $x \in X$。
可以选择仿射开邻域 $U = \Spec(A)$ 作为 $x$ 的邻域，使其所有
不可约分支都经过 $x$。于是 $A \subset A_\mathfrak p$，因为 $A$ 的每个弱伴随素理想
都包含于 $\mathfrak p$，故由《代数》引理
\ref{algebra-lemma-weakly-ass-zero-divisors}，
$A \setminus \mathfrak p$ 中的元素都是非零因子。
容易推出，$A_\mathfrak p$ 的任意非零因子都是 $x$ 的某个
（可能更小的）仿射开邻域上的非零因子的像。这证明了 (3)。
(4) 由 (3) 通过计算茎得出。
\end{proof}

\begin{definition}
\label{divisors-definition-regular-meromorphic-section}
设 $X$ 为局部带环空间。设 $\mathcal{L}$ 为可逆
$\mathcal{O}_X$-模。设 $s$ 是 $\mathcal{L}$ 的一个亚纯截面。若其所诱导的映射
$\mathcal{K}_X \to \mathcal{K}_X(\mathcal{L})$ 为单射，则称其
{\it 正则}。换言之，$s$ 是可逆
$\mathcal{K}_X$-模 $\mathcal{K}_X(\mathcal{L})$ 的正则截面，见定义
\ref{divisors-definition-regular-section}。
\end{definition}

\noindent
下面明确说明何时可以拉回（正则）亚纯截面。

\begin{lemma}
\label{divisors-lemma-meromorphic-sections-pullback}
设 $f : X \to Y$ 为局部带环空间的态射。假设 $f$ 的亚纯函数拉回有定义
（见定义
\ref{divisors-definition-pullback-meromorphic-sections}）。
\begin{enumerate}
\item 设 $\mathcal{F}$ 为 $\mathcal{O}_Y$-模层。存在典范拉回映射
$f^* : \Gamma(Y, \mathcal{K}_Y(\mathcal{F})) \to
\Gamma(X, \mathcal{K}_X(f^*\mathcal{F}))$，用于 $\mathcal{F}$ 的亚纯截面。
\item 设 $\mathcal{L}$ 为可逆 $\mathcal{O}_Y$-模。
设 $s$ 是 $\mathcal{L}$ 的正则亚纯截面。则 $f^*s$ 是
$f^*\mathcal{L}$ 的正则亚纯截面。
\end{enumerate}
\end{lemma}

\begin{proof}
略。
\end{proof}

\begin{lemma}
\label{divisors-lemma-regular-meromorphic-ideal-denominators}
设 $X$ 为概形。设 $\mathcal{L}$ 为可逆 $\mathcal{O}_X$-模。
设 $s$ 为 $\mathcal{L}$ 的正则亚纯截面。记
$\mathcal{I} \subset \mathcal{O}_X$ 为由下述规则定义的理想层：
$$
\mathcal{I}(V)
=
\{f \in \mathcal{O}_X(V) \mid fs \in \mathcal{L}(V)\}.
$$
由于 $\mathcal{L}(V) \subset \mathcal{K}_X(\mathcal{L})(V)$，
该公式有意义。则 $\mathcal{I}$ 是拟凝聚理想层，并且有单射
$$
1 : \mathcal{I} \longrightarrow \mathcal{O}_X,
\quad
s : \mathcal{I} \longrightarrow \mathcal{L}
$$
二者的余核都支撑在 $X$ 的闭无处稠密子集上。
\end{lemma}

\begin{proof}
问题在 $X$ 上是局部的。因此，可以假设 $X = \Spec(A)$ 且
$\mathcal{L} = \mathcal{O}_X$。进一步缩小后，可以假设
$s = a/b$，其中 $a, b \in A$ 在 $A$ 中{\it 都}是非零因子。
置 $I = \{x \in A \mid x(a/b) \in A\}$。

\medskip\noindent
为证明 $\mathcal{I}$ 拟凝聚，需要证明
$I_f = \{x \in A_f \mid x(a/b) \in A_f\}$，对每个 $f \in A$。若
$c/f^n \in A_f$ 且 $(c/f^n)(a/b) \in A_f$，则
$f^mc(a/b) \in A$ 对某个 $m$ 成立，故 $c/f^n \in I_f$。反过来，容易看出
$I_f$ 包含于 $\{x \in A_f \mid x(a/b) \in A_f\}$。
这证明了拟凝聚性。

\medskip\noindent
下面证明最后一个陈述。显然 $(b) \subset I$。因此
$V(I) \subset V(b)$ 是无处稠密子集，因为 $b$ 是非零因子。
所以 $1$ 的余核支撑在某个闭无处稠密集合中。对 $s$ 的余核，
同样的论证成立，因为 $s(b) = (a) \subset sI \subset A$。
\end{proof}

\begin{definition}
\label{divisors-definition-regular-meromorphic-ideal-denominators}
设 $X$ 为概形。设 $\mathcal{L}$ 为可逆 $\mathcal{O}_X$-模。
设 $s$ 为 $\mathcal{L}$ 的正则亚纯截面。引理
\ref{divisors-lemma-regular-meromorphic-ideal-denominators}
中构造的理想层 $\mathcal{I}$ 称为 $s$ 的{\it 分母理想层}。
\end{definition}




\section{亚纯函数与亚纯截面：Noether 情形}
\label{divisors-section-meromorphic-noetherian}

\noindent
对于局部 Noether 概形，可以证明亚纯函数层的一些结果。然而，
\cite{misconceptions} 中有一个例子表明，$\mathcal{K}_X$ 对 Noether 概形
$X$ 也未必拟凝聚。

\begin{lemma}
\label{divisors-lemma-meromorphic-section-restricts-to-zero}
设 $X$ 为拟紧概形。设 $h \in \Gamma(X, \mathcal{O}_X)$ 且
$f \in \Gamma(X, \mathcal{K}_X)$，并假设 $f$ 在 $X_h$ 上的限制为零。
则 $h^n f = 0$ 对某个 $n \gg 0$ 成立。
\end{lemma}

\begin{proof}
可找到 $X$ 的一个由仿射开集 $U$ 构成的覆盖，使得
$f|_U = s^{-1}a$，其中 $a \in \mathcal{O}_X(U)$ 且
$s \in \mathcal{S}(U)$。由于 $X$ 拟紧，可以取这种形式的有限仿射开覆盖。
因此，只需在 $X = \Spec(A)$ 且 $f = s^{-1}a$ 时证明引理。
注意，$s \in A$ 是非零因子，故只需证明 $f = a$ 的情形。条件
$f|_{X_h} = 0$ 蕴含 $a$ 在
$A_h = \mathcal{O}_X(X_h)$ 中的像为零，因为
$\mathcal{O}_X \subset \mathcal{K}_X$。因此，$h^na = 0$ 对某个
$n > 0$ 成立，这正是所需结论。
\end{proof}

\begin{lemma}
\label{divisors-lemma-locally-Noetherian-K}
设 $X$ 为局部 Noether 概形。
\begin{enumerate}
\item 对任意 $x \in X$，$\mathcal{S}_x \subset \mathcal{O}_{X, x}$
是非零因子集合，因而 $\mathcal{K}_{X, x}$ 是
$\mathcal{O}_{X, x}$ 的全商环。
\item 对任意仿射开集 $U \subset X$，环
$\mathcal{K}_X(U)$ 等于 $\mathcal{O}_X(U)$ 的全商环。
\end{enumerate}
\end{lemma}

\begin{proof}
为证明本引理，可以假设 $X$ 是某个 Noether 环 $A$ 的谱。设
$x \in X$ 对应于 $\mathfrak p \subset A$。

\medskip\noindent
证明 (1)。显然，$\mathcal{S}_x$ 包含于
$\mathcal{O}_{X, x} = A_\mathfrak p$ 的非零因子集合。反过来，设
$f, g \in A$、$g \not \in \mathfrak p$，并假设
$f/g$ 是 $A_{\mathfrak p}$ 中的非零因子。令
$I = \{a \in A \mid af = 0\}$。由局部化的正合性可知
$I_{\mathfrak p} = 0$。由于 $A$ Noether，$I$ 有限生成，故
$g'I = 0$，其中存在元素 $g' \in A$，且 $g' \not \in \mathfrak p$。
因此，$f$ 在 $A_{g'}$ 中是非零因子，也就是在
$\mathfrak p$ 的某个 Zariski 开邻域中是非零因子。于是
$f/g$ 是 $\mathcal{S}_x$ 的元素。

\medskip\noindent
证明 (2)。设 $f \in \Gamma(X, \mathcal{K}_X)$ 为亚纯函数。令
$I = \{a \in A \mid af \in A\}$。固定素理想
$\mathfrak p \subset A$，它对应于点 $x \in X$。由 (1)，可以把 $f$
在 $\mathfrak p$ 处茎中的像写为 $a/b$，其中
$a, b \in A_{\mathfrak p}$，且 $b \in A_{\mathfrak p}$ 不是零因子。
写 $b = c/d$，其中 $c, d \in A$、$d \not \in \mathfrak p$。
则 $ad - cf$ 是 $\mathcal{K}_X$ 的一个截面，并在 $x$ 的某个
开邻域上消失。设它在 $D(e)$ 上消失，其中
$e \in A$、$e \not \in \mathfrak p$。由引理
\ref{divisors-lemma-meromorphic-section-restricts-to-zero}，有
$e^n(ad - cf) = 0$，其中 $n \gg 0$ 是某个整数。因此
$e^nc \in I$，且 $e^nc$ 在 $A_{\mathfrak p}$ 中的像是非零因子。令
$\text{Ass}(A) = \{\mathfrak q_1, \ldots, \mathfrak q_t\}$ 为
$A$ 的伴随素理想。考察 $IA_{\mathfrak q_i}$ 并使用《代数》引理
\ref{algebra-lemma-associated-primes-localize}，上述论证表明，
$I \not \subset \mathfrak q_i$ 对每个 $i$ 成立。由《代数》引理
\ref{algebra-lemma-silly}，存在元素 $z \in I$，使得
$z \not \in \bigcup \mathfrak q_i$。由《代数》引理
\ref{algebra-lemma-ass-zero-divisors}，可知 $z$ 不是 $A$ 上的零因子。
因此，$f = (zf)/z$ 是 $A$ 的全商环的元素。这证明了 (2)。
\end{proof}

\begin{lemma}
\label{divisors-lemma-quasi-coherent-K}
设 $X$ 为没有嵌入点的局部 Noether 概形。令 $X^0$ 为 $X$ 的
不可约分支的泛点集合。则有
$$
\mathcal{K}_X =
\bigoplus\nolimits_{\eta \in X^0} j_{\eta, *}\mathcal{O}_{X, \eta} =
\prod\nolimits_{\eta \in X^0} j_{\eta, *}\mathcal{O}_{X, \eta}
$$
其中 $j_\eta : \Spec(\mathcal{O}_{X, \eta}) \to X$ 是《概形》第
\ref{schemes-section-points} 节的典范映射。此外：
\begin{enumerate}
\item $\mathcal{K}_X$ 是拟凝聚 $\mathcal{O}_X$-代数层；
\item 对每个拟凝聚 $\mathcal{O}_X$-模 $\mathcal{F}$，层
$$
\mathcal{K}_X(\mathcal{F}) =
\bigoplus\nolimits_{\eta \in X^0} j_{\eta, *}\mathcal{F}_\eta =
\prod\nolimits_{\eta \in X^0} j_{\eta, *}\mathcal{F}_\eta
$$
即 $\mathcal{F}$ 的亚纯截面层，是拟凝聚的；并且
\item $X$ 的有理函数环就是 $X$ 的亚纯函数环；写成公式即
$R(X) = \Gamma(X, \mathcal{K}_X)$。
\end{enumerate}
\end{lemma}

\begin{proof}
本引理是引理 \ref{divisors-lemma-meromorphic-weakass-finite} 的特例，因为在
局部 Noether 情形，由引理 \ref{divisors-lemma-ass-weakly-ass} 可知，
弱伴随点与伴随点相同。
\end{proof}

\begin{lemma}
\label{divisors-lemma-regular-meromorphic-section-exists-noetherian}
设 $X$ 为没有嵌入点的局部 Noether 概形。设 $\mathcal{L}$ 为可逆
$\mathcal{O}_X$-模。则 $\mathcal{L}$ 有一个正则亚纯截面。
\end{lemma}

\begin{proof}
对每个泛点 $\eta$（作为 $X$ 的泛点），选取生成元 $s_\eta$；它生成秩为
$1$ 的自由模 $\mathcal{L}_\eta$，后者是 Artin 局部环
$\mathcal{O}_{X, \eta}$ 上的模。由引理 \ref{divisors-lemma-quasi-coherent-K} 中对
$\mathcal{K}_X$ 与 $\mathcal{K}_X(\mathcal{L})$ 的描述，立即可知
$s = \prod s_\eta$ 是 $\mathcal{L}$ 的一个正则亚纯截面。
\end{proof}

\begin{lemma}
\label{divisors-lemma-make-maps-regular-section}
假设给定：
\begin{enumerate}
\item 局部 Noether 概形 $X$；
\item 可逆模 $\mathcal{L}$（作为 $\mathcal{O}_X$-模）；
\item 正则亚纯截面 $s$（取值于 $\mathcal{L}$）；以及
\item 凝聚层 $\mathcal{F}$（定义在 $X$ 上），它没有嵌入伴随点且
$\text{Supp}(\mathcal{F}) = X$。
\end{enumerate}
设 $\mathcal{I} \subset \mathcal{O}_X$ 为 $s$ 的分母理想层。
令 $T \subset X$ 为 $\mathcal{O}_X/\mathcal{I}$ 与
$\mathcal{L}/s(\mathcal{I})$ 的支集之并；根据引理
\ref{divisors-lemma-regular-meromorphic-ideal-denominators}，它是无处稠密闭子集
$T \subset X$。则存在典范单射
$$
1 : \mathcal{I}\mathcal{F} \to \mathcal{F}, \quad
s : \mathcal{I}\mathcal{F} \to \mathcal{F} \otimes_{\mathcal{O}_X}\mathcal{L}
$$
二者的余核都支撑在 $T$ 上。
\end{lemma}

\begin{proof}
约化到仿射情形，其中 $\mathcal{L} \cong \mathcal{O}_X$，且
$s = a/b$，这里 $a, b \in A$ 都是非零因子。略去约化步骤的证明。
写 $\mathcal{F} = \widetilde{M}$。令
$I = \{x \in A \mid x(a/b) \in A\}$，于是
$\mathcal{I} = \widetilde{I}$（见引理
\ref{divisors-lemma-regular-meromorphic-ideal-denominators} 的证明）。
注意，$T = V(I) \cup V((a/b)I)$。对任意 $A$-模 $M$，考察映射
$1 : IM \to M$；它给出引理中的映射 $1$。另一方面，考察映射
$\sigma : IM \to M_b, x \mapsto ax/b$。由于 $b$ 不是 $A$ 中的零因子，
而且 $M$ 的支集是 $\Spec(A)$ 且没有嵌入素理想，可知
$b$ 也是 $M$ 上的非零因子。因此，$M \subset M_b$。
由 $I$ 的定义，有 $\sigma(IM) \subset M$；这里把两者看作 $M_b$ 的子模。
因此，得到 $A$-模映射 $s : IM \to M$（即满足
$s(z)/1 = \sigma(z)$ 在 $M_b$ 中对所有 $z \in IM$ 成立的唯一映射）。
它是单射，因为 $a$ 也是非零因子（在 $A$ 与 $M$ 上都是）。
显然，$M/IM$ 被 $I$ 零化，而 $M/s(IM)$ 被 $(a/b)I$ 零化。
引理得证。
\end{proof}




\section{亚纯函数与亚纯截面：约化情形}
\label{divisors-section-meromorphic-reduced}

\noindent
对于约化且局部只有有限多个不可约分支的概形，亚纯函数层是拟凝聚的。

\begin{lemma}
\label{divisors-lemma-reduced-finite-irreducible}
设 $X$ 为约化概形，并且每个拟紧开集都只有有限多个不可约分支。
令 $X^0$ 为 $X$ 的不可约分支的泛点集合。则有
$$
\mathcal{K}_X =
\bigoplus\nolimits_{\eta \in X^0} j_{\eta, *}\kappa(\eta) =
\prod\nolimits_{\eta \in X^0} j_{\eta, *}\kappa(\eta)
$$
其中 $j_\eta : \Spec(\kappa(\eta)) \to X$ 是《概形》第
\ref{schemes-section-points} 节的典范映射。此外：
\begin{enumerate}
\item $\mathcal{K}_X$ 是拟凝聚 $\mathcal{O}_X$-代数层；
\item 对每个拟凝聚 $\mathcal{O}_X$-模 $\mathcal{F}$，层
$$
\mathcal{K}_X(\mathcal{F}) =
\bigoplus\nolimits_{\eta \in X^0} j_{\eta, *}\mathcal{F}_\eta =
\prod\nolimits_{\eta \in X^0} j_{\eta, *}\mathcal{F}_\eta
$$
即 $\mathcal{F}$ 的亚纯截面层，是拟凝聚的；
\item $\mathcal{S}_x \subset \mathcal{O}_{X, x}$ 是非零因子集合，
对任意 $x \in X$；
\item $\mathcal{K}_{X, x}$ 是 $\mathcal{O}_{X, x}$ 的全商环，
对任意 $x \in X$；
\item $\mathcal{K}_X(U)$ 等于 $\mathcal{O}_X(U)$ 的全商环，
对任意仿射开集 $U \subset X$；
\item $X$ 的有理函数环就是 $X$ 的亚纯函数环；写成公式即
$R(X) = \Gamma(X, \mathcal{K}_X)$。
\end{enumerate}
\end{lemma}

\begin{proof}
本引理是引理 \ref{divisors-lemma-meromorphic-weakass-finite} 的特例，因为由引理
\ref{divisors-lemma-weakass-reduced}，约化概形的弱伴随点就是各不可约分支的泛点。
\end{proof}

\begin{lemma}
\label{divisors-lemma-reduced-normalization}
设 $X$ 为概形。假设 $X$ 约化，并且每个拟紧开集
$U \subset X$ 都只有有限多个不可约分支。则正规化态射
$\nu : X^\nu \to X$ 是态射
$$
\underline{\Spec}_X(\mathcal{O}') \longrightarrow X
$$
其中 $\mathcal{O}' \subset \mathcal{K}_X$ 是 $\mathcal{O}_X$
在亚纯函数层中的整闭包。
\end{lemma}

\begin{proof}
将正规化态射 $\nu : X^\nu \to X$ 的定义（见《态射》定义
\ref{morphisms-definition-normalization}）与上面引理
\ref{divisors-lemma-reduced-finite-irreducible} 中对 $\mathcal{K}_X$ 的描述比较即可。
\end{proof}

\begin{lemma}
\label{divisors-lemma-meromorphic-functions-integral-scheme}
设 $X$ 为整概形，泛点为 $\eta$。则：
\begin{enumerate}
\item 亚纯函数层同构于取值为 $X$ 的函数域（见《态射》定义
\ref{morphisms-definition-function-field}）的常值层。
\item 对任意拟凝聚层 $\mathcal{F}$（定义在 $X$ 上），层
$\mathcal{K}_X(\mathcal{F})$ 同构于取值为 $\mathcal{F}_\eta$ 的常值层。
\end{enumerate}
\end{lemma}

\begin{proof}
略。
\end{proof}

\noindent
在某些情形，可以证明正则亚纯截面存在。

\begin{lemma}
\label{divisors-lemma-regular-meromorphic-section-exists}
设 $X$ 为概形。设 $\mathcal{L}$ 为可逆 $\mathcal{O}_X$-模。
在下列每种情形，$\mathcal{L}$ 都有一个正则亚纯截面：
\begin{enumerate}
\item $X$ 整；
\item $X$ 约化，且每个拟紧开集都只有有限多个不可约分支；
\item $X$ 局部 Noether 且没有嵌入点。
\end{enumerate}
\end{lemma}

\begin{proof}
在情形 (1)，令 $\eta \in X$ 为泛点。由引理
\ref{divisors-lemma-meromorphic-functions-integral-scheme} 已知，
$\mathcal{K}_X$ 与 $\mathcal{K}_X(\mathcal{L})$ 分别是取值为
$\kappa(\eta)$ 与 $\mathcal{L}_\eta$ 的常值层。由于
$\dim_{\kappa(\eta)} \mathcal{L}_\eta = 1$，可以选取非零元
$s \in \mathcal{L}_\eta$。显然，$s$ 是 $\mathcal{L}$ 的正则亚纯截面。
在情形 (2)，选取非零元 $s_\eta \in \mathcal{L}_\eta$；这里
$\eta$ 遍历 $X$ 的所有泛点。这是可能的，因为 $\mathcal{L}_\eta$ 是
$1$ 维 $\kappa(\eta)$-向量空间。由引理
\ref{divisors-lemma-reduced-finite-irreducible} 中对
$\mathcal{K}_X$ 与 $\mathcal{K}_X(\mathcal{L})$ 的描述，立即可知
$s = \prod s_\eta$ 是 $\mathcal{L}$ 的正则亚纯截面。情形 (3) 是引理
\ref{divisors-lemma-regular-meromorphic-section-exists-noetherian}。
\end{proof}




\section{Weil 除子}
\label{divisors-section-Weil-divisors}

\noindent
下面对局部 Noether 整概形引入 Weil 除子以及 Weil 除子的有理等价。
由于不假设概形拟紧，在定义 Weil 除子时必须稍加谨慎。必须允许
素除子的无穷和，例如，有理函数可能有无穷多个极点。对于拟紧概形，
这里的 Weil 除子仍像通常一样是有限和。下面这个基本引理常用于证明
闭子概形族局部有限。

\begin{lemma}
\label{divisors-lemma-components-locally-finite}
设 $X$ 为局部 Noether 概形。设 $Z \subset X$ 为闭子概形。
$Z$ 的不可约分支族在 $X$ 中局部有限。
\end{lemma}

\begin{proof}
设 $U \subset X$ 为拟紧开子概形。则 $U$ 是 Noether 概形，因而其
底拓扑空间 Noether（《性质》引理
\ref{properties-lemma-Noetherian-topology}）。因此，每个子空间都
Noether，并且只有有限多个不可约分支（见《拓扑》引理
\ref{topology-lemma-Noetherian}）。
\end{proof}

\noindent
回忆，若 $Z$ 是概形 $X$ 的不可约闭子集，则 $Z$ 在 $X$ 中的余维
等于局部环 $\mathcal{O}_{X, \xi}$ 的维数，其中 $\xi \in Z$ 是泛点。
见《性质》引理 \ref{properties-lemma-codimension-local-ring}。

\begin{definition}
\label{divisors-definition-Weil-divisor}
设 $X$ 为局部 Noether 整概形。
\begin{enumerate}
\item {\it 素除子}是整闭子概形 $Z \subset X$，其余维为 $1$。
\item {\it Weil 除子}是形式和 $D = \sum n_Z Z$，其中求和遍历
$X$ 的素除子，并且族 $\{Z \mid n_Z \not = 0\}$ 局部有限
（《拓扑》定义 \ref{topology-definition-locally-finite}）。
\end{enumerate}
$X$ 上所有 Weil 除子组成的群记作 $\text{Div}(X)$。
\end{definition}

\noindent
下一步是定义与有理函数相伴的 Weil 除子。为此，使用有理函数沿
素除子的消失阶，其定义如下。

\begin{definition}
\label{divisors-definition-order-vanishing}
设 $X$ 为局部 Noether 整概形。设 $f \in R(X)^*$。
对每个素除子 $Z \subset X$，把{\it $f$ 沿 $Z$ 的消失阶}定义为整数
$$
\text{ord}_Z(f) = \text{ord}_{\mathcal{O}_{X, \xi}}(f)
$$
其中右端是《代数》定义 \ref{algebra-definition-ord} 中的概念，
而 $\xi$ 是 $Z$ 的泛点。
\end{definition}

\noindent
注意，对 $f, g \in R(X)^*$，有
$$
\text{ord}_Z(fg) = \text{ord}_Z(f) + \text{ord}_Z(g).
$$
当然可能出现 $\text{ord}_Z(f) < 0$。在这种情形，称 $f$ 沿 $Z$
有一个{\it 极点}，并称 $-\text{ord}_Z(f) > 0$ 为
{\it $f$ 沿 $Z$ 的极点阶}。务必注意，条件
$\text{ord}_Z(f) \geq 0$ 与条件 $f \in \mathcal{O}_{X, \xi}$
并不等价，除非局部环 $\mathcal{O}_{X, \xi}$ 是离散赋值环。

\begin{lemma}
\label{divisors-lemma-divisor-locally-finite}
设 $X$ 为局部 Noether 整概形。设 $f \in R(X)^*$。则族
$$
\{Z \subset X \mid Z\text{ a prime divisor with generic point }\xi
\text{ 且 }f\text{ not in }\mathcal{O}_{X, \xi}\}
$$
与族
$$
\{Z \subset X \mid Z \text{ a prime divisor and }\text{ord}_Z(f) \not = 0\}
$$
在 $X$ 中局部有限。
\end{lemma}

\begin{proof}
存在非空开子概形 $U \subset X$，使得 $f$ 对应于
$\Gamma(U, \mathcal{O}_X^*)$ 的一个截面。因此，引理中两个集合里
可能出现的素除子全都是 $X \setminus U$ 的不可约分支。
于是引理 \ref{divisors-lemma-components-locally-finite} 给出所需结论。
\end{proof}

\noindent
该引理使下述定义有意义。

\begin{definition}
\label{divisors-definition-principal-divisor}
设 $X$ 为局部 Noether 整概形。设 $f \in R(X)^*$。
与 $f$ 相伴的{\it 主 Weil 除子}是 Weil 除子
$$
\text{div}(f) = \text{div}_X(f) = \sum \text{ord}_Z(f) [Z]
$$
其中求和遍历素除子，$\text{ord}_Z(f)$ 如定义
\ref{divisors-definition-order-vanishing} 所述。由引理
\ref{divisors-lemma-divisor-locally-finite}，这是有意义的。
\end{definition}

\begin{lemma}
\label{divisors-lemma-div-additive}
设 $X$ 为局部 Noether 整概形。设 $f, g \in R(X)^*$。则
$$
\text{div}_X(fg) = \text{div}_X(f) + \text{div}_X(g)
$$
作为 $X$ 上的 Weil 除子成立。
\end{lemma}

\begin{proof}
这由 $\text{ord}$ 函数的可加性立即得出。
\end{proof}

\noindent
由上引理可知，主 Weil 除子族构成全体 Weil 除子群的一个子群。
这引出下述定义。

\begin{definition}
\label{divisors-definition-class-group}
设 $X$ 为局部 Noether 整概形。$X$ 的{\it Weil 除子类群}是
Weil 除子群对主 Weil 除子子群的商。记作 $\text{Cl}(X)$。
\end{definition}

\noindent
由构造得到正合复形
\begin{equation}
\label{divisors-equation-Weil-divisor-class}
R(X)^* \xrightarrow{\text{div}} \text{Div}(X) \to \text{Cl}(X) \to 0
\end{equation}
可以把它看作 $\text{Cl}(X)$ 的一个呈示。下一步将把 Weil 除子类群
与 Picard 群联系起来。




\section{与可逆模相伴的 Weil 除子类}
\label{divisors-section-c1}

\noindent
本节将完全沿用第 \ref{divisors-section-Weil-divisors} 节的步骤，为局部 Noether
整概形定义典范映射
$\Pic(X) \to \text{Cl}(X)$。

\medskip\noindent
设 $X$ 为概形。设 $\mathcal{L}$ 为可逆 $\mathcal{O}_X$-模。
设 $\xi \in X$ 为一点。若
$s_\xi, s'_\xi \in \mathcal{L}_\xi$ 生成 $\mathcal{L}_\xi$（这里把它视为
$\mathcal{O}_{X, \xi}$-模），则存在单位 $u \in \mathcal{O}_{X, \xi}^*$，使得
$s_\xi = u s'_\xi$。亚纯截面层 $\mathcal{K}_X(\mathcal{L})$
（即 $\mathcal{L}$ 的亚纯截面层）在 $x$ 处的茎等于
$\mathcal{K}_{X, x} \otimes_{\mathcal{O}_{X, x}} \mathcal{L}_x$。
因此，任意亚纯截面 $s$（取值于 $\mathcal{L}$）在 $x$ 处茎中的像都可写成
$s = fs_\xi$，其中 $f \in \mathcal{K}_{X, x}$。下文将把这简写为
$f = s/s_\xi$。此外，若 $X$ 整，则
$\mathcal{K}_{X, x} = R(X)$ 等于 $X$ 的函数域，所以
$s/s_\xi \in R(X)$。若 $s$ 是正则亚纯截面，则事实上
$s/s_\xi \in R(X)^*$。在整概形上，正则亚纯截面恰好就是非零亚纯截面。
最后可知，$s/s_\xi$ 不依赖于 $s_\xi$ 的选择，至多相差局部环
$\mathcal{O}_{X, x}$ 中一个单位的倍数。综合以上内容，下述定义有意义。

\begin{definition}
\label{divisors-definition-order-vanishing-meromorphic}
设 $X$ 为局部 Noether 整概形。设 $\mathcal{L}$ 为可逆
$\mathcal{O}_X$-模。设
$s \in \Gamma(X, \mathcal{K}_X(\mathcal{L}))$ 为 $\mathcal{L}$ 的
正则亚纯截面。对每个素除子 $Z \subset X$，把
{\it $s$ 沿 $Z$ 的消失阶}定义为整数
$$
\text{ord}_{Z, \mathcal{L}}(s)
= \text{ord}_{\mathcal{O}_{X, \xi}}(s/s_\xi)
$$
其中右端是《代数》定义 \ref{algebra-definition-ord} 中的概念，
$\xi \in Z$ 是泛点，而
$s_\xi \in \mathcal{L}_\xi$ 是一个生成元。
\end{definition}

\noindent
与主除子的情形一样，有下述引理。

\begin{lemma}
\label{divisors-lemma-divisor-meromorphic-locally-finite}
设 $X$ 为局部 Noether 整概形。设 $\mathcal{L}$ 为可逆
$\mathcal{O}_X$-模。设 $s \in \mathcal{K}_X(\mathcal{L})$ 为
$\mathcal{L}$ 的正则（即非零）亚纯截面。则集合
$$
\{Z \subset X \mid Z \text{ a prime divisor with generic point }\xi
\text{ 且 }s\text{ not in }\mathcal{L}_\xi\}
$$
与集合
$$
\{Z \subset X \mid Z \text{ is a prime divisor and }
\text{ord}_{Z, \mathcal{L}}(s) \not = 0\}
$$
在 $X$ 中局部有限。
\end{lemma}

\begin{proof}
存在非空开子概形 $U \subset X$，使得 $s$ 对应于
$\Gamma(U, \mathcal{L})$ 的一个截面，并且该截面生成
$\mathcal{L}$ 在 $U$ 上的限制。因此，引理中的集合里可能出现的素除子全都是
$X \setminus U$ 的不可约分支。于是引理
\ref{divisors-lemma-components-locally-finite} 给出所需结论。
\end{proof}

\begin{lemma}
\label{divisors-lemma-divisor-meromorphic-well-defined}
设 $X$ 为局部 Noether 整概形。设 $\mathcal{L}$ 为可逆
$\mathcal{O}_X$-模。设 $s, s' \in \mathcal{K}_X(\mathcal{L})$ 为
$\mathcal{L}$ 的非零亚纯截面。则 $f = s/s'$ 是 $R(X)^*$ 的元素，并且
$$
\sum \text{ord}_{Z, \mathcal{L}}(s)[Z]
=
\sum \text{ord}_{Z, \mathcal{L}}(s')[Z]
+
\text{div}(f)
$$
作为 Weil 除子成立。
\end{lemma}

\begin{proof}
这由定义立即得出。注意，引理
\ref{divisors-lemma-divisor-meromorphic-locally-finite}
保证这些和确实是 Weil 除子。
\end{proof}

\begin{definition}
\label{divisors-definition-divisor-invertible-sheaf}
设 $X$ 为局部 Noether 整概形。设 $\mathcal{L}$ 为可逆
$\mathcal{O}_X$-模。
\begin{enumerate}
\item 对任意非零亚纯截面 $s$（取值于 $\mathcal{L}$），把
{\it 与 $s$ 相伴的 Weil 除子}定义为
$$
\text{div}_\mathcal{L}(s) =
\sum \text{ord}_{Z, \mathcal{L}}(s) [Z] \in \text{Div}(X)
$$
其中求和遍历素除子。
\item 把{\it 与 $\mathcal{L}$ 相伴的 Weil 除子类}定义为
$\text{div}_\mathcal{L}(s)$ 在 $\text{Cl}(X)$ 中的像，其中
$s$ 是 $\mathcal{L}$ 在 $X$ 上的任意非零亚纯截面。由引理
\ref{divisors-lemma-divisor-meromorphic-well-defined}，这是良定义的。
\end{enumerate}
\end{definition}

\noindent
正如所料，该构造对可逆模可加。

\begin{lemma}
\label{divisors-lemma-c1-additive}
设 $X$ 为局部 Noether 整概形。设 $\mathcal{L}$、$\mathcal{N}$ 为可逆
$\mathcal{O}_X$-模。设 $s$ 与 $t$ 分别为
$\mathcal{L}$ 与 $\mathcal{N}$ 的非零亚纯截面。则 $st$ 是
$\mathcal{L} \otimes \mathcal{N}$ 的非零亚纯截面，并且
$$
\text{div}_{\mathcal{L} \otimes \mathcal{N}}(st)
=
\text{div}_\mathcal{L}(s) + \text{div}_\mathcal{N}(t)
$$
在 $\text{Div}(X)$ 中成立。特别地，
$\mathcal{L} \otimes_{\mathcal{O}_X} \mathcal{N}$ 的 Weil 除子类等于
$\mathcal{L}$ 与 $\mathcal{N}$ 的 Weil 除子类之和。
\end{lemma}

\begin{proof}
设 $s$ 与 $t$ 分别为 $\mathcal{L}$ 与 $\mathcal{N}$ 的非零亚纯截面。
则 $st$ 是 $\mathcal{L} \otimes \mathcal{N}$ 的非零亚纯截面。
设 $Z \subset X$ 为素除子。设 $\xi \in Z$ 为其泛点。选取生成元
$s_\xi \in \mathcal{L}_\xi$ 与 $t_\xi \in \mathcal{N}_\xi$。则
$s_\xi t_\xi$ 是 $(\mathcal{L} \otimes \mathcal{N})_\xi$ 的一个生成元。
所以 $st/(s_\xi t_\xi) = (s/s_\xi)(t/t_\xi)$。因而
$$
\text{ord}_{Z, \mathcal{L} \otimes \mathcal{N}}(st)
=
\text{ord}_{Z, \mathcal{L}}(s) + \text{ord}_{Z, \mathcal{N}}(t)
$$
这由 $\text{ord}_Z$ 函数的可加性得出。
\end{proof}

\noindent
这样得到交换群的同态
\begin{equation}
\label{divisors-equation-c1}
\Pic(X) \longrightarrow \text{Cl}(X)
\end{equation}
它把可逆模映到其 Weil 除子类。

\begin{lemma}
\label{divisors-lemma-normal-c1-injective}
设 $X$ 为局部 Noether 整概形。若 $X$ 正规，则映射
(\ref{divisors-equation-c1}) $\Pic(X) \to \text{Cl}(X)$ 是单射。
\end{lemma}

\begin{proof}
设 $\mathcal{L}$ 为可逆 $\mathcal{O}_X$-模，并且其相伴的 Weil 除子类
平凡。设 $s$ 为 $\mathcal{L}$ 的正则亚纯截面。该假设意味着
$\text{div}_\mathcal{L}(s) = \text{div}(f)$，其中
$f \in R(X)^*$。于是
$t = f^{-1}s$ 是 $\mathcal{L}$ 的正则亚纯截面，且
$\text{div}_\mathcal{L}(t) = 0$，见引理
\ref{divisors-lemma-divisor-meromorphic-well-defined}。下面证明 $t$ 给出
$\mathcal{L}$ 的平凡化，从而完成引理的证明。
为此，可以在 $X$ 上局部工作。因此，可以假设
$X = \Spec(A)$ 仿射且 $\mathcal{L}$ 平凡。于是 $A$ 是 Noether
正规整环，而 $t$ 是其分式域中的元素，并且满足
$\text{ord}_{A_\mathfrak p}(t) = 0$；所取素理想的高度为 $1$，记作
$\mathfrak p$，且它是 $A$ 的素理想。目标是证明 $t$ 是 $A$ 的单位。
由于 $A_\mathfrak p$ 对 $A$ 的高度一素理想而言是离散赋值环
（《代数》引理 \ref{algebra-lemma-criterion-normal}），所以该条件表示
$t \in A_\mathfrak p^*$ 对所有素理想 $\mathfrak p$（高度为 $1$）均成立。
由《代数》引理
\ref{algebra-lemma-normal-domain-intersection-localizations-height-1}，
这蕴含 $t \in A$ 且 $t^{-1} \in A$。证明完成。
\end{proof}

\begin{lemma}
\label{divisors-lemma-local-rings-UFD-c1-bijective}
设 $X$ 为局部 Noether 整概形。考察映射
(\ref{divisors-equation-c1}) $\Pic(X) \to \text{Cl}(X)$。下列条件等价：
\begin{enumerate}
\item $X$ 的局部环都是 UFD；并且
\item $X$ 正规，且 $\Pic(X) \to \text{Cl}(X)$ 为满射。
\end{enumerate}
在这种情形，$\Pic(X) \to \text{Cl}(X)$ 是同构。
\end{lemma}

\begin{proof}
若 (1) 成立，则由《代数》引理 \ref{algebra-lemma-UFD-normal}，
$X$ 正规。因此，映射 (\ref{divisors-equation-c1}) 由引理
\ref{divisors-lemma-normal-c1-injective} 可知为单射。此外，每个素除子 $D \subset X$ 都由引理
\ref{divisors-lemma-weil-divisor-is-cartier-UFD} 成为有效 Cartier 除子。在此情形，
典范截面 $1_D$（它是 $\mathcal{O}_X(D)$ 的截面）
（定义 \ref{divisors-definition-invertible-sheaf-effective-Cartier-divisor}）
恰好沿 $D$ 消失，可知 $D$ 的类是 $\mathcal{O}_X(D)$ 在映射
(\ref{divisors-equation-c1}) 下的像。因此，该映射也为满射。

\medskip\noindent
假设 (2) 成立。选取素除子 $D \subset X$。由于
(\ref{divisors-equation-c1}) 为满射，存在可逆层 $\mathcal{L}$、正则亚纯截面
$s$ 以及 $f \in R(X)^*$，使得
$\text{div}_\mathcal{L}(s) + \text{div}(f) = [D]$。换言之，
$\text{div}_\mathcal{L}(fs) = [D]$。设 $x \in X$，并令
$A = \mathcal{O}_{X, x}$。则 $A$ 是 Noether 局部正规整环，分式域为
$K = R(X)$。每个高度为 $1$ 的素理想（在 $A$ 中）都对应于 $X$ 上的素除子，
而每个可逆 $\mathcal{O}_X$-模限制到 $\Spec(A)$ 后都是平凡可逆模。
由此，对每个高度为 $1$ 的素理想 $\mathfrak p \subset A$，存在元素
$f \in K$，使得 $\text{ord}_{A_\mathfrak p}(f) = 1$，且
$\text{ord}_{A_{\mathfrak p'}}(f) = 0$ 对每个其他高度一素理想
$\mathfrak p'$ 成立。于是由《代数》引理
\ref{algebra-lemma-normal-domain-intersection-localizations-height-1}，
$f \in A$。同理可知，每个元素 $g \in \mathfrak p$ 都可写成
$g = af$，其中 $a \in A$。因此，$A$ 的每个高度一素理想都是主理想，
由《代数》引理 \ref{algebra-lemma-characterize-UFD-height-1}，$A$ 是 UFD。
\end{proof}




\section{关于可逆模的进一步结果}
\label{divisors-section-invertible-modules}

\noindent
本节讨论可逆模的一些性质。

\begin{lemma}
\label{divisors-lemma-in-image-pullback}
设 $\varphi : X \to Y$ 为概形态射。设 $\mathcal{L}$ 为可逆
$\mathcal{O}_X$-模。假设：
\begin{enumerate}
\item $X$ 局部 Noether；
\item $Y$ 局部 Noether、整且正规；
\item $\varphi$ 平坦，且其纤维整（因而非空）；
\item $\varphi$ 或拟紧，或局部有限型；
\item $\mathcal{L}$ 限制到 $\varphi$ 的泛纤维后平凡。
\end{enumerate}
则 $\mathcal{L} \cong \varphi^*\mathcal{N}$，其中右端由某个可逆
$\mathcal{O}_Y$-模 $\mathcal{N}$ 给出。
\end{lemma}

\begin{proof}
设 $\xi \in Y$ 为泛点。设 $X_\xi$ 为 $\varphi$ 在 $\xi$ 上的概形论纤维。
用 $\mathcal{L}_\xi$ 表示 $\mathcal{L}$ 到 $X_\xi$ 的拉回。假设 (5) 意味着
$\mathcal{L}_\xi$ 平凡。选取平凡化截面
$s \in \Gamma(X_\xi, \mathcal{L}_\xi)$。由引理
\ref{divisors-lemma-flat-over-integral-integral-fibre}，$X$ 是整概形。因此可以把
$s$ 看作 $\mathcal{L}$ 的正则亚纯截面。由引理
\ref{divisors-lemma-pullback-meromorphic-sections-defined}，亚纯函数沿
$\varphi$ 的拉回有定义。令 $\mathcal{N} \subset \mathcal{K}_Y$ 为如下
$\mathcal{O}_Y$-模：其在开集 $V \subset Y$ 上的截面，是那些亚纯函数
$g \in \mathcal{K}_Y(V)$，它们满足
$\varphi^*(g)s \in \mathcal{L}(\varphi^{-1}V)$。先验地，$\varphi^*(g)s$ 是
$\mathcal{K}_X(\mathcal{L})$ 在 $\varphi^{-1}V$ 上的截面。我们断言
$\mathcal{N}$ 是可逆 $\mathcal{O}_Y$-模，并且映射
$$
\varphi^*\mathcal{N} \longrightarrow \mathcal{L},\quad
g \longmapsto gs
$$
是同构。

\medskip\noindent
先在下述情形证明该断言：$X$ 与 $Y$ 均仿射，且 $\mathcal{L}$ 平凡。
记 $Y = \Spec(R)$、$X = \Spec(A)$，而 $s$ 由元素
$s \in A \otimes_R K$ 给出，其中 $K$ 是 $R$ 的分式域。可写
$s = a/r$，其中 $r \in R$ 非零且 $a \in A$。由于 $s$ 在泛纤维上生成
$\mathcal{L}$，存在 $s' \in A \otimes_R K$，使得 $ss' = 1$。所以
$s = r'/a'$，其中 $r' \in R$ 非零且 $a' \in A$。设
$\mathfrak p_1, \ldots, \mathfrak p_n \subset R$ 为包含 $rr'$ 的极小素理想。
每个 $R_{\mathfrak p_i}$ 都是离散赋值环（《代数》引理
\ref{algebra-lemma-minimal-over-1} 与
\ref{algebra-lemma-criterion-normal}）。由假设，
$\mathfrak q_i = \mathfrak p_i A$ 是素理想。因此
$\mathfrak q_i A_{\mathfrak q_i}$ 由单个元素生成，进而
$A_{\mathfrak q_i}$ 也是离散赋值环（《代数》引理
\ref{algebra-lemma-characterize-dvr}）。当然，
$R_{\mathfrak p_i} \to A_{\mathfrak q_i}$ 的分歧指数为 $1$。设
$e_i, e'_i \geq 0$ 分别为 $a, a'$ 在 $A_{\mathfrak q_i}$ 中的赋值。
则 $e_i + e'_i$ 是 $rr'$ 在 $R_{\mathfrak p_i}$ 中的赋值。注意，由
《代数》引理
\ref{algebra-lemma-normal-domain-intersection-localizations-height-1}，在
$R$ 中有
$$
\mathfrak p_1^{(e_1 + e'_1)} \cap \ldots \cap \mathfrak p_i^{(e_n + e'_n)} =
(rr')
$$
令
$$
I = \mathfrak p_1^{(e_1)} \cap \ldots \cap \mathfrak p_i^{(e_n)}
\quad\text{且}\quad
I' = \mathfrak p_1^{(e'_1)} \cap \ldots \cap \mathfrak p_i^{(e'_n)}
$$
于是 $II' \subset (rr')$。注意，由《代数》引理
\ref{algebra-lemma-symbolic-power-flat-extension} 与
\ref{algebra-lemma-flat-intersect-ideals}，有
$$
IA =
(\mathfrak p_1^{(e_1)} \cap \ldots \cap \mathfrak p_i^{(e_n)})A =
(\mathfrak p_1A)^{(e_1)} \cap \ldots \cap (\mathfrak p_i A)^{(e_n)}
$$
$I'A$ 也同样如此。因此 $a \in IA$ 且 $a' \in I'A$。由此，
$IA \otimes_A I'A \to rr'A$ 为满射。利用 $R \to A$ 的忠实平坦性，
$I \otimes_R I' \to (rr')$ 也为满射。于是 $II' = (rr')$，并且
$I$ 与 $I'$ 都是秩为 $1$ 的有限局部自由模，见《代数》引理
\ref{algebra-lemma-product-ideals-principal}。因此在 $R$ 上 Zariski 局部地，
可写成 $I = (g)$ 与 $I' = (g')$，其中 $gg' = rr'$。于是
$a = ug$ 且 $a' = u'g'$，其中 $u, u' \in A$。由此可知 $u, u'$ 都是单位。
所以在 $R$ 上 Zariski 局部地，有 $s = ug/r$，该情形的断言得证。

\medskip\noindent
设 $y \in Y$ 为一点。选取 $x \in X$，使它映到 $y$。可以把上一段的结果
应用于 $\Spec(\mathcal{O}_{X, x}) \to \Spec(\mathcal{O}_{Y, y})$。
由此存在元素 $g \in R(Y)^*$；它在相差
$\mathcal{O}_{Y, y}^*$ 中元素之乘法的意义下是良定义的，并且
$\varphi^*(g)s$ 生成 $\mathcal{L}_x$。因而 $\varphi^*(g)s$ 生成
$\mathcal{L}$；这在某个邻域 $U$ 上成立，且 $x$ 属于该邻域。设 $x'$ 是位于
$y$ 上方的另一个点，并设 $g' \in R(Y)^*$ 使得 $\varphi^*(g')s$ 生成
$\mathcal{L}$；这在某个开邻域 $U'$ 上成立，且 $x'$ 属于该邻域。由于纤维不可约，可以选取一点
$x''$，它位于 $U \cap U' \cap \varphi^{-1}(\{y\})$ 中。由环映射
$\mathcal{O}_{Y, y} \to \mathcal{O}_{X, x''}$ 所对应的唯一性，可知
$g$ 与 $g'$ 相差 $\mathcal{O}_{Y, y}^*$ 中元素的乘法。因此可知
$\varphi^*(g)s$ 生成 $\mathcal{L}$；这在 $\varphi^{-1}(y)$ 的某个开邻域上成立。
令 $Z \subset X$ 为如下点的集合：这些点
$z \in X$ 满足 $\varphi^*(g)s$ 不生成 $\mathcal{L}_z$。以上论证表明
$Z$ 是闭集，且有 $Z = \varphi^{-1}(T)$，其中 $T \subset Y$ 是某个子集，并且
$y \not \in T$。若能证明 $T$ 为闭集，则 $g$ 将生成 $\mathcal{N}$（作为
$\mathcal{O}_Y$-模）；这在开邻域 $Y \setminus T$ 上成立，而该邻域包含 $y$，从而完成证明
（略去一些细节）。

\medskip\noindent
若 $\varphi$ 拟紧，则由《态射》引理
\ref{morphisms-lemma-fpqc-quotient-topology}，$T$ 为闭集。若
$\varphi$ 局部有限型，则由《态射》引理
\ref{morphisms-lemma-fppf-open}，$\varphi$ 是开映射。于是
$Y \setminus T$ 作为开集 $X \setminus Z$ 的像是开集。
\end{proof}

\begin{lemma}
\label{divisors-lemma-closure-effective-cartier-divisor}
设 $X$ 为局部 Noether 概形。设 $U \subset X$ 为开集，并设
$D \subset U$ 为有效 Cartier 除子。若 $\mathcal{O}_{X, x}$ 对所有
$x \in X \setminus U$ 都是 UFD，则存在有效 Cartier 除子
$D' \subset X$，使得 $D = U \cap D'$。
\end{lemma}

\begin{proof}
设 $D' \subset X$ 为态射 $D \to X$ 的概形论像。由于 $X$ 局部 Noether，
态射 $D \to X$ 拟紧，见《性质》引理
\ref{properties-lemma-immersion-into-noetherian}。因此由《态射》引理
\ref{morphisms-lemma-quasi-compact-scheme-theoretic-image}，$D'$ 的形成与
取 $X$ 中的开集可交换。故可假设 $X = \Spec(A)$ 仿射。设
$I \subset A$ 为与 $D'$ 对应的理想。设素理想 $\mathfrak p \subset A$
对应于 $X \setminus U$ 中的一点。为完成证明，只需证明
$I_\mathfrak p$ 由单个元素生成，见引理
\ref{divisors-lemma-effective-Cartier-in-points}。因此由《态射》引理
\ref{morphisms-lemma-flat-base-change-scheme-theoretic-image}，可将 $X$ 替换为
$\Spec(A_\mathfrak p)$。换言之，可以假设 $X$ 是局部 UFD $A$ 的谱。
于是 $A$ 的所有局部环都是 UFD。由引理
\ref{divisors-lemma-effective-Cartier-divisor-is-a-sum}，有
$D = \sum a_i D_i$，其中 $D_i \subset U$ 是整的有效 Cartier 除子。
泛点 $\xi_i$（属于 $D_i$）对应于素理想
$\mathfrak p_i \subset A$，其高度为 $1$，见引理
\ref{divisors-lemma-effective-Cartier-codimension-1}。于是有
$\mathfrak p_i = (f_i)$，其中 $f_i \in A$ 是某个素元；从而 $D'$ 由
$\prod f_i^{a_i}$ 截出，正如所需。
\end{proof}

\begin{lemma}
\label{divisors-lemma-extend-invertible-module}
设 $X$ 为局部 Noether 概形。设 $U \subset X$ 为开集，并设
$\mathcal{L}$ 为可逆 $\mathcal{O}_U$-模。若 $\mathcal{O}_{X, x}$ 对所有
$x \in X \setminus U$ 都是 UFD，则存在可逆
$\mathcal{O}_X$-模 $\mathcal{L}'$，使得
$\mathcal{L} \cong \mathcal{L}'|_U$。
\end{lemma}

\begin{proof}
选取 $x \in X$，且 $x \not \in U$。我们将证明存在仿射开邻域
$W \subset X$，使得 $\mathcal{L}|_{W \cap U}$ 可延拓为 $W$ 上的可逆层。
由层的黏合（《层》节 \ref{sheaves-section-glueing-sheaves}），这蕴含可把
$\mathcal{L}$ 延拓到严格更大的开集 $U \cup W$。令
$W = \Spec(A)$ 为一个仿射开邻域。由于 $U \cap W$ 拟仿射，由引理
\ref{divisors-lemma-quasi-projective-Noetherian-pic-effective-Cartier}，可将
$\mathcal{L}|_{W \cap U}$ 写成
$\mathcal{O}(D_1) \otimes \mathcal{O}(D_2)^{\otimes -1}$，其中
$D_1, D_2 \subset W \cap U$ 是某些有效 Cartier 除子。再由引理
\ref{divisors-lemma-closure-effective-cartier-divisor}，$D_1$ 与 $D_2$ 可延拓为
$W$ 上的有效 Cartier 除子，这就给出可逆层的延拓。

\medskip\noindent
若 $X$ Noether（这是实践中最常用的情形），上述结论结合 Noether 归纳法
即完成证明。一般情形如下处理。首先，由于 $U$ 外每一点的局部环都是整环，
且 $X$ 局部 Noether，可知 $U$ 在 $X$ 中的闭包是开集。因此可以假设
$U \subset X$ 稠密且概形论稠密。现在考虑集合 $T$，其元素为三元组
$(U', \mathcal{L}', \alpha)$，其中
$U \subset U' \subset X$ 是开子概形，$\mathcal{L}'$ 是可逆
$\mathcal{O}_{U'}$-模，而
$\alpha : \mathcal{L}'|_U \to \mathcal{L}$ 是同构。在 $T$ 上赋予偏序
$\leq$，定义规则为：
$(U', \mathcal{L}', \alpha) \leq (U'', \mathcal{L}'', \alpha')$
当且仅当 $U' \subset U''$，且存在同构
$\beta : \mathcal{L}''|_{U'} \to \mathcal{L}'$，它与 $\alpha$ 和 $\alpha'$ 相容。
注意，$\beta$（若存在）是唯一的，因为 $U \subset X$ 稠密。证明的第一部分
表明，对任意元素 $t = (U', \mathcal{L}', \alpha)$（属于 $T$），若
$U' \not = X$，则存在 $t' \in T$ 使得 $t' > t$。因此为完成证明，只需说明
Zorn 引理适用。为此，考虑全序子集 $I \subset T$。若 $i \in I$ 对应于
三元组 $(U_i, \mathcal{L}_i, \alpha_i)$，则可如下构造可逆模
$\mathcal{L}'$，它定义在 $U' = \bigcup U_i$ 上。对开且拟紧的
$W \subset U'$，可知 $W \subset U_i$ 对某个 $i$ 成立，并置
$$
\mathcal{L}'(W) = \mathcal{L}_i(W)
$$
转移映射使用各个 $\beta$（它们唯一，因而复合相容）。这就定义了可逆
$\mathcal{O}$-模 $\mathcal{L}'$，定义所在的基是 $U'$ 的拟紧开集；这足以定义一个可逆模
（《层》节 \ref{sheaves-section-bases}）。细节从略。
\end{proof}

\begin{lemma}
\label{divisors-lemma-open-subscheme-UFD}
设 $R$ 为 UFD。下列概形的 Picard 群平凡：
\begin{enumerate}
\item $\Spec(R)$ 及其任意开子概形；
\item $\mathbf{A}^n_R = \Spec(R[x_1, \ldots, x_n])$ 及其任意开子概形。
\end{enumerate}
特别地，仿射 $n$-空间 $\mathbf{A}^n_k$ 在域 $k$ 上的任意开子概形之
Picard 群平凡。
\end{lemma}

\begin{proof}
由《代数》引理 \ref{algebra-lemma-polynomial-ring-UFD}，$R$ 上的任意多项式环
都是 UFD。因此，问题归结为证明 $\Pic(U)$ 平凡，其中
$U \subset \Spec(R)$ 是（非空）开集。

\medskip\noindent
在 Noether 情形，可用引理 \ref{divisors-lemma-extend-invertible-module} 将问题约化为证明
$\Pic(\Spec(R))$ 平凡；而后者已在《代数续篇》引理
\ref{more-algebra-lemma-UFD-Pic-trivial} 中证明。我们建议读者跳过证明的其余部分。

\medskip\noindent
设 $U \subset \Spec(R)$ 为非空开集。选取非零 $f \in R$，使得
$D(f) \subset U$。由于 $f$ 是 $R$ 中素元的乘积，可知只有有限多个素元
$a_1, \ldots, a_n \in R$ 满足 $(a_i) \not \in U$（或等价地，
$U \cap V(a_i) = \emptyset$）。将 $R$ 替换为
$R_{a_1 \ldots a_n}$ 后，可假设 $U$ 包含所有主素理想。（注意，UFD 的局部化
仍是 UFD，例如由《代数》引理
\ref{algebra-lemma-invert-prime-elements} 可知。）

\medskip\noindent
假设 $U$ 包含所有主素理想。写
$U = \Spec(R) \setminus V(I)$，其中 $I \subset R$ 是某个理想。选取非零 $f \in I$，并把
$f = a_1 \ldots a_n$ 写成素元的乘积。由于
$(a_i) \not \in V(I)$，由素理想回避，可选取 $g \in I$，使得
$g \not \in (a_i)$ 对所有 $i$ 均成立。图示为
$$
U' = D(f) \cup D(g) \subset U \subset \Spec(R)
$$
由《代数续篇》引理 \ref{more-algebra-lemma-UFD-Pic-trivial}，$D(f)$ 与
$D(g)$ 的 Picard 群都平凡；而由关于 UFD 的初等论证，映射
$R_f^* \times R_g^* \to R_{fg}^*$，$(u, v) \mapsto uv^{-1}$ 是满射。
由此可知 $\Pic(U')$ 平凡。

\medskip\noindent
因此只需证明：可逆模 $\mathcal{L}$（定义在 $U$ 上）若沿
$j : U' \to U'$ 拉回后平凡，则它平凡。为此，只需证明映射
$\mathcal{L} \to j_*j^*\mathcal{L}$ 是同构。进一步，设
$u \in U \setminus U'$。则
$f, g \in \mathfrak m_u \subset \mathcal{O}_{U, u} = A$ 是局部 UFD 中的元素，
且没有共同素因子。读者容易由考察素因子证明，这意味着
$$
0 \to A \to A_f \times A_g \to A_{fg}
$$
正合。由于 $A = \mathcal{L}_u$（这里使用 $\mathcal{L}$ 可逆），又因为
$(j_*j^*\mathcal{L})_u$ 是 $A_f \times A_g \to A_{fg}$ 的核，结论随即得出。
\end{proof}

\begin{lemma}
\label{divisors-lemma-Pic-projective-space-UFD}
设 $R$ 为 UFD。$\mathbf{P}^n_R$ 的 Picard 群是 $\mathbf{Z}$。更精确地，
存在同构
$$
\mathbf{Z} \longrightarrow \Pic(\mathbf{P}^n_R),\quad
m \longmapsto \mathcal{O}_{\mathbf{P}^n_R}(m)
$$
特别地，射影空间 $\mathbf{P}^n_k$ 在域 $k$ 上的 Picard 群是
$\mathbf{Z}$。
\end{lemma}

\begin{proof}
由引理 \ref{divisors-lemma-open-subscheme-UFD}，开集
$D_+(T_i) \cong \mathbf{A}^n_R$ 的 Picard 群平凡。因此若
$\mathcal{L}$ 是 $\mathbf{P}^n_R$ 上的可逆模，则对于开覆盖
$$
\mathbf{P}^n_R = \bigcup\nolimits_{i = 0, \ldots, n} D_+(T_i)
$$
它由一个 $1$-上闭链给出，该上闭链取值于可逆函子层 $\mathcal{O}^*$。注意，对
$i \not = j$，有
$$
\mathcal{O}^*(D_+(T_i) \cap D_+(T_j)) =
\mathcal{O}^*(D_+(T_iT_j)) = \left(R[T_0, \ldots, T_n]_{(T_iT_j)}\right)^* =
R^* \times (T_i/T_j)^\mathbf{Z}
$$
因此，这样的上闭链 $(g_{ij})$ 由单位
$$
g_{ij} = u_{ij} (T_i/T_j)^{e_{ij}}
$$
给出，其中 $u_{ij} \in R^*$ 且 $e_{ij} \in \mathbf{Z}$，并满足上闭链条件。
$D_+(T_iT_jT_k)$ 上的上闭链条件在 $\#\{i, j, k\} = 3$ 时告诉我们
$$
u_{ik} = u_{ij} u_{jk} \quad\text{且}\quad
e_{ik} = e_{ij} = e_{jk}
$$
于是 $u_{ij} = u_{i1} u_{j1}^{-1}$ 是一个上边界。因此，可逆模的所有同构类
都由取如下上闭链得到：
$$
g_{ij} = (T_i/T_j)^e
$$
其中 $e \in \mathbf{Z}$。由于 $\mathcal{O}(n)$ 有平凡化截面
$T_i^n$（定义在 $D_+(T_i)$ 上），可知 $\mathcal{O}(n)$ 所对应的上闭链是
$(T_i/T_j)^n$。证明完成。
\end{proof}


\section{正规概形上的 Weil 除子}
\label{divisors-section-weil-divisors-normal}

\noindent
首先讨论自反模的性质。

\begin{lemma}
\label{divisors-lemma-reflexive-normal}
设 $X$ 为整的局部 Noether 正规概形。对 $\mathcal{F}$ 与 $\mathcal{G}$（它们是
相干自反 $\mathcal{O}_X$-模），映射
$$
(\SheafHom_{\mathcal{O}_X}(\mathcal{F}, \mathcal{O}_X)
\otimes_{\mathcal{O}_X} \mathcal{G})^{**} \to
\SheafHom_{\mathcal{O}_X}(\mathcal{F}, \mathcal{G})
$$
是同构。规则 $\mathcal{F}, \mathcal{G} \mapsto
(\mathcal{F} \otimes_{\mathcal{O}_X} \mathcal{G})^{**}$
在秩为 $1$ 的相干自反 $\mathcal{O}_X$-模之同构类集合上定义一个交换群律。
\end{lemma}

\begin{proof}
虽非严格必要，但建议先阅读注
\ref{divisors-remark-tensor}，再继续此证明。选取开子概形 $j : U \to X$，使得
$X \setminus U$ 的每个不可约分支之余维均为 $\geq 2$（余维在 $X$ 中取），并且
$j^*\mathcal{F}$ 与 $j^*\mathcal{G}$ 都有限局部自由，见引理
\ref{divisors-lemma-reflexive-over-normal}。映射
$$
\SheafHom_{\mathcal{O}_U}(j^*\mathcal{F}, \mathcal{O}_U)
\otimes_{\mathcal{O}_U} j^*\mathcal{G} \to
\SheafHom_{\mathcal{O}_U}(j^*\mathcal{F}, j^*\mathcal{G})
$$
是同构，因为可在局部检验，而且当这些模有限自由时结论显然。注意，对引理中
显示的箭头应用 $j^*$，所得正是刚刚证明为同构的箭头（略去一个小细节）。由于
$j^*$ 在 $U$ 上的相干自反模与 $X$ 上的相干自反模之间定义范畴等价
（由引理 \ref{divisors-lemma-reflexive-S2-extend} 及 Serre 判别法《性质》引理
\ref{properties-lemma-criterion-normal}），可知引理中的箭头也是同构。
若 $\mathcal{F}$ 的秩为 $1$，则 $j^*\mathcal{F}$ 是可逆
$\mathcal{O}_U$-模，而自反模
$\mathcal{F}^\vee = \SheafHom(\mathcal{F}, \mathcal{O}_X)$
限制后是它的逆。与前面同理可得
$(\mathcal{F} \otimes_{\mathcal{O}_X} \mathcal{F}^\vee)^{**} = \mathcal{O}_X$。
这样便得到引理陈述中群律的逆元。
\end{proof}

\begin{lemma}
\label{divisors-lemma-normal-class-group}
设 $X$ 为整的局部 Noether 正规概形。秩为 $1$ 的相干自反
$\mathcal{O}_X$-模所成的群同构于 Weil 除子类群 $\text{Cl}(X)$（属于 $X$）。
\end{lemma}

\begin{proof}
设 $\mathcal{F}$ 为秩为 $1$ 的相干自反 $\mathcal{O}_X$-模。选取开集
$U \subset X$，使得 $X \setminus U$ 的每个不可约分支之余维均为
$\geq 2$（余维在 $X$ 中取），并且 $\mathcal{F}|_U$ 可逆，见引理
\ref{divisors-lemma-reflexive-over-normal}。注意，$\text{Cl}(U) = \text{Cl}(X)$，因为
$X$ 的 Weil 除子类群只依赖于其有理函数域、余维为 $1$ 的点及其局部环。
因此可以把 $\mathcal{F}$ 的 Weil 除子类定义为 $\mathcal{F}|_U$ 在
$\text{Cl}(U)$ 中的 Weil 除子类。我们略去它不依赖于 $U$ 之选择的验证。

\medskip\noindent
用 $\text{Cl}'(X)$ 表示秩为 $1$ 的相干自反 $\mathcal{O}_X$-模之同构类集合。
上述构造给出群同态
$$
\text{Cl}'(X) \longrightarrow \text{Cl}(X)
$$
这是因为，对任意一对 $\mathcal{F}, \mathcal{G}$（它们属于 $\text{Cl}'(X)$），
可以选取一个对二者都适用的 $U$，而将可逆模映到其 Weil 除子类的赋值
(\ref{divisors-equation-c1}) 是同态。若 $\mathcal{F}$ 属于该映射的核，则由引理
\ref{divisors-lemma-normal-c1-injective}，$\mathcal{F}|_U$ 平凡，继而由引理
\ref{divisors-lemma-reflexive-S2-extend} 及 Serre 判别法《性质》引理
\ref{properties-lemma-criterion-normal}，$\mathcal{F}$ 也平凡。为完成证明，
只需检验该映射为满射。

\medskip\noindent
设 $D = \sum n_Z Z$ 为 $X$ 上的 Weil 除子。我们断言存在开集
$U \subset X$，使得 $X \setminus U$ 的每个不可约分支之余维均为
$\geq 2$（余维在 $X$ 中取），并且 $Z|_U$ 在 $n_Z \not = 0$ 时是有效 Cartier 除子。
为证明该断言，可以假设 $X$ 仿射。于是可以假设
$D = n_1 Z_1 + \ldots + n_r Z_r$ 是有限和，其中
$Z_1, \ldots, Z_r$ 两两不同。删去 $Z_i \cap Z_j$（$i \not = j$）后，
可以假设 $Z_1, \ldots, Z_r$ 两两不交。这把问题约化到单个素除子
$Z$（在 $X$ 上）的情形。由《性质》引理 \ref{properties-lemma-criterion-normal}，
$X$ 满足 $(R_1)$，故局部环
$\mathcal{O}_{X, \xi}$ 是离散赋值环；这里 $\xi$ 是 $Z$ 的泛点。设
$f \in \mathcal{O}_{X, \xi}$ 为一个素元。设 $V \subset X$ 为 $\xi$ 的开邻域，
使得 $f$ 是某个元素 $f \in \mathcal{O}_X(V)$ 的像。缩小 $V$ 后，可假设
$Z \cap V = V(f)$ 概形论地成立，因为这在 $\xi$ 处的局部环中成立。此时取
$$
U = X \setminus (Z \setminus V) = (X \setminus Z) \cup V
$$
便得到所需开集，从而证明该断言。

\medskip\noindent
为证明 $D$ 的除子类在像中，可以写成
$D = \sum_{n_Z < 0} n_Z Z - \sum_{n_Z > 0} (-n_Z) Z$。由上述映射的可加性，
可以而且将假设 $n_Z \leq 0$ 对所有素除子 $Z$ 均成立
（若读者愿意，可以避开这一步）。设 $U \subset X$ 如上述断言所述。
若 $U$ 拟紧，则写
$D|_U = -n_1 Z_1 - \ldots - n_r Z_r$，其中 $Z_i$ 是两两不同的素除子且
$n_i > 0$；并考虑可逆 $\mathcal{O}_U$-模
$$
\mathcal{L} =
\mathcal{I}_1^{n_1} \ldots \mathcal{I}_r^{n_r} \subset \mathcal{O}_U
$$
其中 $\mathcal{I}_i$ 是 $Z_i$ 的理想层。由 $U$ 的选取及引理
\ref{divisors-lemma-sum-effective-Cartier-divisors}，这是可逆的。此外，
$\text{div}_\mathcal{L}(1) = D|_U$。由引理
\ref{divisors-lemma-reflexive-S2-extend}，有 $\mathcal{L} = \mathcal{F}|_U$，其中存在某个
秩为 $1$ 的相干自反 $\mathcal{O}_X$-模 $\mathcal{F}$；因而
$D$ 在所述映射的像中。

\medskip\noindent
若 $U$ 不拟紧，则按上面的显示公式局部定义
$\mathcal{L} \subset \mathcal{O}_U$。读者可验证该构造能够黏合，随后与前面
完全相同地完成证明。细节从略。
\end{proof}

\begin{lemma}
\label{divisors-lemma-structure-sheaf-Xs}
设 $X$ 为整的局部 Noether 正规概形。设 $\mathcal{F}$ 为秩为 1 的相干自反
$\mathcal{O}_X$-模。设 $s \in \Gamma(X, \mathcal{F})$。令
$$
U = \{x \in X \mid s : \mathcal{O}_{X, x} \to \mathcal{F}_x
\text{ 为同构}\}
$$
则 $j : U \to X$ 是 $X$ 的开子概形，并且
$$
j_*\mathcal{O}_U =
\colim (\mathcal{O}_X \xrightarrow{s} \mathcal{F}
\xrightarrow{s} \mathcal{F}^{[2]}
\xrightarrow{s} \mathcal{F}^{[3]}
\xrightarrow{s} \ldots)
$$
其中 $\mathcal{F}^{[1]} = \mathcal{F}$，并且归纳地定义
$\mathcal{F}^{[n + 1]} =
(\mathcal{F} \otimes_{\mathcal{O}_X} \mathcal{F}^{[n]})^{**}$。
\end{lemma}

\begin{proof}
由《模》引理 \ref{modules-lemma-finite-type-surjective-on-stalk} 与
\ref{modules-lemma-finite-type-to-coherent-injective-on-stalk}，集合 $U$ 是开集。
注意，由《性质》引理 \ref{properties-lemma-immersion-into-noetherian}，$j$ 拟紧。
为证明最后的陈述，只需证明：对每个拟紧开集 $W \subset X$，存在同构
$$
\colim \Gamma(W, \mathcal{F}^{[n]})
\longrightarrow
\Gamma(U \cap W, \mathcal{O}_U)
$$
它是 $\mathcal{O}_X(W)$-模同构，且与限制映射相容。我们略去相容性的验证。
将 $X$ 替换为 $W$，并用 Hom 改写上式后，可知只需构造同构
$$
\colim \Hom_{\mathcal{O}_X}(\mathcal{O}_X, \mathcal{F}^{[n]})
\longrightarrow
\Hom_{\mathcal{O}_U}(\mathcal{O}_U, \mathcal{O}_U)
$$
选取开集 $V \subset X$，使得 $X \setminus V$ 的每个不可约分支之余维均为
$\geq 2$（余维在 $X$ 中取），并且 $\mathcal{F}|_V$ 可逆，见引理
\ref{divisors-lemma-reflexive-over-normal}。于是，限制定义了秩为 $1$ 的相干自反模范畴
之间的等价（这两个范畴分别在 $X$ 与 $V$ 上），也定义了秩为 $1$ 的相干
自反模范畴之间的等价（这两个范畴分别在 $U$ 与 $V \cap U$ 上）。见引理
\ref{divisors-lemma-reflexive-S2-extend} 及 Serre 判别法
《性质》引理 \ref{properties-lemma-criterion-normal}。因此只需构造同构
$$
\colim \Gamma(V, (\mathcal{F}|_V)^{\otimes n}) \longrightarrow
\Gamma(V \cap U, \mathcal{O}_U)
$$
由于 $\mathcal{F}|_V$ 可逆，而且 $U \cap V$ 正是 $s|_V$ 生成这个可逆模的
点所成的集合，这是《性质》引理
\ref{properties-lemma-invert-s-sections} 的特殊情形（该映射也有显式公式）。
\end{proof}

\begin{lemma}
\label{divisors-lemma-Xs-codim-complement}
假设与记号同引理 \ref{divisors-lemma-structure-sheaf-Xs}。若 $s$ 非零，则
$X \setminus U$ 的每个不可约分支之余维均为 $1$（余维在 $X$ 中取）。
\end{lemma}

\begin{proof}
设 $\xi \in X$ 是一个不可约分支 $Z$ 的泛点，而该分支属于 $X \setminus U$。将 $X$
替换为 $\xi$ 的一个开邻域后，可以假设 $Z = X \setminus U$ 不可约。由于
$s : \mathcal{O}_U \to \mathcal{F}|_U$ 是同构，若 $Z$ 在 $X$ 中的余维为
$\geq 2$，则由引理 \ref{divisors-lemma-reflexive-S2-extend} 及 Serre 判别法《性质》引理
\ref{properties-lemma-criterion-normal}，
$s : \mathcal{O}_X \to \mathcal{F}$ 是同构。这将意味着
$Z = \emptyset$，矛盾。
\end{proof}

\begin{remark}
\label{divisors-remark-structure-sheaf-Xs}
设 $A$ 为 Noether 正规整环。设 $M$ 为秩为 $1$ 的有限自反 $A$-模。设
$s \in M$ 非零。设 $\mathfrak p_1, \ldots, \mathfrak p_r$ 为高度为
$1$ 的素理想（属于 $A$），它们位于 $M/As$ 的支集中。由引理
\ref{divisors-lemma-structure-sheaf-Xs}，开集 $U$ 为
$$
U = \Spec(A) \setminus
\left(V(\mathfrak p_1) \cup \ldots \cup V(\mathfrak p_r)\right)
$$
这是引理 \ref{divisors-lemma-Xs-codim-complement} 的结论。此外，若 $M^{[n]}$ 表示
$M \otimes_A \ldots \otimes_A M$（$n$ 个因子）的自反包，则由引理
\ref{divisors-lemma-structure-sheaf-Xs}，
$$
\Gamma(U, \mathcal{O}_U) = \colim M^{[n]}
$$
\end{remark}

\begin{lemma}
\label{divisors-lemma-affine-Xs}
假设与记号同引理 \ref{divisors-lemma-structure-sheaf-Xs}。下列条件等价：
\begin{enumerate}
\item 包含态射 $j : U \to X$ 是仿射的；并且
\item 对每个 $x \in X \setminus U$，存在 $n > 0$，使得
$s^n \in \mathfrak m_x \mathcal{F}^{[n]}_x$。
\end{enumerate}
\end{lemma}

\begin{proof}
假设 (1)。对 $x \in X \setminus U$，逆像 $U_x$（它是 $U$ 在典范态射
$f_x : \Spec(\mathcal{O}_{X, x}) \to X$ 下的逆像）是仿射的，且不包含
$x$。因此 $\mathfrak m_x \Gamma(U_x, \mathcal{O}_{U_x})$ 是单位理想。特别地，
可写成
$$
1 = \sum f_i g_i
$$
其中 $f_i \in \mathfrak m_x$ 且 $g_i \in \Gamma(U_x, \mathcal{O}_{U_x})$。
由引理 \ref{divisors-lemma-structure-sheaf-Xs}，有
$\Gamma(U_x, \mathcal{O}_{U_x}) = \colim \mathcal{F}^{[n]}_x$，其转移映射由乘以
$s$ 给出。因此对某个 $n > 0$，有
$$
s^n = \sum f_i t_i
$$
其中某些 $t_i = s^ng_i \in \mathcal{F}^{[n]}_x$。故 (2) 成立。

\medskip\noindent
反之，假设 (2) 成立。$j$ 是否仿射是 $X$ 上的局部问题，见《态射》引理
\ref{morphisms-lemma-characterize-affine}。因此可以而且将假设 $X$ 仿射。
目标是证明 $U$ 仿射。由《概形上同调》引理
\ref{coherent-lemma-affine-if-quasi-affine}，只需证明
$H^p(U, \mathcal{O}_U) = 0$ 对 $p > 0$ 成立。由于
$H^p(U, \mathcal{O}_U) = H^0(X, R^pj_*\mathcal{O}_U)$
（《概形上同调》引理
\ref{coherent-lemma-quasi-coherence-higher-direct-images-application}），并且
$R^pj_*\mathcal{O}_U$ 拟相干（《概形上同调》引理
\ref{coherent-lemma-quasi-coherence-higher-direct-images}），只需证明茎
$(R^pj_*\mathcal{O}_U)_x$ 在点 $x \in X$ 处为零。考虑基变换图
$$
\xymatrix{
U_x \ar[d]_{j_x} \ar[r] & U \ar[d]^j \\
\Spec(\mathcal{O}_{X, x}) \ar[r] & X
}
$$
由《概形上同调》引理
\ref{coherent-lemma-flat-base-change-cohomology}，有
$(R^pj_*\mathcal{O}_U)_x = R^pj_{x, *}\mathcal{O}_{U_x}$。因此可假设 $X$ 是以
$x$ 为闭点的局部概形，并且必须证明 $U$ 仿射（因为由上面给出的引用，这与
所需消失等价）。特别地，$d = \dim(X)$ 有限（《代数》命题
\ref{algebra-proposition-dimension}）。若 $x \in U$，则 $U = X$，结论显然。
若 $d = 0$ 且 $x \not \in U$，则 $U = \emptyset$，结论也显然。现在假设
$d > 0$ 且 $x \not \in U$。由于
$j_*\mathcal{O}_U = \colim \mathcal{F}^{[n]}$，假设意味着可写成
$$
1 = \sum f_i g_i
$$
其中对某个 $n > 0$，有 $f_i \in \mathfrak m_x$ 且
$g_i \in \mathcal{O}(U)$。由对 $d$ 的归纳可知，
$D(f_i) \cap U$ 对所有 $i$ 都仿射：把刚才的整个论证中的 $X$ 替换为 $D(f_i)$，最终会
得到维数严格小于 $d$ 的 Noether 局部环。因此由《性质》引理
\ref{properties-lemma-characterize-affine}，$U$ 仿射，正如所需。
\end{proof}


\section{相对 Proj}
\label{divisors-section-relative-proj}

\noindent
这里给出相对 Proj 的一些结果。先从若干非常基本的结果开始。回忆，相对 Proj
总是在基底上分离，见《构造》引理
\ref{constructions-lemma-relative-proj-separated}。

\begin{lemma}
\label{divisors-lemma-relative-proj-quasi-compact}
设 $S$ 为概形。设 $\mathcal{A}$ 为拟相干分次 $\mathcal{O}_S$-代数。设
$p : X = \underline{\text{Proj}}_S(\mathcal{A}) \to S$ 为
$\mathcal{A}$ 的相对 Proj。若下列条件之一成立：
\begin{enumerate}
\item $\mathcal{A}$ 作为 $\mathcal{A}_0$-代数层是有限型的；
\item $\mathcal{A}$ 由 $\mathcal{A}_1$ 生成（作为 $\mathcal{A}_0$-代数），且
$\mathcal{A}_1$ 是有限型 $\mathcal{A}_0$-模；
\item 存在有限型拟相干 $\mathcal{A}_0$-子模
$\mathcal{F} \subset \mathcal{A}_{+}$，使得
$\mathcal{A}_{+}/\mathcal{F}\mathcal{A}$ 是
$\mathcal{A}/\mathcal{F}\mathcal{A}$ 的局部幂零理想层；
\end{enumerate}
则 $p$ 拟紧。
\end{lemma}

\begin{proof}
该问题在基底上是局部的，见《概形》引理
\ref{schemes-lemma-quasi-compact-affine}。因此可假设 $S$ 仿射。记
$S = \Spec(R)$，且 $\mathcal{A}$ 对应于分次 $R$-代数 $A$。于是
$X = \text{Proj}(A)$，见《构造》节
\ref{constructions-section-relative-proj-via-glueing}。在情形 (1)，进一步局部化后
可假设 $A$ 由齐次元素 $f_1, \ldots, f_n \in A_{+}$ 生成（在 $A_0$ 上）。于是由
《代数》引理 \ref{algebra-lemma-S-plus-generated}，有
$A_{+} = (f_1, \ldots, f_n)$。在情形 (3)，可知
$\mathcal{F} = \widetilde{M}$，其中 $A_0$-模
$M \subset A_{+}$ 是有限型的。记
$M = \sum A_0f_i$。把 $f_i = \sum f_{i, j}$ 写成齐次分量之分解。条件 (3)
表示 $A_{+} \subset \sqrt{(f_{i, j})}$。因此在这两种情形，由《构造》引理
\ref{constructions-lemma-proj-quasi-compact}，可知 $\text{Proj}(A)$ 拟紧。
最后，(2) 由 (1) 得出。
\end{proof}

\begin{lemma}
\label{divisors-lemma-relative-proj-finite-type}
设 $S$ 为概形。设 $\mathcal{A}$ 为拟相干分次 $\mathcal{O}_S$-代数。设
$p : X = \underline{\text{Proj}}_S(\mathcal{A}) \to S$ 为
$\mathcal{A}$ 的相对 Proj。若 $\mathcal{A}$ 作为 $\mathcal{O}_S$-代数层是
有限型的，则 $p$ 是有限型态射，且 $\mathcal{O}_X(d)$ 是有限型
$\mathcal{O}_X$-模。
\end{lemma}

\begin{proof}
该假设蕴含 $p$ 拟紧，见引理 \ref{divisors-lemma-relative-proj-quasi-compact}。因此只需
证明 $p$ 局部有限型。故该问题在基底和目标上是局部的，见《态射》引理
\ref{morphisms-lemma-locally-finite-type-characterize}。记 $S = \Spec(R)$，且
$\mathcal{A}$ 对应于分次 $R$-代数 $A$。在 $S$ 上进一步局部化后，可假设
$A$ 是有限型 $R$-代数。概形 $X$ 由环 $A_{(f)}$ 的谱黏合而成，其中
$f \in A_{+}$ 齐次。由《代数》引理
\ref{algebra-lemma-dehomogenize-finite-type} 的 (1)，每个这样的环在 $R$ 上都是
有限型的。因此 $\text{Proj}(A)$ 在 $R$ 上有限型。关于
$\mathcal{O}_X(d)$ 的陈述使用《代数》引理
\ref{algebra-lemma-dehomogenize-finite-type} 的 (2)。
\end{proof}

\begin{lemma}
\label{divisors-lemma-relative-proj-universally-closed}
设 $S$ 为概形。设 $\mathcal{A}$ 为拟相干分次 $\mathcal{O}_S$-代数。设
$p : X = \underline{\text{Proj}}_S(\mathcal{A}) \to S$ 为
$\mathcal{A}$ 的相对 Proj。若 $\mathcal{O}_S \to \mathcal{A}_0$ 是整代数映射
\footnote{换言之，$\mathcal{O}_S$ 在 $\mathcal{A}_0$ 中的整闭包（见《态射》定义
\ref{morphisms-definition-integral-closure}）等于 $\mathcal{A}_0$。}，并且
$\mathcal{A}$ 作为 $\mathcal{A}_0$-代数有限型，则 $p$ 万有闭。
\end{lemma}

\begin{proof}
该问题在基底上是局部的。因此可假设 $X = \Spec(R)$ 仿射。设
$\mathcal{A}$ 为拟相干 $\mathcal{O}_X$-代数，它与分次 $R$-代数 $A$ 相伴。
假设是 $R \to A_0$ 为整映射，且 $A$ 在 $A_0$ 上有限型。把
$X \to \Spec(R)$ 写成复合 $X \to \Spec(A_0) \to \Spec(R)$。由于
$R \to A_0$ 是整环映射，由《态射》引理
\ref{morphisms-lemma-integral-universally-closed}，可知
$\Spec(A_0) \to \Spec(R)$ 万有闭。拟紧态射（见《构造》引理
\ref{constructions-lemma-proj-quasi-compact}）
$$
X = \text{Proj}(A) \to \Spec(A_0)
$$
由《构造》引理 \ref{constructions-lemma-proj-valuative-criterion} 满足赋值判别法的
存在性部分，因而由《概形》命题
\ref{schemes-proposition-characterize-universally-closed} 万有闭。所以
$X \to \Spec(R)$ 作为万有闭态射的复合是万有闭的。
\end{proof}

\begin{lemma}
\label{divisors-lemma-relative-proj-proper}
设 $S$ 为概形。设 $\mathcal{A}$ 为拟相干分次 $\mathcal{O}_S$-代数。设
$p : X = \underline{\text{Proj}}_S(\mathcal{A}) \to S$ 为
$\mathcal{A}$ 的相对 Proj。下列条件等价：
\begin{enumerate}
\item $\mathcal{A}_0$ 是有限型 $\mathcal{O}_S$-模，且 $\mathcal{A}$ 作为
$\mathcal{A}_0$-代数有限型；
\item $\mathcal{A}_0$ 是有限型 $\mathcal{O}_S$-模，且 $\mathcal{A}$ 作为
$\mathcal{O}_S$-代数有限型。
\end{enumerate}
若这些条件成立，则 $p$ 局部射影，特别地它是固有态射。
\end{lemma}

\begin{proof}
假设 $\mathcal{A}_0$ 是有限型 $\mathcal{O}_S$-模。选取仿射开集
$U = \Spec(R) \subset X$，使得 $\mathcal{A}$ 对应于分次 $R$-代数 $A$，且
$A_0$ 是有限 $R$-模。条件 (1) 意味着（在 $S$ 上进一步局部化后）$A$ 是
有限型 $A_0$-代数，而条件 (2) 意味着（在 $S$ 上进一步局部化后）$A$ 是
有限型 $R$-代数。因此由《代数》引理
\ref{algebra-lemma-compose-finite-type}，这两个条件相互蕴含。

\medskip\noindent
局部射影态射是固有的，见《态射》引理
\ref{morphisms-lemma-locally-projective-proper}。因此现在可假设
$S = \Spec(R)$ 且 $X = \text{Proj}(A)$，并且 $A_0$ 在 $R$ 上有限而 $A$ 在
$R$ 上有限型。我们将证明 $X = \text{Proj}(A) \to \Spec(R)$ 是射影态射。
我们建议读者自己证明这一点：直接构造 $X$ 到 $R$ 上射影空间中的闭浸入，
而不要阅读下面给出的论证。

\medskip\noindent
由引理 \ref{divisors-lemma-relative-proj-finite-type}，可知 $X$ 在 $\Spec(R)$ 上有限型。
《构造》引理 \ref{constructions-lemma-ample-on-proj} 告诉我们，
$\mathcal{O}_X(d)$ 在 $X$ 上对某个 $d \geq 1$ 丰沛（见《性质》节
\ref{properties-section-ample}）。因此 $X \to \Spec(R)$ 是拟射影态射（由
《态射》定义 \ref{morphisms-definition-quasi-projective}）。由《态射》引理
\ref{morphisms-lemma-quasi-projective-open-projective}，可知 $X$ 同构于某个在
$\Spec(R)$ 上射影的概形之开子概形。因此，为完成证明，只需证明
$X \to \Spec(R)$ 万有闭（使用《态射》引理
\ref{morphisms-lemma-image-proper-scheme-closed}）。这由引理
\ref{divisors-lemma-relative-proj-universally-closed} 得出。
\end{proof}

\begin{lemma}
\label{divisors-lemma-relative-proj-projective}
设 $S$ 为概形。设 $\mathcal{A}$ 为拟相干分次 $\mathcal{O}_S$-代数。设
$p : X = \underline{\text{Proj}}_S(\mathcal{A}) \to S$ 为
$\mathcal{A}$ 的相对 Proj。若 $\mathcal{A}$ 由 $\mathcal{A}_1$ 生成（在
$\mathcal{A}_0$ 上），且 $\mathcal{A}_1$ 是有限型 $\mathcal{O}_S$-模，则
$p$ 是射影态射。
\end{lemma}

\begin{proof}
具体而言，与分次 $\mathcal{O}_S$-代数映射
$$
\text{Sym}_{\mathcal{O}_X}^*(\mathcal{A}_1)
\longrightarrow
\mathcal{A}
$$
相伴的态射是闭浸入 $X \to \mathbf{P}(\mathcal{A}_1)$，见《构造》引理
\ref{constructions-lemma-surjective-generated-degree-1-map-relative-proj}。
\end{proof}

\begin{lemma}
\label{divisors-lemma-relative-proj-flat}
设 $S$ 为概形。设 $\mathcal{A}$ 为拟相干分次 $\mathcal{O}_S$-代数。设
$p : X = \underline{\text{Proj}}_S(\mathcal{A}) \to S$ 为
$\mathcal{A}$ 的相对 Proj。若 $\mathcal{A}_d$ 是平坦 $\mathcal{O}_S$-模
（当 $d \gg 0$ 时），则 $p$ 平坦，且 $\mathcal{O}_X(d)$ 在 $S$ 上平坦。
\end{lemma}

\begin{proof}
在仿射局部，$X$ 在 $S$ 上的平坦性约化为下述陈述：设 $R$ 为环，设 $A$ 为
分次 $R$-代数，且 $A_d$ 在 $R$ 上平坦（当 $d \gg 0$ 时）；设
$f \in A_d$，其中 $d > 0$。则 $A_{(f)}$ 在 $R$ 上平坦。由于
$A_{(f)} = \colim A_{nd}$，且转移映射由乘以 $f$ 给出，该结论由《代数》引理
\ref{algebra-lemma-colimit-flat} 得出。用同样的论证可得
$\mathcal{O}_X(d)$ 在 $S$ 上的平坦性。
\end{proof}

\begin{lemma}
\label{divisors-lemma-relative-proj-finite-presentation}
设 $S$ 为概形。设 $\mathcal{A}$ 为拟相干分次 $\mathcal{O}_S$-代数。设
$p : X = \underline{\text{Proj}}_S(\mathcal{A}) \to S$ 为
$\mathcal{A}$ 的相对 Proj。若 $\mathcal{A}$ 是有限呈示
$\mathcal{O}_S$-代数，则 $p$ 是有限呈示态射，且 $\mathcal{O}_X(d)$ 是有限呈示
$\mathcal{O}_X$-模。
\end{lemma}

\begin{proof}
在仿射局部，这约化为下述陈述：设 $R$ 为环，且设 $A$ 为有限呈示分次
$R$-代数。则 $\text{Proj}(A) \to \Spec(R)$ 是有限呈示态射，而且
$\mathcal{O}_{\text{Proj}(A)}(d)$ 是有限呈示
$\mathcal{O}_{\text{Proj}(A)}$-模。有限呈示条件蕴含可以选取呈示
$$
A = R[X_1, \ldots, X_n]/(F_1, \ldots, F_m)
$$
其中 $R[X_1, \ldots, X_n]$ 是如下分次多项式环：赋予权 $d_i$ 给 $X_i$；而
$F_1, \ldots, F_m$ 是次数为 $e_j$ 的齐次多项式。设 $R_0 \subset R$ 为由
多项式 $F_1, \ldots, F_m$ 的系数生成的子环。然后置
$A_0 = R_0[X_1, \ldots, X_n]/(F_1, \ldots, F_m)$。按构造，
$A = A_0 \otimes_{R_0} R$。因此由《构造》引理
\ref{constructions-lemma-base-change-map-proj}，只需对
$X_0 = \text{Proj}(A_0)$（在 $R_0$ 上）证明该结果。由引理
\ref{divisors-lemma-relative-proj-finite-type}，已知 $X_0$ 在 $R_0$ 上有限型，且
$\mathcal{O}_{X_0}(d)$ 是有限型拟相干 $\mathcal{O}_{X_0}$-模。由于 $R_0$
是 Noether 环（它是有限生成 $\mathbf{Z}$-代数），可知 $X_0$ 在 $R_0$ 上有限呈示
（《态射》引理 \ref{morphisms-lemma-noetherian-finite-type-finite-presentation}），
而由《概形上同调》引理 \ref{coherent-lemma-coherent-Noetherian}，
$\mathcal{O}_{X_0}(d)$ 是有限呈示模。证明完成。
\end{proof}


\section{相对 Proj 的闭子概形}
\label{divisors-section-closed-in-relative-proj}

\noindent
这里给出关于相对 Proj 的闭子概形的若干辅助引理。

\begin{lemma}
\label{divisors-lemma-closed-subscheme-proj}
设 $S$ 为概形。设 $\mathcal{A}$ 为拟相干分次
$\mathcal{O}_S$-代数。设
$p : X = \underline{\text{Proj}}_S(\mathcal{A}) \to S$ 为
$\mathcal{A}$ 的相对 Proj。设 $i : Z \to X$ 为闭子概形。记
$\mathcal{I} \subset \mathcal{A}$；它是下述典范映射的核：
$$
\mathcal{A}
\longrightarrow
\bigoplus\nolimits_{d \geq 0} p_*\left((i_*\mathcal{O}_Z)(d)\right).
$$
若 $p$ 拟紧，则有同构
$Z = \underline{\text{Proj}}_S(\mathcal{A}/\mathcal{I})$。
\end{lemma}

\begin{proof}
由《构造》引理 \ref{constructions-lemma-relative-proj-separated}，态射 $p$
是分离的。由于 $p$ 拟紧，$p_*$ 把拟相干模变为拟相干模，见《概形》引理
\ref{schemes-lemma-push-forward-quasi-coherent}。因此 $\mathcal{I}$ 是拟相干
$\mathcal{O}_S$-模。特别地，$\mathcal{B} = \mathcal{A}/\mathcal{I}$ 是拟相干
分次 $\mathcal{O}_S$-代数。由函子性得到的态射
$Z' = \underline{\text{Proj}}_S(\mathcal{B}) \to
\underline{\text{Proj}}_S(\mathcal{A})$ 处处有定义且为闭浸入，见《构造》引理
\ref{constructions-lemma-surjective-graded-rings-map-relative-proj}。
因此只需证明 $Z = Z'$ 作为 $X$ 的闭子概形相等。

\medskip\noindent
由此，该问题在基底上是局部的，故可假设
$S = \Spec(R)$ 且 $X = \text{Proj}(A)$，其中存在一个在 $R$ 上的分次代数
$A$。假设 $\mathcal{I} = \widetilde{I}$，其中
$I \subset A$ 是分次理想。由《构造》引理
\ref{constructions-lemma-proj-quasi-compact}，存在
$f_0, \ldots, f_n \in A_{+}$，使得
$A_{+} \subset \sqrt{(f_0, \ldots, f_n)}$；换言之，
$X = \bigcup D_{+}(f_i)$。因此只需验证
$Z \cap D_{+}(f_i) = Z' \cap D_{+}(f_i)$ 对每个 $i$ 都成立。重新编号后可假设
$i = 0$。设 $Z \cap D_{+}(f_0)$ 与 $Z' \cap D_{+}(f_0)$
分别由理想 $J$ 与 $J'$ 截出，这两个理想属于 $A_{(f_0)}$。

\medskip\noindent
先证包含关系 $J' \subset J$。设 $d$ 是
$\deg(f_0), \ldots, \deg(f_n)$ 的最小公倍数。注意，每个扭曲
$\mathcal{O}_X(nd)$ 都可逆，并且由 $f_i^{nd/\deg(f_i)}$ 在
$D_{+}(f_i)$ 上将其平凡化；对于任意拟相干模 $\mathcal{F}$（在 $X$ 上），乘法映射
$\mathcal{O}_X(nd) \otimes_{\mathcal{O}_X} \mathcal{F}(m)
\to \mathcal{F}(nd + m)$
都是同构，见《构造》引理
\ref{constructions-lemma-when-invertible}。注意，$J'$ 是由元素
$g/f_0^e$ 生成的理想，其中 $g \in I$ 是次数为
$e\deg(f_0)$ 的齐次元素（见《构造》引理
\ref{constructions-lemma-surjective-graded-rings-map-proj} 的证明）。
当然，把 $g$ 替换为 $f_0^lg$（这里 $l$ 取适当值），则总可假设
$d | e$。由于 $g$ 作为 $\mathcal{O}_X(e\deg(f_0))$ 的截面限制到
$Z$ 后消失，可知 $g/f_0^d$ 是 $J$ 的元素。因此 $J' \subset J$。

\medskip\noindent
反之，假设 $g/f_0^e \in J$。仍可假设 $d | e$。选取
$i \in \{1, \ldots, n\}$。此时 $Z \cap D_{+}(f_i)$ 由某个理想
$J_i \subset A_{(f_i)}$ 截出，并且
$$
J \cdot A_{(f_0f_i)} = J_i \cdot A_{(f_0f_i)}.
$$
右端是 $J_i$ 关于
$f_0^{\deg(f_i)}/f_i^{\deg(f_0)}$ 的局部化。因此，对某个充分可除的
$e_i \gg 0$，有
$$
f_0^{e_i}g/f_i^{(e_i + e)\deg(f_0)/\deg(f_i)} \in J_i
$$
这证明 $f_0^{\max(e_i)}g$ 是 $I$ 的元素，因为它在每个仿射开集
$D_{+}(f_i)$ 上的限制都在闭子概形
$Z \cap D_{+}(f_i)$ 上消失。所以 $g/f_0^e \in J'$，从而如所需有
$J \subset J'$。
\end{proof}

\begin{example}
\label{divisors-example-closed-subscheme-of-proj}
设 $A$ 为分次环。设 $X = \text{Proj}(A)$ 且 $S = \Spec(A_0)$。
给定分次理想 $I \subset A$，由《构造》引理
\ref{constructions-lemma-surjective-graded-rings-map-proj} 得到闭子概形
$V_+(I) = \text{Proj}(A/I) \to X$。把引理
\ref{divisors-lemma-closed-subscheme-proj} 的结果翻译到这一情形，可知若 $X$ 拟紧，
则任意闭子概形 $Z$ 都具有形式 $V_+(I(Z))$，其中分次理想
$I(Z) \subset A$ 由下式给出：
$$
I(Z) = \Ker(A \longrightarrow
\bigoplus\nolimits_{n \geq 0} \Gamma(Z, \mathcal{O}_Z(n)))
$$
于是可以提出以下两个自然问题：
\begin{enumerate}
\item 哪些理想 $I$ 具有形式 $I(Z)$？
\item 如何描述运算 $I \mapsto I(V_+(I))$？
\end{enumerate}
我们将在 $A$ 为 Noether 环时回答这些问题。

\medskip\noindent
首先假设 $A$ 由 $A_1$ 生成（在 $A_0$ 上）。在此情形，对任意理想
$I \subset A$，映射
$A/I \to \bigoplus \Gamma(\text{Proj}(A/I), \mathcal{O})$
的核是 $A/I$ 的挠元素集合，见《概形上同调》命题
\ref{coherent-proposition-coherent-modules-on-proj}。因此
$$
I(V_+(I)) = \{x \in A \mid A_n x \subset I\text{ 对某个 }n \geq 0\}
$$
右端的理想有时称为 $I$ 的饱和化。这回答了 (2)；对 (1) 的回答是：
一个理想具有形式 $I(Z)$，当且仅当它是饱和的，即等于自身的饱和化。

\medskip\noindent
若 $A$ 是一般的 Noether 分次环，则使用《概形上同调》命题
\ref{coherent-proposition-coherent-modules-on-proj-general}。因此，若 $d$
等于 $A$ 在 $A_0$ 上一组生成元的次数之最小公倍数，则
$$
I(V_+(I)) = \{x \in A \mid (Ax)_{nd} \subset I\text{ 对所有 }n \gg 0\}
$$
这可能不同于 $I$ 的饱和化；当 $d \not = 1$ 时会出现这种情形。例如，假设
$A = \mathbf{Q}[x, y]$，其中 $\deg(x) = 2$ 且 $\deg(y) = 3$。
于是 $d = 6$。令 $I = (y^2)$。可知 $y \in I(V_+(I))$，因为对于任意
齐次 $f \in A$，若 $6 | \deg(fy)$，则 $y | f$，故 $fy \in I$。
由此 $I(V_+(I)) = (y)$，但 $x^n y \not \in I$ 对所有 $n$
都成立；因此 $I(V_+(I))$ 不等于该饱和化。
\end{example}

\begin{lemma}
\label{divisors-lemma-equation-codim-1-in-projective-space}
设 $R$ 为 Noether UFD。设 $Z \subset \mathbf{P}^n_R$ 为闭子概形，
它不含嵌入点，并且 $Z$ 的每个不可约分支的余维均为 $1$（在
$\mathbf{P}^n_R$ 中）。则理想
$I(Z) \subset R[T_0, \ldots, T_n]$（它对应于 $Z$）为主理想。
\end{lemma}

\begin{proof}
注意，$X = \mathbf{P}^n_R$ 的局部环都是 UFD，因为 $X$ 被同构于
$\mathbf{A}^n_R$ 的仿射片覆盖，而且 $R[x_1, \ldots, x_n]$ 是 UFD
（《代数》引理 \ref{algebra-lemma-polynomial-ring-UFD}）。因此由引理
\ref{divisors-lemma-codimension-1-is-effective-Cartier}，$Z$ 是有效 Cartier 除子。
设 $\mathcal{I} \subset \mathcal{O}_X$ 为与 $Z$ 对应的拟相干理想层。
选取同构 $\mathcal{O}(m) \to \mathcal{I}$，其中
$m \in \mathbf{Z}$，见引理
\ref{divisors-lemma-Pic-projective-space-UFD}。于是复合
$$
\mathcal{O}_X(m) \to \mathcal{I} \to \mathcal{O}_X
$$
非零。可知 $m \leq 0$，且 $\mathcal{O}_X(m)^{\otimes -1} = \mathcal{O}_X(-m)$
的相应截面由某个
$F \in R[T_0, \ldots, T_n]$ 给出，其次数为 $-m$，见《概形上同调》引理
\ref{coherent-lemma-cohomology-projective-space-over-ring}。因此，在第
$i$ 个标准开集 $U_i = D_+(T_i)$ 上，闭子概形
$Z \cap U_i$ 由理想
$$
(F(T_0/T_i, \ldots, T_n/T_i)) \subset R[T_0/T_i, \ldots, T_n/T_i]
$$
截出。因此分次理想
$I(Z) = \Ker(R[T_0, \ldots, T_n] \to \bigoplus \Gamma(\mathcal{O}_Z(m)))$
的齐次元素是满足下式的齐次多项式 $G$：
$$
G(T_0/T_i, \ldots, T_n/T_i) \in (F(T_0/T_i, \ldots, T_n/T_i))
$$
其中 $i = 0, \ldots, n$。清去分母可知，存在 $e_i \geq 0$，使得
$$
T_i^{e_i}G \in (F)
$$
对 $i = 0, \ldots, n$ 成立。由于 $R$ 是 UFD，
$R[T_0, \ldots, T_n]$ 也是 UFD。因此
$F | T_0^{e_0}G$ 与 $F | T_1^{e_1}G$ 蕴含 $F | G$，因为
$T_0^{e_0}$ 与 $T_1^{e_1}$ 没有公因子。故 $I(Z) = (F)$。
\end{proof}

\noindent
当闭子概形在局部由有限多个方程截出时，可以用 $\mathcal{A}$ 的有限型理想层
定义它。

\begin{lemma}
\label{divisors-lemma-closed-subscheme-proj-finite}
设 $S$ 为拟紧拟分离概形。设 $\mathcal{A}$ 为拟相干分次
$\mathcal{O}_S$-代数。设
$p : X = \underline{\text{Proj}}_S(\mathcal{A}) \to S$ 为
$\mathcal{A}$ 的相对 Proj。设 $i : Z \to X$ 为闭子概形。
若 $p$ 拟紧且 $i$ 是有限呈示态射，则存在 $d > 0$ 和有限型拟相干
$\mathcal{O}_S$-子模 $\mathcal{F} \subset \mathcal{A}_d$，使得
$Z = \underline{\text{Proj}}_S(\mathcal{A}/\mathcal{F}\mathcal{A})$。
\end{lemma}

\begin{proof}
由引理 \ref{divisors-lemma-closed-subscheme-proj}，已知存在拟相干分次理想层
$\mathcal{I} \subset \mathcal{A}$，使得
$Z = \underline{\text{Proj}}(\mathcal{A}/\mathcal{I})$。由于 $S$ 拟紧，
可以选取有限仿射开覆盖 $S = U_1 \cup \ldots \cup U_n$。记
$U_i = \Spec(R_i)$。设 $\mathcal{A}|_{U_i}$ 对应于分次
$R_i$-代数 $A_i$，而 $\mathcal{I}|_{U_i}$ 对应于分次理想
$I_i \subset A_i$。注意，$p^{-1}(U_i) = \text{Proj}(A_i)$ 是
$R_i$ 上的概形。由于 $p$ 拟紧，可以选取有限多个齐次元素
$f_{i, j} \in A_{i, +}$，使得
$p^{-1}(U_i) = \bigcup_j D_{+}(f_{i, j})$。关于
$Z \to X$ 的条件意味着 $Z$ 的理想层（作为 $\mathcal{O}_X$ 的理想）有限型，见
《态射》引理 \ref{morphisms-lemma-closed-immersion-finite-presentation}。
因此可以找到有限多个齐次元素
$h_{i, j, k} \in I_i \cap A_{i, +}$，使得
$Z \cap D_{+}(f_{i, j})$ 的理想由元素
$h_{i, j, k}/f_{i, j}^{e_{i, j, k}}$ 生成。选取 $d > 0$，使其为所有整数
$\deg(f_{i, j})$ 与 $\deg(h_{i, j, k})$ 的公倍数。由《性质》引理
\ref{properties-lemma-quasi-coherent-colimit-finite-type}，存在有限型拟相干子模
$\mathcal{F} \subset \mathcal{I}_d$，使得所有局部截面
$$
h_{i, j, k}f_{i, j}^{(d - \deg(h_{i, j, k}))/\deg(f_{i, j})}
$$
都是 $\mathcal{F}$ 的截面。按构造，$\mathcal{F}$ 即为所求。
\end{proof}

\noindent
引理 \ref{divisors-lemma-closed-subscheme-proj-finite} 的下列版本将在引理
\ref{divisors-lemma-composition-admissible-blowups} 的证明中使用。

\begin{lemma}
\label{divisors-lemma-closed-subscheme-proj-finite-type}
设 $S$ 为拟紧拟分离概形。设 $\mathcal{A}$ 为拟相干分次
$\mathcal{O}_S$-代数。设
$p : X = \underline{\text{Proj}}_S(\mathcal{A}) \to S$ 为
$\mathcal{A}$ 的相对 Proj。设 $i : Z \to X$ 为闭子概形。
设 $U \subset S$ 为开集。假设：
\begin{enumerate}
\item $p$ 拟紧；
\item $i$ 是有限呈示态射；
\item $U \cap p(i(Z)) = \emptyset$；
\item $U$ 拟紧；
\item $\mathcal{A}_n$ 都是有限型 $\mathcal{O}_S$-模（对所有 $n$）。
\end{enumerate}
则存在 $d > 0$ 和有限型拟相干 $\mathcal{O}_S$-子模
$\mathcal{F} \subset \mathcal{A}_d$，使得 (a)
$Z = \underline{\text{Proj}}_S(\mathcal{A}/\mathcal{F}\mathcal{A})$，且 (b)
$\mathcal{A}_d/\mathcal{F}$ 的支集与 $U$ 不交。
\end{lemma}

\begin{proof}
设 $\mathcal{I} \subset \mathcal{A}$ 为引理
\ref{divisors-lemma-closed-subscheme-proj} 中构造的拟相干分次理想层。令
$U_i$、$R_i$、$A_i$、$I_i$、$f_{i, j}$、$h_{i, j, k}$ 与 $d$
如引理 \ref{divisors-lemma-closed-subscheme-proj-finite} 的证明中所构造。由于
$U \cap p(i(Z)) = \emptyset$，可知
$\mathcal{I}_d|_U = \mathcal{A}_d|_U$（由我们把 $\mathcal{I}$
构造为核的方式）。由于 $U$ 拟紧，可选取有限仿射开覆盖
$U = W_1 \cup \ldots \cup W_m$。由于 $\mathcal{A}_d$ 有限型，可以找到
有限多个截面 $g_{t, s} \in \mathcal{A}_d(W_t)$，它们生成
$\mathcal{A}_d|_{W_t} = \mathcal{I}_d|_{W_t}$（作为
$\mathcal{O}_{W_t}$-模）。为完成证明，注意由《性质》引理
\ref{properties-lemma-quasi-coherent-colimit-finite-type}，存在有限型子模
$\mathcal{F} \subset \mathcal{I}_d$，使得所有局部截面
$$
h_{i, j, k}f_{i, j}^{(d - \deg(h_{i, j, k}))/\deg(f_{i, j})}
\quad\text{且}\quad
g_{t, s}
$$
都是 $\mathcal{F}$ 的截面。按构造，$\mathcal{F}$ 即为所求。
\end{proof}

\begin{lemma}
\label{divisors-lemma-conormal-sheaf-section-projective-bundle}
设 $X$ 为概形。设 $\mathcal{E}$ 为拟相干
$\mathcal{O}_X$-模。存在双射
$$
\left\{
\begin{matrix}
\text{sections }\sigma\text{ of the } \\
\text{morphism } \mathbf{P}(\mathcal{E}) \to X
\end{matrix}
\right\}
\leftrightarrow
\left\{
\begin{matrix}
\text{surjections }\mathcal{E} \to \mathcal{L}\text{ 其中} \\
\mathcal{L}\text{ is an invertible }\mathcal{O}_X\text{-模}
\end{matrix}
\right\}
$$
在此情形，$\sigma$ 是闭浸入，并且存在典范同构
$$
\Ker(\mathcal{E} \to \mathcal{L})
\otimes_{\mathcal{O}_X} \mathcal{L}^{\otimes -1}
\longrightarrow
\mathcal{C}_{\sigma(X)/\mathbf{P}(\mathcal{E})}
$$
该双射与同构都和基变换相容。
\end{lemma}

\begin{proof}
回忆，$\pi : \mathbf{P}(\mathcal{E}) \to X$ 是 $\mathcal{E}$
上的对称代数之相对 Proj，见《构造》定义
\ref{constructions-definition-projective-bundle}。因此立即得到截面 $\sigma$
的描述；这来自《构造》引理 \ref{constructions-lemma-apply-relative} 中对
$\mathbf{P}(\mathcal{E})$ 的点函子的描述。
由于 $\pi$ 分离，任意截面都是闭浸入（《构造》引理
\ref{constructions-lemma-relative-proj-separated} 和《概形》引理
\ref{schemes-lemma-section-immersion}）。
设 $U \subset X$ 为仿射开集，并设
$k \in \mathcal{E}(U)$ 与 $s \in \mathcal{E}(U)$ 为局部截面，使得 $k$
在 $\mathcal{L}$ 中映为零，而 $s$ 映为一个生成元 $\overline{s}$
（属于 $\mathcal{L}$）。于是 $f = k/s$ 是
$\mathcal{O}_{\mathbf{P}(\mathcal{E})}$ 的截面，它定义在开邻域
$D_+(s)$ 上；这是 $s(U)$ 在 $\pi^{-1}(U)$ 中的一个开邻域。此外，由于 $k$
在 $\mathcal{L}$ 中映为零，可知 $f$ 是 $s(U)$ 在 $\pi^{-1}(U)$ 中的
理想层的截面。因此可以取像 $\overline{f}$，即 $f$ 在
$\mathcal{C}_{\sigma(X)/\mathbf{P}(\mathcal{E})}(U)$ 中的像。
我们断言：(1) 像 $\overline{f}$ 只依赖于截面 $k$
与 $\overline{s}$，而不依赖于 $s$ 的选取；(2) 以这种方式在 $U$
上得到一个同构（见下文）。一旦 (1) 与 (2) 成立，便可看出该构造与
$U' \to U$ 给出的基变换相容，其中 $U'$ 仿射；这就证明这些局部映射可以
黏合，并且与任意基变换相容。

\medskip\noindent
为证明 (1) 与 (2)，我们把具体情形写出来。即设
$U = \Spec(A)$，并设 $\mathcal{E} \to \mathcal{L}$ 对应于
$A$-模映射 $M \to N$。于是
$k \in K = \Ker(M \to N)$，并且 $s \in M$ 映为生成元
$\overline{s}$（属于 $N$）。因此 $M = K \oplus A s$。于是
$$
\text{Sym}(M) = \text{Sym}(K)[s]
$$
考虑等同 $\text{Sym}(K) \to \text{Sym}(M)_{(s)}$，它由规则
$g \mapsto g/s^n$ 给出，其中 $g \in \text{Sym}^n(K)$。这给出同构
$D_+(s) = \Spec(\text{Sym}(K))$，使得 $\sigma$ 对应于环映射
$\text{Sym}(K) \to A$，它把 $K$ 映为零。通过该同构可见，与 $K$
对应的拟相干模被等同于
$\mathcal{C}_{\sigma(U)/D_+(s)}$，这就证明了 (2)。
最后，假设 $s' = k' + s$，其中 $k' \in K$。则
$$
k/s' = (k/s) (s/s') = (k/s) (s'/s)^{-1} = (k/s) (1 + k'/s)^{-1}
$$
在 $\sigma(U)$ 的一个开邻域内成立，该邻域位于 $D_+(s)$ 中。因此可见
$s'/s$ 的限制等于 $1$；这里限制在 $\sigma(U)$ 上。并且 $k/s'$ 与 $k/s$ 在余法层中
映为同一元素，从而证明了 (1)。
\end{proof}

\section{爆破}
\label{divisors-section-blowing-up}

\noindent
爆破是代数几何中的重要工具。

\begin{definition}
\label{divisors-definition-blow-up}
设 $X$ 为概形。设 $\mathcal{I} \subset \mathcal{O}_X$ 为拟相干理想层，
并设 $Z \subset X$ 为与 $\mathcal{I}$ 对应的闭子概形，见《概形》定义
\ref{schemes-definition-immersion}。所谓{\it $X$ 沿 $Z$ 的爆破}，或
所谓{\it $X$ 在理想层 $\mathcal{I}$ 处的爆破}，是态射
$$
b :
\underline{\text{Proj}}_X
\left(\bigoplus\nolimits_{n \geq 0} \mathcal{I}^n\right)
\longrightarrow
X
$$
爆破的{\it 例外除子}是逆像 $b^{-1}(Z)$。有时把 $Z$ 称为爆破的
{\it 中心}。
\end{definition}

\noindent
我们稍后将看到，例外除子是有效 Cartier 除子。此外，爆破可刻画为
$X$ 上使 $Z$ 的逆像成为有效 Cartier 除子的“最小”概形。

\medskip\noindent
若 $b : X' \to X$ 是 $X$ 在 $Z$ 处的爆破，则常以
$\mathcal{O}_{X'}(n)$ 表示结构层的扭曲。注意，这些都是可逆
$\mathcal{O}_{X'}$-模，并且
$\mathcal{O}_{X'}(n) = \mathcal{O}_{X'}(1)^{\otimes n}$，因为 $X'$
是一个拟相干分次 $\mathcal{O}_X$-代数的相对 Proj，而该代数由次数
$1$ 的部分生成，见《构造》引理
\ref{constructions-lemma-apply-relative}。注意，$\mathcal{O}_{X'}(1)$
相对于 $b$ 甚丰沛；尽管 $b$ 未必有限型，甚至未必拟紧，这是因为 $X'$
自带到 $\mathbf{P}(\mathcal{I})$ 中的闭浸入，见《态射》例
\ref{morphisms-example-very-ample}。

\begin{lemma}
\label{divisors-lemma-blowing-up-affine}
设 $X$ 为概形。设 $\mathcal{I} \subset \mathcal{O}_X$ 为拟相干理想层。
设 $U = \Spec(A)$ 为 $X$ 的仿射开子概形，并设 $I \subset A$ 为与
$\mathcal{I}|_U$ 对应的理想。若 $b : X' \to X$ 是 $X$ 在
$\mathcal{I}$ 处的爆破，则存在典范同构
$$
b^{-1}(U) = \text{Proj}(\bigoplus\nolimits_{d \geq 0} I^d)
$$
把 $b^{-1}(U)$ 等同于 $I$ 在 $A$ 中的 Rees 代数之齐次谱。此外，
$b^{-1}(U)$ 具有由仿射爆破代数 $A[\frac{I}{a}]$ 的谱组成的仿射开覆盖。
\end{lemma}

\begin{proof}
第一条陈述由相对 Proj 的黏合构造直接得出，见《构造》节
\ref{constructions-section-relative-proj-via-glueing}。对 $a \in I$，记
$a^{(1)}$ 为 $a$ 视作次数 $1$ 的元素（在 Rees 代数
$\bigoplus_{n \geq 0} I^n$ 中）。由于这些元素在 $A$ 上生成
Rees 代数，可知 $\text{Proj}(\bigoplus_{d \geq 0} I^d)$ 被仿射开集
$D_{+}(a^{(1)})$ 覆盖。仿射概形 $D_{+}(a^{(1)})$ 是仿射爆破代数
$A' = A[\frac{I}{a}]$ 的谱，见《代数》定义
\ref{algebra-definition-blow-up}。证明完成。
\end{proof}

\begin{lemma}
\label{divisors-lemma-flat-base-change-blowing-up}
\begin{slogan}
爆破与平坦基变换可交换。
\end{slogan}
设 $X_1 \to X_2$ 为概形的平坦态射。设 $Z_2 \subset X_2$ 为闭子概形。
设 $Z_1$ 为 $Z_2$ 在 $X_1$ 中的逆像。设 $X'_i$ 为以 $Z_i$ 为中心对
$X_i$ 的爆破。则存在概形的笛卡尔图
$$
\xymatrix{
X_1' \ar[r] \ar[d] & X_2' \ar[d] \\
X_1 \ar[r] & X_2
}
$$
\end{lemma}

\begin{proof}
设 $\mathcal{I}_2$ 为 $Z_2$ 在 $X_2$ 中的理想层。记给定态射为
$g : X_1 \to X_2$。则理想层 $\mathcal{I}_1$（对应于 $Z_1$）是映射
$g^*\mathcal{I}_2 \to \mathcal{O}_{X_1}$ 的像（由逆像的定义，见
《概形》定义 \ref{schemes-definition-inverse-image-closed-subscheme}）。
由《构造》引理 \ref{constructions-lemma-relative-proj-base-change}，
$X_1 \times_{X_2} X_2'$ 是
$\bigoplus_{n \geq 0} g^*\mathcal{I}_2^n$ 的相对 Proj。由于 $g$ 平坦，
映射 $g^*\mathcal{I}_2^n \to \mathcal{O}_{X_1}$ 为单射，其像为
$\mathcal{I}_1^n$。因此 $X_1 \times_{X_2} X_2' = X_1'$。
\end{proof}

\begin{lemma}
\label{divisors-lemma-blowing-up-gives-effective-Cartier-divisor}
设 $X$ 为概形。设 $Z \subset X$ 为闭子概形。爆破
$b : X' \to X$（即以 $Z$ 为中心对 $X$ 的爆破）具有下列性质：
\begin{enumerate}
\item $b|_{b^{-1}(X \setminus Z)} : b^{-1}(X \setminus Z) \to X \setminus Z$
是同构；
\item 例外除子 $E = b^{-1}(Z)$ 是 $X'$ 上的有效 Cartier 除子；
\item 存在典范同构
$\mathcal{O}_{X'}(-1) = \mathcal{O}_{X'}(E)$。
\end{enumerate}
\end{lemma}

\begin{proof}
由于爆破与到开子概形的限制可交换（引理
\ref{divisors-lemma-flat-base-change-blowing-up}），第一条陈述只不过是说：
有 $X' = X$；这里 $Z = \emptyset$。在此情形，我们在理想层
$\mathcal{I} = \mathcal{O}_X$ 处爆破，结论由《构造》例
\ref{constructions-example-trivial-proj} 得出。

\medskip\noindent
第二条陈述在 $X$ 上是局部的，故可假设 $X$ 仿射。记
$X = \Spec(A)$ 且 $Z = \Spec(A/I)$。由引理
\ref{divisors-lemma-blowing-up-affine}，$X'$ 被仿射爆破代数
$A' = A[\frac{I}{a}]$ 的谱覆盖。此时 $IA' = aA'$，并且按《代数》引理
\ref{algebra-lemma-affine-blowup}，$a$ 在 $A'$ 中的像是非零因子。这证明了
引理，因为 $Z$ 在 $\Spec(A')$ 中的逆像对应于
$\Spec(A'/IA') \subset \Spec(A')$。

\medskip\noindent
考虑典范映射
$\psi_{univ, 1} : b^*\mathcal{I} \to \mathcal{O}_{X'}(1)$，见《构造》定义
\ref{constructions-definition-relative-proj} 后面的讨论。我们断言，它分解为同构
$\mathcal{I}_E \to \mathcal{O}_{X'}(1)$（这证明最后一条断言）。
具体而言，在上面提到的、与爆破代数
$A' = A[\frac{I}{a}]$ 对应的仿射开集上，$\psi_{univ, 1}$ 对应于
$A'$-模映射
$$
I \otimes_A A'
\longrightarrow
\left(\Big(\bigoplus\nolimits_{d \geq 0} I^d\Big)_{a^{(1)}}\right)_1
$$
其中 $a^{(1)}$ 如《代数》定义 \ref{algebra-definition-blow-up} 所述。
我们略去验证它正是映射 $I \otimes_A A' \to IA' = aA'$。
\end{proof}

\begin{lemma}[爆破的泛性质]
\label{divisors-lemma-universal-property-blowing-up}
设 $X$ 为概形。设 $Z \subset X$ 为闭子概形。设 $\mathcal{C}$ 为
$(\Sch/X)$ 的如下全子范畴：其对象是满足下列条件的 $Y \to X$，即
$Z$ 在 $Y$ 中的逆像是有效 Cartier 除子。则爆破
$b : X' \to X$（即以 $Z$ 为中心对 $X$ 作爆破）是 $\mathcal{C}$ 的终对象。
\end{lemma}

\begin{proof}
由引理 \ref{divisors-lemma-blowing-up-gives-effective-Cartier-divisor}，
$b : X' \to X$ 是 $\mathcal{C}$ 的对象。设 $f : Y \to X$ 是
$\mathcal{C}$ 的对象。必须证明存在唯一态射 $Y \to X'$，且该态射在 $X$ 上。令
$D = f^{-1}(Z)$。设 $\mathcal{I} \subset \mathcal{O}_X$ 为 $Z$ 的理想层，
并设 $\mathcal{I}_D$ 为 $D$ 的理想层。于是
$f^*\mathcal{I} \to \mathcal{I}_D$ 是到一个可逆
$\mathcal{O}_Y$-模的满射。它扩张为映射
$\psi : \bigoplus f^*\mathcal{I}^d \to \bigoplus \mathcal{I}_D^d$，
这是分次 $\mathcal{O}_Y$-代数映射。（注意，有
$\mathcal{I}_D^d = \mathcal{I}_D^{\otimes d}$，因为 $D$ 是有效 Cartier
除子。）由《构造》节
\ref{constructions-section-relative-proj} 的内容，三元组
$(1, f : Y \to X, \psi)$ 定义了态射 $Y \to X'$，它在 $X$ 上。其限制
$$
Y \setminus D \longrightarrow X' \setminus b^{-1}(Z) = X \setminus Z
$$
是唯一的。由引理 \ref{divisors-lemma-complement-effective-Cartier-divisor}，开集
$Y \setminus D$ 在 $Y$ 中概形论稠密。因此由《态射》引理
\ref{morphisms-lemma-equality-of-morphisms}，态射 $Y \to X'$ 唯一
（并且由《构造》引理
\ref{constructions-lemma-relative-proj-separated}，$b$ 也分离）。
\end{proof}

\begin{lemma}
\label{divisors-lemma-characterize-affine-blowup}
设 $b : X' \to X$ 为概形 $X$ 沿闭子概形 $Z$ 的爆破。设
$U = \Spec(A)$ 为 $X$ 的仿射开集，并设 $I \subset A$ 为与
$Z \cap U$ 对应的理想。设 $a \in I$，并设 $x' \in X'$ 为映到
$U$ 中某点的点。则 $x'$ 属于仿射开集
$U' = \Spec(A[\frac{I}{a}])$，当且仅当 $a$ 在
$\mathcal{O}_{X', x'}$ 中的像截出例外除子。
\end{lemma}

\begin{proof}
由于 $U'$ 上的例外除子由 $a$ 在
$A' = A[\frac{I}{a}]$ 中的像截出，一个方向是清楚的。反之，假设
$a$ 在 $\mathcal{O}_{X', x'}$ 中的像截出 $E$。由于 $I$ 的每个元素在
定义 $E$ 的理想中，这是在 $b^{-1}(U)$ 上而言；可知 $I$ 的元素都可被
$a$ 整除（在 $\mathcal{O}_{X', x'}$ 中）。因此对 $f \in I^n$，可写成
$f = \psi(f) a^n$，其中 $\psi(f) \in \mathcal{O}_{X', x'}$ 为某个元素。
注意，由于 $a$ 映到 $\mathcal{O}_{X', x'}$ 的非零因子，元素
$\psi(f)$ 由此唯一确定。于是定义
$$
A' \longrightarrow \mathcal{O}_{X', x'},\quad
f/a^n \longmapsto \psi(f)
$$
这里使用《代数》定义 \ref{divisors-definition-blow-up} 后给出的爆破代数描述。
上面提到的唯一性表明，这是 $A$-代数同态。它给出态射
$\Spec(\mathcal{O}_{X', x"}) \to \Spec(A') = U'$。由爆破的泛性质
（引理 \ref{divisors-lemma-universal-property-blowing-up}），这是 $X'$ 上的态射，
当然这蕴含 $x' \in U'$。
\end{proof}

\begin{lemma}
\label{divisors-lemma-blow-up-effective-Cartier-divisor}
设 $X$ 为概形。设 $Z \subset X$ 为有效 Cartier 除子。$X$ 在 $Z$
处的爆破是 $X$ 的恒等态射。
\end{lemma}

\begin{proof}
由爆破的泛性质（引理 \ref{divisors-lemma-universal-property-blowing-up}）立即得出。
\end{proof}

\begin{lemma}
\label{divisors-lemma-blow-up-reduced-scheme}
设 $X$ 为概形。设 $\mathcal{I} \subset \mathcal{O}_X$ 为拟相干理想层。
若 $X$ 约化，则爆破 $X'$（即 $X$ 在 $\mathcal{I}$ 处的爆破）约化。
\end{lemma}

\begin{proof}
结合引理 \ref{divisors-lemma-blowing-up-affine} 与《代数》引理
\ref{algebra-lemma-blowup-reduced}。
\end{proof}

\begin{lemma}
\label{divisors-lemma-blow-up-integral-scheme}
设 $X$ 为概形。设 $\mathcal{I} \subset \mathcal{O}_X$ 为非零拟相干理想层。
若 $X$ 整，则爆破 $X'$（即 $X$ 在 $\mathcal{I}$ 处的爆破）整。
\end{lemma}

\begin{proof}
结合引理 \ref{divisors-lemma-blowing-up-affine} 与《代数》引理
\ref{algebra-lemma-blowup-domain}。
\end{proof}

\begin{lemma}
\label{divisors-lemma-blow-up-and-irreducible-components}
设 $X$ 为概形。设 $Z \subset X$ 为闭子概形。设 $b : X' \to X$ 为
$X$ 沿 $Z$ 的爆破。则 $b$ 诱导双射：从 $X'$ 的不可约分支的泛点集合，
到 $X$ 的不可约分支中那些不属于 $Z$ 的泛点集合。
\end{lemma}

\begin{proof}
例外除子 $E \subset X'$ 是有效 Cartier 除子，并且
$X' \setminus E \to X \setminus Z$ 是同构，见引理
\ref{divisors-lemma-blowing-up-gives-effective-Cartier-divisor}。因此只需证明下述事实：
给定有效 Cartier 除子 $D \subset S$，其中 $S$ 是概形；$S$ 的不可约分支的泛点
没有一个属于 $D$。为此，可以用仿射开覆盖的成员替换 $S$。所以由引理
\ref{divisors-lemma-characterize-effective-Cartier-divisor}，可假设
$S = \Spec(A)$ 且 $D = V(f)$，其中 $f \in A$ 为非零因子。于是必须证明
$f$ 不属于任何极小素理想 $\mathfrak p \subset A$。否则，$f$ 将映为
$R_\mathfrak p$ 的极大理想中所含的非零因子，这与《代数》引理
\ref{algebra-lemma-minimal-prime-reduced-ring} 矛盾。
\end{proof}

\begin{lemma}
\label{divisors-lemma-blow-up-pullback-effective-Cartier}
设 $X$ 为概形。设 $b : X' \to X$ 为 $X$ 沿某闭子概形的爆破。拉回
$b^{-1}D$ 对所有有效 Cartier 除子 $D \subset X$ 都有定义，并且亚纯函数沿
$b$ 的拉回也有定义（定义
\ref{divisors-definition-pullback-effective-Cartier-divisor} 与
\ref{divisors-definition-pullback-meromorphic-sections}）。
\end{lemma}

\begin{proof}
由引理 \ref{divisors-lemma-blowing-up-affine} 与
\ref{divisors-lemma-characterize-effective-Cartier-divisor}，这约化为下述代数事实：
设 $A$ 为环，$I \subset A$ 为理想，$a \in I$，且 $x \in A$ 为非零因子。
则 $x$ 在 $A[\frac{I}{a}]$ 中的像是非零因子。事实上，假设
$x (y/a^n) = 0$ 在 $A[\frac{I}{a}]$ 中成立。则
$a^mxy = 0$ 在 $A$ 中成立，其中 $m$ 为某个整数。因此 $a^my = 0$，
因为 $x$ 是非零因子。
所以如所需，$y/a^n$ 在 $A[\frac{I}{a}]$ 中为零。
\end{proof}

\begin{lemma}
\label{divisors-lemma-blowing-up-two-ideals}
设 $X$ 为概形。设 $\mathcal{I}, \mathcal{J} \subset \mathcal{O}_X$
为拟相干理想层。设 $b : X' \to X$ 为 $X$ 在 $\mathcal{I}$ 处的爆破。
设 $b' : X'' \to X'$ 为 $X'$ 在
$b^{-1}\mathcal{J} \mathcal{O}_{X'}$ 处的爆破。则 $X'' \to X$
典范同构于 $X$ 在 $\mathcal{I}\mathcal{J}$ 处的爆破。
\end{lemma}

\begin{proof}
设 $E \subset X'$ 为 $b$ 的例外除子；由引理
\ref{divisors-lemma-blowing-up-gives-effective-Cartier-divisor}，它是有效 Cartier 除子。
由引理 \ref{divisors-lemma-blow-up-pullback-effective-Cartier}，
$(b')^{-1}E$ 是 $X''$ 上的有效 Cartier 除子。设 $E' \subset X''$ 为
$b'$ 的例外除子（它也是有效 Cartier 除子）。考虑有效 Cartier 除子
$E'' = E' + (b')^{-1}E$。按构造，$E''$ 的理想为
$(b \circ b')^{-1}\mathcal{I} (b \circ b')^{-1}\mathcal{J} \mathcal{O}_{X''}$。
因此由引理 \ref{divisors-lemma-universal-property-blowing-up}，存在从 $X''$ 到爆破
$c : Y \to X$ 的典范态射，其中后者是 $X$ 在
$\mathcal{I}\mathcal{J}$ 处的爆破。反之，由于
$\mathcal{I}\mathcal{J}$ 拉回为可逆理想，可知
$c^{-1}\mathcal{I}\mathcal{O}_Y$ 定义了有效 Cartier 除子，见引理
\ref{divisors-lemma-sum-closed-subschemes-effective-Cartier}。所以由引理
\ref{divisors-lemma-universal-property-blowing-up}，存在态射
$c' : Y \to X'$，它在 $X$ 上。于是
$(c')^{-1}b^{-1}\mathcal{J}\mathcal{O}_Y = c^{-1}\mathcal{J}\mathcal{O}_Y$，
它也定义了有效 Cartier 除子。因此存在态射
$c'' : Y \to X''$，它在 $X'$ 上。我们略去验证该态射与先前构造的态射
$X'' \to Y$ 互为逆态射。
\end{proof}

\begin{lemma}
\label{divisors-lemma-blowing-up-projective}
设 $X$ 为概形。设 $\mathcal{I} \subset \mathcal{O}_X$ 为拟相干理想层。
设 $b : X' \to X$ 为 $X$ 在理想层 $\mathcal{I}$ 处的爆破。若
$\mathcal{I}$ 有限型，则
\begin{enumerate}
\item $b : X' \to X$ 是射影态射；
\item $\mathcal{O}_{X'}(1)$ 是相对于 $b$ 丰沛的可逆层。
\end{enumerate}
\end{lemma}

\begin{proof}
分次 $\mathcal{O}_X$-代数的满射
$$
\text{Sym}_{\mathcal{O}_X}^*(\mathcal{I})
\longrightarrow
\bigoplus\nolimits_{d \geq 0} \mathcal{I}^d
$$
由《构造》引理
\ref{constructions-lemma-surjective-generated-degree-1-map-relative-proj}
定义闭浸入
$$
X' = \underline{\text{Proj}}_X (\bigoplus\nolimits_{d \geq 0} \mathcal{I}^d)
\longrightarrow
\mathbf{P}(\mathcal{I}).
$$
因此由《态射》定义 \ref{morphisms-definition-projective}，$b$ 是射影态射。
第二条陈述例如由《态射》引理
\ref{morphisms-lemma-characterize-relatively-ample} 中对相对丰沛可逆层的刻画
得出。略去若干细节。
\end{proof}

\begin{lemma}
\label{divisors-lemma-composition-finite-type-blowups}
\begin{slogan}
爆破的复合仍为爆破。
\end{slogan}
设 $X$ 为拟紧拟分离概形。设 $Z \subset X$ 为有限呈示闭子概形。
设 $b : X' \to X$ 为以 $Z$ 为中心的爆破。设 $Z' \subset X'$ 为有限呈示
闭子概形。设 $X'' \to X'$ 为以 $Z'$ 为中心的爆破。则存在有限呈示闭子概形
$Y \subset X$，使得
\begin{enumerate}
\item 在集合论意义下，$Y = Z \cup b(Z')$；
\item 复合 $X'' \to X$ 同构于 $X$ 在 $Y$ 处的爆破。
\end{enumerate}
\end{lemma}

\begin{proof}
条件 $Z \to X$ 有限呈示意味着 $Z$ 由有限型拟相干理想层
$\mathcal{I} \subset \mathcal{O}_X$ 截出，见《态射》引理
\ref{morphisms-lemma-closed-immersion-finite-presentation}。写成
$\mathcal{A} = \bigoplus_{n \geq 0} \mathcal{I}^n$，于是
$X' = \underline{\text{Proj}}(\mathcal{A})$。注意，由《性质》引理
\ref{properties-lemma-quasi-coherent-finite-type-ideals}，
$X \setminus Z$ 是 $X$ 的拟紧开集。由于
$b^{-1}(X \setminus Z) \to X \setminus Z$ 是同构（引理
\ref{divisors-lemma-blowing-up-gives-effective-Cartier-divisor}），同一结果表明
$b^{-1}(X \setminus Z) \setminus Z'$ 是 $X'$ 的拟紧开集。因此
$U = X \setminus (Z \cup b(Z'))$ 是 $X$ 的拟紧开集。由引理
\ref{divisors-lemma-closed-subscheme-proj-finite-type}，存在 $d > 0$ 和有限型
$\mathcal{O}_X$-子模 $\mathcal{F} \subset \mathcal{I}^d$，使得
$Z' = \underline{\text{Proj}}(\mathcal{A}/\mathcal{F}\mathcal{A})$，
并且 $\mathcal{I}^d/\mathcal{F}$ 的支集包含在 $X \setminus U$ 中。

\medskip\noindent
由于 $\mathcal{F} \subset \mathcal{I}^d$ 是 $\mathcal{O}_X$-子模，
可以把 $\mathcal{F} \subset \mathcal{I}^d \subset \mathcal{O}_X$
视为 $X$ 上的有限型拟相干理想层。为避免混淆，将其记作
$\mathcal{J} \subset \mathcal{O}_X$。由于
$\mathcal{I}^d / \mathcal{J}$ 与 $\mathcal{O}/\mathcal{I}^d$
的支集都在 $X \setminus U$ 上，可知 $V(\mathcal{J})$ 包含在
$X \setminus U$ 中。反之，由于 $\mathcal{J} \subset \mathcal{I}^d$，
可知 $Z \subset V(\mathcal{J})$。在
$X \setminus Z \cong X' \setminus b^{-1}(Z)$ 上，理想层
$\mathcal{J}$ 截出 $Z'$（见下面的显示公式）。所以
$V(\mathcal{J})$ 等于 $Z \cup b(Z')$。从而在集合论意义下，也有
$V(\mathcal{I}\mathcal{J}) = Z \cup b(Z')$。此外，
$\mathcal{I}\mathcal{J}$ 作为两个有限型理想之积仍有限型。我们断言
$X'' \to X$ 同构于 $X$ 在 $\mathcal{I}\mathcal{J}$ 处的爆破；令
$Y = V(\mathcal{I}\mathcal{J})$ 即完成引理的证明。

\medskip\noindent
首先回忆，$X$ 在 $\mathcal{I}\mathcal{J}$ 处的爆破与 $X'$ 在
$b^{-1}\mathcal{J} \mathcal{O}_{X'}$ 处的爆破相同，见引理
\ref{divisors-lemma-blowing-up-two-ideals}。因此只需证明 $X'$ 在
$b^{-1}\mathcal{J} \mathcal{O}_{X'}$ 处的爆破与 $X'$ 在 $Z'$ 处的爆破
一致。我们将证明
$$
b^{-1}\mathcal{J} \mathcal{O}_{X'} = \mathcal{I}_E^d \mathcal{I}_{Z'}
$$
作为 $X''$ 上的理想层成立。这将证明所需结论，因为
$\mathcal{I}_E^d$ 截出有效 Cartier 除子 $dE$，并且可以使用引理
\ref{divisors-lemma-blow-up-effective-Cartier-divisor} 与
\ref{divisors-lemma-blowing-up-two-ideals}。

\medskip\noindent
为验证所显示的理想等式，可以作局部计算。采用引理
\ref{divisors-lemma-blowing-up-affine} 中的记号 $A$、$I$、$a \in I$，可知
$\mathcal{F}$ 对应于 $R$-子模 $M \subset I^d$，后者同构地映到理想
$J \subset R$。条件
$Z' = \underline{\text{Proj}}(\mathcal{A}/\mathcal{F}\mathcal{A})$
意味着 $Z' \cap \Spec(A[\frac{I}{a}])$ 由元素
$m/a^d$（其中 $m \in M$）生成的理想截出。设元素 $m \in M$ 对应于函数
$f \in J$。则在仿射爆破代数 $A' = A[\frac{I}{a}]$ 中，有
$f = (a^dm)/a^d = a^d (m/a^d)$。因此等式成立。
\end{proof}

\section{严格变换}
\label{divisors-section-strict-transform}

\noindent
本节简要讨论爆破下的严格变换。设 $S$ 为概形，设 $Z \subset S$ 为闭子概形。
设 $b : S' \to S$ 为 $S$ 在 $Z$ 处的爆破，并记 $E \subset S'$ 为例外除子
$E = b^{-1}Z$。下面常常考虑概形 $X$（它在 $S$ 上），并作笛卡尔图
$$
\xymatrix{
\text{pr}_{S'}^{-1}E \ar[r] \ar[d] &
X \times_S S' \ar[r]_-{\text{pr}_X} \ar[d]_{\text{pr}_{S'}} &
X \ar[d]^f \\
E \ar[r] & S' \ar[r] & S
}
$$
由于 $E$ 是有效 Cartier 除子（引理
\ref{divisors-lemma-blowing-up-gives-effective-Cartier-divisor}），可知
$\text{pr}_{S'}^{-1}E \subset X \times_S S'$ 局部为主理想所定义
（引理 \ref{divisors-lemma-pullback-locally-principal}）。因此
$\text{pr}_{S'}^{-1}E$ 在 $X \times_S S'$ 中的补集是逆紧的
（引理 \ref{divisors-lemma-complement-locally-principal-closed-subscheme}）。
所以，对拟相干 $\mathcal{O}_{X \times_S S'}$-模 $\mathcal{G}$，
由支集在 $\text{pr}_{S'}^{-1}E$ 上的截面组成的子层是拟相干子模，见《性质》引理
\ref{properties-lemma-sections-supported-on-closed-subset}。若 $\mathcal{G}$
是拟相干代数层，例如 $\mathcal{G} = \mathcal{O}_{X \times_S S'}$，
则该子层是 $\mathcal{G}$ 的理想。

\begin{definition}
\label{divisors-definition-strict-transform}
沿用上面的 $Z \subset S$ 与 $f : X \to S$。
\begin{enumerate}
\item 给定拟相干 $\mathcal{O}_X$-模 $\mathcal{F}$，$\mathcal{F}$
关于 $S$ 在 $Z$ 处之爆破的{\it 严格变换} $\mathcal{F}'$，是
$\text{pr}_X^*\mathcal{F}$ 对支集在 $\text{pr}_{S'}^{-1}E$ 上的截面子模
取商。
\item $X$ 的{\it 严格变换}是闭子概形
$X' \subset X \times_S S'$，它由 $\mathcal{O}_{X \times_S S'}$
中支集在 $\text{pr}_{S'}^{-1}E$ 上的截面组成的拟相干理想截出。
\end{enumerate}
\end{definition}

\noindent
注意，沿一次爆破取严格变换依赖于用来作爆破的闭子概形
（而不只依赖态射 $S' \to S$）。这个概念常用于 $S$ 的闭子概形。
事实表明，$X$ 的严格变换是 $X$ 的一次爆破。

\begin{lemma}
\label{divisors-lemma-strict-transform}
在定义 \ref{divisors-definition-strict-transform} 的情形中：
\begin{enumerate}
\item 严格变换 $X'$（属于 $X$）是 $X$ 的爆破，其中心是闭子概形
$f^{-1}Z$（位于 $X$ 中）；
\item 对拟相干 $\mathcal{O}_X$-模 $\mathcal{F}$，严格变换
$\mathcal{F}'$ 典范同构于沿 $X' \to X \times_S S'$ 的下述正像：
$\mathcal{F}$ 关于爆破 $X' \to X$ 的严格变换。
\end{enumerate}
\end{lemma}

\begin{proof}
设 $X'' \to X$ 为 $X$ 在 $f^{-1}Z$ 处的爆破。由爆破的泛性质
（引理 \ref{divisors-lemma-universal-property-blowing-up}），存在交换图
$$
\xymatrix{
X'' \ar[r] \ar[d] & X \ar[d] \\
S' \ar[r] & S
}
$$
从而有态射 $X'' \to X \times_S S'$。所以第一条断言是说，该态射是以
$X'$ 为像的闭浸入。问题在 $X$ 上是局部的，故可假设 $X$ 与 $S$ 都仿射。
记 $S = \Spec(A)$、$X = \Spec(B)$，并设 $Z$ 由理想
$I \subset A$ 截出。置 $J = IB$。映射
$B \otimes_A \bigoplus_{n \geq 0} I^n \to \bigoplus_{n \geq 0} J^n$
定义了闭浸入 $X'' \to X \times_S S'$，见《构造》引理
\ref{constructions-lemma-base-change-map-proj} 与
\ref{constructions-lemma-surjective-graded-rings-generated-degree-1-map-proj}。
我们略去验证该态射与上面由泛性质构造的态射相同。选取 $a \in I$，它对应于
仿射开集 $\Spec(A[\frac{I}{a}]) \subset S'$，见引理
\ref{divisors-lemma-blowing-up-affine}。$\Spec(A[\frac{I}{a}])$ 在严格变换
$X'$ 中的逆像（这里严格变换属于 $X$）是下述环的谱：
$$
B' = (B \otimes_A A[\textstyle{\frac{I}{a}}])/a\text{-power-torsion}
$$
见《性质》引理
\ref{properties-lemma-sections-supported-on-closed-subset}。另一方面，令
$b \in J$ 为 $a$ 的像，则 $\Spec(B[\frac{J}{b}])$ 是
$\Spec(A[\frac{I}{a}])$ 在 $X''$ 中的逆像。由《代数》引理
\ref{algebra-lemma-blowup-base-change}，开集
$\Spec(B[\frac{J}{b}])$ 同构地映到开子概形
$\text{pr}_{S'}^{-1}(\Spec(A[\frac{I}{a}]))$；后者位于 $X'$ 中。因此
$X'' \to X'$ 是同构。

\medskip\noindent
在上述记号下，设 $\mathcal{F}$ 对应于 $B$-模 $N$。$\mathcal{F}$
的严格变换对应于 $B \otimes_A A[\frac{I}{a}]$-模
$$
N' = (N \otimes_A A[\textstyle{\frac{I}{a}}])/a\text{-power-torsion}
$$
见《性质》引理
\ref{properties-lemma-sections-supported-on-closed-subset}。
$\mathcal{F}$ 关于 $X$ 在 $f^{-1}Z$ 处之爆破的严格变换，对应于
$B[\frac{J}{b}]$-模
$N \otimes_B B[\frac{J}{b}]/b\text{-power-torsion}$。完全仿照上面即可证明
这两个模同构。略去细节。
\end{proof}

\begin{lemma}
\label{divisors-lemma-strict-transform-flat}
在定义 \ref{divisors-definition-strict-transform} 的情形中：
\begin{enumerate}
\item 如果 $X$ 平坦于 $S$，至少在映到 $Z$ 的每一点处如此，则
$X$ 的严格变换等于基变换 $X \times_S S'$。
\item 设 $\mathcal{F}$ 为拟相干 $\mathcal{O}_X$-模。
如果 $\mathcal{F}$ 平坦于 $S$，至少在映到 $Z$ 的每一点处如此，则
严格变换 $\mathcal{F}'$（属于 $\mathcal{F}$）等于拉回
$\text{pr}_X^*\mathcal{F}$。
\end{enumerate}
\end{lemma}

\begin{proof}
我们证明第 (2) 条；由 $S$ 上概形的严格变换之定义，它蕴含第 (1) 条。
问题在 $X$ 上是局部的。因此可假设 $S = \Spec(A)$、$X = \Spec(B)$，并且
$\mathcal{F}$ 对应于 $B$-模 $N$。于是，$\mathcal{F}'$ 在开集
$\Spec(B \otimes_A A[\frac{I}{a}])$ 上（该开集位于 $X \times_S S'$ 中）对应于模
$$
N' = (N \otimes_A A[\textstyle{\frac{I}{a}}])/a\text{-power-torsion}
$$
见《性质》引理
\ref{properties-lemma-sections-supported-on-closed-subset}。
所以只须证明 $a$-幂挠在 $N \otimes_A A[\frac{I}{a}]$ 中为零。
设 $y \in N \otimes_A A[\frac{I}{a}]$ 且 $a^n y = 0$。如果
$\mathfrak q \subset B$ 是素理想且 $a \not \in \mathfrak q$，那么 $y$ 在
$(N \otimes_A A[\frac{I}{a}])_\mathfrak q$ 中的像为零。另一方面，如果
$a \in \mathfrak q$，则 $N_\mathfrak q$ 是平坦 $A$-模，因而
$N_\mathfrak q \otimes_A A[\frac{I}{a}]
=(N \otimes_A A[\frac{I}{a}])_\mathfrak q$
没有 $a$-幂挠（因为 $A[\frac{I}{a}]$ 没有）。故 $y$ 在这个局部化中的像
也为零。由《代数》引理 \ref{algebra-lemma-characterize-zero-local}，可知
$y$ 为零。
\end{proof}

\begin{lemma}
\label{divisors-lemma-strict-transform-affine}
设 $S$ 为概形。设 $Z \subset S$ 为闭子概形。设 $b : S' \to S$ 为
沿 $Z$ 对 $S$ 作的爆破。设 $g : X \to Y$ 为 $S$ 上概形之间的仿射态射。
设 $\mathcal{F}$ 为 $X$ 上的拟相干层。设
$g' : X \times_S S' \to Y \times_S S'$ 为 $g$ 的基变换。
设 $\mathcal{F}'$ 为 $\mathcal{F}$ 关于 $b$ 的严格变换。那么
$g'_*\mathcal{F}'$ 是 $g_*\mathcal{F}$ 的严格变换。
\end{lemma}

\begin{proof}
由《概形上同调》引理 \ref{coherent-lemma-affine-base-change}，有
$g'_*\text{pr}_X^*\mathcal{F} = \text{pr}_Y^*g_*\mathcal{F}$。
设 $\mathcal{K} \subset \text{pr}_X^*\mathcal{F}$ 为由支集在
$Z$ 的逆像上的截面组成的子层；该逆像位于 $X \times_S S'$ 中。由《性质》引理
\ref{properties-lemma-push-sections-supported-on-closed-subset}，正像
$g'_*\mathcal{K}$ 是 $\text{pr}_Y^*g_*\mathcal{F}$ 的截面子层，
其支集在 $Z$ 的逆像上；该逆像位于 $Y \times_S S'$ 中。由于 $g'$ 是仿射的
（《态射》引理 \ref{morphisms-lemma-base-change-affine}），函子 $g'_*$ 正合，
结论随即成立。
\end{proof}

\begin{lemma}
\label{divisors-lemma-strict-transform-different-centers}
设 $S$ 为概形。设 $Z \subset S$ 为闭子概形。设 $D \subset S$ 为
有效 Cartier 除子。设 $Z' \subset S$ 为由 $Z$ 与 $D$ 的理想层之积
截出的闭子概形。设 $S' \to S$ 为 $S$ 沿 $Z$ 的爆破。
\begin{enumerate}
\item $S$ 沿 $Z'$ 的爆破同构于 $S' \to S$。
\item 设 $f : X \to S$ 为概形态射，设 $\mathcal{F}$ 为拟相干
$\mathcal{O}_X$-模。如果 $\mathcal{F}$ 没有支集在 $f^{-1}D$ 上的非零
局部截面，那么 $\mathcal{F}$ 关于沿 $Z$ 爆破的严格变换，与
$\mathcal{F}$ 关于 $S$ 沿 $Z'$ 爆破的严格变换相同。
\end{enumerate}
\end{lemma}

\begin{proof}
第一条由引理 \ref{divisors-lemma-blowing-up-two-ideals} 和
\ref{divisors-lemma-blow-up-effective-Cartier-divisor} 合并即得。
利用引理 \ref{divisors-lemma-blowing-up-affine}，第二条转化为下面的代数问题。
设 $A$ 为环，$I \subset A$ 为理想，$x \in A$ 为非零因子，且
$a \in I$。设 $M$ 为 $A$-模，其 $x$-挠为零。要证：
$a$-幂挠在 $M \otimes_A A[\frac{I}{a}]$ 中等于 $xa$-幂挠。
理由如下：映射 $A \to A[\frac{I}{a}]$ 的核与余核都是 $a$-幂挠，
所以将 $a$ 可逆化之后，该映射成为同构。因此
$M \to M \otimes_A A[\frac{I}{a}]$ 的核与余核也都是 $a$-幂挠。
这就蕴含所求结论。
\end{proof}

\begin{lemma}
\label{divisors-lemma-strict-transform-composition-blowups}
设 $S$ 为概形。设 $Z \subset S$ 为闭子概形。设 $b : S' \to S$ 为
中心是 $Z$ 的爆破。设 $Z' \subset S'$ 为闭子概形。设 $S'' \to S'$ 为
中心是 $Z'$ 的爆破。设 $Y \subset S$ 为闭子概形，使得集合论上
$Y = Z \cup b(Z')$，并且复合 $S'' \to S$ 同构于 $S$ 沿 $Y$ 的爆破。
在这种情形中，给定任意概形 $X$（定义在 $S$ 上）以及
$\mathcal{F} \in \QCoh(\mathcal{O}_X)$，有：
\begin{enumerate}
\item $\mathcal{F}$ 关于 $S$ 沿 $Y$ 爆破的严格变换，等于下述严格变换：
对于 $S'' \to S'$ 沿 $Z'$ 的爆破，取 $\mathcal{F}$ 关于
$S' \to S$（即 $S$ 沿 $Z$ 的爆破）的严格变换之严格变换；并且
\item $X$ 关于 $S$ 沿 $Y$ 爆破的严格变换，等于下述严格变换：
对于 $S'' \to S'$ 沿 $Z'$ 的爆破，取 $X$ 关于
$S' \to S$（即 $S$ 沿 $Z$ 的爆破）的严格变换之严格变换。
\end{enumerate}
\end{lemma}

\begin{proof}
设 $\mathcal{F}'$ 为 $\mathcal{F}$ 关于 $S' \to S$（即 $S$ 沿 $Z$ 的爆破）
的严格变换。设 $\mathcal{F}''$ 为 $\mathcal{F}'$ 关于
$S'' \to S'$（即 $S'$ 沿 $Z'$ 的爆破）的严格变换。
设 $\mathcal{G}$ 为 $\mathcal{F}$ 关于 $S'' \to S$（即 $S$ 沿 $Y$ 的爆破）
的严格变换。我们还如下标记各态射：
$$
\xymatrix{
X \times_S S'' \ar[r]_q \ar[d]^{f''} &
X \times_S S' \ar[r]_p \ar[d]^{f'} &
X \ar[d]^f \\
S'' \ar[r] & S' \ar[r] & S
}
$$
由定义有满射 $p^*\mathcal{F} \to \mathcal{F}'$ 和满射
$q^*\mathcal{F}' \to \mathcal{F}''$；利用 $q^*$ 的右正合性，可将二者合成满射
$(p \circ q)^*\mathcal{F} \to \mathcal{F}''$。另有满射
$(p \circ q)^*\mathcal{F} \to \mathcal{G}$。所以只须证明这两个满射的核相同。

\medskip\noindent
满射 $p^*\mathcal{F} \to \mathcal{F}'$ 的核支集在
$(f \circ p)^{-1}Z$ 上，故这个映射在其补集的各点处是同构。因此
$q^*p^*\mathcal{F} \to q^*\mathcal{F}'$ 的核支集在
$(f \circ p \circ q)^{-1}Z$ 上。映射
$q^*\mathcal{F}' \to \mathcal{F}''$ 的核支集在 $(f' \circ q)^{-1}Z'$ 上。
合在一起可知，$(p \circ q)^*\mathcal{F} \to \mathcal{F}''$ 的核支集在
$(f \circ p \circ q)^{-1}Z \cup (f' \circ q)^{-1}Z' =
(f \circ p \circ q)^{-1}Y$ 上。由 $\mathcal{G}$ 的构造可得分解
$(p \circ q)^*\mathcal{F} \to \mathcal{F}'' \to \mathcal{G}$。
为完成证明，只须说明 $\mathcal{F}''$ 没有支集在
$(f \circ p \circ q)^{-1}(Y) =
(f \circ p \circ q)^{-1}Z \cup (f' \circ q)^{-1}Z'$ 上的非零（局部）截面。
把引理 \ref{divisors-lemma-strict-transform-different-centers} 应用于
$\mathcal{F}'$（它在 $X \times_S S'$ 上，而后者定义在 $S'$ 上）、
闭子概形 $Z'$ 以及有效 Cartier 除子 $b^{-1}Z$，即得此结论。
\end{proof}

\begin{lemma}
\label{divisors-lemma-strict-transform-universally-injective}
在定义 \ref{divisors-definition-strict-transform} 的情形中。假设
$$
0 \to \mathcal{F}_1 \to \mathcal{F}_2 \to \mathcal{F}_3 \to 0
$$
是 $X$ 上拟相干层的正合列，并且在任意基变换 $T \to S$ 后仍正合。
那么，严格变换 $\mathcal{F}_i'$（关于任意爆破 $S' \to S$）也组成短正合列
$0 \to \mathcal{F}'_1 \to \mathcal{F}'_2 \to \mathcal{F}'_3 \to 0$。
\end{lemma}

\begin{proof}
可在 $S$ 与 $X$ 上局部化，并假设二者都仿射。然后可把
$\mathcal{F}_i$ 正像到 $S$，见引理 \ref{divisors-lemma-strict-transform-affine}。
可以假设我们的爆破是态射 $1 : S \to S$，它与有效 Cartier 除子
$D \subset S$ 相伴。此时，代数表述如下：设 $A$ 为环，并设
$0 \to M_1 \to M_2 \to M_3 \to 0$ 是 $A$-模的泛正合列。
令 $a\in A$。则序列
$$
0 \to
M_1/a\text{-幂挠子} \to
M_2/a\text{-幂挠子} \to
M_3/a\text{-幂挠子} \to 0
$$
也正合。事实上，最后一个映射的满射性与第一个映射的单射性是立即的。
问题在于中间位置的正合性。设 $x \in M_2$ 在
$M_3/a\text{-幂挠子}$ 中的像为零。则
$y = a^n x \in M_1$，且这对某个 $n$ 成立。于是 $y$ 在 $M_2/a^nM_2$ 中的像为零。
由于 $M_1 \to M_2$ 是泛单射，可知 $y$ 在 $M_1/a^nM_1$ 中的像为零。
因而有 $y = a^n z$，其中 $z \in M_1$ 是某个元素。所以
$a^n(x - y) = 0$。故 $y$ 映到 $x$ 的类；该类位于
$M_2/a\text{-幂挠子}$ 中，
这正是所需结论。
\end{proof}

\section{容许爆破}
\label{divisors-section-admissible-blowups}

\noindent
为了对所用的爆破有更多控制，我们引入下面的标准术语。

\begin{definition}
\label{divisors-definition-admissible-blowup}
设 $X$ 为概形。设 $U \subset X$ 为开子概形。称态射
$X' \to X$ 为一个{\it $U$-容许爆破}，如果存在有限呈示的闭浸入
$Z \to X$，其中 $Z$ 与 $U$ 不交，并且 $X'$ 同构于 $X$ 沿 $Z$ 的爆破。
\end{definition}

\noindent
回忆，$Z \to X$ 有限呈示，当且仅当理想层
$\mathcal{I}_Z \subset \mathcal{O}_X$ 为有限型，见《态射》引理
\ref{morphisms-lemma-closed-immersion-finite-presentation}。
特别地，$U$-容许爆破是射影态射，见引理
\ref{divisors-lemma-blowing-up-projective}。
注意，产生同一态射的中心可以不止一个。因此，要求仅仅是存在某个与
$U$ 不交的中心，它产生 $X'$。最后，由于态射 $b : X' \to X$ 在 $U$ 上
是同构（见引理 \ref{divisors-lemma-blowing-up-gives-effective-Cartier-divisor}），
我们常常滥用记号，也把 $U$ 看作 $X'$ 的开子概形。

\begin{lemma}
\label{divisors-lemma-composition-admissible-blowups}
\begin{slogan}
容许爆破对复合稳定。
\end{slogan}
设 $X$ 为拟紧且拟分离概形。设 $U \subset X$ 为拟紧开子概形。
设 $b : X' \to X$ 为 $U$-容许爆破。设 $X'' \to X'$ 为
$U$-容许爆破。那么复合 $X'' \to X$ 是 $U$-容许爆破。
\end{lemma}

\begin{proof}
这是更精确的引理 \ref{divisors-lemma-composition-finite-type-blowups} 的直接推论。
\end{proof}

\begin{lemma}
\label{divisors-lemma-extend-admissible-blowups}
设 $X$ 为拟紧且拟分离概形。设 $U, V \subset X$ 为拟紧开子概形。
设 $b : V' \to V$ 为 $U \cap V$-容许爆破。则存在 $U$-容许爆破
$X' \to X$，它在 $V$ 上的限制为 $V'$。
\end{lemma}

\begin{proof}
设 $\mathcal{I} \subset \mathcal{O}_V$ 为有限型拟相干理想层，满足
$V(\mathcal{I})$ 与 $U \cap V$ 不交，并使 $V'$ 同构于 $V$ 沿
$\mathcal{I}$ 的爆破。设
$\mathcal{I}' \subset \mathcal{O}_{U \cup V}$ 为拟相干理想层，它在
$U$ 上的限制是 $\mathcal{O}_U$，在 $V$ 上的限制是 $\mathcal{I}$
（见《层》\ref{sheaves-section-glueing-sheaves} 节）。由《性质》引理
\ref{properties-lemma-extend}，存在有限型拟相干理想层
$\mathcal{J} \subset \mathcal{O}_X$，它在 $U \cup V$ 上的限制是
$\mathcal{I}'$。引理得证。
\end{proof}

\begin{lemma}
\label{divisors-lemma-dominate-admissible-blowups}
设 $X$ 为拟紧且拟分离概形。设 $U \subset X$ 为拟紧开子概形。
设 $b_i : X_i \to X$（$i = 1, \ldots, n$）为 $U$-容许爆破。
存在 $U$-容许爆破 $b : X' \to X$，使得
(a) $b$ 分解为 $X' \to X_i \to X$，其中 $i = 1, \ldots, n$；且
(b) 每个态射 $X' \to X_i$ 都是 $U$-容许爆破。
\end{lemma}

\begin{proof}
设 $\mathcal{I}_i \subset \mathcal{O}_X$ 为有限型拟相干理想层，满足
$V(\mathcal{I}_i)$ 与 $U$ 不交，并使 $X_i$ 同构于 $X$ 沿
$\mathcal{I}_i$ 的爆破。置
$\mathcal{I} = \mathcal{I}_1 \cdot \ldots \cdot \mathcal{I}_n$，并令
$X'$ 为 $X$ 沿 $\mathcal{I}$ 的爆破。由引理
\ref{divisors-lemma-blowing-up-two-ideals}，$X' \to X$ 通过 $b_i$ 分解。
\end{proof}

\begin{lemma}
\label{divisors-lemma-separate-disjoint-opens-by-blowing-up}
\begin{slogan}
通过爆破分离不可约分支。
\end{slogan}
设 $X$ 为拟紧且拟分离概形。设 $U, V$ 为 $X$ 的两个拟紧且不交的
开子概形。则存在一个 $U \cup V$-容许爆破 $b : X' \to X$，使得
$X'$ 是开子概形的不交并 $X' = X'_1 \amalg X'_2$，并且
$b^{-1}(U) \subset X'_1$ 及 $b^{-1}(V) \subset X'_2$。
\end{lemma}

\begin{proof}
选取有限型拟相干理想层 $\mathcal{I}$ 以及相应的 $\mathcal{J}$，使得
$X \setminus U = V(\mathcal{I})$ 以及相应的
$X \setminus V = V(\mathcal{J})$，见《性质》引理
\ref{properties-lemma-quasi-coherent-finite-type-ideals}。
于是集合论上 $V(\mathcal{I}\mathcal{J}) = X$，因而
$\mathcal{I}\mathcal{J}$ 是局部幂零理想层。由于 $\mathcal{I}$ 与
$\mathcal{J}$ 都是有限型的，并且 $X$ 拟紧，存在 $n > 0$ 使得
$\mathcal{I}^n \mathcal{J}^n = 0$。我们可以而且确实把 $\mathcal{I}$
换成 $\mathcal{I}^n$，把 $\mathcal{J}$ 换成 $\mathcal{J}^n$。于是
$\mathcal{I} \mathcal{J} = 0$。设 $b : X' \to X$ 为沿
$\mathcal{I} + \mathcal{J}$ 的爆破。它是 $U \cup V$-容许的，因为
$V(\mathcal{I} + \mathcal{J}) = X \setminus U \cup V$。我们将证明
$X'$ 是开子概形的不交并 $X' = X'_1 \amalg X'_2$，并且
$b^{-1}\mathcal{I}|_{X'_2} = 0$ 与 $b^{-1}\mathcal{J}|_{X'_1} = 0$；
这就证明了引理。

\medskip\noindent
我们将使用引理 \ref{divisors-lemma-blowing-up-affine} 中对爆破的描述。
假设 $U = \Spec(A) \subset X$ 是仿射开集，使得
$\mathcal{I}|_U$ 以及相应的 $\mathcal{J}|_U$，分别对应于有限生成理想
$I \subset A$ 以及相应的 $J \subset A$。于是
$$
b^{-1}(U) = \text{Proj}(A \oplus (I + J) \oplus (I + J)^2 \oplus \ldots)
$$
它由仿射开子集 $A[\frac{I + J}{x}]$ 与
$A[\frac{I + J}{y}]$ 覆盖，其中 $x \in I$ 且 $y \in J$。由于
$x \in I$ 在 $A[\frac{I + J}{x}]$ 中是非零因子，并且 $IJ = 0$，可知
$J A[\frac{I + J}{x}] = 0$。由于 $y \in J$ 在
$A[\frac{I + J}{y}]$ 中是非零因子，并且 $IJ = 0$，可知
$I A[\frac{I + J}{y}] = 0$。此外，
$$
\Spec(A[\textstyle{\frac{I + J}{x}}]) \cap
\Spec(A[\textstyle{\frac{I + J}{y}}]) =
\Spec(A[\textstyle{\frac{I + J}{xy}}]) = \emptyset
$$
因为 $xy$ 既是非零因子又为零。所以 $b^{-1}(U)$ 是下列两个开子概形的
不交并：$U_1$ 定义为标准开集
$\Spec(A[\frac{I + J}{x}])$（$x \in I$）之并；而 $U_2$ 是仿射开集
$\Spec(A[\frac{I + J}{y}])$（$y \in J$）之并。我们已经看到，
$b^{-1}\mathcal{I}\mathcal{O}_{X'}$ 在 $U_2$ 上的限制为零，并且
$b^{-1}\mathcal{I}\mathcal{O}_{X'}$ 在 $U_1$ 上的限制为零。
我们略去验证这些开子概形可以黏合成全局开子概形 $X'_1$ 与 $X'_2$。
\end{proof}

\begin{lemma}
\label{divisors-lemma-blowing-up-denominators}
设 $X$ 为局部 Noether 概形。设 $\mathcal{L}$ 为可逆
$\mathcal{O}_X$-模。设 $s$ 为 $\mathcal{L}$ 的正则亚纯截面。设
$U \subset X$ 为最大的开子概形，使得 $s$ 对应于 $\mathcal{L}$ 在 $U$ 上的
一个截面。爆破 $b : X' \to X$ 是沿 $s$ 的分母理想作的；它是
$U$-容许的。存在有效 Cartier 除子 $D \subset X'$ 以及同构
$$
b^*\mathcal{L} = \mathcal{O}_{X'}(D - E),
$$
其中 $E \subset X'$ 是例外除子，并且经由这个同构，亚纯截面
$b^*s$ 对应于亚纯截面 $1_D \otimes (1_E)^{-1}$。
\end{lemma}

\begin{proof}
由定义 \ref{divisors-definition-regular-meromorphic-ideal-denominators} 中分母理想的定义，
立即可见 $b$ 是 $U$-容许爆破。关于记号 $1_D$、$1_E$ 和
$\mathcal{O}_{X'}(D - E)$，请见定义
\ref{divisors-definition-invertible-sheaf-effective-Cartier-divisor}。
拉回 $b^*s$ 由引理 \ref{divisors-lemma-blow-up-pullback-effective-Cartier} 和
\ref{divisors-lemma-meromorphic-sections-pullback} 定义。因此引理的陈述是有意义的。
可把最后一个断言重新解释为：$b^*s$ 是
$b^*\mathcal{L}(E)$ 的整体正则截面，其零点概形为 $D$。
这唯一确定了 $D$，所以为了证明引理，可以在 $X$ 与 $X'$ 上仿射局部地工作。
假设 $X = \Spec(A)$ 为仿射概形，并且
$\mathcal{L} = \mathcal{O}_X$。于是 $s$ 是正则亚纯函数；再作缩小后，
可假设 $s = a'/a$，其中 $a', a \in A$ 都是非零因子。
于是 $s$ 的分母理想对应于理想
$I = \{x \in A \mid xa' \in aA\}$。回忆，$X'$ 由仿射爆破代数的谱覆盖，
这些代数为 $A' = A[\frac{I}{x}]$，其中 $x \in I$
（引理 \ref{divisors-lemma-blowing-up-affine}）。固定 $x \in I$，并写成
$xa' = a a''$，其中某个 $a'' \in A$。除子 $E \subset X'$ 由
$x \in A'$ 截出（这是在 $A'$ 的谱上），因而 $1/x$ 是
$\mathcal{O}_{X'}(E)$ 在 $\Spec(A')$ 上的生成元。最后，在 $A'$ 的全商环中，
有 $a'/a = a''/x$。所以 $b^*s = a'/a$ 限制为
$\mathcal{O}_{X'}(E)$ 的正则截面；在 $\Spec(A')$ 上，这个截面由
$a''/x$ 给出。这就完成了证明。
（除子 $D \cap \Spec(A')$ 由 $a''$ 在 $A'$ 中的像截出。）
\end{proof}

\section{爆破与平坦性}
\label{divisors-section-blowup-flat}

\noindent
我们继续《代数进阶》\ref{more-algebra-section-blowup-flat} 节中开始的讨论。
我们还将在《平坦性进阶》\ref{flat-section-blowup-flat} 节中证明进一步的结果。


\begin{lemma}
\label{divisors-lemma-strict-transform-blowup-fitting-ideal}
设 $S$ 为概形。设 $\mathcal{F}$ 为有限型拟相干
$\mathcal{O}_S$-模。设 $Z_k \subset S$ 为由
$\text{Fit}_k(\mathcal{F})$ 截出的闭子概形，见
\ref{divisors-section-fitting-ideals} 节。设 $S' \to S$ 为 $S$ 沿 $Z_k$ 的爆破，
并设 $\mathcal{F}'$ 为 $\mathcal{F}$ 的严格变换。那么
$\mathcal{F}'$ 局部可由 $\leq k$ 个截面生成。
\end{lemma}

\begin{proof}
回忆，$\mathcal{F}'$ 局部可由 $\leq k$ 个截面生成，当且仅当
$\text{Fit}_k(\mathcal{F}') = \mathcal{O}_{S'}$，见引理
\ref{divisors-lemma-fitting-ideal-generate-locally}。所以本引理是《代数进阶》引理
\ref{more-algebra-lemma-blowup-fitting-ideal} 的翻译。
\end{proof}

\begin{lemma}
\label{divisors-lemma-strict-transform-blowup-fitting-ideal-locally-free}
设 $S$ 为概形。设 $\mathcal{F}$ 为有限型拟相干
$\mathcal{O}_S$-模。设 $Z_k \subset S$ 为由
$\text{Fit}_k(\mathcal{F})$ 截出的闭子概形，见
\ref{divisors-section-fitting-ideals} 节。假设 $\mathcal{F}$ 局部自由且秩为 $k$，
这是在 $S \setminus Z_k$ 上。设 $S' \to S$ 为 $S$ 沿 $Z_k$ 的爆破，并设
$\mathcal{F}'$ 为 $\mathcal{F}$ 的严格变换。那么 $\mathcal{F}'$ 局部自由且
秩为 $k$。
\end{lemma}

\begin{proof}
这是《代数进阶》引理
\ref{more-algebra-lemma-blowup-fitting-ideal-locally-free} 的翻译。
\end{proof}

\begin{lemma}
\label{divisors-lemma-blowup-fitting-ideal}
设 $X$ 为概形。设 $\mathcal{F}$ 为有限呈示的
$\mathcal{O}_X$-模。设 $U \subset X$ 为概形论稠密开集，使得
$\mathcal{F}|_U$ 为常秩 $r$ 的有限局部自由模。则：
\begin{enumerate}
\item 爆破 $b : X' \to X$ 是对 $X$ 沿第 $r$ 个 Fitting 理想
（属于 $\mathcal{F}$）作的，并且是 $U$-容许的；
\item 严格变换 $\mathcal{F}'$（取自 $\mathcal{F}$，关于 $b$）局部自由且秩为 $r$；
\item 核 $\mathcal{K}$（对应于满射 $b^*\mathcal{F} \to \mathcal{F}'$）
有限呈示，并且 $\mathcal{K}|_U = 0$；
\item $b^*\mathcal{F}$ 与 $\mathcal{K}$ 都是完美
$\mathcal{O}_{X'}$-模，其 Tor 维数 $\leq 1$。
\end{enumerate}
\end{lemma}

\begin{proof}
由引理 \ref{divisors-lemma-fitting-ideal-of-finitely-presented}，理想
$\text{Fit}_r(\mathcal{F})$ 是有限型的；由引理
\ref{divisors-lemma-fitting-ideal-finite-locally-free}，它在 $U$ 上的限制等于
$\mathcal{O}_U$。所以 $b : X' \to X$ 是 $U$-容许的，见定义
\ref{divisors-definition-admissible-blowup}。

\medskip\noindent
由引理 \ref{divisors-lemma-fitting-ideal-finite-locally-free}，
$\text{Fit}_{r - 1}(\mathcal{F})$ 在 $U$ 上的限制为零；由于 $U$ 概形论稠密，
可知 $\text{Fit}_{r - 1}(\mathcal{F}) = 0$ 在整个 $X$ 上成立。因此，再由引理
\ref{divisors-lemma-fitting-ideal-finite-locally-free}，$\mathcal{F}$ 局部自由且秩为 $r$，
这是在下述子概形的补集上：该子概形由第 $r$ 个 Fitting 理想
（属于 $\mathcal{F}$）截出（这个补集可能大于 $U$，这正是论证中必须完成这一步的原因）。
所以，由引理 \ref{divisors-lemma-strict-transform-blowup-fitting-ideal-locally-free}，
严格变换
$$
b^*\mathcal{F} \longrightarrow \mathcal{F}'
$$
局部自由且秩为 $r$。这个映射的核 $\mathcal{K}$ 支集在爆破 $b$ 的
例外除子上，因而 $\mathcal{K}|_U = 0$。最后，由于
$\mathcal{F}'$ 有限局部自由，并且所示箭头是满射，可以在 $X'$ 上局部地把
$b^*\mathcal{F}$ 写成 $\mathcal{K}$ 与 $\mathcal{F}'$ 的直和。
由于 $b^*\mathcal{F}'$ 是有限呈示的
（《模》引理 \ref{modules-lemma-pullback-finite-presentation}），
$\mathcal{K}$ 也如此。

\medskip\noindent
关于 Tor 维数的陈述由《代数进阶》引理
\ref{more-algebra-lemma-fitting-ideals-and-pd1} 得出。
\end{proof}

\section{修改}
\label{divisors-section-modifications}

\noindent
本节将汇集如下类型的结果：作一次修改之后，某某性质成立。以后将看到，
修改可由一次爆破支配（《平坦性进阶》引理
\ref{flat-lemma-dominate-modification-by-blowup}）。

\begin{lemma}
\label{divisors-lemma-filter-after-modification}
设 $X$ 为整概形。设 $\mathcal{E}$ 为有限局部自由
$\mathcal{O}_X$-模。存在修改 $f : X' \to X$，使得
$f^*\mathcal{E}$ 有一个滤过，其逐次商都是可逆 $\mathcal{O}_{X'}$-模。
\end{lemma}

\begin{proof}
关于秩 $r$ 作归纳证明，其中这是 $\mathcal{E}$ 的秩。如果 $r = 1$ 或 $r = 0$，
引理显然成立。假设 $r > 1$。令
$P = \mathbf{P}(\mathcal{E})$，其结构态射为 $\pi : P \to X$，见《构造》
\ref{constructions-section-projective-bundle} 节。于是 $\pi$ 是固有的
（引理 \ref{divisors-lemma-relative-proj-proper}）。有典范满射
$$
\pi^*\mathcal{E} \to \mathcal{O}_P(1)
$$
其核有限局部自由且秩为 $r - 1$。选取非空开子概形
$U \subset X$，使得 $\mathcal{E}|_U \cong \mathcal{O}_U^{\oplus r}$。
于是 $P_U = \pi^{-1}(U)$ 同构于 $\mathbf{P}^{r - 1}_U$。
特别地，存在截面 $s : U \to P_U$（属于 $\pi$）。设
$X' \subset P$ 为态射 $U \to P_U \to P$ 的概形论像。那么
$X'$ 是整概形（《态射》引理
\ref{morphisms-lemma-scheme-theoretic-image-reduced}），态射
$f = \pi|_{X'} : X' \to X$ 是固有的（《态射》引理
\ref{morphisms-lemma-closed-immersion-proper} 和
\ref{morphisms-lemma-composition-proper}），并且
$f^{-1}(U) \to U$ 是同构。所以 $f$ 是修改
（《态射》定义 \ref{morphisms-definition-modification}）。
由构造，拉回 $f^*\mathcal{E}$ 有一个二步滤过；它的商可逆，因为该商等于
$\mathcal{O}_P(1)|_{X'}$；它的子模 $\mathcal{E}'$ 局部自由且秩为
$r - 1$。由归纳假设，可以找到修改 $g : X'' \to X'$，使得
$g^*\mathcal{E}'$ 有如引理陈述中的滤过。因此
$f \circ g : X'' \to X$ 就是所求的修改。
\end{proof}

\begin{lemma}
\label{divisors-lemma-extend-rational-map-after-modification}
设 $S$ 为概形。设 $X$、$Y$ 为 $S$ 上概形。假设 $X$ 是 Noether 的，
并且 $Y$ 在 $S$ 上固有。给定 $S$-有理映射
$f : U \to Y$，其源概形为 $X$、目标为 $Y$，则存在态射 $p : X' \to X$ 与
$S$-态射 $f' : X' \to Y$，使得：
\begin{enumerate}
\item $p$ 是固有的，并且 $p^{-1}(U) \to U$ 是同构；
\item $f'|_{p^{-1}(U)}$ 等于 $f \circ p|_{p^{-1}(U)}$。
\end{enumerate}
\end{lemma}

\begin{proof}
记包含态射 $j : U \to X$。设 $X' \subset Y \times_S X$ 为
$(f, j) : U \to Y \times_S X$ 的概形论像（《态射》定义
\ref{morphisms-definition-scheme-theoretic-image}）。投影
$g : Y \times_S X \to X$ 是固有的（《态射》引理
\ref{morphisms-lemma-base-change-proper}）。态射 $p : X' \to X$ 是
$X' \to Y \times_S X$ 与 $g$ 的复合；它是固有的（《态射》引理
\ref{morphisms-lemma-closed-immersion-proper} 和
\ref{morphisms-lemma-composition-proper}）。由于 $g$ 是分离的，并且
$U \subset X$ 是逆紧的（因为 $X$ 是 Noether 的），由《态射》引理
\ref{morphisms-lemma-scheme-theoretic-image-of-partial-section} 可知
$p^{-1}(U) \to U$ 是同构。另一方面，$f' : X' \to Y$ 是
$X' \to Y \times_S X$ 与投影 $Y \times_S X \to Y$ 的复合；它与
$f$ 一致，这是在 $p^{-1}(U)$ 上。
\end{proof}
